%% manual-kindle.tex -- the same manual, at Kindle screen geometry.
%%
%% Build with `make kindle'; the result is manual-kindle.pdf, which is copied
%% to the device as-is.  There is no second copy of the text: this file sets
%% three switches and then reads manual.tex, which carries two \ifdefined
%% \VNBkindle blocks (page geometry, and the verbatim font) and is otherwise
%% shared with the ordinary build.
%%
%% Why not post-process manual.pdf with k2pdfopt or a crop-and-scale script:
%% those reflow the page, and the manual is full of verbatim S-expressions
%% whose line breaks and indentation carry the meaning.  Re-typesetting at the
%% target size keeps them intact, and keeps the text as vector glyphs.
%%
%% Page: 105 x 140 mm is the 3:4 of a Paperwhite-class screen (1236 x 1648 px),
%% so the device scales it without letterboxing.  For a different device, edit
%% the geometry in manual.tex part 1 and \VNBkindlecode below.
%%
%% \VNBkindlecode is sized so that the manual's widest verbatim line -- 87
%% columns, in ch-proofs -- still fits the 97 mm text block.  Raise it and the
%% code blocks start running into the margin; the LaTeX log reports each one as
%% an Overfull \hbox, so the check is mechanical.

\PassOptionsToClass{9pt,oneside,openany}{amsbook}
\def\VNBkindle{}
\def\VNBkindlecode{\fontsize{5.3pt}{6.4pt}\selectfont}
\def\VNBkindletab{\fontsize{5.3pt}{6.4pt}\selectfont}

\documentclass{amsbook}[openany]
\renewcommand\thesection{\thechapter.\arabic{section}}
\usepackage[T1]{fontenc}
\usepackage{amsmath,amssymb}
\usepackage{textcomp}
\usepackage{fancyvrb}
\usepackage{upquote}
%% ---- Kindle build, part 1 of 2: page geometry -----------------------------
%% manual-kindle.tex defines \VNBkindle and then \input's this file, so both
%% kindle blocks are inert in the ordinary build.  geometry goes here, BEFORE
%% hyperref, which is the supported order.  Part 2 (the code font) is below,
%% after the verbatim environments it overrides.
\ifdefined\VNBkindle
  \usepackage[paperwidth=105mm,paperheight=140mm,
              top=5mm,bottom=5mm,left=4mm,right=4mm]{geometry}
\fi
%% INDEX.  No \usepackage{makeidx}: amsbook already provides \see, \seealso,
%% the theindex environment and \printindex (amsbook.cls:1329,1190,1331), and
%% makeidx re-\newcommand's \see -- "Command \see already defined", a fatal
%% error.  \makeindex itself is in the LaTeX kernel, so this line is all that
%% is needed.  Entries come from the generators; see \printindex at the end.
\makeindex
\usepackage{hyperref}
\usepackage[protrusion=true,expansion=false]{microtype}
\usepackage{array}
\usepackage{booktabs}
\usepackage{longtable}
\usepackage{enumitem}
%% List style, set ONCE here rather than per-list.  The manual had accumulated
%% six different labels for the same thing -- $\circ$ (36 lists, in three
%% spellings), $-$ (13), \textendash (4), $\ast$ (3), $\checkmark$ (1) -- each
%% carrying its own bracketed option block, plus two bare lists that fell back
%% to the amsbook bullet.  Levels are now distinguished by INDENTATION, with a
%% dash rather than a bullet at both: \textendash and not $-$, because $-$ is a
%% math minus, set in math italic at math spacing.
\setlist[itemize,1]{label=\textendash, leftmargin=0.85cm, itemsep=2mm}
\setlist[itemize,2]{label=\textendash, leftmargin=0.85cm, itemsep=1mm,
                    topsep=1mm}
\newtheorem{thm}{Theorem}[chapter]
\newtheorem{prop}[thm]{Proposition}
\newtheorem{cor}[thm]{Corollary}
\newtheorem{lem}[thm]{Lemma}
\newtheorem{claim}{Claim}

\theoremstyle{definition}
\newtheorem{dfn}{Definition}[chapter]
\newtheorem{example}{Example}[chapter]
\theoremstyle{remark}
\newtheorem{remark}{Remark}[chapter]
\theoremstyle{remark}
\pdfglyphtounicode{visiblespace}{0020}
\pdfgentounicode=1

%% Code blocks are set off by INDENTATION, not by a frame.  VNBcode used to be
%% Verbatim with frame=single (91 boxes) and VNBsnip the identical environment
%% without the frame (49 uses); nothing in the text distinguished them, so the
%% choice between box and no-box across 140 blocks was editorial noise.  They
%% are now one style.  The indent matches the lists' leftmargin (0.85cm), so
%% code and list items hang on the same left edge.
\DefineVerbatimEnvironment{VNBcode}{Verbatim}{xleftmargin=0.85cm,fontsize=\footnotesize}
\DefineVerbatimEnvironment{VNBsnip}{Verbatim}{xleftmargin=0.85cm,fontsize=\footnotesize}

%% ---- Kindle build, part 2 of 2: the code font -----------------------------
%% The widest verbatim line in the manual is 87 columns, and a code block is
%% precisely what a reflow tool (k2pdfopt and friends) mangles.  So the code
%% font is shrunk until 87 columns fit the text block instead; \VNBkindlecode
%% is set in manual-kindle.tex, where it can be retuned without touching this
%% file.  Prose keeps the ordinary body size.
\ifdefined\VNBkindle
  %% Same unframed style as the ordinary build, at the Kindle code size and a
  %% narrower indent -- 0.85cm of a 97 mm measure is a lot to give away.
  \DefineVerbatimEnvironment{VNBcode}{Verbatim}
    {xleftmargin=3mm,fontsize=\VNBkindlecode}
  \DefineVerbatimEnvironment{VNBsnip}{Verbatim}
    {xleftmargin=3mm,fontsize=\VNBkindlecode}
  \setlist[itemize,1]{label=\textendash, leftmargin=5mm, itemsep=1mm}
  \setlist[itemize,2]{label=\textendash, leftmargin=5mm, itemsep=0.5mm,
                      topsep=0.5mm}
  %% A 97 mm measure is narrower than any identifier this manual likes to name.
  %% `hyphenat[htt]' lets a long \texttt token break; the emergency stretch lets
  %% TeX buy a loose line rather than run into the margin; and the tables, which
  %% were laid out for a 345 pt measure, are set at \VNBkindletab.
  \usepackage[htt]{hyphenat}
  \usepackage{etoolbox}
  \setlength{\emergencystretch}{3em}
  \AtBeginEnvironment{tabular}{\VNBkindletab}
  %% BEFORE, not AtBegin, so the size covers the table's header rows too.
  \BeforeBeginEnvironment{longtable}{\VNBkindletab}
  \setlength{\tabcolsep}{2pt}
  %% Displayed formulas: the manual's are long enough that a dozen of them run
  %% clean off a 105 mm page -- not merely into the margin, off the paper, with
  %% the right-hand side of the equation gone.  So \[ ... \] is boxed, measured,
  %% and scaled down only when it does not fit.  A display that fits is centred
  %% and untouched, so nothing changes but the offenders.
  \usepackage{graphicx}
  \newsavebox{\VNBdispbox}
  \renewcommand{\[}{\par\smallskip\begin{lrbox}{\VNBdispbox}$\displaystyle}
  \renewcommand{\]}{$\end{lrbox}%
    \noindent
    \ifdim\wd\VNBdispbox>\linewidth
      \resizebox{\linewidth}{!}{\usebox{\VNBdispbox}}%
    \else
      \hbox to \linewidth{\hfil\usebox{\VNBdispbox}\hfil}%
    \fi
    \par\smallskip}
\fi

\title{\textbf{VNB Proof Checker}\\[4pt]
       \large User Manual}
\author{F. J. Thayer}
\date{}

\begin{document}
\maketitle
\tableofcontents
\newpage

%% =========================================================
\chapter{Introduction}

\begin{sloppypar}The VNB Proof Checker is an interactive proof assistant implemented in
MIT~Scheme.  Its logical foundation is a first-order set theory in the
style of von~Neumann--Bernays (VNB): a single-sorted universe of
\emph{classes}, with \textsc{set} as a distinguished predicate and
membership ($\in$) as the only non-logical primitive.\end{sloppypar}
[[NOTE: References needed. NO HALLUCINATIONS please]]

The proof architecture follows the sequent-calculus approach of the
IMPS~2.0 system~\cite{IMPS}: proofs are \emph{deduction graphs} whose
nodes are sequents and whose edges are inference steps.  Every
inference step is validated by a small, auditable \emph{primitive
inference kernel}; higher-level tactics and a form of conditional
rewrite rules, called macetes in the IMPS system, are built on top of
the inference kernel and cannot introduce unsoundness.
The deduction-graph proof structure and the use of local contexts in
macete application are due to L.~G.~Monk~\cite{Monk1988}, whose work
was a direct influence on IMPS.
[[Reference to IMPS and an online Portuguese language dictionary
(e.g.\ \url{https://dicionario.acad-ciencias.pt/pesquisa/macete/} has
an entry)]]

This project is, among other things, a proof of concept: a rigorous
mathematical proof system built entirely by AI
[[Note: we can reference documentation for Claude]],
with architectural direction from F.~J.~Thayer.  It is inspired by
ideas from IMPS~\cite{IMPS}.

%% =========================================================
\section{VNB: primitives and reductions}
\label{sec:primitives}

\emph{[References to VNB in the published literature will be added here
after a literature search.  No content will be added without accurate
citations.]}

Pure VNB set theory is a spare formalism.  The natural numbers, the
reals, functions, and ordinals can all be \emph{constructed} from the
axioms of set theory---but those constructions are long, intricate, and
bear little resemblance to how mathematics is actually written or thought
about.  A proof that $\sin(\pi/2) = 1$ should not need to unfold the
Dedekind-cut definition of the reals.

This system therefore starts from a broad primitive base.  VNB set
theory supplies the logical foundation, but many mathematical objects
are taken as \emph{primitive constants} and \emph{primitive
constructors}, governed by axioms rather than by explicit
set-theoretic construction:

\begin{itemize}[label = $\circ$, leftmargin= 0.9cm]
\item \textbf{Number systems} \texttt{NN}, \texttt{ZZ}, \texttt{QQ},
  \texttt{RR}, \texttt{CC} are primitive constants denoting sets.  Their
  elements and operations are characterised by ring, field, and order
  axioms, not by Dedekind cuts, Cauchy sequences, or ordinal arithmetic.
\item \textbf{Ordinals} \texttt{ORD} form a primitive proper class with
  its own ordering predicate \texttt{<=\_ORD}, independent of the
  set-membership relation.
\item \textbf{Function spaces} \texttt{fun($A$, $B$)} are primitive
  set constructors.  Functions are not encoded as sets of ordered
  pairs; application \texttt{$f$($x$)} is a primitive operation.
\item \textbf{Mathematical structures} such as groups, vector spaces,
  and Banach spaces are declared with \texttt{def-structure}, which
  installs accessor operations and a typing predicate.
\end{itemize}

In principle every one of these primitives could be reduced to the bare
VNB axioms---ordinals by the von~Neumann construction, reals by Dedekind
cuts, functions by sets of ordered pairs, and so on.  This approach
simply takes these reductions as \emph{existence theorems} that
have been proved (or could be proved) elsewhere, and works with the
objects directly at their natural level of abstraction.  The resulting
proofs are shorter, more readable, and correspond more closely to
informal mathematical practice.

A further departure from standard set theory is the treatment of
\emph{partial application}.  Any syntactic expression \texttt{f(x)}
is a well-formed term, but it may be \emph{undefined} if \texttt{f} is
not a function applicable to \texttt{x}.  The system provides both
strict equality \texttt{a = b} (false if either side is undefined) and
quasi-equality \texttt{a == b} (true when both sides are undefined in
the same way), following the approach of IMPS~\cite{IMPS}.

A consequence of this broad primitive base is that membership in a
primitive constant has \emph{indeterminate truth value} in general.
The formula \texttt{x in 7} is syntactically well-formed --- VNB is
single-sorted and \texttt{in} accepts any two class expressions ---
but the system makes no claim about whether it holds. Please note that
this does not mean that the formula is undefined: whether it is true
or not depends on how $7$ is interpreted in VNB. This system
takes $7$ as a primitive number constant and makes no identification
of it with any set-theoretic object.  Membership statements of this
kind are grammatically meaningful but proof-theoretically inert.

%% =========================================================
\section{Architecture}
\label{sec:architecture}

\subsection{Proof architecture}

\begin{sloppypar}
The proof architecture follows the sequent-calculus approach of the
IMPS~2.0 system~\cite{IMPS}: proofs are \emph{deduction graphs} whose
nodes are sequents and whose edges are inference steps.  Every inference
step is validated by a small, auditable \emph{primitive inference kernel};
higher-level tactics and macetes are built on top of it and cannot
introduce unsoundness.
\end{sloppypar}

\subsection{Theory architecture}

In contrast to the \emph{little theories} methodology of
IMPS~\cite{IMPS}---where each mathematical domain inhabits its own
separate mini-theory---VNB works with a \emph{single, growing current
theory}.  Mathematical structures are declared directly in that one theory
using \texttt{def-structure}, and their axioms accumulate there alongside
the base VNB axioms and the primitive number-system axioms.  This reflects
the view of mathematics as one unified body of knowledge rather than a
federation of isolated theories.

The system does permit \emph{minimal} or even \emph{odd-ball theories}
for the cases where a genuinely separate context becomes necessary.
These theories can be created with \texttt{make-theory} and carry
the primitive class constants but \emph{no structural axioms}.  Such
theories are expected to be exceptional; all routine mathematical work
takes place in the single current theory.

\begin{remark}\label{remark:27137-61828-986360}
%
For the potential user interested in \emph{constructivist} mathematics,
this system is not appropriate. These potential users should consider
alternatives such as Coq, which has a very well established library of
mathematics.
%
\end{remark}

\subsection{Design goals}

\begin{itemize}[ label = $\circ$, leftmargin = 1cm]
\item \textbf{Enriched set-theoretic foundation.}  VNB set theory with
  primitive constants for standard mathematical objects; no type
  hierarchy; partial function application with undefinedness.
\item \textbf{Transparent kernel.}  All proof steps reduce to a small
number of
  primitive inferences; the kernel is under 300 lines of Scheme.
\item \textbf{Composable rewriting.}  A macete mechanism (borrowed from
  IMPS) allows theorems to be automatically converted into named rewrite
  strategies that can be sequenced and combined.
\item \textbf{N-ary constructors.}  Ordered tuples \texttt{[$e_1$,
  \ldots, $e_n$]} and Cartesian products \texttt{cartesian($A_1$, \ldots,
  $A_n$)} are first-class, supporting the definition of mathematical
  structures with multiple carriers and operations.
\end{itemize}

\subsection{Why MIT Scheme?}

MIT~Scheme was chosen for several practical reasons: it is small (the
proof kernel is under 300 lines of idiomatic Scheme), easy to learn
and write code, free and readily available on all major platforms, and has
a long track record of correctness and stability.  Since all code is
open source, porting to another language with AI assistance should be
straightforward.

%% =========================================================

%% =========================================================
\section{Quick start}
\label{sec:quickstart-chap}

\subsection{Installation and startup}

\subsubsection{Requirements}

MIT~Scheme~11.2 or later, and for best interactive use one of
GNU/Emacs or XEmacs.  With minor modifications, the native Scheme text
editor Edwin also works.  No other dependencies.

\subsubsection{Directory layout}

\begin{VNBsnip}
~/prover/
  prover                 executable script
  load.scm               loader (sources all files in order)
  expressions.scm        expression language
  wff.scm                <wff> record type, wff-child, wff-in-theory, wff-equiv?
  sequents.scm           sequents over <wff>; context operations
  deduction-graphs.scm   proof graph data structures
  primitive-inferences.scm  inference kernel
  macetes.scm            macete mechanism
  theory.scm             theory management + VNB base axioms
  structures.scm         mathematical structure definitions
  number-systems.scm     NN, ZZ, QQ, RR, CC constants and axioms
  contexts.scm           local context management, make-wff, unravel
  arith-eval.scm         ground arithmetic decision procedure
  proof-commands.scm     high-level proof commands
  interactive.scm        interactive short-form commands (sp, di, ai, ...)
  vnb.el                 GNU Emacs / XEmacs interface
\end{VNBsnip}

\subsection{Starting the checker}

\begin{VNBsnip}
~/prover/prover              # interactive REPL
~/prover/prover proof.scm    # run a proof script and exit
~/prover/prover -i proof.scm # run a script, then stay interactive
\end{VNBsnip}

The interactive session is a standard MIT~Scheme REPL with the proof
checker pre-loaded.
\subsection{The Emacs interface}
\label{sec:emacs}

The file \texttt{vnb.el} provides a GNU~Emacs and XEmacs interface for
the proof checker.  It requires no external packages beyond
\texttt{comint}, which ships with every Emacs distribution.

\subsubsection{Installation}

Add the following to your Emacs init file (\texttt{\textasciitilde/.emacs}
or \texttt{\textasciitilde/.xemacs/init.el}):

\begin{verbatim}
(add-to-list 'load-path "~/prover")
(require 'vnb)
\end{verbatim}

\noindent
\textbf{Do not byte-compile \texttt{vnb.el}.}  GNU~Emacs and XEmacs use
incompatible byte-code formats; loading the source file interpreted is
the correct approach for both.

\subsubsection{Starting the interface}

\begin{VNBsnip}
M-x vnb
\end{VNBsnip}

This starts the prover subprocess (if not already running), and opens a
two-buffer layout:
\begin{itemize}[label=\textendash, leftmargin=1cm]
\item \textbf{Left window --- \texttt{*VNB*} REPL.}  A \texttt{comint}
  buffer running MIT~Scheme with the proof checker pre-loaded.  Type any
  Scheme expression and press \texttt{RET} to evaluate it.
\item \textbf{Right window --- \texttt{*VNB State*}.}  A read-only
  display buffer that is updated automatically after every command that
  calls \texttt{(show)}.  It always shows the current proof state: open
  goals, assumptions, and focus.
\end{itemize}

The two buffers are kept synchronised by a comint output filter that
watches for the sentinel lines \texttt{;;VNB-STATE-BEGIN} /
\texttt{;;VNB-STATE-END} emitted by \texttt{(show)} and routes the
enclosed text to the state buffer.

\subsubsection{Key bindings in the \texttt{*VNB*} REPL buffer}

\begin{itemize}[label = $\circ$, leftmargin=1cm]
\item \texttt{RET} Submit expression to the prover (standard comint).
\item \texttt{C-c C-s} Send \texttt{(show)} to refresh the state display.
\item \texttt{C-c C-p} Switch to the \texttt{*VNB State*} buffer.
\end{itemize}

\subsubsection{Key bindings in the \texttt{*VNB State*} buffer}

\begin{itemize}[label = $\circ$, leftmargin=1cm]
\item \texttt{q} Return to the \texttt{*VNB*} REPL buffer.
\item \texttt{C-c C-s} Send \texttt{(show)} to refresh the state display.
\end{itemize}

\subsubsection{The \texttt{*VNB Commands*} scratch buffer}

\begin{VNBsnip}
M-x vnb-command-buffer
\end{VNBsnip}

Opens a \texttt{*VNB Commands*} buffer in \texttt{vnb-command-mode}.
This buffer works exactly like the Emacs \texttt{*scratch*} buffer (which
uses \texttt{lisp-interaction-mode}), except that \texttt{C-j} sends
expressions to the \emph{VNB prover} instead of evaluating them as Emacs
Lisp.

Write proof commands as S-expressions in the buffer.  Position point
immediately after any complete expression and press \texttt{C-j}:
the expression is sent to the running VNB process, and the return value
is inserted on a new line below point.  The buffer accumulates a
readable session transcript.

\begin{VNBcode}
(sp (make-wff-from-string "n in NN and m in NN implies m in NN and n in NN"))
;;VNB-STATE-BEGIN
1 open goal(s).
Focus:
  [0] ()  =>  (((n in nn) and (m in nn)) implies ((m in nn) and (n in nn)))
;;VNB-STATE-END

(di)       C-j
;;VNB-STATE-BEGIN
2 open goal(s).
Focus:
  [1] ((n in nn) and (m in nn))  =>  ((m in nn) and (n in nn))
Other open goals:
  [0] ()  =>  (((n in nn) and (m in nn)) implies ((m in nn) and (n in nn)))
;;VNB-STATE-END
\end{VNBcode}

\paragraph{Key bindings in \texttt{vnb-command-mode}.}

\begin{itemize}[label = $\circ$, leftmargin=1cm]
\item \texttt{C-j} Send S-expression before point to VNB; insert result below.
\item \texttt{C-c C-c} Send entire buffer to VNB; append result.
\item \texttt{TAB} Complete VNB command name at point.
\item \texttt{C-c C-d} Describe VNB command: show signature and documentation.
\item \texttt{C-c C-s} Refresh \texttt{*VNB State*} proof-state buffer.
\end{itemize}

\subsubsection{Command reference and completion}

\begin{sloppypar}\texttt{vnb.el} maintains the list \texttt{vnb-commands} of all proof
commands, each with a usage signature and a one-line description.  The
derived alist \texttt{vnb-commands-alist} maps short names to signatures
and is used by \texttt{TAB} completion in \texttt{vnb-command-mode}.\end{sloppypar}

\begin{VNBsnip}
M-x vnb-describe-command   RET   di   RET
; Echo area: (di)
;   Direct inference: decompose current goal by its top connective. ...
\end{VNBsnip}

\subsubsection{Programmatic evaluation: \texttt{vnb-eval-string}}

For scripting or integration with other Emacs packages:

\begin{VNBsnip}
; str is a Scheme expression as a string
(vnb-eval-string str)           ; => Scheme return value as string
(vnb-eval-string str timeout)   ; => same, with explicit timeout (seconds)
\end{VNBsnip}

\texttt{(vnb-eval-string STR \&optional TIMEOUT)} sends the string
\texttt{STR} to the running prover, waits up to \texttt{TIMEOUT} seconds
(default 30) for the next REPL prompt, and returns the Scheme return
value as an Emacs string.  It signals an error if the prover is not
running.

%% =========================================================
\subsection{A complete proof}
\label{sec:quickstart}

\begin{sloppypar}
After installing the Emacs interface (Section~\ref{sec:emacs}), launch
the prover with \texttt{M-x vnb}.  Two buffers appear: \texttt{*VNB*}
(the REPL) and \texttt{*VNB State*} (the live proof state).  Type the
following five commands at the \texttt{*VNB*} prompt, pressing
\texttt{C-j} (or RET) after each:
\end{sloppypar}

\begin{VNBsnip}
(sp (make-wff-from-string "n in NN and m in NN implies m in NN and n in NN"))
(di)                                   ; assume both, prove (m in nn) and (n in nn)
(ai "n in NN and m in NN")            ; split into (n in nn) and (m in nn)
(di)                                   ; split goal: prove (m in nn), then (n in nn)
(ass) (ass)                            ; both are in the assumption list
\end{VNBsnip}

The \texttt{*VNB State*} buffer updates after each command.  The final
state shows \texttt{Proof complete.}

\paragraph{What just happened.}
\texttt{(make-wff-from-string \ldots)} parses a string in the VNB string
syntax and wraps the result as a typed \texttt{wff} in the current
theory.  \texttt{(sp \ldots)} starts a new proof.  \texttt{(di)} is
direct-inference: it decomposes the goal by its outermost connective.
\texttt{(ai \ldots)} is antecedent-inference: it decomposes a named
assumption; the argument is a formula string (or a raw S-expression for
low-level use) that identifies the assumption by alpha-equivalence.
\texttt{(ass)} discharges a goal that already appears in the assumption
list.

The same proof at the Scheme level (no short forms) reads:
\begin{VNBsnip}
(let* ((wic (make-wff-from-string
              "n in NN and m in NN implies m in NN and n in NN"))
       (ps  (start-proof wic))
       (ps  (cmd-direct-inference ps))
       (ps  (cmd-antecedent-inference ps (parse-string "n in NN and m in NN")))
       (ps  (cmd-direct-inference ps))
       (ps  (cmd-assumption ps))
       (ps  (cmd-assumption ps)))
  (print-proof-state ps))
\end{VNBsnip}
The interactive short forms in \texttt{interactive.scm} are simple
wrappers: each mutates a global \texttt{*ps*} and calls \texttt{(show)}
to reprint the state.  All examples in this manual use the short forms;
the long form is available when scripting non-interactively.

\paragraph{What \texttt{make-wff} checks.}
\texttt{make-wff} (and \texttt{make-wff-from-string}, which calls it)
performs a position-typed walk of its argument and signals an error
rather than wrap a malformed expression.  The checks are:
\begin{sloppypar}
\begin{itemize}[label=$\ast$, leftmargin=1cm]
\item \textbf{Numbers in wff positions} are rejected:
  \texttt{(make-wff-from-string "1 and 35")} $\Rightarrow$ ``number in wff
  position 1''.
\item \textbf{Term-forming heads in wff positions} are rejected:
  \texttt{(make-wff-from-string "choice(A)")} as a top-level wff fails because
  \texttt{choice} produces a term, not a wff.
\item \textbf{Wff-forming heads in term positions} are rejected:
  \texttt{(make-wff-from-string "(A and B) in C")} fails because \texttt{A and B}
  is a wff and the first argument of \texttt{in} must be a term.
\item \textbf{Bare symbols in wff positions} are rejected.  Only
  \texttt{truth} and \texttt{falsity} are valid bare-symbol wffs;
  anything else (e.g.\ \texttt{A implies B}) is an error.
\item \textbf{Free-symbol role conflicts}: a symbol that appears free
  in both wff and term positions in the same formula is rejected.
\item \textbf{Wff-position binders}: a binder whose body uses the
  bound variable in wff position is rejected within its scope.  This
  catches \texttt{forall([p], p implies p)}: \texttt{p} is bound as a
  class variable but the body uses it as a proposition.
\item \textbf{Bound vs.\ free same-name use} is \emph{accepted with a
  warning}.  In \texttt{forall([a], a in nn) implies a in zz} the
  inner \texttt{a} is bound and the outer \texttt{a} is free.  The
  validator emits a \texttt{;; warning} but does not reject.
\item \textbf{Arities and bound-variable shapes} are checked: every
  connective and quantifier gets the right number of arguments;
  \texttt{forall}, \texttt{forsome}, \texttt{iota},
  \texttt{lambda}, \texttt{lambdoid}, and \texttt{sep} bound variables must be
  symbols.
\end{itemize}
\end{sloppypar}
%
\begin{remark}\label{remark:27135-34761-138906}
%
The set-theoretic \texttt{power(A)} (power set) and arithmetic
\texttt{x \^{} n} (exponentiation) share the same underlying symbol because
MIT~Scheme reads identifiers case-insensitively; \texttt{make-wff}
distinguishes them by arity (1 vs.\ 2 arguments).  All downstream
operations (\texttt{free-vars}, \texttt{subst-free},
\texttt{alpha-equiv?}, macete pattern matching) honour the same
arity-based dispatch.

%
\end{remark}

%% =========================================================
\subsection{A macete example}
\label{sec:quickstart-macete}

Macetes are named rewrite rules derived from installed theorems.
To apply a macete outside a proof, install a theorem and call
\texttt{apply-macete}:

\begin{VNBcode}
(install-theorem! 'eq-symm
  (make-wff-from-string "forall([x, y], x = y iff y = x)"))

(apply-macete 'eq-symm (make-wff-from-string "3 + 5 = 8"))
; => wff for "8 = 3 + 5"

(apply-macete 'eq-symm (make-wff-from-string "n = m implies truth"))
; => wff for "m = n implies truth"

(apply-macete 'eq-symm (make-wff-from-string "3 in nn"))
; => #f          -- pattern x = y does not match; no rewrite fires
\end{VNBcode}

\texttt{apply-macete} accepts and returns a \texttt{wff} object; it
returns \texttt{\#f} when the pattern matches nowhere.
To apply a macete as a proof step, use
\texttt{(mac \textit{name})} inside a proof (see
Section~\ref{sec:macetes}).

%% =========================================================


%% =========================================================


%% =========================================================
\chapter{Classification of results}
\label{ch:classification}

Every named result in the current theory carries a \emph{classification}.
It answers two independent questions, and keeping them apart is the whole
point of this chapter:

\begin{itemize}
\item \textbf{What justifies the statement?} --- is it a primitive
  assumed as a base, a definition true by construction, a theorem with a
  machine proof, or a result trusted on a stated warrant without proof?
\item \textbf{May the proof engine activate it unbidden?} --- whether the
  prover is allowed to apply the result as a rewrite rule during a proof.
\end{itemize}

\noindent
These two questions are \emph{orthogonal}.  The first sorts results into
the four kinds below; the second is a separate matter, discussed in
Section~\ref{sec:firing}.  A common confusion --- ``is this proved?'' ---
conflates them, and most of what follows is aimed at pulling them apart.

%% =========================================================
\section{The four kinds}
\label{sec:four-kinds}

\begin{itemize}
\item \textbf{Axiom.} A primitive of VNB --- the base we build on.  The
  axiom set is small and essentially fixed; see
  Section~\ref{sec:primitives}.  The only reason to add one is the
  discovery that the existing primitives do not capture the full power of
  VNB.  An axiom is assumed, not proved.

\item \textbf{Definition.} True \emph{by construction}: a defining
  equation, or an unfolding of one.  A definition is \emph{not a theorem}
  --- there is nothing to prove, since it holds by the meaning of the
  symbol it introduces.  Consequently it is never submitted to the VNB
  test and is \emph{not} a member of the Proof Support Set.  It does,
  however, still \textbf{fire as a rewrite} (Section~\ref{sec:firing}):
  the defining equation is exactly the rule that unfolds the symbol.

\item \textbf{Theorem.} Carries a machine proof.  Its proof script runs
  end-to-end and is re-certified by the VNB test
  (Section~\ref{sec:vnb-test}).  This is the only kind whose truth rests
  on a checked derivation inside the system.

\item \textbf{Proof Support Set (PSS).} A theorem \emph{trusted on a
  warrant, without a machine proof} --- and so \emph{excused} from the VNB
  test.  Every PSS member carries credentials: a plain statement of what
  it says plus a short justification of why it is accepted
  (Section~\ref{sec:warrants}).  The PSS is the curated foundational
  engine of obvious-but-tedious facts (finite-sum identities, ordinal and
  set commonplaces) that are not in doubt and not worth the cost of
  mechanising.  A second motive is reuse: a proof typically calls on many
  small, well-established intermediate facts that recur across unrelated
  developments, and collecting them in the PSS makes each available
  everywhere at once rather than re-derived context after context.
\end{itemize}

\begin{remark}
The PSS exists for a practical reason.  Pushing into new
mathematical territory, a proof constantly bumps into basic --- and
genuinely important --- facts whose proofs are well known but tedious and
time-consuming to mechanise.  Stopping to grind each one out would stall
the advance.  The PSS is the mechanism for brushing them aside: assert the
fact on a warrant, keep moving.  That is exactly why it was introduced.
\end{remark}

\begin{remark}
``Asserted'' is not a fifth kind.  A result installed without a proof is
debt against the fixed axiom base: it is destined either to be discharged
(promoted to a \emph{theorem} with a proof) or recognised as a genuine
primitive gap (an \emph{axiom}), or --- if it is an obvious support fact
--- credentialed into the \emph{PSS}.  An asserted result carrying no
warrant at all is a backlog item, not a stable category.
\end{remark}

%% =========================================================
\section{Activation versus justification}
\label{sec:firing}

Whether the prover may apply a result as a rewrite rule is
\emph{independent} of the four-way classification above.  Axioms,
definitions, theorems, and PSS members can \emph{all} be rewrite rules,
and any of them can equally be inert (cited only by name).  Activation is
the business of the \emph{macete} mechanism, a separate registry described
in Section~\ref{sec:macetes}: installing a result as a macete with
\texttt{install-theorem!} is what lets the engine fire it, and that is a
decision about \emph{use}, not about what justifies the statement.

The clearest case is the definition: it is not a PSS member, yet its
defining equation is precisely the rewrite that unfolds the defined
symbol.  Conversely, a fully proved theorem may be left inert if it is too
specialised to be worth firing automatically.

\begin{center}
\begin{tabular}{@{}llll@{}}
\toprule
\textbf{Kind} & \textbf{Believed because} & \textbf{In VNB test?} & \textbf{May fire?}\\
\midrule
Axiom      & assumed (primitive)   & no  & yes\\
Definition & true by construction  & no  & yes\\
Theorem    & machine proof         & yes & yes\\
PSS        & a stated warrant      & no (excused) & yes\\
\bottomrule
\end{tabular}
\end{center}

\noindent
The last column is the same ``yes'' in every row: firing rights cut across
the classification entirely.

%% =========================================================
\section{Warrants}
\label{sec:warrants}

A definition needs no warrant: it is true by construction, so there is
nothing to justify.  A PSS member, by contrast, is an asserted theorem ---
something that \emph{could} be false and is not being proved here --- and
so it must carry credentials.  A \emph{warrant} is that justification: a
comprehensible statement of the fact together with at least a couple of
lines saying why we accept it without a machine proof.

A warrant is attached with

\begin{VNBsnip}
(warrant! 'name 'kind "free-text justification")
\end{VNBsnip}

\noindent
where \texttt{kind} is one of the recognised grades of justification (the
set \texttt{*warrant-kinds*}), weakest to strongest:

\begin{itemize}
\item \texttt{hand-wave} --- a heuristic or plausibility argument;
  convincing, not rigorous.
\item \texttt{well-known} --- a standard textbook fact, asserted without
  argument.
\item \texttt{reference} --- cites a specific source, named in the text.
\item \texttt{informal} --- a rigorous paper-proof exists (often sketched
  in the file comment or the text), just not mechanised in VNB.
\item \texttt{proof} --- a machine-checked VNB proof exists.
\end{itemize}

\begin{remark}
A warrant is a promissory note, not a permanent verdict: any weaker kind
may later be discharged into \texttt{proof}.  Nothing logical depends on
it --- it is metadata for review and triage.  A result with no warrant at
all ranks below even \texttt{hand-wave}.
\end{remark}

\begin{remark}
Warrants propagate.  When a biconditional theorem is installed, its
automatically generated \texttt{-rev} companion inherits the same warrant,
so the reverse-direction rewrite is credentialed identically to the
forward one.
\end{remark}

%% =========================================================
\section{Proof debt: bills and trust tiers}
\label{sec:proof-debt}

A machine proof is only as unconditional as the facts it cites.  When a
proof reaches \texttt{QED} the system computes its \emph{bill}: the set of
\emph{asserted} facts it ultimately rests on.  The bill is transitive.  A
cited \emph{proven} lemma contributes its own bill; a \emph{primitive}
axiom or a \emph{definition} contributes nothing (it is trusted base); an
\emph{asserted} fact --- a PSS member, or an as-yet-unwarranted assertion
--- contributes itself.  Each fact in the bill is a \emph{leaf}.

The bill is printed inline at \texttt{QED} and collected in the ledger
\texttt{reference/PROOF-DEBT.md}:

\begin{itemize}
\item \texttt{proven modulo 0} --- the \emph{empty bill}.  The proof rests
  only on the trusted base: the fixed primitive axioms, the kernel
  inference rules, and definitions.  Nothing asserted stands between the
  theorem and the kernel.  This is the strongest report a proof can make,
  and it carries \emph{no} trust tier.
\item \texttt{proven modulo \{a, b, \ldots\}} --- unconditional
  \emph{given} $a, b, \ldots$.  Discharging every leaf (proving it, or
  turning it into a definition) would collapse the bill to $0$.
\end{itemize}

\noindent
A non-empty bill is stamped with a \emph{trust tier}
\texttt{[trust:\ \textit{K}]}, and $K$ is the warrant of the bill's
\emph{weakest} leaf: a single shaky citation caps the whole proof, however
solid the rest.  The tiers are the warrant kinds of
Section~\ref{sec:warrants} with one addition below the weakest ---
\texttt{none}, an asserted leaf carrying no warrant at all --- so, weakest
to strongest:
\[
  \texttt{none} \;<\; \texttt{hand-wave} \;<\; \texttt{well-known}
  \;<\; \texttt{reference} \;<\; \texttt{informal} \;<\; \texttt{proof}.
\]
Thus \texttt{[trust:\ none]} is the \emph{weakest} thing a proof can
report --- some leaf justifies nothing --- and \texttt{modulo 0}, which
prints no tier at all, is the strongest.  The two are opposite ends of the
same scale, and are easily misread for one another; the ledger's header
carries this same legend for exactly that reason.

\begin{remark}
The tier is a moving target by design.  Discharging a \texttt{none} or
\texttt{hand-wave} leaf into a warranted or proven fact lifts the tier of
every proof that cites it.  The ledger's reverse index ranks the asserted
leaves by how many proofs each one caps, so the highest-leverage debts to
retire are read off the top.
\end{remark}

%% =========================================================
\section{Circular dependency}
\label{sec:proof-cycle}

The bill of Section~\ref{sec:proof-debt} is computed from a \emph{citation
graph}: at each \texttt{QED} the system records which named results the finished
proof cited.  That graph carries a soundness obligation of its own, independent
of the trust tiers: \textbf{no proven theorem may depend --- transitively,
through its citations --- on itself.}

The failure it guards against is a circle that no single step betrays.  Assert
$X$; prove $Y$ citing $X$; then ``prove'' $X$ citing $Y$.  Each proof is locally
valid, yet together they rest on nothing --- $X$ is believed because of $Y$, and
$Y$ because of $X$.  A bill, a frozen snapshot taken at \texttt{QED}, cannot
reveal such a loop, so the check walks the \emph{live} citation graph instead.

\texttt{(proof-cycle-check)} returns every dependency cycle among the proven
theorems --- the empty list when there are none.  \texttt{load.scm} runs it as a
load-time gate and, on any surviving cycle, \textbf{aborts the load} rather than
merely warn: a genuine circular proof is a soundness bug and is treated as one.

Two kinds of citation edge are deliberately \emph{not} dependencies, and are
excluded when the graph is recorded.  Without the exclusion an ordinary
well-founded induction would read as a cycle:

\begin{itemize}
\item \textbf{A proof's citation of its own name.}  A proof by induction applies
  its induction hypothesis --- the theorem's own statement at a smaller
  instance --- as a rewrite; likewise a structure predicate is unfolded by a
  macete bearing the very name the finished theorem will carry.  Neither is a
  dependence of the theorem on itself.
\item \textbf{A citation that is a definition when it is recorded.}  Unfolding a
  definition adds no debt (Section~\ref{sec:proof-debt}) and no dependency; a
  definitional macete that happens to share a later theorem's name is excluded
  on the same ground.
\end{itemize}

\noindent
Excluding these leaves every genuine circular dependency intact --- such a loop
runs through \emph{proven} citations, which are always kept --- so the gate
neither cries wolf on sound induction nor lets a real circle through.

%% =========================================================
\section{The VNB test}
\label{sec:vnb-test}

The \emph{VNB test} is the standing examination that certifies the
\textbf{theorem} bucket.  It re-proves, in sequence, every result that has
a runnable proof script, and reports whether each still reaches
\texttt{QED}.  Definitions and PSS members are excused by construction:
the former have nothing to prove, the latter are accepted on a warrant.
Run it from the prover root:

\begin{VNBsnip}
./vnb-test
\end{VNBsnip}

\noindent
It harvests the proofs that run at load time and the standalone proof
scripts, runs each, and writes a \emph{certification stamp}
(\texttt{reference/certification.scm}) recording the date and which results
passed, together with a human-readable report
(\texttt{reference/VNB-TEST.md}).  The next time the prover loads, the
stamp is read back, so a result's status can report ``certified by the VNB
test on \textit{date}'' rather than merely ``has a proof script.''

\begin{remark}
The VNB test is \emph{not} the unit test suite.
\texttt{test-suite.scm} is a fast feature-and-regression suite of several
hundred cases run constantly during development; it checks that the
machinery behaves, not that the library's theorems still derive.  The VNB
test is the slower, periodic exam that re-runs the mathematics
itself --- re-deriving every proof in the library, as any serious proof
system must.  Keep both: they answer different questions.
\end{remark}

%% =========================================================
\section{Inspecting a result}
\label{sec:status-query}

Three commands expose the classification at the REPL.

\begin{itemize}
\item \texttt{(status 'name)} prints, on one screen, a result's
  provenance, whether it is a PSS member, any warrant it carries, and its
  certification state (``certified by the VNB test on \textit{date}'' or
  ``not certified since the last run'').

\item \texttt{(status-audit)} writes \texttt{reference/STATUS-AUDIT.md}: a
  census of the whole library by kind, flagging hygiene issues --- asserted
  results with no warrant, theorem-shaped names that have drifted out of
  the theorem bucket, and the like.

\item \texttt{(pss-letter 'name)} drafts a plain-English ``recommendation
  letter'' for a candidate PSS member: an English rendering of what it
  says, its current standing, and an assessment against the admission
  gates (true; a sound, terminating rewrite shape; general rather than
  bespoke; non-redundant).  It is an aid to deciding whether a result
  belongs in the curated support set, not an automatic admission.
\end{itemize}


%% =========================================================
\chapter{The expression language}
\label{sec:language}


The language has two syntactic layers.

The \emph{internal layer} represents every expression as a plain
Lisp S-expression.  This representation is time-tested and deliberately
simple: it requires no parser, supports straightforward recursive
processing, and is directly inspectable and constructible at the Scheme
REPL.  All prover machinery---inference rules, substitution, pattern
matching, macetes, arithmetic evaluation---operates on this layer.
Users are not required to engage with it directly; however, in the
tradition of early Lisp, they are free to do so.  Constructing or
manipulating S-expressions by hand, loading axioms as quoted lists,
or writing Scheme code that builds formulas programmatically are all
legitimate and supported uses.

The \emph{string syntax} is the surface notation accepted by
\texttt{make-wff-from-string} and emitted by \texttt{wff->string}
/ \texttt{pp}: binary operators are written infix, multi-argument
terms use functional notation, quantifiers use the binding-list form,
and set expressions use brace notation.
Section~\ref{sec:parser-printer} describes both layers in full.
The examples in this manual use string syntax for user-facing
formulas; internal S-expression forms appear only in theory and
structure definitions.

MIT~Scheme reads symbols case-insensitively (folding to lower case),
so \texttt{FORALL}, \texttt{Forall}, and \texttt{forall} are the same
symbol at the Scheme level.

As should be the case in any formal system, we distingish between
syntax and semantics. In what follows, we distinguish two distinct
types of syntactic categories, each with its own semantics: formulas
and terms. These syntactic categories can be viewed in either the
internal syntax or the string syntax. In the presentation below, we
adhere to the string syntax.


\section{Formulas}

Semantically, a formula is a syntactic expression that can have only
two values: TRUTH or FALSEHOOD. 
%
\begin{remark}\label{remark:27137-52766-301830}
%
Note that ``UNDEFINED'' is not a possible value for a
formula. However, the value of any particular formula may not be
known. In fact, since we are silent about how elements in our system
may be interpreted in a pure von Neumann Bernays theory, different
interpretations may yield different values.
%
\end{remark}

Formulas can be produced from other ones by basic logical operators,
including quantification over variables that range over objects.
%
In the following table, the symbols $p$, $q$ are not variables that
range over formulas, but rather placeholders that can instantiated
with well-formed formulas.

\begin{itemize}[label=\textendash, leftmargin=1cm]
\item \texttt{truth} propositional truth $\top$
\item \texttt{falsity} propositional falsity $\bot$
\item \texttt{not($p$)} negation $\lnot p$
\item \texttt{$p$ and $q$} conjunction $p \land q$
\item \texttt{$p$ or $q$} disjunction $p \lor q$
\item \texttt{$p$ implies $q$} implication $p \Rightarrow q$
\item \texttt{$p$ iff $q$} biconditional $p \Leftrightarrow q$
\item \texttt{forall([$b_1$, \ldots, $b_n$], $p$)} universal quantification; see Section~\ref{sec:multiquant}
\item \texttt{forsome([$b_1$, \ldots, $b_n$], $p$)} existential quantification; see Section~\ref{sec:multiquant}
\item \texttt{$a$ = $b$} strict equality: both defined and equal
\item \texttt{$a$ == $b$} quasi-equality: both defined and equal, or both undefined
\item \texttt{$a$ in $b$} membership $a \in b$
\item \texttt{$A$ subset $B$} subclass $A \subseteq B$; see Section~\ref{sec:class-set-discipline}
\end{itemize}

\noindent
The quantifiers \texttt{forall} and \texttt{forsome} range over \emph{all
classes} (the system is single-sorted).  There is no restriction to sets.

\subsection{Quantification}
\label{sec:multiquant}

Both \texttt{forall} and \texttt{forsome} require a bracket-enclosed
\emph{binding list}, even for a single variable:

\begin{itemize}[label = $\circ$, leftmargin=1cm]
\item \texttt{forall([binding, ...], body)}

\item \texttt{forsome([binding, ...], body)}
\end{itemize}

\noindent
\texttt{forall(x, x = x)} is a syntax error; the correct form is
\texttt{forall([x], x = x)}.

\noindent
Each \emph{binding} in the string syntax is one of:

\begin{itemize}[label = $\circ$, leftmargin=1cm]
\item \texttt{$x$} unrestricted variable ranging over all classes
\item \texttt{$x$ in $A$} variable $x$ restricted to class $A$
\item \texttt{[$n_1$, \ldots, $n_k$] in $A$} tuple destructuring: fresh $t \in A$, components bound to \texttt{nth(1,~$t$)}, \ldots, \texttt{nth($k$,~$t$)}
\end{itemize}

\noindent
In the \emph{S-expression syntax} a binding can additionally group
several unrestricted variables into a single spec:
\[
  \texttt{(FORALL ((}x_1\ x_2\ \ldots\ x_n\texttt{) ...) body)}
\]
The grouped binding \texttt{($x_1$ $x_2$ \ldots\ $x_n$)} is exactly
equivalent to the sequence of single-variable bindings
\texttt{($x_1$) ($x_2$) \ldots\ ($x_n$)}.
The utility functions \texttt{vnb-unwind} and \texttt{vnb-wind}
convert between the two forms:

\begin{VNBsnip}
(vnb-unwind '(FORALL ((x y z)) body))
;=> (FORALL ((x)) (FORALL ((y)) (FORALL ((z)) body)))

(vnb-wind '(FORALL ((x)) (FORALL ((y)) (FORALL ((z)) body))))
;=> (FORALL ((x) (y) (z)) body)
\end{VNBsnip}

\noindent
Bindings are processed left to right, outermost first.  Thus in the
\emph{internal syntax} they expand to fully nested single-variable form
as follows:

\begin{VNBsnip}
forall([x, y in NN, z], y >= 0)
; => (FORALL x (FORALL y (IMPLIES (IN y NN) (FORALL z (>= y 0)))))

forall([[E, +, *] in METRIC], IS-METRIC(E, +, *))
; => (FORALL t_0 (IMPLIES (IN t_0 METRIC)
;                         (IS-METRIC (NTH 1 t_0) (NTH 2 t_0) (NTH 3 t_0))))
\end{VNBsnip}

\paragraph{Display.}
The printer reconstructs binding-list form from the internal nested
representation.  Consecutive same-type quantifiers are collapsed into
a single binding list; a change of quantifier type stops the collapsing.
For example, \texttt{forall([x in NN], forall([y], p))} prints as
\texttt{forall([x in NN, y], p)}.

\paragraph{Variable shadowing.}
Bound variable names may be reused within nested quantifiers.
The inner binding shadows the outer; \texttt{fresh-var} in
\texttt{subst-free} ensures substitution never captures a free variable.
Such formulas are valid but confusing; users should avoid them.

\subsection{Equality and quasi-equality}

Terms may be \emph{undefined} (e.g.\ an application \texttt{f(x)} when
\texttt{f} is not a function applicable to \texttt{x}).  The system
provides two equality predicates:

\paragraph{Strict equality \texttt{$a$ = $b$}.}
True iff both \texttt{a} and \texttt{b} are defined and equal.
False if either is undefined.  In particular:
\begin{itemize}[label=$-$, leftmargin=1cm]
\item \texttt{george(s) = 75} is \emph{false} when \texttt{george}
  is not a function.
\item \texttt{george(s) = george(s)} is \emph{false} for the same
  reason.
\item A term $t$ is defined iff $t = t$:
  $\mathit{DEF}(t) \;\equiv\; (t = t)$.
\end{itemize}

\paragraph{Quasi-equality \texttt{$a$ == $b$}.}
True iff both \texttt{a} and \texttt{b} are defined and equal, \emph{or}
both are undefined.  Key properties:
\begin{itemize}[label=$-$, leftmargin=1cm]
\item \texttt{t == t} is \emph{always true}, for any term $t$,
  defined or not.
\item \texttt{george(s) == george(s)} is \emph{true}.
\item \texttt{$a$ = $b$} implies \texttt{$a$ == $b$}, but not vice versa.
\item \texttt{$a$ == $b$} and $\mathit{DEF}(a)$ together imply
  \texttt{$a$ = $b$}.
\end{itemize}

\noindent
Undefinedness propagates through all operations: if \texttt{george(s)}
is undefined then \texttt{george(s) + 1} is also undefined,
and \texttt{george(s) + 1 = 2} is false.

Variables and class constants (\texttt{ORD}, \texttt{RR},
\texttt{EMPTY-SET}, \ldots) are always defined.  The application
\texttt{f(x)} is defined when \texttt{f in fun(A, B)} and
\texttt{x in A} are provable for some $A$, $B$.

\paragraph{Reflexivity in proofs.}
\texttt{cmd-quasi-reflexivity} closes any goal of the form
\texttt{a == a} unconditionally.
\texttt{cmd-reflexivity} closes \texttt{a = a} only when \texttt{a} is
demonstrably defined.

\section{Class constructors}
\label{sec:class-set-discipline}

Proper VNB has two kinds of object: \emph{sets} (elements of
$\mathit{SET}$) and \emph{proper classes} (collections that can have
members but are not themselves members of anything).  We retain that
distinction and add a third kind, \emph{functoids}.

A \textbf{functoid} is an operation specified by a \emph{signature}
\[
  A_1 \times A_2 \times \cdots \times A_n \;\longrightarrow\; B
\]
where $A_1, \ldots, A_n$ are domain classes and $B$ is the codomain
class.  None of them need be sets.

\begin{remark}
The motivation comes from category theory.  A functor is conceptually a
mapping between collections of mathematical structures, but its domain
is typically a proper class (e.g.\ the class of all groups, or all
topological spaces).  In a reductionist treatment one could introduce a
distinguished class $\mathit{FUNCTOID} \subseteq \mathit{CLASS}$,
whose members are ordered pairs encoding signature and rule.  This
system instead treats functoids as a new kind of object, coordinate
with sets and classes rather than reducible to them.  The principal
consequence is that the formal language does not quantify over
functoids: ``all functoids $F$ satisfy\ldots'' is metalanguage, not an
object-language formula.  Functoids are introduced from the predefined
list below or via \texttt{def-functoid} (Section~\ref{sec:def-functoid}).
\end{remark}

\texttt{SET} is a primitive class constant.  The formula
\texttt{$a$ in SET} expresses ``$a$ is a set''; there is no separate
predicate form \texttt{SET($a$)}.
The key axiom distinguishing VNB from typed systems:
\[
  (a \in b) \;\Longrightarrow\; a \in \mathit{SET}
\]
Only sets can be members of anything: membership forces the \emph{left}
argument to be a set.  The right argument---the containing class---may be
a proper class.  ($3 \in \mathit{ORD}$ is well-formed even though the
class of all ordinals is a proper class.)

Function application is written $f(x_1, \ldots, x_n)$ in both cases:
for a function $f$ this abbreviates $\texttt{apply-function}(f, x_1, \ldots, x_n)$,
and for a functoid it abbreviates $\texttt{apply-functoid}(f, x_1, \ldots, x_n)$.

The subset relation
\[
  A \subseteq B \;\equiv\; \forall x.\; x \in A \Rightarrow x \in B
\]
is a primitive binary predicate written \texttt{A subset B}
(\texttt{subset-def}).  An immediate corollary of
\texttt{membership-implies-sethood} is that \emph{every class is a
subclass of $\mathit{SET}$} (axiom \texttt{subset-set}): only sets can
appear on the left of~$\in$.

The constructors below are \emph{total} functoids: each is defined for
\emph{any} class arguments and always produces a well-defined class.
Whether the result is a set depends on a separate \emph{sethood-closure}
condition on the arguments; those conditions are listed in
Section~\ref{sec:axioms} and Section~\ref{sec:funapp}.
Class comprehension $\{x \mid p\}$ is the sole exception: it is a
\emph{schema}, yielding one instance per syntactically correct formula
$p$, each taking $0$ arguments.

\begin{remark}
Totality is the VNB tradition.  \texttt{fun(BONGO, HARP)} is a
perfectly well-defined class---the class of all total functions from
\texttt{BONGO} to \texttt{HARP}---even when neither is a set.  The
axiom \texttt{fun-set-iff} then says that this class is itself a set
if and only if both \texttt{BONGO} and \texttt{HARP} are sets.
Definability and sethood are separate questions.
\end{remark}

\begin{remark}
A related discipline: when we add axioms about a constructor we state
them at the abstract level---how \texttt{fun}, \texttt{pair},
\texttt{tuples}, \texttt{list}, etc.\ \emph{behave}---never in terms
of how they might be concretely encoded as sets at the vanilla VNB
level.  So even if some particular encoding happened to make the list
\texttt{[2,3]} literally equal to the integer \texttt{23}, no theorem
of VNB+ may distinguish such coincidences.  The user should never
have to know whether ordered pairs are Kuratowski or Wiener, or
whether \texttt{nn} is built from von Neumann or Zermelo numerals.
\end{remark}

\begin{remark}[Variadic functoids]
Several of the constructors that follow---\texttt{union},
\texttt{intersection}, \texttt{cartesian}, and the type constructor
\texttt{fun}---are \emph{variadic functoids}.  The number of arguments
varies between invocations but is fixed when the expression is written:
any \emph{finite} number is allowed.

This is distinct from variadic \emph{functions} of type
$\mathit{FUN}(\mathit{TUPLES}(A), B)$, which take a single set-level
tuple as their one argument (see the remark on $\mathit{TUPLES}$ in
Section~\ref{sec:constructor-examples}), and from fixed-arity functoids
such as ring addition
$\mathit{ADD}(r) \in \mathit{FUN}(\mathit{CARTESIAN}(A(r), A(r)), A(r))$,
which have a single argument slot of fixed Cartesian type.
\end{remark}

\medskip\noindent The remaining constructors of this section are listed
below.  Each entry states the defining membership property and points to
the corresponding proof commands in
Section~\ref{sec:nary-proof-commands}.

\begin{itemize}[label = $\circ$, leftmargin=1cm]

\item \texttt{union($A_1$,\ldots,$A_n$)}
  $x \in C$ iff $x \in A_i$ for some $i$.
  Axiom \texttt{union-membership} states this IFF for $n=2$.
  At any arity, the kernel rule \texttt{pi-union-intro!} reduces a goal
  $x \in \text{union}(A_1,\ldots,A_n)$ to the single subgoal $x \in A_k$
  for a chosen index $k$; \texttt{pi-union-elim!} consumes a hypothesis
  $x \in \text{union}(A_1,\ldots,A_n)$ and case-splits it into $n$
  branches, each with $x \in A_i$ assumed.  See
  Section~\ref{sec:nary-proof-commands}.
  The macete \texttt{union-decompose} rewrites
  $x \in \text{union}(A_1,\ldots,A_n)$ to the disjunction
  $(x \in A_1) \lor \cdots \lor (x \in A_n)$ in one step.
  Result is a set when all $A_i \in \mathit{SET}$
  (\texttt{union-set-closure}; n-ary by induction).

\item \texttt{intersection($A_1$,\ldots,$A_n$)}
  $x \in C$ iff $x \in A_i$ for all $i$.
  Axiom \texttt{intersection-membership} states this IFF for $n=2$.
  At any arity, the kernel rule \texttt{pi-intersection-intro!} reduces
  a goal $x \in \text{intersection}(A_1,\ldots,A_n)$ to $n$ subgoals
  $x \in A_i$; \texttt{pi-intersection-elim!} extracts $x \in A_k$ from
  a hypothesis $x \in \text{intersection}(A_1,\ldots,A_n)$ for a chosen
  index $k$.  See Section~\ref{sec:nary-proof-commands}.
  The macete \texttt{intersection-decompose} rewrites
  $x \in \text{intersection}(A_1,\ldots,A_n)$ to the conjunction
  $(x \in A_1) \land \cdots \land (x \in A_n)$ in one step.
  Result is a set when any $A_i \in \mathit{SET}$
  (\texttt{intersection-set-closure}).

\item \texttt{complement-in($A$, $B$)}
  Relative complement $A \setminus B$: $x \in C$ iff $x \in A$ and
  $x \notin B$ (\texttt{complement-in-membership}).
  $B$ need not be a set; result is a set when $A \in \mathit{SET}$
  (\texttt{complement-in-set-closure}).

\item \texttt{cartesian($A_1$,\ldots,$A_n$)}
  $x \in C$ iff $x = [a_1,\ldots,a_n]$ with $a_i \in A_i$.
  Result is a set iff all $A_i \in \mathit{SET}$ (\texttt{cartesian-set-iff}).

\item \texttt{tuples($A$)}
  $x \in C$ iff $x = [a_1,\ldots,a_n]$ with all $a_i \in A$,
  for some $n \ge 0$.
  Result is a set when $A \in \mathit{SET}$ (\texttt{tuples-sethood}).

\item \texttt{fun($A$)}, \texttt{fun($A$, $B$)}
  The unary $\mathit{FUN}(A)$ is the class of all total functions with
  domain $A$ (no codomain restriction); proper class in general.  The
  binary $\mathit{FUN}(A,B)$ is the restriction
  $\{f \in \mathit{FUN}(A) : \forall x \in A.\;f(x) \in B\}$.
  $\mathit{FUN}(A,B) \in \mathit{SET}$ iff $A, B \in \mathit{SET}$
  (\texttt{fun-set-iff}); $\mathit{FUN}(A) \in \mathit{SET}$ never holds
  in general because the codomain is unconstrained.
  See Section~\ref{sec:funapp}.

  The strengthened \texttt{fun-domain-apply-def} now reads
  $f \in \mathit{FUN}(A) \Rightarrow ((f\,x) = (f\,x) \Leftrightarrow x \in A)$.
  Recall VNB equality is partial: $(t = t)$ is the definedness predicate
  for $t$, so the iff says \emph{$f(x)$ is defined exactly when $x \in A$}.
  Consequence: the witness $A$ is unique, justifying ``the'' domain of $f$.

\item \texttt{is-fun($f$)}
  The predicate ``$f$ is a function'': there exists a set $A$ with
  $f \in \mathit{FUN}(A)$.  Axiom \texttt{is-fun-def}:
  $\mathit{IS\text{-}FUN}(f) \Leftrightarrow \exists A \in \mathit{SET}.\;
   f \in \mathit{FUN}(A)$.

\item \texttt{dom($f$)}
  The \emph{domain} of $f$ as a class.  Characterised by membership
  (\texttt{dom-membership}):
  $x \in \mathit{DOM}(f) \Leftrightarrow x \in \mathit{SET}
   \land (f\,x) = (f\,x)$, i.e.\ $x$ is a set on which $f$ is defined.
  When $f \in \mathit{FUN}(A)$, $\mathit{DOM}(f)$ and $A$ agree on
  members (\texttt{dom-fun-membership}).  When additionally
  $A \in \mathit{SET}$, $\mathit{DOM}(f) \in \mathit{SET}$ and
  $\mathit{DOM}(f) = A$ as sets, by extensionality
  (\texttt{dom-of-fun}).

\item \texttt{res($f$, $B$)}
  The \emph{restriction} of $f$ to $B$.  Semantically
  $\mathit{RES}(f, B) = \lambda x \in B.\;f(x)$.  Total as a constructor
  (defined for all $f$ and $B$ -- bongo-style); typed by:
  \begin{itemize}[label=$-$, leftmargin=1cm]
  \item \texttt{res-typing}:
    $f \in \mathit{FUN}(A) \land B \subseteq A
    \Rightarrow \mathit{RES}(f, B) \in \mathit{FUN}(B)$.
  \item \texttt{res-apply}:
    $f \in \mathit{FUN}(A) \land B \subseteq A \land x \in B
    \Rightarrow (\mathit{RES}(f, B))(x) = f(x)$.
  \item \texttt{res-codomain} (derived): codomain inherits from
    $f \in \mathit{FUN}(A, C)$.
  \end{itemize}
  Outside the case $B \subseteq A$ the term exists but carries no
  $\mathit{FUN}$-typing guarantees.

\item \texttt{partial-fun($A$)}, \texttt{partial-fun($A$, $C$)}
  Partial functions on $A$.  Members are total functions whose domain
  is some subset of $A$:
  $\mathit{PARTIAL\text{-}FUN}(A) = \{f \mid \exists B \in \mathit{POWER}(A).
   \; f \in \mathit{FUN}(B)\}$
  (\texttt{partial-fun-membership}, generally a proper class because
  the codomain is unconstrained).  The binary form fixes a codomain:
  $\mathit{PARTIAL\text{-}FUN}(A, C) = \{f \mid \exists B \in \mathit{POWER}(A).
   \; f \in \mathit{FUN}(B, C)\}$
  (\texttt{partial-fun-binary-membership}).  Sethood
  (\texttt{partial-fun-binary-sethood}):
  $A, C \in \mathit{SET} \Rightarrow \mathit{PARTIAL\text{-}FUN}(A, C)
   \in \mathit{SET}$ (stated as axiom; provable once a class-union
  principle over a set-indexed family is in place).

\item \texttt{\{$x$ in $A$: $p$\}} / \texttt{\{$x$ in $A$ | $p$\}}
  Separation $\mathit{SEP}(x, A, p) = \{x \in A \mid p\}$: $x \in C$
  iff $x \in A$ and $p[x]$ holds.  Because $p$ is genuinely a schema
  (a formula, not a first-order term), the characterization is given
  by kernel rules rather than a single FOL axiom:
  \texttt{pi-sep-sethood!} ($\mathit{SEP}(x, A, p) \in \mathit{SET}$
  reduces to $A \in \mathit{SET}$);
  \texttt{pi-sep-mem-intro!} ($y \in \mathit{SEP}(x, A, p)$ splits
  into $y \in A$ and $p[x:=y]$);
  \texttt{pi-sep-mem-elim!} (assumption $y \in \mathit{SEP}(x, A, p)$
  yields both $y \in A$ and $p[x:=y]$ in context).
  Result is a set when $A \in \mathit{SET}$.

\item \texttt{\{$x$ | $p$\}}
  Class comprehension $\mathit{COMP}(x, p) = \{x \mid p\}$: $y \in C$
  iff $y \in \mathit{SET}$ and $p[x:=y]$ (members of any class are
  sets).  Same schema treatment as separation:
  \texttt{pi-comp-mem-intro!} and \texttt{pi-comp-mem-elim!}.
  Proper class in general unless provably bounded.

\item \texttt{power($A$)}, \texttt{power($A$, $B$)}
  Unary \texttt{power($A$)} is the power set: $x \in C$ iff
  $x \in \mathit{SET}$ and $x \subseteq A$ (\texttt{power-set}).
  Binary \texttt{power($A$, $B$)} is the function-space exponentiation
  $A^B$, axiomatized by \texttt{power-exp}:
  \[
     \mathit{POWER}(A, B) \;=\; \mathit{FUN}(B, A)
  \]
  i.e.\ functions $B \to A$.  This convention matches the arithmetic
  evaluator: $\texttt{power}(2, 3) = 2^3 = 8$ counts functions from
  a 3-set to a 2-set.  Sethood and membership inherit from
  \texttt{fun-set-iff} / \texttt{fun-codomain-iff}.

\end{itemize}

\begin{remark}
\texttt{union($A_1$, \ldots, $A_n$)} takes a finite listed union.
The \emph{big-union} $\bigcup A = \{x \mid \exists a \in A.\; x \in a\}$
of an arbitrary class $A$ of sets is a separate constructor
\texttt{UNION-SET}, currently pending (Section~\ref{sec:axioms}).
\end{remark}

\section{Term constructors}
\label{sec:choice-iota-lambda}

\emph{Term constructors} build syntactic objects whose intended
representation are objects in the VNB universe.

\begin{itemize}[label = $\circ$, leftmargin=1cm]
\item \texttt{[$e_1$, \ldots, $e_n$]} ordered $n$-tuple $(e_1, \ldots, e_n)$; always a set
\item \texttt{nth($k$, $e$)} $k$-th component of tuple $e$ (1-indexed); $k$ may be any term;
  \texttt{nth($k$, [$e_1$,\ldots,$e_n$])} reduces to $e_k$ via
  \texttt{cmd-nth-reduce} (Section~\ref{sec:commands})
\item \texttt{length($L$)} number of elements in tuple $L$; a natural number:
  \begin{itemize}[label=$-$, leftmargin=1cm]
  \item \texttt{length-of-empty}: $\texttt{length}([\,]) = 0$
  \item \texttt{length-in-nn}: $L \in \mathit{TUPLES}(\mathit{SET}) \Rightarrow
    \texttt{length}(L) \in \mathit{NN}$
  \item \texttt{nth-in-range}: $L \in \mathit{TUPLES}(A),\; i \in \mathit{NN},\;
    1 \le i \le \texttt{length}(L) \Rightarrow \texttt{nth}(i,L) \in A$
  \end{itemize}
  A full recursive characterisation requires a \texttt{prepend} constructor
  and the \texttt{tuples-induction} schema, both pending
  (Section~\ref{sec:pending-axioms}).
\item \texttt{EMPTY-SET} the empty set $\emptyset$; always a set; $= \texttt{make-set}([\,])$
\item \texttt{\{$a$\}} singleton $\{a\}$; always a set; $= \texttt{make-set}([$a$])$
\item \texttt{make-set($L$)} unordered finite set from tuple $L \in \mathit{TUPLES}(A)$,
  $A \in \mathit{SET}$; always a set:
  \begin{itemize}[label=$-$, leftmargin=1cm]
  \item \texttt{make-set-membership}: $x \in \texttt{make-set}(L)
    \Longleftrightarrow \exists\,i \in \mathit{NN}.\;
    1 \le i \le \texttt{length}(L) \land \texttt{nth}(i,L) = x$
    (the $1 \le i \le \texttt{length}(L)$ bounds matter:
    $\texttt{nth}$ out-of-range is unspecified)
  \item \texttt{make-set-sethood}: $L \in \mathit{TUPLES}(A),\; A \in \mathit{SET}
    \Rightarrow \texttt{make-set}(L) \in \mathit{SET}$
  \item \texttt{make-set-empty}: $\texttt{make-set}([\,]) = \emptyset$
  \end{itemize}
  The notation $\{a,b,c,\ldots\}$ is sugar for $\texttt{make-set}([a,b,c,\ldots])$;
  repetitions are harmless.
\item \texttt{choice($A$)} primitive choice operator: an element of nonempty class~$A$
  (undefined when $A$ is empty).  Axiom \texttt{choice-axiom}:
  \[
    (\exists x.\; x \in A) \;\Longrightarrow\; \texttt{choice}(A) \in A
  \]
  This is the Bourbaki-style global choice: a single canonical operator
  defined uniformly on all nonempty classes, not just sets.
\item \texttt{iota($x$, $p$)} definite description: the unique object satisfying~$p$;
  a binder.  Defined only when $\exists! x.\; p(x)$ is provable.  The
  characterization is the kernel rule \texttt{pi-iota-def!}, invoked
  with the iota term as argument; it posts two subgoals:
  \begin{enumerate}[label=(\arabic*),leftmargin=1.2cm]
  \item the existence-and-uniqueness obligation
    $\exists x.\;p \land (\forall y.\; p[x:=y] \Rightarrow x = y)$, and
  \item the original goal with the defining property
    $p[x := \texttt{iota}(x, p)]$ added as an assumption.
  \end{enumerate}
  A planned \texttt{define-object} form will wrap this pattern:
\begin{VNBsnip}
(define-object sqtwo (IOTA x (= (* x x) 2)))
; requires a proved unique-existence theorem; in this language,
; spelt out as (FORSOME x (AND (= (* x x) 2)
;                              (FORALL y (IMPLIES (= (* y y) 2) (= x y)))))
; (a FORSOME-UNIQUE macro abbreviation is planned but not yet defined)
; installs sqtwo satisfying: sqtwo * sqtwo = 2
\end{VNBsnip}
\item \texttt{lambda([$x_1$ in $A_1$, \ldots, $x_n$ in $A_n$], $\mathit{body}$)}
  anonymous multi-argument function; a binder.
  Each variable has its own independent domain, exactly as in quantification.
  All domains $A_i$ must belong to \texttt{SET}; the result is an element
  of $\mathit{FUN}$.  Two underlying representations:
  \begin{itemize}[label = $-$, leftmargin=1cm]
  \item \emph{record form} (\texttt{<functoid>}): produced by the parser
    when the lambda is given explicit binder domains
    \texttt{lambda([x in A], body)}.  Reduced by
    \texttt{cmd-functoid-beta} (short form \texttt{(beta)}).
  \item \emph{symbolic form} (\texttt{VNB-LAMBDA}): a raw S-expression
    binder \texttt{(VNB-LAMBDA $x$ $\mathit{body}$)} or
    \texttt{(VNB-LAMBDA (LIST $x_1\,\ldots\,x_n$) $\mathit{body}$)}.
    Used in axioms written directly as S-expressions
    (e.g.\ \texttt{RING-PROD}, \texttt{CC-RING}, \texttt{sum}).
    Typing rule \texttt{pi-lambda-type!} (single-binder shape):
    goal $\mathit{VNB\text{-}LAMBDA}(x, \mathit{body}) \in
    \mathit{FUN}(A, B)$ reduces to the subgoal
    $\forall x.\; x \in A \Rightarrow \mathit{body} \in B$.
    $\beta$-reduction \texttt{pi-lambda-beta!} (short form
    \texttt{(lam-b)}) rewrites
    $\bigl(\mathit{VNB\text{-}LAMBDA}(\overline x, \mathit{body})\bigr)
    (\overline v) = \mathit{body}[\overline x := \overline v]$
    (parallel substitution).
  \end{itemize}
  Cardinal-exponent typing schema (record-form, summarized; multi-binder
  derives from the single-binder rule plus uncurrying):
  \begin{multline*}
    \bigl(\forall x_i \in A_i.\;\mathit{body}(x_1,\ldots,x_n) \in B\bigr)
    \;\Longrightarrow\; \\
    \texttt{lambda}([x_1 \in A_1,\ldots,x_n \in A_n],\,\mathit{body})
      \in \mathit{FUN}(A_1 \times \cdots \times A_n,\,B)
  \end{multline*}

\item \texttt{lambdoid([$x_1$ in $A_1$, \ldots, $x_n$ in $A_n$], $\mathit{body}$)}
  like \texttt{lambda}, but the domains $A_i$ may be arbitrary classes
  (including proper classes such as \texttt{ORD}).  The result is tagged
  as a \emph{functoid} at the meta-level; it does not necessarily belong
  to any \texttt{FUN} class but still supports $\beta$-reduction via
  \texttt{(beta)}.
\end{itemize}

\begin{remark}[Functoids as meta-level objects]
Both \texttt{lambda} and \texttt{lambdoid} produce \emph{functoid records} at
the implementation level --- tagged Scheme structures distinct from plain
S-expressions and from VNB sets or classes.  This design follows the IMPS
principle that every structured entity carries its own recognisable tag; it
prevents the expression-traversal code from ever confusing a functoid with an
ordinary compound term.

Internally, the application
$\texttt{lambda}([x_1 \in A_1,\ldots], \mathit{body})(v_1,\ldots)$
is represented as
$(\texttt{apply-functoid}\;\langle\textit{record}\rangle\;v_1\;\cdots)$,
a mixed structure whose \texttt{car} is the symbol \texttt{apply-functoid}
and whose \texttt{cadr} is the record itself.  The surface syntax
$f(x)$ is used for \emph{all} application; the parser emits
\texttt{apply-functoid} only when $f$ is already tagged as a functoid.
\end{remark}

\begin{remark}
Tuples $[e_1,\ldots,e_n]$ are primitive here.  One could in principle reduce
them to iterated pairs (Kuratowski style), but we have no interest in doing so:
$[\ ]$ is perfectly well-defined as a term constructor of arity zero, and that
is all we need.
\end{remark}

\begin{remark}
The index $k$ in $\texttt{nth}(k, e)$ may be any term (not only a literal
integer).  The primitive inference \texttt{cmd-nth-reduce} does require a
literal integer to reduce $\texttt{nth}(k, [e_1,\ldots,e_n])$ to $e_k$, but
$\texttt{nth}(i, L)$ with a variable index $i$ is a perfectly valid term and
appears in \texttt{make-set-membership} and \texttt{nth-in-range}.
For example, $\texttt{nth}(3,\,[a,b,c,c,e])$ reduces to $c$.
\end{remark}

\begin{remark}[Indexing convention: lists are 1-indexed]
Tuples / lists in VNB are \emph{1-indexed}: \texttt{nth}$(i, L)$ is valid
for $1 \le i \le \texttt{length}(L)$, and $\texttt{nth}(1,\,[a, b, c]) = a$.
Lists and functions on initial segments of $\mathit{NN}$ are genuinely
different things, and we do \emph{not} reduce one to the other.  When an
interval constructor is added, the canonical domain for an $n$-element list
is $\mathit{INTERVAL}(1, n)$, \emph{not} $\mathit{INTERVAL}(0, n-1)$.
The 1-indexing convention matches IMPS and is built into every axiom
that mentions \texttt{nth} or \texttt{length} (\texttt{make-set-membership},
\texttt{nth-in-range}, \texttt{length-of-empty}).
\end{remark}

\begin{remark}
\texttt{choice}, \texttt{iota}, \texttt{lambda}, and \texttt{lambdoid} all
participate correctly in \texttt{free-vars}, \texttt{subst-free},
\texttt{alpha-equiv?}, and macete pattern-matching.  Both \texttt{lambda} and
\texttt{lambdoid} produce \emph{functoid records} --- Scheme tagged structures
separate from S-expressions --- so they are never accidentally mistaken for
ordinary terms by the expression machinery.
\end{remark}

\section{Function application}
\label{sec:funapp}

Function application is written as a compound term: \texttt{$f$($x$)} for
any object $f$ and any argument $x$.  Examples:
\begin{VNBsnip}
a + b               ; addition applied to a, b
NORM(E, v)          ; NORM of Banach space E applied to vector v
GEORGE(s)           ; syntactically valid, but may be undefined
\end{VNBsnip}

\texttt{fun}($A$, $B$) (with $A, B \in \mathit{SET}$) is a primitive set constructor: its members are
the total functions from $A$ to $B$.  Since $A, B \in \mathit{SET}$,
\texttt{fun}($A$,$B$) is itself a set (\texttt{fun-set-iff}), and each
$f \in \mathit{FUN}(A,B)$ is a set (\texttt{membership-implies-sethood}).

\begin{remark}
The soundness of $\mathit{FUN}(A,B)$ is confirmed by the standard
set-theoretic construction: a function $f : A \to B$ is a set of pairs
$(a, b) \in A \times B$ that is functional (each $a$ has at most one $b$)
and total (every $a \in A$ appears).  Since $A \times B$ is a set, so
is any subset of it, making $f$ a set.
\end{remark}


\begin{remark}
Category theory motivates a natural extension of $\mathit{FUN}$.  A
functor between large categories---such as the category of all groups or
all topological spaces---maps a proper class of objects to a proper class
of objects.  Since $\mathit{FUN}$ requires a set-valued domain, functors
fall outside $\mathit{FUN}$ entirely.

The planned treatment introduces \emph{operators} as a third kind of
entity alongside sets and classes.  There is no class $\mathit{OPER}(A,B)$
(no such collection can exist without violating the set/class dichotomy),
but there will be a predicate
\[
  \mathit{ISOPER}(A,\, B,\, f)
\]
asserting that $f$ behaves as a total function from $A$ to $B$ even when
$A$ or $B$ is a proper class.  An accompanying term constructor
$\texttt{apply-operator}(f, x)$ then gives the value of $f$ at $x$, for
any $x \in A$, under the hypothesis $\mathit{ISOPER}(A,B,f)$.

Operators are not members of anything---they do not appear on the left of
$\in$---so they are consistent with
\texttt{membership-implies-sethood}.  The three-level ontology (sets,
classes, operators) keeps the VNB foundation intact while accommodating
the full generality of categorical constructions.
\end{remark}

%
The base theory installs the following axioms.  The primary form of
\texttt{FUN} is the unary $\mathit{FUN}(A)$ -- all total functions whose
domain is $A$, with no codomain restriction; the binary $\mathit{FUN}(A,B)$
is the derived form $\{f \in \mathit{FUN}(A) : \forall x \in A.\;f(x) \in B\}$.

\begin{itemize}[label = $\circ$, leftmargin=1cm]
\item \texttt{fun-domain-apply-def}
  $f \in \mathit{FUN}(A) \land x \in A \;\Rightarrow\; f(x) = f(x)$
\item \texttt{fun-domain-extensionality}
  $f,g \in \mathit{FUN}(A) \land (\forall x \in A.\;f(x) = g(x)) \;\Rightarrow\; f = g$
\item \texttt{fun-codomain-iff}
  $f \in \mathit{FUN}(A,B) \;\Leftrightarrow\;
   (f \in \mathit{FUN}(A) \land \forall x \in A.\;f(x) \in B)$
\item \texttt{fun-apply-type}
  $f \in \mathit{FUN}(A,B) \land x \in A \;\Rightarrow\; f(x) \in B$
  (in \texttt{axioms.scm}; derivable from \texttt{fun-codomain-iff})
\item \texttt{fun-set-iff}
  $\mathit{FUN}(A,B) \in \mathit{SET} \;\Leftrightarrow\; (A \in \mathit{SET}) \land (B \in \mathit{SET})$
\item \texttt{fun-elements-are-sets}
  $f \in \mathit{FUN}(A,B) \land A \in \mathit{SET} \;\Rightarrow\; f \in \mathit{SET}$
\item \texttt{fun-domain-elements-are-sets}
  $f \in \mathit{FUN}(A) \land A \in \mathit{SET} \;\Rightarrow\; f \in \mathit{SET}$
\item \texttt{cartesian-set-iff}
  $A \times B \in \mathit{SET} \;\Leftrightarrow\; (A \in \mathit{SET}) \land (B \in \mathit{SET})$
\item \texttt{tuples-sethood}
  $A \in \mathit{SET} \;\Rightarrow\; \mathit{TUPLES}(A) \in \mathit{SET}$
\end{itemize}

\begin{remark}
The constructor $\mathit{TUPLES}(A)$ is the class of all finite sequences
$[a_1, \ldots, a_n]$ with each $a_i \in A$, for any $n \ge 0$ (including
the empty sequence $[\,]$).  It fills a gap that $\mathit{CARTESIAN}$ does
not cover: $\mathit{CARTESIAN}(A_1, \ldots, A_n)$ fixes the arity $n$ at
the time the expression is written, whereas $\mathit{TUPLES}(A)$ lets $n$
range over all natural numbers.

This supports \emph{variadic} functions whose single argument is a tuple
of any length: an element of $\mathit{FUN}(\mathit{TUPLES}(A), B)$ takes
one tuple and returns a value in $B$.  This is \emph{not} the same as a
fixed-arity $n$-ary function, which has type
$\mathit{FUN}(\mathit{CARTESIAN}(A_1, \ldots, A_n), B)$ and is bound to
exactly $n$ argument slots by its type.  The variadic flavour is the
right one when an operator must handle inputs of any length (e.g.\ a sum
over an indexed sequence); the fixed-arity flavour is the right one when
the arity is intrinsic to the operator (e.g.\ ring addition, which is
honestly binary).  Since $\mathit{TUPLES}(A) \in \mathit{SET}$ when
$A \in \mathit{SET}$, the variadic case is an ordinary instance of
$\mathit{FUN}(A', B)$ with $A' = \mathit{TUPLES}(A)$.
\end{remark}

\section{Constructor examples}
\label{sec:constructor-examples}

Quick reference for every class and term constructor.  Full descriptions
appear in Sections~\ref{sec:class-set-discipline}
and~\ref{sec:choice-iota-lambda}.

\paragraph{Functoids (class constructors).}  Each is defined for any class
arguments; the result is always a well-defined class.  Each row gives the
condition on $x$ for membership in the constructed class $C$, and the
sethood-closure condition on the arguments (when there is one).

\begin{center}
\footnotesize
\begin{tabular}{@{}p{4.5cm}p{4.8cm}p{3.0cm}@{}}
\toprule
Expression & $x \in C$ iff \ldots & $C \in \mathit{SET}$ when \\
\midrule
\texttt{union($A_1$,\ldots,$A_n$)} & $x \in A_i$ for some $i$ & all $A_i \in \mathit{SET}$ \\[2pt]
\texttt{intersection($A_1$,\ldots,$A_n$)} & $x \in A_i$ for all $i$ & any $A_i \in \mathit{SET}$ \\[2pt]
\texttt{complement-in($A$,$B$)} & $x \in A$ and $x \notin B$ & $A \in \mathit{SET}$ \\[2pt]
\texttt{cartesian($A_1$,\ldots,$A_n$)} & $x = [a_1,\ldots,a_n]$, $a_i \in A_i$ & all $A_i \in \mathit{SET}$ \\[2pt]
\texttt{tuples($A$)} & $x = [a_1,\ldots,a_n]$, $a_i \in A$, $n \ge 0$ & $A \in \mathit{SET}$ \\[2pt]
\texttt{fun($A$,$B$)} & $x$ is a total function from $A$ to $B$ & $A, B \in \mathit{SET}$ \\[2pt]
\texttt{\{$x$ in $A$: $p$\}} & $x \in A$ and $p(x)$ holds & $A \in \mathit{SET}$ \\[2pt]
\texttt{\{$x$ | $p$\}} & $p(x)$ holds & rarely \\[2pt]
\texttt{power($A$)} & $x \in \mathit{SET}$ and $x \subseteq A$ & $A \in \mathit{SET}$ \\
\bottomrule
\end{tabular}
\end{center}

\paragraph{Term constructors.}

\begin{center}
\footnotesize
\begin{tabular}{@{}p{5.2cm}p{7.1cm}@{}}
\toprule
Expression & Key property \\
\midrule
\texttt{EMPTY-SET} & empty set $\emptyset$; always a set; no members \\[2pt]
\texttt{\{$a$\}} & singleton $\{a\}$; always a set \\[2pt]
\texttt{\{$a$, $b$, \ldots\}} & finite set via \texttt{make-set}; repetitions harmless \\[2pt]
\texttt{[$e_1$, \ldots, $e_n$]} & ordered $n$-tuple; always a set \\[2pt]
\texttt{nth($k$, $e$)} & $k$-th component (1-indexed); reduces via \texttt{cmd-nth-reduce} \\[2pt]
\texttt{choice($A$)} & $(\exists x.\;x\in A) \Rightarrow \texttt{choice}(A) \in A$ \\[2pt]
\texttt{iota($x$, $p$)} & unique object satisfying $p$; requires $\exists!\,x.\;p(x)$ \\[2pt]
\texttt{lambda([$x_1$ in $A_1$,\ldots], $b$)} & multi-arg function; domains in \texttt{SET}; $\beta$-reduces via \texttt{(beta)} \\[2pt]
\texttt{lambdoid([$x_1$ in $A_1$,\ldots], $b$)} & like \texttt{lambda} but domains may be proper classes \\
\bottomrule
\end{tabular}
\end{center}

%% =========================================================

\section{Destructuring quantification}

The syntax extends the bounded-variable binding form
\texttt{[$x$ in $A$]} to allow the bound variable to be a
\emph{tuple pattern}:
\[
  \texttt{[[$n_1$, \ldots, $n_k$] in $A$]}
\]
Here $A$ is \emph{any class} --- a Cartesian product, a number system,
a separation class, a structure predicate, or any user-defined class.
There is no requirement that $A$ be a mathematical structure in the sense
of \texttt{def-structure}.

The binding introduces an anonymous implicit tuple variable~$t_0$, adds
$t_0 \in A$ to the assumptions, and binds each name $n_i$
to the $i$-th component of~$t_0$:
\[
  n_i \;\longmapsto\; \texttt{nth}(i,\; t_0)
\]
The variable~$t_0$ never appears to the user.  The scope of
$n_1, \ldots, n_k$ is \emph{lexical}: they have no meaning outside the
quantifier body.

\paragraph{Examples.}

Ranging over a Cartesian product (no structure involved):
\begin{VNBsnip}
forall([[x, y] in cartesian(RR, RR)], x + y = y + x)

; Expands to (internal form):
(FORALL t_0
  (IMPLIES (IN t_0 (CARTESIAN RR RR))
    (= (+ (NTH 1 t_0) (NTH 2 t_0))
       (+ (NTH 2 t_0) (NTH 1 t_0)))))
\end{VNBsnip}

Ranging over a named structure (the structure case is not special):
\begin{VNBsnip}
forall([[A, +, *, f] in NORMEDSPACE], v in A implies is-continuous(f, A, RR))

; Expands to (internal form):
(FORALL t_0
  (IMPLIES (IN t_0 NORMEDSPACE)
    (IMPLIES (IN v (NTH 1 t_0))
      (is-continuous (NTH 4 t_0) (NTH 1 t_0) RR))))
\end{VNBsnip}

Destructuring bindings compose freely with plain and bounded bindings
(Section~\ref{sec:multiquant}):
\begin{VNBsnip}
forall([[A, +, *, f] in NORMEDSPACE, v in A], is-continuous(f, A, RR))
\end{VNBsnip}


\section{Parser and printer}
\label{sec:parser-printer}

Four entry points provide string-level access to VNB formulas:

\begin{VNBsnip}
(make-wff-from-string str)  ; string -> <wff>  (validates via make-wff)
(parse-string str)          ; string -> raw S-expression (no validation)
(wff->string w)             ; <wff>  -> string
(pp w)                      ; display wff->string then newline
\end{VNBsnip}

\noindent
\texttt{parse-string} converts a string to a raw S-expression without
validation; it is used internally by the interactive commands and is
available for low-level \texttt{cmd-*} calls.  For normal use, pass
formula strings directly to \texttt{ai}, \texttt{bc}, \texttt{inst},
\texttt{cut}, \texttt{wk}, and \texttt{vnb-create-num-theorem}.
\texttt{make-wff-from-string} is the right choice when constructing a
formula to prove (for \texttt{sp}).

\paragraph{Operator precedence (high to low).}
\begin{center}
\small
\begin{tabular}{lll}
\toprule
Level & Operators & Associativity \\
\midrule
primary & atoms, \texttt{f(a,b)}, \texttt{[a,b]}, \texttt{(expr)} & --- \\
power & \texttt{\textasciicircum} & right \\
unary & \texttt{-x} & right \\
mul & \texttt{*}, \texttt{/} & left (n-ary \texttt{*}) \\
add & \texttt{+}, \texttt{-} & left (n-ary \texttt{+}) \\
cmp & \texttt{in}, \texttt{=}, \texttt{==}, \texttt{<=} & binary \\
and & \texttt{and} & left (n-ary) \\
or & \texttt{or} & left (n-ary) \\
iff & \texttt{iff} & left (n-ary) \\
implies & \texttt{implies} & right \\
\bottomrule
\end{tabular}
\end{center}

\paragraph{Division and negation sugar.}
\begin{VNBsnip}
x / y    =>  (* x (recip y))   ; division is syntactic sugar
-x       =>  (- x)             ; unary minus; lower prec than ^
-x ^ 2   =>  (- (power x 2))  ; i.e. -(x^2), not (-x)^2
\end{VNBsnip}

\paragraph{Quantifier syntax.}
Quantifiers use the binding-list form:
\begin{VNBsnip}
forall([x], body)              ; bare variable (all classes)
forall([x in A], body)         ; guarded variable
forall([[a, b, c] in A], body) ; tuple destructuring
forall([x in A, y, z in B], body)  ; mixed bindings
forsome([x in NN], body)       ; existential; same forms
\end{VNBsnip}

\noindent
The bracket syntax after the function name distinguishes quantifiers
from ordinary function calls.

\paragraph{Display conventions (\texttt{wff->string} / \texttt{pp}).}
\begin{itemize}[label=\textendash, leftmargin=1cm]
\item \begin{sloppypar}\textbf{Quantifiers} print in binding-list form:
  the primitive expanded form is reverse-engineered.
  \texttt{forall([x in nn, y], p(x, y))}\end{sloppypar}
\item \textbf{Logical connectives} \texttt{and}, \texttt{or}, \texttt{iff}
  print n-ary infix with surrounding parentheses:
  \texttt{((p and q) and r)}.
  \texttt{implies} prints binary infix: \texttt{(p implies q)}.
\item \textbf{Comparisons} \texttt{in}, \texttt{=}, \texttt{==},
  \texttt{<=} print binary infix: \texttt{(x in nn)}, \texttt{(x = y)}.
\item \textbf{Arithmetic} \texttt{+}, \texttt{*} print n-ary infix:
  \texttt{(x + y + z)}.  Binary \texttt{-} prints as \texttt{(x - y)};
  unary as \texttt{-x}.  Exponentiation as \texttt{(x \textasciicircum{} y)}.
\item \textbf{Lists} \texttt{[a, b, c]} in both input and output: the
  list syntax is the same in string form (internally stored as \texttt{LIST}).
\item \textbf{All other compound terms} use functional notation
  \texttt{f(x, y, z)}: \texttt{nth}, \texttt{cartesian}, \texttt{fun},
  \texttt{power} (unary set power), \texttt{not}, and user-defined heads.
  Nullary symbols print without parentheses.
\end{itemize}

\paragraph{Example round-trips.}
\begin{VNBsnip}
(pp (make-wff-from-string "forall([x in NN, y in NN], x + y = y + x)"))
; => forall([x in nn, y in nn], ((x + y) = (y + x)))

(pp (make-wff-from-string "not(P) implies falsity"))
; => (not(p) implies falsity)

(pp (make-wff-from-string "2 ^ 10 = 1024"))
; => ((2 ^ 10) = 1024)

(expression->string '(forall x (implies (in x nn)
                       (forall y (implies (in y nn) (= (+ x y) (+ y x)))))))
; => "forall([x in nn, y in nn], ((x + y) = (y + x)))"
\end{VNBsnip}

\noindent
All identifiers are case-folded to lowercase by MIT~Scheme, so
\texttt{NN} and \texttt{nn} are the same symbol; \texttt{FORALL} and
\texttt{forall} are also identical.  The printer always emits lowercase.

%% =========================================================

%% =========================================================


%% =========================================================
%% This file was a short chapter of its own until 2026-09-20.  It is now the closing part of the
%% chapter on the expression language, so that the chapter on proofs follows that chapter directly.
\section{Formulas as typed objects: \texttt{wff}}
\label{chap:syntax-contexts}

A \emph{well-formed formula} in this system is not a bare S-expression
but a typed object: a formula that has passed the well-formedness check.  In the REPL it prints as:
\begin{VNBsnip}
#[wff N "formula string"]
\end{VNBsnip}
where \texttt{N} is an internal identifier and the quoted string is the
pretty-printed formula.  Every formula in a sequent is a \texttt{wff}.

\texttt{(make-wff expr)} constructs such an object, from a string in
surface syntax or from an S-expression.  \texttt{(free-vars p)} returns
the variables that are free in it.

Two procedures expose the internals of a \texttt{wff}:
\begin{VNBsnip}
(wff->sexp p)        ; => raw S-expression (no output)
(display-contents p) ; prints kind, library name, formula
\end{VNBsnip}
\texttt{wff->sexp} is defined in \texttt{wff.scm};
\texttt{display-contents} is defined in \texttt{contexts.scm}.
Both are available in interactive sessions.

Example (the base library is \texttt{VNB-LIBRARY}):
\begin{VNBsnip}
(define p (make-wff "forall([x in A], 0 <= funny-name)"))
; => #[wff 12 "forall([x in a], 0 <= funny-name)"]

(free-vars p)
; => (a funny-name)   ; x is bound by the forall
\end{VNBsnip}


%% =========================================================
\chapter{Proofs}
\label{ch:proofs}

\section{Sequents and deduction graphs}\label{section:sequents}

A \emph{sequent} is a pair $(L, B)$, where $L$ is a list of formulas
called \emph{assumptions} and $B$ is a single formula called the
\emph{assertion}.  A sequent is represented as
\[
  A_1,\, A_2,\, \ldots,\, A_n \;\Longrightarrow\; B
\]
where $A_1, \ldots, A_n$ are the assumptions (or hypotheses or the context) and $B$
is the assertion. In the course of a proof, many sequents may be created.

A \emph{deduction graph} is a directed acyclic graph of
\emph{sequent nodes} connected by \emph{inference nodes}.  An inference
node records which rule was applied, which sequent nodes are its
hypotheses, and which sequent node is its conclusion.

A sequent node is \emph{grounded} when it has an inference node all of
whose hypothesis nodes are themselves grounded.  Leaf inferences (with no
hypothesis nodes) are immediately grounded.  Grounding propagates upward.
A proof is \emph{complete} when the root sequent node is grounded.

\begin{remark}
A \emph{proof state} (\texttt{<proof-state>}) is not the same thing as a
deduction graph.  It is a thin wrapper that bundles three items together: the
deduction graph itself, the \emph{root} sequent node (the original goal,
fixed at \texttt{start-proof} time and checked by \texttt{qed}), and the
\emph{focus} sequent node (the subgoal currently being worked on).  The
deduction graph is the mathematical object that grows as rules are applied;
the proof state is the navigational view into it that the interactive session
and the \texttt{cmd-*} commands operate on.  In the interactive session,
\texttt{*ps*} always holds the current proof state.
\end{remark}

%% =========================================================
\section{Kernel operations}\label{sec:kernel-boundary}

A proof in VNB is carried out on a deduction graph, and the graph is
changed only by \emph{kernel operations}.  A kernel operation is applied
to one sequent node $N$.  It adds one inference node: the conclusion is
$N$, the rule recorded is the name of the operation, and the hypotheses
are sequent nodes that the operation builds from $N$ in a way fixed by
the operation.  Nothing else changes the graph.

There are \VNBkernelopcount{} kernel operations.  The list below is
produced from the list the system itself checks each time it starts
(Section~\ref{sec:kernel-checked}).

\VNBkerneloptable

The last \VNBkernelopcountoracle{} are \emph{oracles}.  An oracle runs a
decision procedure, written in Scheme, on the assertion of $N$ and some
of its assumptions.  If the procedure succeeds, the operation adds an
inference node, with no hypotheses except in the case of
\texttt{arith-simplify}, which posts one.  The graph records the name of
the oracle; for \texttt{ineq} and \texttt{sos} it also records the
certificate the procedure found.  The inference is written only if the
oracle's checker accepts it (Section~\ref{sec:inference-checked}).

The kernel operations, together with the axioms
(Section~\ref{sec:proof-debt}), are everything VNB assumes.  A result
reported \texttt{modulo 0} rests on these and on nothing else.

Every other proof command is a \emph{tactic}.  A tactic is a Scheme
procedure that calls kernel operations, possibly many, possibly after a
search.  It has no other access to the graph.  A tactic can fail to find
a proof.  It cannot add an inference that no kernel operation would add.
Section~\ref{sec:tactics-ref} describes the tactics, and
Section~\ref{sec:command-uses} lists, for each command, the kernel
operations it can reach.

\begin{remark}
One command may record different operations according to the form of the
assertion: \texttt{di} records \texttt{and-intro}, \texttt{implies-intro},
\texttt{forall-intro} and so on.  Two commands may record the same
operation.  The name recorded in the graph is that of the operation, not
of the command.
\end{remark}

\subsection{What the operations do}\label{sec:kernel-ops-detail}

%% Inference figures.  \VNBfig{operation}{hypotheses}{conclusion}, and
%% \VNBfigS with the hypotheses stacked one per row (separated by \\); the side
%% conditions follow in a VNBcond list.  Defined here because this is the only
%% section that uses them.
\providecommand{\VNBfig}[3]{\[\frac{\displaystyle #2}{\displaystyle #3}\;\;\mbox{\texttt{#1}}\]}
\providecommand{\VNBnohyp}{\phantom{\Gamma}}
\providecommand{\FV}{\mathrm{FV}}
\providecommand{\aeq}{=_{\alpha}}
\providecommand{\VNBfigS}[3]{\[\frac{\displaystyle\begin{array}{c}#2\end{array}}{\displaystyle #3}\;\;\mbox{\texttt{#1}}\]}
\newenvironment{VNBcond}{\sloppy\begin{description}[font=\normalfont\itshape,leftmargin=0.85cm,itemsep=0.5mm,topsep=0.5mm]}{\end{description}}

In this section $N$ is the node the operation is applied to, with
sequent $\Gamma \Longrightarrow G$.  ``Hypotheses'' are the sequents of
the new hypothesis nodes.  $B[x{:=}t]$ is $B$ with $t$
substituted for the free occurrences of $x$, bound variables being
renamed where needed to avoid capture.  $\FV(E)$ is the set of variables
free in $E$, and $\FV(\Gamma)$ the union over the formulas of $\Gamma$.
An operation that \emph{declines} leaves the graph unchanged.  Posting a
sequent that the graph already contains, up to renaming of bound
variables, returns the existing node, so a new hypothesis may already be
grounded.

An operation that takes an index or a pattern records it beside its
name: the graph holds \texttt{(cartesian-elim 2)}, \texttt{(macete
\textit{name} $L$ $R$)}, \texttt{(ineq \textit{certificate})}.  The name under which the inference is counted, and under which
it appears in the list above, is the leading symbol.  The arguments are
data about one application.

\emph{Definedness.}  VNB terms may fail to denote, and the sign of
denotation is $t = t$: strict equality holds only between defined terms
(Chapter~\ref{sec:language}).  Two operations, \texttt{forall-elim}
and \texttt{reflexivity}, therefore consult a \emph{definedness
certificate}, one test applied to a term and the assumptions of $N$.
$\Gamma$ certifies $t$ when $t$ is a variable or an atomic constant, a
numeral or a ground arithmetic term, a term that $\Gamma$ types
($t \in X$) or equates ($t = u$), a term occurring outside every binder
in an assumption of the form $\in$, $=$, $\leq$ or $<$, a total class
constructor applied to certified arguments, a separation, comprehension
or indexed union whose domain is certified, an accessor of a structure
$\Gamma$ carries, an applied structure operation whose arguments
$\Gamma$ types in its declared domain, or $f(a)$ with $f \in
\mathrm{FUN}(D, \cdot)$ and $a \in D$ in $\Gamma$.  A term whose head is
\textsc{choice} or \textsc{iota}, and an application of an untyped
function, are never certified.  The policy is
\texttt{docs/definedness-instantiation-2026-09-18.md}.

\subsubsection*{Reading the figures}

Each operation is stated below as an inference figure.  The sequents above
the line are the hypotheses, the sequent below it is the sequent of $N$,
and the labelled conditions under the figure are the side conditions.  A
figure is the relation that the operation's \emph{checker} tests before
the inference is written (Section~\ref{sec:inference-checked}): an
inference of that form that meets the conditions is accepted, whichever
procedure produced it, and any other is refused.  The operation itself
produces one instance of its figure.  Where it has a choice to make
(which occurrences to rewrite, what to call a new variable), the prose
after the figure says which choice it makes.  A condition labelled
\emph{Not checked} is imposed by the operation but not tested by its
checker.

The figures share five conventions.
%% the conventions: header of rule-checkers-logic.scm, lines 14-27;
%% chk-set=? :82, chk-ctx+? :100, chk-ctx-+? :103, chk-arity? :133,
%% chk-match-under :167, chk-adds-nothing? :126
\begin{itemize}
\item \emph{Contexts are sets.}  $A \aeq B$ means that $A$ and $B$ are
  the same formula up to renaming of bound variables.  Two contexts are
  equal when every formula of each is $\aeq$ to a formula of the other.
  $\Gamma, A$ is $\Gamma \cup \{A\}$, so it is $\Gamma$ itself when
  $\Gamma$ already contains $A$; $\Gamma - F$ is $\Gamma$ with every
  formula $\aeq F$ removed.  The order of a context carries no meaning.
\item \emph{Arity.}  The number of hypotheses is the number drawn.  A
  figure with two sequents above the line refuses an inference with one
  or with three.  Where the number depends on the formula (one hypothesis
  for each component of a tuple), the figure says so.
\item \emph{Assumption.}  ``$F \in \Gamma$'' means that some member of
  $\Gamma$ is $\aeq F$.  $F$ is a member of the context, never a
  subformula of one.  The operation is told which member by its argument;
  the checker is not, and searches $\Gamma$ for one that fits.
\item \emph{Matching.}  A term that occurs in a hypothesis but not in the
  conclusion (the $t$ of \texttt{forall-elim}, the witness of
  \texttt{forsome-intro}, the eigenvariables of \texttt{forall-intro}
  and \texttt{forsome-elim}) is not recorded.  The checker recovers it by
  one-way matching: the conclusion's formula, with a hole in place of the
  variable, is matched against the hypothesis' formula.  A hole is never
  filled by a term that contains a variable bound inside the pattern.
\item \emph{Weakening case.}  \texttt{forall-elim}, \texttt{forsome-elim},
  \texttt{theorem-assumption}, \texttt{if-true} and \texttt{if-false}
  also accept a hypothesis whose assertion is $G$ and whose context is
  included in $\Gamma$.  That is what the operation produces when the
  formula it adds is already in $\Gamma$, and the inference is then an
  instance of \texttt{weakening}.
\end{itemize}

\subsection{Logic, equality, induction and the structure eliminators}

\subsubsection*{Direct inference (command \texttt{di})}

The command records one of five operations, according to the form of $G$
(\texttt{pi-direct-inference!}).  On any other form it declines.
%% pi-direct-inference!, primitive-inferences.scm:75

\begin{center}
\begin{tabular}{@{}ll@{}}
\toprule
$G$ & operation \\
\midrule
$A \wedge B$ & \texttt{and-intro} \\
$A \Rightarrow B$ & \texttt{implies-intro} \\
$A \Leftrightarrow B$ & \texttt{iff-intro} \\
$\neg A$ & \texttt{not-intro} \\
$\forall x.\, B$ & \texttt{forall-intro} \\
\bottomrule
\end{tabular}
\end{center}

%% rule-checkers-logic.scm:347 and-intro, :386 implies-intro, :412 iff-intro, :399 not-intro
\VNBfig{and-intro}{\Gamma \Longrightarrow A \qquad \Gamma \Longrightarrow B}{\Gamma \Longrightarrow A \wedge B}
\VNBfig{implies-intro}{\Gamma, A \Longrightarrow B}{\Gamma \Longrightarrow A \Rightarrow B}
\VNBfig{iff-intro}{\Gamma, A \Longrightarrow B \qquad \Gamma, B \Longrightarrow A}{\Gamma \Longrightarrow A \Leftrightarrow B}
\VNBfig{not-intro}{\Gamma, A \Longrightarrow \textsc{falsity}}{\Gamma \Longrightarrow \neg A}

To prove $\forall x.\, B$ one proves $B$ for an $x$ about which nothing
is assumed.  In a sequent that means a variable that does not occur free
in any assumption: if $\Gamma$ said something about $x$, a proof of $B$
from $\Gamma$ would be a proof about that particular $x$, not about every
$x$.  A variable introduced under this condition is called an
\emph{eigenvariable}\index{eigenvariable}, and the condition is the \emph{eigenvariable
condition}.  \texttt{forall-intro} removes a chain of leading universal
quantifiers in one step, and collects on the way the guards $y \in A$
that restrict them.  Write the assertion as $G_0 = G$ and, for
$1 \leq i \leq k$, $G_{i-1} = \forall x_i.\, C_i$.
%% chk-forall-strip :441, chk-forall-eigen-ok? :491, chk-forall-intro-at :514,
%% forall-intro checker :555, all rule-checkers-logic.scm
\VNBfig{forall-intro}{\Gamma, g_1, \ldots, g_m \Longrightarrow G_k}{\Gamma \Longrightarrow G_0}
\begin{VNBcond}
\item[Chain] $1 \leq k$, and for each $i$ the formula $C_i[x_i{:=}y_i]$
  is either $(y_i \in A_i) \Rightarrow G_i$, in which case $y_i \in A_i$
  is \emph{collected}, or else it is $G_i$.  $g_1, \ldots, g_m$ are the
  collected formulas.
\item[Eigenvariable] Each $y_i$ is a variable, and
  $y_i \notin \FV(\Gamma)$, $y_i \notin \FV(G_{i-1})$ and
  $y_i \notin V_{i-1}$, where $V_0 = \emptyset$ and $V_i$ is
  $V_{i-1} \cup \FV(y_i \in A_i)$ when step $i$ collected a guard and
  $V_{i-1} \cup \{y_i\}$ otherwise.
\item[Matching] The $y_i$ are read off the hypothesis.  The checker
  tries $k$ equal to the whole length of the chain first and then each
  shorter length, since a shorter chain is also a sound inference.
\end{VNBcond}
In words: an eigenvariable is free neither in the context, nor in the
formula still to be generalised at that step, nor in a guard collected
before it (domain included), and the eigenvariables are pairwise
distinct.

The operation (\texttt{peel-foralls-raw}) always takes the whole chain.
It keeps $y_i = x_i$ when $x_i$ is free neither in $\Gamma$ nor in a
guard already collected, and otherwise takes a new variable, spelled as
$x_i$ followed by an underscore and a number.
%% peel-foralls-raw, primitive-inferences.scm:41
A guard is collected only when it has the form $y_i \in A_i$ for the
variable just bound.  Any other antecedent stays in the assertion:
$\forall x.\, P(x) \Rightarrow C$ yields the hypothesis
$\Gamma \Longrightarrow P(x) \Rightarrow C$, and a second direct inference
gives $\Gamma, P(x) \Longrightarrow C$.

\subsubsection*{\texttt{or-intro-left}, \texttt{or-intro-right}
  (commands \texttt{oi-l}, \texttt{oi-r})}

%% rule-checkers-logic.scm:362, :374
\VNBfig{or-intro-left}{\Gamma \Longrightarrow A}{\Gamma \Longrightarrow A \vee B}
\VNBfig{or-intro-right}{\Gamma \Longrightarrow B}{\Gamma \Longrightarrow A \vee B}

\subsubsection*{\texttt{forsome-intro}: supplying a witness (command \texttt{ew})}

%% rule-checkers-logic.scm:574; pi-exists-witness!, primitive-inferences.scm:357
\VNBfig{forsome-intro}{\Gamma \Longrightarrow B[x{:=}t]}{\Gamma \Longrightarrow \exists x.\, B}
\begin{VNBcond}
\item[Matching] $t$ is read off the hypothesis.
\item[Not checked] Well-formedness of $B[x{:=}t]$.  The operation
  checks it and raises an error, recording nothing, on a malformed
  witness.
\end{VNBcond}
The argument of the operation is $t$.  No definedness obligation is
posted for $t$, so the operation is not the mirror of
\texttt{forall-elim} in that respect.

\subsubsection*{\texttt{truth-intro}}

%% rule-checkers-logic.scm:588
\VNBfig{truth-intro}{\VNBnohyp}{\Gamma \Longrightarrow \textsc{truth}}

This operation is unreachable.  No proof command calls the procedure
that records it, and a \textsc{truth} assertion is already closed by
\texttt{assumption} and by \texttt{reflexivity}, both of which accept it.
It is listed because the code is present and a future command could
reach it.

\subsubsection*{Antecedent inference (command \texttt{ai})}

The argument is an assumption $F$ of $N$.  The command records one of
five operations according to the form of $F$
(\texttt{pi-antecedent-inference!}).
%% pi-antecedent-inference!, primitive-inferences.scm:132
\begin{center}
\begin{tabular}{@{}ll@{}}
\toprule
$F$ & operation \\
\midrule
$A \wedge B$ & \texttt{and-elim} \\
$A \vee B$ & \texttt{or-elim} \\
$\neg A$, with $A$ in $\Gamma$ & \texttt{not-elim} \\
$\exists x.\, B$ & \texttt{forsome-elim} \\
$A \Leftrightarrow B$ & \texttt{iff-elim} \\
\bottomrule
\end{tabular}
\end{center}
The command declines when $F$ is not an assumption of $N$, when $F$ is
$\neg A$ and $A$ is not in $\Gamma$, and when $F$ has any other form: an
implication, a universal formula or an atomic formula.  An implication in
the context is used by \texttt{detach} and \texttt{backchain}, a universal
formula by \texttt{forall-elim}.

%% rule-checkers-logic.scm:597 and-elim, :615 or-elim, :633 not-elim, :645 iff-elim
\VNBfig{and-elim}{(\Gamma - F), A, B \Longrightarrow G}{\Gamma \Longrightarrow G}
\begin{VNBcond}
\item[Assumption] $F = A \wedge B$, $F \in \Gamma$.
\end{VNBcond}
\VNBfig{or-elim}{(\Gamma - F), A \Longrightarrow G \qquad (\Gamma - F), B \Longrightarrow G}{\Gamma \Longrightarrow G}
\begin{VNBcond}
\item[Assumption] $F = A \vee B$, $F \in \Gamma$.
\end{VNBcond}
\VNBfig{not-elim}{\VNBnohyp}{\Gamma \Longrightarrow G}
\begin{VNBcond}
\item[Assumption] $\neg A \in \Gamma$ and $A \in \Gamma$.
\end{VNBcond}
\VNBfig{iff-elim}{(\Gamma - F), A \Rightarrow B, B \Rightarrow A \Longrightarrow G}{\Gamma \Longrightarrow G}
\begin{VNBcond}
\item[Assumption] $F = A \Leftrightarrow B$, $F \in \Gamma$.
\end{VNBcond}

%% forsome-elim checker, rule-checkers-logic.scm:665
\VNBfig{forsome-elim}{(\Gamma - F), B[x{:=}y] \Longrightarrow G}{\Gamma \Longrightarrow G}
\begin{VNBcond}
\item[Assumption] $F = \exists x.\, B$, $F \in \Gamma$.
\item[Eigenvariable] $y$ is a variable, and $y \notin \FV(G)$,
  $y \notin \FV(\Gamma - F)$, $y \notin \FV(F)$.  When $x$ is not free
  in $B$ there is no $y$ to test.
\item[Matching] $y$ is read off the hypothesis.
\end{VNBcond}
Nothing else is required: in particular nothing about the other
assumptions of $N$.  The operation names $y$ by the counter
(\texttt{fresh-var}): $x$ followed by an underscore and $n$, with $n$ the
next value for which that variable is free in none of $B$, $G$ and
$\Gamma - F$.  It never reuses $x$ itself.  The weakening case of the
conventions arises here only when $x$ is not free in $B$ and $B$ is
already in $\Gamma - F$.
%% pi-antecedent-inference! FORSOME arm, primitive-inferences.scm:169;
%% fresh-var, expressions.scm:701

\subsubsection*{\texttt{forall-elim}: instantiation of a universal assumption
  (commands \texttt{inst}, \texttt{inst+}; used by \texttt{fact})}

%% forall-elim checker, rule-checkers-logic.scm:702; pi-instantiate!, primitive-inferences.scm:282
\VNBfig{forall-elim}{\Gamma, B[x{:=}t] \Longrightarrow G \qquad \bigl(\Gamma \Longrightarrow t = t\bigr)}{\Gamma \Longrightarrow G}
\begin{VNBcond}
\item[Assumption] $\forall x.\, B \in \Gamma$.  It stays in the context.
\item[Matching] $t$ is read off the first hypothesis.
\item[Definedness] The second hypothesis is present exactly when
  $\Gamma$ does not certify $t$ (the certificate at the head of this
  section).
\item[Not checked] The checker accepts the second hypothesis present or
  absent.  When it is present, the checker requires it to be
  $\Gamma \Longrightarrow u = u$ for some term $u$, and does not require
  $u$ to be $t$.  Definedness of $t$ is therefore enforced by the
  operation alone.
\end{VNBcond}
The arguments are the universal assumption and $t$.  The second
hypothesis is there because a term may be undefined ($f(a)$ for $a$
outside the domain of $f$, a description with no referent), and a
universal formula speaks only of defined values.  $B[x{:=}t]$ is checked
for well-formedness; a malformed term raises an error and leaves the
graph unchanged.  The same operation is recorded, once per term, by the
procedure behind \texttt{apply-thm}, which cites a theorem and then
instantiates it at a list of terms (\texttt{pi-spec!}).

\subsubsection*{\texttt{assumption}: the goal is already in the context
  (command \texttt{ass})}

%% rule-checkers-logic.scm:738
\VNBfig{assumption}{\VNBnohyp}{\Gamma \Longrightarrow G}
\begin{VNBcond}
\item[Assumption] $G \in \Gamma$, or $G$ is \textsc{truth}.
\end{VNBcond}

\subsubsection*{\texttt{theorem-assumption}: citing a theorem or an axiom
  (command \texttt{ta}; used by \texttt{fact})}

%% rule-checkers-logic.scm:750; chk-installed-theorem-name :329
\VNBfig{theorem-assumption}{\Gamma, S \Longrightarrow G}{\Gamma \Longrightarrow G}
\begin{VNBcond}
\item[Installed] $S \aeq$ the formula installed under some name in the
  theorem table.
\end{VNBcond}
The argument of the operation is a name, and $S$ is the formula installed
under it; the operation declines when there is none.  The name is not
recorded with the inference, so the checker looks the statement up in
the table.

This is the only way a formula enters a proof without being proved in
it.  $S$ is an axiom (there are \VNBaxiomcount{} of them, listed in the
catalog), a definition, a proven theorem, or an asserted support.  The
bill printed when a theorem is installed
(Section~\ref{sec:proof-debt}) lists the asserted supports reached
through this operation, directly or through the theorems cited.

\subsubsection*{\texttt{cut}: introducing a lemma (commands \texttt{cut},
  \texttt{have!})}

%% rule-checkers-logic.scm:767; pi-cut!, primitive-inferences.scm:223
\VNBfig{cut}{\Gamma \Longrightarrow A \qquad \Gamma, A \Longrightarrow G}{\Gamma \Longrightarrow G}

The argument is $A$.  It is checked for well-formedness first, and a
malformed lemma raises an error before anything is recorded.

\begin{remark}
The operation never declines.  When $\Gamma$ already contains $A$, the
second hypothesis is $\Gamma \Longrightarrow G$, which is the sequent of
$N$ itself, so the inference has $N$ among its own hypotheses and
grounds nothing.  The same happens when $A$ is $G$.  A command that cuts
a formula already present therefore appears to succeed and leaves the
proof where it was.
\end{remark}

\subsubsection*{\texttt{weakening}: dropping assumptions (command \texttt{wk})}

%% rule-checkers-logic.scm:781; pi-weaken!, primitive-inferences.scm:208
\VNBfig{weakening}{\Delta \Longrightarrow G}{\Gamma \Longrightarrow G}
\begin{VNBcond}
\item[Inclusion] $\Delta \subseteq \Gamma$.
\end{VNBcond}
The operation takes one assumption $F$ of $N$ and posts
$\Delta = \Gamma - F$; it declines when $F$ is not an assumption of $N$.

\subsubsection*{\texttt{detach}: forward modus ponens (command \texttt{detach!})}

%% rule-checkers-logic.scm:790; pi-detach!, primitive-inferences.scm:442
\VNBfig{detach}{\Gamma, B \Longrightarrow G}{\Gamma \Longrightarrow G}
\begin{VNBcond}
\item[Assumption] $A \Rightarrow B \in \Gamma$ and $A \in \Gamma$.
\end{VNBcond}
The argument is the implication.  Both it and its antecedent stay.

\subsubsection*{\texttt{backchain}: reducing the goal to an antecedent
  (commands \texttt{bc}, \texttt{bc*})}

%% rule-checkers-logic.scm:807; pi-backchain!, primitive-inferences.scm:419
\VNBfig{backchain}{\Gamma \Longrightarrow A}{\Gamma \Longrightarrow G}
\begin{VNBcond}
\item[Assumption] $A \Rightarrow G \in \Gamma$.
\end{VNBcond}

\subsubsection*{\texttt{proof-by-contradiction} (command \texttt{pbc})}

%% rule-checkers-logic.scm:824
\VNBfig{proof-by-contradiction}{\Gamma, \neg G \Longrightarrow \textsc{falsity}}{\Gamma \Longrightarrow G}

There are no arguments, and the operation never declines.

\subsubsection*{\texttt{contraposition}}

%% rule-checkers-logic.scm:835
\VNBfig{contraposition}{\Gamma \Longrightarrow \neg B}{\Gamma \Longrightarrow \neg A}
\begin{VNBcond}
\item[Assumption] $A \Rightarrow B \in \Gamma$.
\end{VNBcond}
This operation is unreachable: no proof command calls the procedure that
records it.

\subsubsection*{\texttt{eq-subst}: rewriting the assertion by a context
  equation (command \texttt{subst})}

%% rule-checkers-logic.scm:857; chk-replaced? :241
\VNBfig{eq-subst}{\Gamma \Longrightarrow G'}{\Gamma \Longrightarrow G}
\begin{VNBcond}
\item[Assumption] $\Gamma$ contains $s = t$, $t = s$, $s \equiv t$ or
  $t \equiv s$.
\item[Replacement] $G' \not\aeq G$, and $G'$ is $G$ with some
  occurrences of $s$ replaced by $t$, or with some occurrences of $t$
  replaced by $s$.  A binder whose bound variable the replacement renamed
  is matched with its renamed form.
\item[Not checked] That no replaced occurrence lies in the scope of a
  binder whose variable is free in $s$ or in $t$.  The checker compares
  $G$ and $G'$ position by position and does not test it.
\end{VNBcond}
The argument is the equation $s = t$ (or $s \equiv t$).  The operation
replaces \emph{every} occurrence of $s$ by $t$, in operator position as
well as in argument position, and declines when nothing changes.  It
leaves alone a subterm below a binder whose bound variable occurs free in
$s$ or in $t$, since rewriting there would capture; when $s$ is a
variable it replaces the free occurrences only.  $G'$ is checked for
well-formedness.
%% pi-eq-subst! :563, replace-term :496, primitive-inferences.scm

\begin{remark}
The direction of the rewrite is fixed by the argument, not by the
orientation in which $\Gamma$ happens to hold the equation.  To use a
context equation from right to left, name it the other way round.
\end{remark}

\subsubsection*{\texttt{reflexivity} and \texttt{quasi-reflexivity}
  (commands \texttt{rfl}, \texttt{qrfl})}

%% rule-checkers-logic.scm:879, :890
\VNBfig{reflexivity}{\VNBnohyp}{\Gamma \Longrightarrow t = t}
\begin{VNBcond}
\item[Definedness] $\Gamma$ certifies $t$.  The operation also accepts
  the assertion \textsc{truth}.
\item[Not checked] The definedness condition.  The checker accepts
  $\Gamma \Longrightarrow u = u$ for any term $u$.
\end{VNBcond}
\VNBfig{quasi-reflexivity}{\VNBnohyp}{\Gamma \Longrightarrow t \equiv t}

\texttt{quasi-reflexivity} asks nothing further: quasi-equality holds
between two undefined terms as well.

\subsubsection*{\texttt{if-true}, \texttt{if-false}
  (commands \texttt{if-true}, \texttt{if-false})}

%% chk-if-rule, rule-checkers-logic.scm:900-933
\VNBfig{if-true}{\Gamma \Longrightarrow p \qquad \Gamma, \mathrm{IF}(p,a,b) = a \Longrightarrow G}{\Gamma \Longrightarrow G}
\VNBfig{if-false}{\Gamma \Longrightarrow \neg p \qquad \Gamma, \mathrm{IF}(p,a,b) = b \Longrightarrow G}{\Gamma \Longrightarrow G}

The argument is the conditional term $\mathrm{IF}(p, a, b)$.  It need not
occur in $G$.  The assertion is not rewritten: the branch equation
arrives as an assumption, and it takes an \texttt{eq-subst} to use it.
An argument that is not a conditional term raises an error rather than
declining.

\subsubsection*{\texttt{cartesian-intro}, \texttt{cartesian-elim}
  (commands \texttt{ci}, \texttt{ce})}

%% rule-checkers-logic.scm:939, :963
\VNBfig{cartesian-intro}{\Gamma \Longrightarrow a_1 \in A_1 \quad \cdots \quad \Gamma \Longrightarrow a_n \in A_n}{\Gamma \Longrightarrow [a_1, \ldots, a_n] \in A_1 \times \cdots \times A_n}
\begin{VNBcond}
\item[Arity] The literal tuple and the literal product have the same
  length $n$, and there is one hypothesis per component, in order.
\end{VNBcond}
\VNBfig{cartesian-elim $k$}{\Gamma, a_k \in A_k \Longrightarrow G}{\Gamma \Longrightarrow G}
\begin{VNBcond}
\item[Assumption] $[a_1, \ldots, a_n] \in A_1 \times \cdots \times A_n
  \in \Gamma$, with tuple and product of the same length.  It stays.
\item[Index] $k$ is recorded with the name, and $1 \leq k \leq n$.
\end{VNBcond}

\subsubsection*{\texttt{tuples-intro}, \texttt{tuples-elim}
  (commands \texttt{ti}, \texttt{te})}

$\mathrm{TUPLES}(A)$ is the class of finite tuples over $A$.
%% rule-checkers-logic.scm:991, :1012
\VNBfig{tuples-intro}{\Gamma \Longrightarrow a_1 \in A \quad \cdots \quad \Gamma \Longrightarrow a_n \in A}{\Gamma \Longrightarrow [a_1, \ldots, a_n] \in \mathrm{TUPLES}(A)}
\begin{VNBcond}
\item[Arity] One hypothesis per component; none for the empty tuple.
\end{VNBcond}
\VNBfig{tuples-elim $k$}{\Gamma, a_k \in A \Longrightarrow G}{\Gamma \Longrightarrow G}
\begin{VNBcond}
\item[Assumption] $[a_1, \ldots, a_n] \in \mathrm{TUPLES}(A) \in \Gamma$.
\item[Index] $k$ is recorded with the name, and $1 \leq k \leq n$.
\end{VNBcond}

\subsubsection*{\texttt{union-intro}, \texttt{union-elim}
  (commands \texttt{ui}, \texttt{ue})}

$\mathrm{UNION}$ takes $n$ arguments.
%% rule-checkers-logic.scm:1038, :1060
\VNBfig{union-intro $k$}{\Gamma \Longrightarrow x \in A_k}{\Gamma \Longrightarrow x \in A_1 \cup \cdots \cup A_n}
\begin{VNBcond}
\item[Index] $k$ is recorded with the name, and $1 \leq k \leq n$.
\end{VNBcond}
\VNBfig{union-elim}{(\Gamma - F), x \in A_1 \Longrightarrow G \quad \cdots \quad (\Gamma - F), x \in A_n \Longrightarrow G}{\Gamma \Longrightarrow G}
\begin{VNBcond}
\item[Assumption] $F = x \in A_1 \cup \cdots \cup A_n$, $F \in \Gamma$.
  It is consumed.
\item[Arity] One hypothesis per set, in order.
\end{VNBcond}

\subsubsection*{\texttt{intersection-intro}, \texttt{intersection-elim}
  (commands \texttt{ii}, \texttt{ie})}

%% rule-checkers-logic.scm:1079, :1098
\VNBfig{intersection-intro}{\Gamma \Longrightarrow x \in A_1 \quad \cdots \quad \Gamma \Longrightarrow x \in A_n}{\Gamma \Longrightarrow x \in A_1 \cap \cdots \cap A_n}
\VNBfig{intersection-elim $k$}{\Gamma, x \in A_k \Longrightarrow G}{\Gamma \Longrightarrow G}
\begin{VNBcond}
\item[Assumption] $x \in A_1 \cap \cdots \cap A_n \in \Gamma$.  It stays.
\item[Index] $k$ is recorded with the name, and $1 \leq k \leq n$.
\end{VNBcond}

\subsubsection*{\texttt{nth-reduce}, \texttt{length-reduce},
  \texttt{functoid-beta} (commands \texttt{nth-r}, \texttt{len-r},
  \texttt{beta})}

The three have the same figure and differ in the contraction.
%% chk-contracts? rule-checkers-logic.scm:284; checkers :1137, :1153, :1174
\VNBfig{}{\Gamma \Longrightarrow G'}{\Gamma \Longrightarrow G}
\begin{VNBcond}
\item[Contraction] $G' \not\aeq G$, and $G'$ is obtained from $G$ by
  contracting redexes, where a redex that a contraction creates may be
  contracted in its turn.  The redexes are, respectively:
  $\mathrm{NTH}(k, [t_1, \ldots, t_n])$ with $1 \leq k \leq n$, which
  contracts to $t_k$; $\mathrm{LENGTH}([t_1, \ldots, t_n])$, which
  contracts to $n$; and a functoid applied to as many arguments as it
  binds, which contracts to its body with the arguments substituted in
  parallel for the bound variables.
\end{VNBcond}
None takes an argument.  Each operation contracts every redex of its
kind in $G$ and declines when $G'$ is $G$.  \texttt{nth-reduce} raises an
error on an index outside $1 \leq k \leq n$.  Nothing is asked of the
entries: a literal tuple has $n$ places whether or not the terms in them
denote.

\subsubsection*{\texttt{nn-induction} (command \texttt{ni})}

%% rule-checkers-logic.scm:1191
\VNBfig{nn-induction}{\Gamma \Longrightarrow B[n{:=}0] \qquad \Gamma \Longrightarrow \forall n.\, n \in \mathbb{N} \Rightarrow (B \Rightarrow B[n{:=}\mathrm{succ}\,n])}{\Gamma \Longrightarrow \forall n.\, n \in \mathbb{N} \Rightarrow B}
\begin{VNBcond}
\item[Shape] The guard is $n \in \mathbb{N}$ on the variable of the
  outermost quantifier.
\end{VNBcond}
The operation has no argument and declines on any other shape.  In
particular the induction variable must be the outermost one: a statement
that quantifies over something else first is not of this shape.

\subsubsection*{\texttt{transfinite-induction} (command \texttt{tfi})}

%% rule-checkers-logic.scm:1247; chk-tfi-beta-ok? :1241
\VNBfig{transfinite-induction}{\Gamma \Longrightarrow \forall \alpha.\, \bigl(\alpha \in \mathrm{ORD} \wedge \forall \beta.\, \beta <_{\mathrm{ORD}} \alpha \Rightarrow P[\alpha{:=}\beta]\bigr) \Rightarrow P}{\Gamma \Longrightarrow \forall \alpha.\, \alpha \in \mathrm{ORD} \Rightarrow P}
\begin{VNBcond}
\item[Freshness] $\beta$ is a variable, $\beta \neq \alpha$,
  $\beta \notin \FV(P)$, $\beta \notin \FV(\Gamma)$.
\item[Matching] $\beta$ is read off the hypothesis.
\end{VNBcond}
This is strong induction: the hypothesis of the step is $P$ at every
smaller ordinal.

\begin{remark}
The operation \texttt{transfinite-induction} and the axiom of the same
name in \texttt{structure-library/ordinals.scm} are two different
objects.  One is code in the kernel, the other a formula in the catalog.
\end{remark}

\subsubsection*{\texttt{transfinite-induction-3cases} (command \texttt{tfi3})}

%% rule-checkers-logic.scm:1274
\[
\frac{\displaystyle
  \begin{array}{l}
  \Gamma \Longrightarrow P[\alpha{:=}0] \\[2pt]
  \Gamma \Longrightarrow \forall \alpha.\,
    (\alpha \in \mathrm{ORD} \wedge P) \Rightarrow P[\alpha{:=}\mathrm{succ}_{\mathrm{ORD}}\,\alpha] \\[2pt]
  \Gamma \Longrightarrow \forall \alpha.\,
    \bigl(\mathrm{LIMIT\text{-}ORD}(\alpha) \wedge
          \forall \beta.\, \beta <_{\mathrm{ORD}} \alpha \Rightarrow P[\alpha{:=}\beta]\bigr)
    \Rightarrow P
  \end{array}}
 {\displaystyle \Gamma \Longrightarrow \forall \alpha.\, \alpha \in \mathrm{ORD} \Rightarrow P}
\]
\begin{VNBcond}
\item[Freshness] As for \texttt{transfinite-induction}, with $\beta$ read
  off the third hypothesis.
\end{VNBcond}
The three hypotheses are the base, successor and limit cases, in that
order.

\subsection{Set formation, descriptions and functions}

The three binders of this group carry a formula or a class expression in
their body, with a variable of the binder free in it.  A first-order
axiom would need a variable ranging over formulas, which VNB does not
have, so each is characterised by operations instead.
$\mathrm{SEP}(x, A, p)$ is $\{x \in A \mid p\}$,
$\mathrm{COMP}(x, p)$ is $\{x \mid p\}$, and
$\mathrm{BIG\text{-}UNION}(z, A, b)$ is $\bigcup_{z \in A} b$.

\subsubsection*{\texttt{sep-sethood}, \texttt{sep-mem-intro},
  \texttt{sep-mem-elim}
  (commands \texttt{sep-set}, \texttt{sep-mi}, \texttt{sep-me})}

%% rcs-sep-sethood :107, rcs-sep-mem-intro :160, rcs-sep-mem-elim :167, rule-checkers-schema.scm
\VNBfig{sep-sethood}{\Gamma \Longrightarrow A \in \mathrm{SET}}{\Gamma \Longrightarrow \mathrm{SEP}(x,A,p) \in \mathrm{SET}}
\VNBfig{sep-mem-intro}{\Gamma \Longrightarrow y \in A \qquad \Gamma \Longrightarrow p[x{:=}y]}{\Gamma \Longrightarrow y \in \mathrm{SEP}(x,A,p)}
\VNBfig{sep-mem-elim}{(\Gamma - F), y \in A, p[x{:=}y] \Longrightarrow G}{\Gamma \Longrightarrow G}
\begin{VNBcond}
\item[Assumption] $F = y \in \mathrm{SEP}(x,A,p)$, $F \in \Gamma$.  It
  is consumed.
\end{VNBcond}

\subsubsection*{\texttt{comp-mem-intro}, \texttt{comp-mem-elim}
  (commands \texttt{comp-mi}, \texttt{comp-me})}

A comprehension is in general a proper class, and its members are sets,
so the domain condition of separation is replaced by membership in
$\mathrm{SET}$.
%% rcs-comp-mem-intro :177, rcs-comp-mem-elim :182, rule-checkers-schema.scm
\VNBfig{comp-mem-intro}{\Gamma \Longrightarrow y \in \mathrm{SET} \qquad \Gamma \Longrightarrow p[x{:=}y]}{\Gamma \Longrightarrow y \in \mathrm{COMP}(x,p)}
\VNBfig{comp-mem-elim}{(\Gamma - F), y \in \mathrm{SET}, p[x{:=}y] \Longrightarrow G}{\Gamma \Longrightarrow G}
\begin{VNBcond}
\item[Assumption] $F = y \in \mathrm{COMP}(x,p)$, $F \in \Gamma$.  It is
  consumed.
\end{VNBcond}
There is no sethood operation for comprehension.  A comprehension is a
set only when it is separately shown to be bounded by one.

\subsubsection*{\texttt{big-union-sethood}, \texttt{big-union-mem-intro},
  \texttt{big-union-mem-elim}
  (commands \texttt{bu-set}, \texttt{bu-mi}, \texttt{bu-me})}

Write $U$ for $\mathrm{BIG\text{-}UNION}(z, A, b)$.
%% rcs-big-union-sethood :194, rcs-big-union-mem-intro :234,
%% rcs-big-union-mem-elim :260, rule-checkers-schema.scm
\VNBfigS{big-union-sethood}{\Gamma \Longrightarrow A \in \mathrm{SET} \\ \Gamma \Longrightarrow \forall v.\, v \in A \Rightarrow b[z{:=}v] \in \mathrm{SET}}{\Gamma \Longrightarrow U \in \mathrm{SET}}
\begin{VNBcond}
\item[Freshness] $v = z$, or $v \notin \FV(A)$.
\item[Not checked] That $v \notin \FV(b)$ when $v \neq z$.  Without it
  the quantifier of the second hypothesis captures a free $v$ of $b$.
  The operation takes $v$ new, free in none of $b$, $A$, $G$ and
  $\Gamma$ (\texttt{pi-big-union-sethood!}, primitive-inferences.scm).
\end{VNBcond}
%% pi-big-union-sethood!, primitive-inferences.scm:1651
\VNBfig{big-union-mem-intro}{\Gamma \Longrightarrow w \in A \qquad \Gamma \Longrightarrow x \in b[z{:=}w]}{\Gamma \Longrightarrow x \in U}
\begin{VNBcond}
\item[Matching] $w$ is read off the first hypothesis.
\end{VNBcond}
\VNBfig{big-union-mem-elim}{(\Gamma - F), e \in A, x \in b[z{:=}e] \Longrightarrow G}{\Gamma \Longrightarrow G}
\begin{VNBcond}
\item[Assumption] $F = x \in U$, $F \in \Gamma$.  It is consumed.
\item[Eigenvariable] $e$ is a variable free in none of $G$, $A$, $b$,
  $x$ and $\Gamma - F$.
\end{VNBcond}
\texttt{big-union-mem-intro} takes the witness $w$ as its argument, and
checks $x \in b[z{:=}w]$ for well-formedness.

\subsubsection*{\texttt{iota-def}, \texttt{iota-in-elim}
  (commands \texttt{iota-d}, \texttt{iota-e})}

$\mathrm{IOTA}(x, p)$ is the unique $x$ satisfying $p$, and denotes
nothing when there is no such $x$ or more than one.  Write $\iota$ for
$\mathrm{IOTA}(x, p)$.  Both operations take the term as their argument;
it need not occur in $G$.  Both grant the defining property
$p[x{:=}\iota]$ as an assumption, and they differ in what pays for it.
%% rcs-iota-def :308, rcs-iota-in-elim :372, rcs--context-defines? :350, rule-checkers-schema.scm
\VNBfigS{iota-def}{\Gamma \Longrightarrow \exists x.\, \bigl(p \wedge \forall y.\, p[x{:=}y] \Rightarrow x = y\bigr) \\ \Gamma, p[x{:=}\iota] \Longrightarrow G}{\Gamma \Longrightarrow G}
\begin{VNBcond}
\item[Freshness] $y$ is a variable, $y \neq x$, $y \notin \FV(p)$.
\end{VNBcond}
\VNBfig{iota-in-elim}{\Gamma, p[x{:=}\iota] \Longrightarrow G}{\Gamma \Longrightarrow G}
\begin{VNBcond}
\item[Denotation] $\Gamma$ contains $\iota \in X$ for some $X$, or an
  equation $\iota = u$ or $u = \iota$.
\end{VNBcond}
\texttt{iota-def} pays with a proof of existence and uniqueness.
\texttt{iota-in-elim} pays with the context: the two forms of the
Denotation condition are the first two clauses of the definedness
certificate, and the operation declines without one of them.

\subsubsection*{\texttt{lambda-type} (command \texttt{lam-t})}

$\mathrm{VNB\text{-}LAMBDA}(\sigma, A, b)$ is the set of pairs
$\{(x, b) : x \in A\}$, with $\sigma$ the binder, a variable or a list of
variables.  The term carries its domain.
%% rcs-lambda-type :437, rcs--lambda-typing-ok? :414, rule-checkers-schema.scm
\VNBfig{lambda-type}{\Gamma \Longrightarrow \forall v.\, v \in A \Rightarrow b[x{:=}v] \in B \qquad \Gamma \Longrightarrow A \in \mathrm{SET}}{\Gamma \Longrightarrow \mathrm{VNB\text{-}LAMBDA}(x, A, b) \in \mathrm{FUN}(A, B)}
\begin{VNBcond}
\item[Domain] The domain the term declares is $\aeq$ the domain of the
  $\mathrm{FUN}$.
\item[Freshness] $v = x$, or $v \notin \FV(b) \cup \FV(B)$.
\item[Binder list] For $\sigma = [x_1, \ldots, x_n]$ the domain is
  $A_1 \times \cdots \times A_n$ with the same $n$, and the first
  hypothesis is the $n$-fold guarded universal
  $\forall v_1.\, v_1 \in A_1 \Rightarrow \cdots \Rightarrow
   \forall v_n.\, v_n \in A_n \Rightarrow b[x_1{:=}v_1, \ldots, x_n{:=}v_n] \in B$,
  each $v_i$ fresh as above and the $v_i$ pairwise distinct.
\end{VNBcond}
The operation has no argument, and declines unless the Domain and Binder
list conditions hold.  Neither condition is decoration.  Without the
first, one term types into $\mathrm{FUN}(A, B)$ for every $A$.  Without
the sethood hypothesis, a function on a proper class, which is itself a
proper class, would be certified into a $\mathrm{FUN}$.

\subsubsection*{\texttt{lambda-beta}, \texttt{lambda-beta-hyp}
  (commands \texttt{lam-b}, \texttt{lam-b-h})}

A \emph{redex} is a term
\[
  \mathrm{VNB\text{-}LAMBDA}(\sigma, A, b)(u_1, \ldots, u_n),
\]
with $n = 1$ for a single binder and $n$ the length of $\sigma$ for a
binder list.  It contracts to $b$ with the $u_i$ substituted in parallel
for the variables of $\sigma$.  Write $o_1, \ldots, o_m$ for the
\emph{obligations}, $m \geq 0$.
%% rcs-lambda-beta :854, rcs-lambda-beta-hyp :881, rcs--beta-match/fuel :701,
%% rcs--beta-licences-ok? :816, rcs--beta-licensed? :775, rule-checkers-schema.scm
\VNBfigS{lambda-beta}{\Gamma \Longrightarrow G' \\ \Gamma \Longrightarrow o_1 \quad \cdots \quad \Gamma \Longrightarrow o_m}{\Gamma \Longrightarrow G}
\VNBfigS{lambda-beta-hyp}{(\Gamma - F), F' \Longrightarrow G \\ \Gamma \Longrightarrow o_1 \quad \cdots \quad \Gamma \Longrightarrow o_m}{\Gamma \Longrightarrow G}
\begin{VNBcond}
\item[Contraction] $G'$ (respectively $F'$, for some $F \in \Gamma$) is
  reached from $G$ (respectively $F$) by a sequence of contractions,
  each taken at a position of the formula reached so far.  A redex that
  a contraction creates may be contracted in its turn.
\item[Licence] Each redex contracted at a position $\pi$ is licensed
  there or owed.  It is \emph{licensed} when $u \in A$ (for $n = 1$) or
  $[u_1, \ldots, u_n] \in A$, or, when $A$ is a product
  $A_1 \times \cdots \times A_n$, every $u_i \in A_i$, is among the
  formulas known at $\pi$: the assumptions of $\Gamma$ in which no
  variable bound between the root and $\pi$ occurs free, together with
  the guards of the guarded universals that enclose $\pi$ and the domain
  memberships of the $\mathrm{VNB\text{-}LAMBDA}$, $\mathrm{SEP}$ and
  $\mathrm{BIG\text{-}UNION}$ binders that enclose it.  It is
  \emph{owed} when that membership is one of the $o_j$ and contains no
  variable bound between the root and $\pi$.  The argument may be taken
  as written or with the redexes inside it contracted; in the second
  case those inner redexes count as contracted and need their own
  licence.
\item[Obligations] Every $o_j$ is owed by some redex.
\item[Bound] The checker's search takes at most 400 contractions, and
  refuses when that runs out.
\end{VNBcond}
\texttt{lambda-beta} takes no argument; \texttt{lambda-beta-hyp} takes
the assumption $F$.  Each reduces every redex it can license or owe, in
one bottom-up pass, and declines when the reduction changes nothing.
The obligations are posted after the reduced formula, one per redex
whose licence is not evident.  A member of $\mathrm{FUN}(A, B)$ is
defined exactly on $A$, so reducing off the domain without the
obligation would make a term denote that the library says does not.

\begin{remark}
The obligation is posted in the context of $N$.  A redex under a binder
that has not yet been introduced therefore owes a membership in which the
bound variable is free, and that leaf cannot be closed.  Introducing the
binder and typing the argument first is not a preference but the
condition under which the obligation is provable.
\end{remark}

\subsection{Rewriting}

\subsubsection*{\texttt{macete}: rewriting the assertion by a theorem
  (commands \texttt{mac}, \texttt{macm})}

A \emph{macete} is a rewrite rule read off an installed theorem of the
form
\[
  \forall x_1 \ldots x_n.\; C_1 \Rightarrow (\cdots (C_k \Rightarrow L = R)),
\]
where $=$ may also be quasi-equality or $\Leftrightarrow$, and the
conditions $C_i$ may be absent.  A name may also denote a rule that is
not a theorem: a functoid $F$ with parameters $\vec p$ and body $e$
denotes $F(\vec p) \to e$, a structure accessor with slot $k$ denotes
$\mathrm{ACC}(s) \to \mathrm{NTH}(k, s)$, and a functor projection
\texttt{F@ACC} denotes the corresponding component of $F$'s body.
%% rkw--macete-definition/compute, rule-checkers-rewrite.scm:593
The \emph{local context} of a position in a formula is $\Gamma$, with
conjunctions split into their conjuncts, extended on the way down:
entering $B$ in $A \Rightarrow B$ adds $A$; entering either conjunct of a
conjunction adds the other; entering either disjunct of $A \vee B$ adds
the negation of the other; entering the scope of a binder
($\forall$, $\exists$, $\mathrm{IOTA}$, $\mathrm{COMP}$, the body of
$\mathrm{SEP}$, $\mathrm{BIG\text{-}UNION}$ or $\mathrm{VNB\text{-}LAMBDA}$,
a functoid) removes every formula in which the bound variable occurs
free (Monk~\cite{Monk1988}).
%% rkw--descend, rule-checkers-rewrite.scm:471
%% rkw--check-macete, rule-checkers-rewrite.scm:687; rkw--condition-discharged? :390
\VNBfig{macete $N$ $L$ $R$}{\Gamma \Longrightarrow G' \qquad \Gamma \Longrightarrow c_1 \quad \cdots \quad \Gamma \Longrightarrow c_m}{\Gamma \Longrightarrow G}
\begin{VNBcond}
\item[Licence] $N$ names an installed theorem, a functoid, an accessor
  or a functor projection, and denotes the rule
  $(C_1, \ldots, C_k;\ L_0 \to R_0)$.
\item[Tag] The recorded $L$ and $R$ are $\sigma L_0$ and $\sigma R_0$
  for one substitution $\sigma$.  (Not tested for the two variadic rules,
  \texttt{union-decompose} and \texttt{intersection-decompose}.)
\item[Rewrite] $G'$ is $G$ with some subterms of the form $\tau L_0$
  replaced by $\tau R_0$, $\tau$ a substitution of the rule's variables
  found by matching.  A pattern $\mathrm{succ}(P)$ matches a positive
  integer literal $m$ by matching $P$ against $m-1$.
\item[Conditions] At each rewritten position every $\tau C_i$ is
  \textsc{truth}, or is in the local context of the position, or is a
  closed arithmetic sentence that evaluates to true, or is one of the
  $c_j$.
\end{VNBcond}
The operation replaces each subterm that matches $L_0$ and whose
conditions hold in the local context; a match whose conditions do not
all hold is left alone, and the rewriter descends into its subterms
instead.  It declines when $G'$ is $G$.  Its local context is the one
above, except that inside a conjunction $A \wedge B$ it extends only the
context of $B$, by $A$.  Hypotheses $c_j$ appear only when the command
posts the conditions it could not discharge as subgoals.  The name of
the theorem and the two patterns are recorded with the operation, so
that the inference can be checked against the rule it claims to apply
(Section~\ref{sec:inference-checked}).

\subsubsection*{\texttt{macete-hyp}: rewriting an assumption
  (commands \texttt{mac-h}, \texttt{mac-h*})}

%% rkw--check-macete-hyp, rule-checkers-rewrite.scm:741
\VNBfig{macete-hyp $N$ $L$ $R$}{(\Gamma - H), H' \Longrightarrow G \qquad \Gamma \Longrightarrow c_1 \quad \cdots \quad \Gamma \Longrightarrow c_m}{\Gamma \Longrightarrow G}
\begin{VNBcond}
\item[Assumption] $H \in \Gamma$.  Exactly one formula of $\Gamma$ is
  missing from the first hypothesis, and at most one formula is new in
  it (none when $H'$ was already in $\Gamma$).
\item[Licence, Tag, Rewrite, Conditions] As for \texttt{macete}, with
  $H$ and $H'$ in place of $G$ and $G'$, and the local context started
  from all of $\Gamma$, $H$ included.
\item[Not checked] The two conditions below on which the operation
  declines.  Neither is needed for soundness: every rule a name denotes
  replaces a term by an equal or quasi-equal one, or a formula by an
  equivalent one, under its conditions.
\end{VNBcond}
The arguments are the name and $H$.  The rewrite is made whether or not
the conditions hold, and one hypothesis $\Gamma \Longrightarrow c_j$ is
posted for each instantiated condition that the local context does not
contain.  Besides declining when the rewrite changes nothing, the
operation declines on two conditions read off the named theorem.  The
theorem's core must be an equivalence, that is an $\Leftrightarrow$ or
one of the two equalities, since replacing an assumption by a
consequence of it would lose information.  And no variable of $R$ or of
a condition may be left undetermined by the match, since such a variable
would be conjured free into the context.

\subsubsection*{\texttt{cartesian-decompose} (command \texttt{mac})}

%% rkw--check-cartesian-decompose :870, rkw--cartesian-shape? :819, rule-checkers-rewrite.scm
\VNBfig{cartesian-decompose}{\Gamma \Longrightarrow G'}{\Gamma \Longrightarrow G}
\begin{VNBcond}
\item[Rewrite] $G'$ is $G$ with some subformulas
  $x \in A_1 \times \cdots \times A_n$, $n \geq 1$, replaced by
  \[
    \exists v_1.\, v_1 \in A_1 \wedge
    \bigl(\cdots \wedge
    \exists v_n.\, v_n \in A_n \wedge x = [v_1, \ldots, v_n] \bigr).
  \]
\item[Freshness] The $v_i$ are pairwise distinct, and none is free in
  $x$ or in any $A_j$.
\end{VNBcond}
No argument beyond the name.  The operation rewrites every such
subformula, with new $v_i$, and declines when $G'$ is $G$.  This is the
membership condition of the product, and it is not a formula in the
catalog: the fresh witnesses cannot be written as a fixed template, so it
is code.  It is on the primitive shelf, and contributes nothing to a
bill.

\subsubsection*{\texttt{tuple-equality-decompose} (command \texttt{mac})}

%% rkw--check-tuple-equality-decompose :902, rkw--tuple-conjunction :883, rule-checkers-rewrite.scm
\VNBfig{tuple-equality-decompose}{\Gamma \Longrightarrow G'}{\Gamma \Longrightarrow G}
\begin{VNBcond}
\item[Rewrite] $G'$ is $G$ with some subformulas
  $[a_1, \ldots, a_n] = [b_1, \ldots, b_n]$, literal tuples of the same
  length, replaced by $a_1 = b_1 \wedge (a_2 = b_2 \wedge \cdots)$; by
  $a_1 = b_1$ when $n = 1$, and by \textsc{truth} when $n = 0$.
\end{VNBcond}
No argument beyond the name.  The operation rewrites every such
subformula and declines when $G'$ is $G$.  Tuples of different lengths
are not rewritten.  That two of them cannot be equal is a separate fact
and not this operation's business.

\subsection{Oracles}

An oracle decides a class of assertions by computation.  Six of the
seven, on success, record an inference with no hypotheses, so $N$ is
grounded at once; \texttt{arith-simplify} posts one hypothesis.
\texttt{ineq} and \texttt{sos} record, with the name, the certificate the
computation produced, and print it; the others record the name alone.
Each inference is checked before it is written, as described in
Section~\ref{sec:inference-checked}.  The figures below state what the
checker of each oracle tests.
%% ineq tag: pi-ineq!, structure-library/ineq-oracle.scm:308;
%% sos tag: structure-library/sos-oracle.scm:120; arith-eval.scm:382-388

\subsubsection*{\texttt{arith-ground}, \texttt{arith-forsome},
  \texttt{arith-simplify} (command \texttt{arith})}

One procedure records these three, trying them in this order.
%% rco--check-arith-ground :242, rco--check-arith-forsome :265,
%% rco--check-arith-simplify :298, rule-checkers-oracle.scm
\VNBfig{arith-ground}{\VNBnohyp}{\Gamma \Longrightarrow G}
\begin{VNBcond}
\item[Evaluation] $G$ is closed, and exact evaluation gives true.
\end{VNBcond}
\VNBfig{arith-forsome}{\VNBnohyp}{\Gamma \Longrightarrow \exists x.\, x \in C \wedge e = x}
\begin{VNBcond}
\item[Evaluation] $e$ evaluates exactly to a number $v$, and the body
  evaluates to true at $x = v$.  The equation may also be written
  $x = e$.
\end{VNBcond}
\VNBfig{arith-simplify}{\Gamma \Longrightarrow G'}{\Gamma \Longrightarrow G}
\begin{VNBcond}
\item[Evaluation] $G'$ is $G$ with some ground subterms replaced by
  the numbers to which they evaluate exactly.
\end{VNBcond}

\texttt{arith-ground} fires when $G$ evaluates to true.  The evaluator
reads $=$ and $\leq$ between terms it can reduce to exact numbers,
membership in $\mathbb{N}$, $\mathbb{Z}$, $\mathbb{Q}$, $\mathbb{R}$ and
$\mathbb{C}$, and the connectives $\neg$, $\wedge$, $\vee$,
$\Rightarrow$.  Everything else is undecided: a free variable, a
quantifier, a biconditional, a strict $<$, or a value that is not exact.
Arithmetic is exact throughout; a computed floating-point value is
refused rather than rounded.

\texttt{arith-forsome} fires on the shape of its figure when $e$ reduces
to a number and that number lies in $C$.

\texttt{arith-simplify} is the fallback.  When replacing the ground
numeric subterms of $G$ by their values changes $G$, it posts the single
hypothesis $\Gamma \Longrightarrow G'$.

The command declines when none of the three applies.

\subsubsection*{\texttt{ring-simplify}: identities of rings
  (commands \texttt{rs}, \texttt{simp})}

%% rco--check-rs :939, rco--crs-number-ok? :881, rco--num-gens-ok? :846, rule-checkers-oracle.scm
\VNBfig{ring-simplify}{\VNBnohyp}{\Gamma \Longrightarrow \forall x_1.\, x_1 \in D_1 \Rightarrow \cdots \forall x_k.\, x_k \in D_k \Rightarrow e_1 = e_2}
\begin{VNBcond}
\item[Normal form] $e_1$ and $e_2$ have the same normal form in the
  free associative ring $\mathbb{Z}[\textit{generators}]$.
\item[Typing] Each $D_i$ is one of $\mathbb{N}$, $\mathbb{Z}$,
  $\mathbb{Q}$, $\mathbb{R}$, $\mathbb{C}$, and every generator of
  $e_1$ or $e_2$ is one of the $x_i$ or is typed in one of these domains
  by an assumption $g \in D$ of $\Gamma$.
\end{VNBcond}
No arguments.  A \emph{generator} is a maximal subterm that ground
evaluation cannot reduce to a number, and multiplication does not
commute.  The operation declines unless both conditions hold.  The
generators counted are those of the two sides as written, including any
that cancel.  Otherwise $x - x = 0$ would close for an $x$ of which
nothing is known, and in VNB a membership is also what says that a term
denotes.

\subsubsection*{\texttt{comm-ring-simplify}: identities of commutative rings
  (command \texttt{crs})}

%% rco--check-crs :929, rco--crs-structure-ok? :894, rule-checkers-oracle.scm
\VNBfig{comm-ring-simplify}{\VNBnohyp}{\Gamma \Longrightarrow \forall x_1.\, x_1 \in D_1 \Rightarrow \cdots \forall x_k.\, x_k \in D_k \Rightarrow e_1 = e_2}
\begin{VNBcond}
\item[Normal form] $e_1$ and $e_2$ have the same normal form in the
  free commutative ring $\mathbb{Z}[\textit{generators}]$, literal
  powers being expanded first.
\item[Typing] Either as for \texttt{ring-simplify}; or the operations
  are those of a structure $R$, $(\mathrm{ADD}\ R)$, $(\mathrm{MUL}\ R)$,
  $(\mathrm{NEG}\ R)$, $(\mathrm{ZERO}\ R)$, $(\mathrm{ONE}\ R)$,
  $\mathrm{IS\text{-}COMMUTATIVE\text{-}RING}(R)$ is a peeled guard or is
  in $\Gamma$, and every generator is a variable bound by a guard
  $x_i \in \mathrm{CARR}(R)$ or is typed in $\mathrm{CARR}(R)$ by an
  assumption.
\end{VNBcond}
No arguments.  The operation tries the number reading first and the
structure reading second.

An existential assertion is outside this operation and every other
oracle except \texttt{arith-forsome}'s one shape: none of them supplies a
witness.  For $\exists a \in \mathbb{Z}.\, (x+y)^3 = x^3 + y \cdot a$ the
witness is given with \texttt{forsome-intro} (command \texttt{ew}), and
the equation that remains is an identity within reach of
\texttt{comm-ring-simplify}.

\subsubsection*{\texttt{ineq}: an oracle for linear arithmetic over the reals
  (command \texttt{ineq})}

%% rco--check-ineq :556, rco--farkas-ok? :515, rco--ineq-constraint :494, rule-checkers-oracle.scm
\VNBfig{ineq $c$}{\VNBnohyp}{\Gamma \Longrightarrow \forall x_1.\, x_1 \in \mathbb{R} \Rightarrow \cdots \forall x_k.\, x_k \in \mathbb{R} \Rightarrow G_0}
\begin{VNBcond}
\item[Certificate] $c$ is a list of pairs $(\mathit{id}\ .\ m)$,
  each $m$ an exact nonnegative rational, where $\mathit{id}$ is
  \texttt{goal}, the negation of $G_0$; or \texttt{h}$I$, the $I$-th
  assumption of $N$ counted from one; or \texttt{h}$I$\texttt{-rev}, the
  $I$-th assumption, an equation, read from right to left.  When $G_0$
  is an equation, $c$ is a pair of two such lists, one for each
  direction.
\item[Premises] Each assumption the certificate names is an equation or
  a $\leq$ or $<$ inequality.
\item[Absurdity] The sum of the named constraints, each written
  $\ell \leq 0$ or $\ell < 0$ and multiplied by its $m$, has every atom's
  coefficient zero, and its constant is $> 0$, or $\geq 0$ when a strict
  constraint enters with positive multiplier.
\item[Typing] Every atom of a constraint that enters with a positive
  multiplier is certified real: $\Gamma$ contains $a \in \mathbb{R}$, or $a$ is one of the
  $x_i$, or $a$ is $|t|$ with every atom of $t$ certified.
\item[Eigenvariable] No $x_i$ is free in a premise the certificate
  names.
\end{VNBcond}
The arguments of the operation are the positions in $\Gamma$ of the
assumptions to be used as premises, counted from one.  When the
procedure succeeds $N$ is grounded, and otherwise the operation
declines.

The assertion and each named premise are read as linear equations or
inequalities in their \emph{atoms}, the maximal subterms that are not
sums, differences, or products by a numeral.  A named premise that does
not read this way is left out, which can only make the procedure prove
less; an index outside the range of $\Gamma$ is an error and abandons the
call.  The operation requires the Typing condition of every atom, an
atom that cancels included, and declines when a variable bound by one of
the leading quantifiers occurs free in \emph{any} assumption of $N$.
Fourier--Motzkin elimination is run on the premises and the negation of
the assertion.  The procedure succeeds when the elimination reaches
$0 < 0$.  The nonnegative combination of premises that produces it, a
\emph{Farkas certificate}, is printed and is recorded in the graph with
the name of the operation.

\begin{remark}
The condition on bound variables is the eigenvariable condition of
\texttt{forall-intro}.  Until 2026-09-20 the oracle did not impose it,
and from $x \in \mathbb{R},\; x \leq 0$ it closed
$\forall x \in \mathbb{R}.\; x \leq 0$.  The defect was found when the
checker of Section~\ref{sec:inference-checked} was written for this
operation, and no proof in the library had used it.
\end{remark}

\subsubsection*{\texttt{sos}: nonnegativity by a sum of squares
  (command \texttt{sos})}

%% rco--check-sos :956, rule-checkers-oracle.scm; tag built at structure-library/sos-oracle.scm:120
\VNBfig{sos $c$}{\VNBnohyp}{\Gamma \Longrightarrow \forall x_1.\, x_1 \in \mathbb{R} \Rightarrow \cdots \forall x_k.\, x_k \in \mathbb{R} \Rightarrow a \leq b}
\begin{VNBcond}
\item[Certificate] $c$ is a list of pairs $(c_i\ .\ \lambda_i)$, each
  $\lambda_i$ an exact nonnegative rational.
\item[Identity] $b - a = \lambda_1 c_1^2 + \cdots + \lambda_n c_n^2$ as
  polynomials in the free commutative ring over the generators.
\item[Typing] Every generator of $a$, $b$ and the $c_i$ is certified
  real, as for \texttt{ineq}.
\end{VNBcond}
The assertion may also be written $b \geq a$.  The arguments of the
operation are the terms $c_1, \ldots, c_n$ to be squared.  An exact
linear program decides whether weights $\lambda_i \geq 0$ exist; on
success the pairs $(c_i\ .\ \lambda_i)$ are printed and recorded in the
graph with the name.

The oracle verifies a certificate, it does not search for one: the
squares are supplied by the caller.  A strict inequality is not accepted,
since a square may be zero.

\subsection{Which command reaches which operation}\label{sec:command-uses}

A tactic has no access to the graph except through the operations above,
so the operations one tactic can reach are the bound on what it can do.
The table below gives that bound for each proof command: the operations
recorded by any kernel procedure reachable from the command's own code,
whether or not the library happens to exercise it.  It is produced from
\texttt{reference/kernel-map-static.sexp}, the result of reading the
source of every loaded file with the Scheme reader.  One thing is left
out of every row: after any command that changes a proof, the system tries
to close the definedness obligations that instantiation has left
(Section~\ref{sec:kernel-ops-detail}, \texttt{forall-elim}), and that step
can reach twenty kernel procedures of its own.  It is the same for every
command and is stated once in \texttt{reference/KERNEL-MAP.md}.

A command listed as \texttt{none} records nothing.  Those are the
navigation, search and session commands.

%% The table and its counts, generated by gen-tactic-uses.py from
%% reference/kernel-map-static.sexp and the command registry.
%% GENERATED by gen-tactic-uses.py from reference/kernel-map-static.sexp and
%% tactics-help.scm (*tactic-kind*).  Do not edit.
\newcommand{\VNBtacticusescount}{105}
\newcommand{\VNBtacticusesmax}{64}
\newcommand{\VNBtacticusestable}{%
\begin{longtable}{@{}p{0.17\textwidth}p{0.79\textwidth}@{}}
\toprule
command & operations it can reach \\
\midrule
\endfirsthead
\toprule
command & operations it can reach \\
\midrule
\endhead
\texttt{ai} & \texttt{and-elim} \texttt{forsome-elim} \texttt{iff-elim} \texttt{not-elim} \texttt{or-elim}\\
\texttt{apply-thm} & \texttt{forall-elim} \texttt{theorem-assumption}\\
\texttt{arith} & \texttt{arith-forsome} \texttt{arith-ground} \texttt{arith-simplify} \texttt{cartesian-decompose} \texttt{eq-subst} \texttt{macete} \texttt{tuple-equality-decompose}\\
\texttt{arith--decide} & \texttt{arith-forsome} \texttt{arith-ground} \texttt{arith-simplify}\\
\texttt{ass} & \texttt{assumption}\\
\texttt{ass-all} & \texttt{assumption}\\
\texttt{backup-one} & none\\
\texttt{bc} & \texttt{backchain}\\
\texttt{bc*} & \texttt{assumption} \texttt{backchain} \texttt{cut} \texttt{forall-elim} \texttt{theorem-assumption}\\
\texttt{bc*--attempt} & \texttt{assumption} \texttt{backchain} \texttt{cut} \texttt{forall-elim} \texttt{theorem-assumption}\\
\texttt{beta} & \texttt{functoid-beta}\\
\texttt{bu-me} & \texttt{big-union-mem-elim}\\
\texttt{bu-mi} & \texttt{big-union-mem-intro}\\
\texttt{bu-set} & \texttt{big-union-sethood}\\
\texttt{calc} & \emph{all 64 reachable operations}\\
\texttt{ce} & \texttt{cartesian-elim}\\
\texttt{choose} & \texttt{and-elim} \texttt{and-intro} \texttt{arith-forsome} \texttt{arith-ground} \texttt{arith-simplify} \texttt{assumption} \texttt{cartesian-decompose} \texttt{cut} \texttt{detach} \texttt{eq-subst} \texttt{forall-elim} \texttt{forall-intro} \texttt{forsome-elim} \texttt{forsome-intro} \texttt{iff-elim} \texttt{iff-intro} \texttt{implies-intro} \texttt{macete} \texttt{not-elim} \texttt{not-intro} \texttt{or-elim} \texttt{theorem-assumption} \texttt{tuple-equality-decompose}\\
\texttt{choose-pos} & \texttt{and-elim} \texttt{and-intro} \texttt{arith-forsome} \texttt{arith-ground} \texttt{arith-simplify} \texttt{assumption} \texttt{cartesian-decompose} \texttt{cut} \texttt{detach} \texttt{eq-subst} \texttt{forall-elim} \texttt{forall-intro} \texttt{forsome-elim} \texttt{forsome-intro} \texttt{iff-elim} \texttt{iff-intro} \texttt{implies-intro} \texttt{macete} \texttt{not-elim} \texttt{not-intro} \texttt{or-elim} \texttt{theorem-assumption} \texttt{tuple-equality-decompose}\\
\texttt{ci} & \texttt{cartesian-intro}\\
\texttt{comp-me} & \texttt{comp-mem-elim}\\
\texttt{comp-mi} & \texttt{comp-mem-intro}\\
\texttt{contra} & \texttt{and-elim} \texttt{and-intro} \texttt{arith-forsome} \texttt{arith-ground} \texttt{arith-simplify} \texttt{assumption} \texttt{cartesian-decompose} \texttt{cut} \texttt{detach} \texttt{eq-subst} \texttt{forall-elim} \texttt{forall-intro} \texttt{forsome-elim} \texttt{iff-elim} \texttt{iff-intro} \texttt{implies-intro} \texttt{ineq} \texttt{macete} \texttt{not-elim} \texttt{not-intro} \texttt{or-elim} \texttt{theorem-assumption} \texttt{tuple-equality-decompose}\\
\texttt{crs} & \texttt{comm-ring-simplify}\\
\texttt{cut} & \texttt{cut}\\
\texttt{detach!} & \texttt{detach}\\
\texttt{di} & \texttt{and-intro} \texttt{forall-intro} \texttt{iff-intro} \texttt{implies-intro} \texttt{not-intro}\\
\texttt{dial} & none\\
\texttt{dial-wff} & none\\
\texttt{dk-focus!} & none\\
\texttt{eps-part} & \texttt{and-elim} \texttt{and-intro} \texttt{arith-forsome} \texttt{arith-ground} \texttt{arith-simplify} \texttt{assumption} \texttt{cartesian-decompose} \texttt{cut} \texttt{detach} \texttt{eq-subst} \texttt{forall-elim} \texttt{forall-intro} \texttt{forsome-elim} \texttt{iff-elim} \texttt{iff-intro} \texttt{implies-intro} \texttt{macete} \texttt{not-elim} \texttt{not-intro} \texttt{or-elim} \texttt{theorem-assumption} \texttt{tuple-equality-decompose}\\
\texttt{ew} & \texttt{forsome-intro}\\
\texttt{fact} & \texttt{arith-forsome} \texttt{arith-ground} \texttt{arith-simplify} \texttt{cut} \texttt{detach} \texttt{forall-elim} \texttt{theorem-assumption}\\
\texttt{focus} & none\\
\texttt{focus-id} & none\\
\texttt{goal-status} & none\\
\texttt{grind} & \texttt{and-elim} \texttt{and-intro} \texttt{forall-intro} \texttt{forsome-elim} \texttt{iff-elim} \texttt{iff-intro} \texttt{implies-intro} \texttt{macete-hyp} \texttt{not-elim} \texttt{not-intro} \texttt{or-elim}\\
\texttt{grind-and-mp} & \texttt{and-elim} \texttt{and-intro} \texttt{arith-forsome} \texttt{arith-ground} \texttt{arith-simplify} \texttt{assumption} \texttt{cut} \texttt{detach} \texttt{forall-elim} \texttt{forall-intro} \texttt{forsome-elim} \texttt{iff-elim} \texttt{iff-intro} \texttt{implies-intro} \texttt{macete-hyp} \texttt{not-elim} \texttt{not-intro} \texttt{or-elim}\\
\texttt{have!} & \texttt{and-intro} \texttt{arith-forsome} \texttt{arith-ground} \texttt{arith-simplify} \texttt{assumption} \texttt{cartesian-decompose} \texttt{cut} \texttt{detach} \texttt{eq-subst} \texttt{forall-elim} \texttt{forall-intro} \texttt{iff-intro} \texttt{implies-intro} \texttt{macete} \texttt{not-intro} \texttt{theorem-assumption} \texttt{tuple-equality-decompose}\\
\texttt{ie} & \texttt{intersection-elim}\\
\texttt{if-false} & \texttt{if-false}\\
\texttt{if-true} & \texttt{if-true}\\
\texttt{ii} & \texttt{intersection-intro}\\
\texttt{in-rr} & \texttt{and-intro} \texttt{arith-forsome} \texttt{arith-ground} \texttt{arith-simplify} \texttt{assumption} \texttt{cartesian-decompose} \texttt{cartesian-intro} \texttt{cut} \texttt{detach} \texttt{eq-subst} \texttt{forall-elim} \texttt{forall-intro} \texttt{iff-intro} \texttt{implies-intro} \texttt{macete} \texttt{not-intro} \texttt{theorem-assumption} \texttt{tuple-equality-decompose}\\
\texttt{in-rr--refocus!} & none\\
\texttt{ineq} & \texttt{ineq}\\
\texttt{inst} & \texttt{forall-elim}\\
\texttt{inst+} & \texttt{arith-forsome} \texttt{arith-ground} \texttt{arith-simplify} \texttt{cut} \texttt{detach} \texttt{forall-elim}\\
\texttt{iota-d} & \texttt{iota-def}\\
\texttt{iota-e} & \texttt{iota-in-elim}\\
\texttt{lam-b} & \texttt{lambda-beta}\\
\texttt{lam-b-h} & \texttt{lambda-beta-hyp}\\
\texttt{lam-t} & \texttt{lambda-type}\\
\texttt{len-r} & \texttt{length-reduce}\\
\texttt{mac} & \texttt{cartesian-decompose} \texttt{macete} \texttt{tuple-equality-decompose}\\
\texttt{mac-h} & \texttt{macete-hyp}\\
\texttt{mac-h*} & \texttt{and-elim} \texttt{forsome-elim} \texttt{iff-elim} \texttt{macete-hyp} \texttt{not-elim} \texttt{or-elim}\\
\texttt{macm} & \texttt{cartesian-decompose} \texttt{macete} \texttt{tuple-equality-decompose}\\
\texttt{minimize!} & \texttt{and-elim} \texttt{and-intro} \texttt{arith-forsome} \texttt{arith-ground} \texttt{arith-simplify} \texttt{assumption} \texttt{cartesian-decompose} \texttt{cut} \texttt{detach} \texttt{eq-subst} \texttt{forall-elim} \texttt{forall-intro} \texttt{forsome-elim} \texttt{forsome-intro} \texttt{iff-elim} \texttt{iff-intro} \texttt{implies-intro} \texttt{macete} \texttt{not-elim} \texttt{not-intro} \texttt{or-elim} \texttt{reflexivity} \texttt{sep-mem-elim} \texttt{sep-mem-intro} \texttt{theorem-assumption} \texttt{tuple-equality-decompose}\\
\texttt{mp} & \texttt{arith-forsome} \texttt{arith-ground} \texttt{arith-simplify} \texttt{assumption} \texttt{cut} \texttt{detach} \texttt{forall-elim}\\
\texttt{ni} & \texttt{nn-induction}\\
\texttt{nth-r} & \texttt{nth-reduce}\\
\texttt{obtain} & \emph{all 64 reachable operations}\\
\texttt{obtain-at} & \texttt{and-elim} \texttt{and-intro} \texttt{arith-forsome} \texttt{arith-ground} \texttt{arith-simplify} \texttt{assumption} \texttt{cartesian-decompose} \texttt{cut} \texttt{detach} \texttt{eq-subst} \texttt{forall-elim} \texttt{forall-intro} \texttt{forsome-elim} \texttt{iff-elim} \texttt{iff-intro} \texttt{implies-intro} \texttt{macete} \texttt{not-elim} \texttt{not-intro} \texttt{or-elim} \texttt{theorem-assumption} \texttt{tuple-equality-decompose}\\
\texttt{oi-l} & \texttt{or-intro-left}\\
\texttt{oi-r} & \texttt{or-intro-right}\\
\texttt{orelse} & none\\
\texttt{pbc} & \texttt{proof-by-contradiction}\\
\texttt{prep} & none\\
\texttt{prop} & \texttt{and-elim} \texttt{and-intro} \texttt{arith-forsome} \texttt{arith-ground} \texttt{arith-simplify} \texttt{assumption} \texttt{cartesian-decompose} \texttt{cut} \texttt{detach} \texttt{eq-subst} \texttt{forall-elim} \texttt{forall-intro} \texttt{forsome-elim} \texttt{iff-elim} \texttt{iff-intro} \texttt{implies-intro} \texttt{macete} \texttt{not-elim} \texttt{not-intro} \texttt{or-elim} \texttt{or-intro-left} \texttt{or-intro-right} \texttt{proof-by-contradiction} \texttt{theorem-assumption} \texttt{tuple-equality-decompose}\\
\texttt{push-not-h} & \texttt{and-elim} \texttt{and-intro} \texttt{arith-forsome} \texttt{arith-ground} \texttt{arith-simplify} \texttt{assumption} \texttt{cartesian-decompose} \texttt{cut} \texttt{detach} \texttt{eq-subst} \texttt{forall-elim} \texttt{forall-intro} \texttt{forsome-elim} \texttt{forsome-intro} \texttt{iff-elim} \texttt{iff-intro} \texttt{implies-intro} \texttt{macete} \texttt{not-elim} \texttt{not-intro} \texttt{or-elim} \texttt{or-intro-left} \texttt{or-intro-right} \texttt{proof-by-contradiction} \texttt{theorem-assumption} \texttt{tuple-equality-decompose}\\
\texttt{qed} & none\\
\texttt{qrfl} & \texttt{quasi-reflexivity}\\
\texttt{quietly} & none\\
\texttt{repeat} & none\\
\texttt{replay-proof} & \emph{all 64 reachable operations}\\
\texttt{rfl} & \texttt{reflexivity}\\
\texttt{rs} & \texttt{ring-simplify}\\
\texttt{save-proof} & none\\
\texttt{scout} & \emph{all 64 reachable operations}\\
\texttt{scout-run} & none\\
\texttt{scout-show} & \emph{all 64 reachable operations}\\
\texttt{sep-me} & \texttt{sep-mem-elim}\\
\texttt{sep-mi} & \texttt{sep-mem-intro}\\
\texttt{sep-set} & \texttt{sep-sethood}\\
\texttt{show} & none\\
\texttt{simp} & \texttt{comm-ring-simplify} \texttt{cut} \texttt{eq-subst}\\
\texttt{slot} & \texttt{cartesian-decompose} \texttt{macete} \texttt{tuple-equality-decompose}\\
\texttt{slot-h} & \texttt{macete-hyp}\\
\texttt{sos} & \texttt{sos}\\
\texttt{sp} & none\\
\texttt{subst} & \texttt{eq-subst}\\
\texttt{supply} & \emph{all 64 reachable operations}\\
\texttt{ta} & \texttt{theorem-assumption}\\
\texttt{te} & \texttt{tuples-elim}\\
\texttt{tfi} & \texttt{transfinite-induction}\\
\texttt{tfi3} & \texttt{transfinite-induction-3cases}\\
\texttt{ti} & \texttt{tuples-intro}\\
\texttt{ue} & \texttt{union-elim}\\
\texttt{ui} & \texttt{union-intro}\\
\texttt{undo} & none\\
\texttt{use-at} & \texttt{and-intro} \texttt{arith-forsome} \texttt{arith-ground} \texttt{arith-simplify} \texttt{assumption} \texttt{cartesian-decompose} \texttt{cut} \texttt{detach} \texttt{eq-subst} \texttt{forall-elim} \texttt{forall-intro} \texttt{iff-intro} \texttt{implies-intro} \texttt{macete} \texttt{not-intro} \texttt{theorem-assumption} \texttt{tuple-equality-decompose}\\
\texttt{use-em} & \texttt{and-elim} \texttt{and-intro} \texttt{arith-forsome} \texttt{arith-ground} \texttt{arith-simplify} \texttt{assumption} \texttt{cartesian-decompose} \texttt{cut} \texttt{detach} \texttt{eq-subst} \texttt{forall-elim} \texttt{forall-intro} \texttt{forsome-elim} \texttt{iff-elim} \texttt{iff-intro} \texttt{implies-intro} \texttt{macete} \texttt{not-elim} \texttt{not-intro} \texttt{or-elim} \texttt{or-intro-left} \texttt{or-intro-right} \texttt{proof-by-contradiction} \texttt{theorem-assumption} \texttt{tuple-equality-decompose}\\
\texttt{vlet} & \texttt{and-elim} \texttt{cut} \texttt{forsome-elim} \texttt{iff-elim} \texttt{not-elim} \texttt{or-elim}\\
\texttt{wbc} & \texttt{and-elim} \texttt{arith-forsome} \texttt{arith-ground} \texttt{arith-simplify} \texttt{cut} \texttt{detach} \texttt{forall-elim} \texttt{forsome-elim} \texttt{iff-elim} \texttt{not-elim} \texttt{or-elim} \texttt{theorem-assumption}\\
\texttt{wk} & \texttt{weakening}\\
\bottomrule
\end{longtable}
\noindent\footnotesize \textsuperscript{r}: read from the command registry rather
than from the kernel map, which predates the command.\normalsize}

\VNBtacticusestable

\subsection{How an inference is recorded, and checked}\label{sec:inference-checked}

Every kernel operation ends by calling one procedure,
\texttt{dg-apply-rule!}\ (\texttt{deduction-graphs.scm}).  It is handed
the rule, the hypothesis sequents and the node $N$.  It checks the
inference (below); then it posts each hypothesis sequent as a node,
reusing a node the graph already has for the same sequent; then it
creates the inference node and writes the arrows.  No other procedure
writes to the graph.
%% dg-apply-rule!, deduction-graphs.scm:375; make-inference-node :59
The inference node records exactly three things: the rule, the list of
hypothesis nodes, and the conclusion node $N$.  The rule is the name of
the operation, or a list headed by that name when the operation carries
data about the application: an index (\texttt{(cartesian-elim 2)}), the
name and the patterns of a macete (\texttt{(macete \textit{name} $L$
$R$)}), or the certificate of an oracle (\texttt{(ineq $c$)},
\texttt{(sos $c$)}).

\subsubsection*{The construction and the check}

An operation has two parts.  The \emph{builder} is a procedure
\texttt{pi-\ldots!}\ (in \texttt{primitive-inferences.scm},
\texttt{macetes.scm} and the oracle files).  It reads $N$ and the
operation's argument and \emph{constructs} the hypothesis sequents: for
\texttt{and-intro}, \texttt{pi-direct-inference!} takes the two conjuncts
of $G$ and makes $\Gamma \Longrightarrow A$ and
$\Gamma \Longrightarrow B$; for \texttt{forall-intro} it also chooses the
eigenvariables; for an oracle it runs the decision procedure.  The
\emph{checker} is a separate procedure, registered for the operation's
name in one of four files (\texttt{rule-checkers-logic.scm},
\texttt{-schema.scm}, \texttt{-rewrite.scm}, \texttt{-oracle.scm}).
\texttt{dg-apply-rule!}\ hands it the finished result: the rule, the
hypothesis sequents and the sequent of $N$.  The checker does not call
the builder and does not compute hypotheses from $N$.  It decides whether
the sequents it is given stand in the relation of the operation's figure
(Section~\ref{sec:kernel-ops-detail}), and answers yes, or no with a
reason.

The checkers of the logic and schema operations use four things from the
expression layer: substitution, the free-variable computation, equality
of formulas up to renaming of bound variables ($\aeq$), and one-way
matching.  The rewrite checkers add two: a theorem is read as a rewrite
rule by the same procedure the rewriter uses
(\texttt{extract-rewrite-patterns}), because that reading is what the
name of a macete denotes; and a closed arithmetic condition is evaluated
by the arithmetic evaluator.  The oracle checkers carry their own exact
arithmetic, linear forms and polynomials, written separately from the
oracles.
%% headers of rule-checkers-logic.scm (1-13), rule-checkers-rewrite.scm (18-26),
%% rule-checkers-oracle.scm (1-14)

Two of the relations, written out as the checker tests them:
%% and-intro checker, rule-checkers-logic.scm:347; forall-elim checker :702
\begin{itemize}
\item \texttt{and-intro} accepts hypotheses $H_1, H_2$ and conclusion
  $\Gamma \Longrightarrow G$ exactly when $G$ is $A \wedge B$, there are
  two hypotheses, the context of each equals $\Gamma$ as a set, the
  assertion of $H_1$ is $\aeq A$ and that of $H_2$ is $\aeq B$.
\item \texttt{forall-elim} accepts exactly when there are one or two
  hypotheses, the assertion of $H_1$ is $\aeq G$, and either the context
  of $H_1$ is included in $\Gamma$, or $\Gamma$ contains a formula
  $\forall x.\, B$ and the context of $H_1$ is $\Gamma \cup \{D\}$ for a
  formula $D$ that matching shows to be $B[x{:=}t]$ for some term $t$;
  and, if there is a second hypothesis, it is
  $\Gamma \Longrightarrow u = u$ for some term $u$.
\end{itemize}

The eigenvariable condition of \texttt{forall-intro} is tested as
follows.  The checker counts the length of
the chain of leading universal quantifiers of $G$, passing through each
guard $x \in A$ on the variable just bound, and tries that length $k$
first and then each smaller one.  For a given $k$ it reads the
eigenvariables $y_1, \ldots, y_k$ off the hypothesis
$\Gamma' \Longrightarrow B'$ by matching: $G$ with its $k$ bound
variables replaced by holes is matched, guard by guard, against the
formulas of $\Gamma'$ that are not in $\Gamma$, and its body against
$B'$.  A hole that matching does not reach belongs to a variable absent
from the body, and is given the bound variable's own name.  It then
walks the chain forwards with a set $V$ of excluded variables, which
starts as $\FV(\Gamma)$.  At step $i$ the formula is
$\forall x_i.\, C_i$, and $y_i$ must be a variable, not in $V$, and not
free in $\forall x_i.\, C_i$.  If $C_i[x_i{:=}y_i]$ is
$(y_i \in A_i) \Rightarrow D$, the guard $y_i \in A_i$ is collected, its
free variables ($y_i$ and those of $A_i$) are added to $V$, and the walk
continues with $D$; otherwise $y_i$ is added to $V$ and the walk
continues with $C_i[x_i{:=}y_i]$.  The inference is accepted when the
formula reached after $k$ steps is $\aeq B'$ and $\Gamma'$ is $\Gamma$
together with the collected guards.
%% chk-forall-intro-at, rule-checkers-logic.scm:514; chk-forall-eigen-ok? :491;
%% chk-forall-fill :477; chk-vacuous-name :540; the k-loop in the checker at :555

If no checker is registered for the name, or the checker refuses,
\texttt{dg-apply-rule!}\ raises an error and the graph is unchanged.  A
refusal is printed even when output is suppressed, and the load ends by
reporting the number of inferences verified and the number refused.
%% dg-check-inference!, dg-check-report-refusal!, deduction-graphs.scm

An error in a builder goes unnoticed only if the checker's relation also
admits the inference the builder produced.  Where a figure of
Section~\ref{sec:kernel-ops-detail} carries a condition labelled
\emph{Not checked}, the relation is weaker than the operation, and at
that point the operation is the only safeguard.  There are four such
places: definedness in \texttt{forall-elim} and \texttt{reflexivity}, by
policy, since the certificate decides when an obligation may be omitted
and the checker accepts the inference either way; capture in
\texttt{eq-subst}; the variable of the second hypothesis of
\texttt{big-union-sethood}; and the two refusals of
\texttt{macete-hyp}, which are not needed for soundness.
%% policy: header of rule-checkers-logic.scm, lines 29-43

\subsubsection*{Oracles: a certificate check or a second decision}

An oracle has no proof to replay, so what its checker can establish
depends on what the oracle records.  There are two cases.
%% header of rule-checkers-oracle.scm, lines 22-62

\emph{Certificate check} (\texttt{ineq}, \texttt{sos}).  The oracle
records, with its name, the object that shows the assertion true, and
the checker verifies that object by arithmetic of its own, without
elimination, search or any code of the oracle.  For \texttt{ineq} the
object is a Farkas certificate: the checker reads the assertion and the
named premises as linear forms, multiplies each by its multiplier, adds,
and confirms that every atom cancels and that the constant left is an
absurd inequality.  For \texttt{sos} it is the list of squares and
weights: the checker expands $\lambda_1 c_1^2 + \cdots + \lambda_n c_n^2$
and $b - a$ as polynomials and compares them monomial by monomial.  Either
the arithmetic confirms the certificate, or the inference is refused.
%% rco--farkas-ok? :515, rco--check-sos :956, rule-checkers-oracle.scm

\begin{sloppypar}
\emph{Second decision} (\texttt{arith-ground}, \texttt{arith-forsome},
\texttt{arith-simplify}, \texttt{comm-ring-simplify},
\texttt{ring-simplify}).  These oracles record their name only, so the
checker decides the same question again with a second implementation of
the same mathematics: exact evaluation for \texttt{arith}, expansion of
both sides to a sum of monomials for the two ring oracles.  The two
implementations share no code beyond the expression layer.  A second
decision catches an error in how the oracle reads its input or reports
its answer: an atom misread, a typing it failed to require, a premise
dropped, a case it mishandles that the second implementation handles
correctly.  It does not catch an error that both implementations make
the same way, for instance a mistake in the mathematics of the normal
form itself.
\end{sloppypar}
%% rco--check-arith-ground :242, rco--check-crs :929, rco--check-rs :939

Neither kind of check adds proving power.  A checker only judges an
inference that an oracle has already produced; an assertion the oracle
declines is not proved by its checker.

\subsubsection*{Where a certificate is kept, and how it is shown}

The certificate of \texttt{ineq} is part of the rule recorded in the
inference node, beside the name of the operation.  The rule has the form
\texttt{(ineq $c$)}, where $c$ is a list of pairs
\texttt{($\mathit{id}$ .\ $m$)}: $m$ is an exact nonnegative rational,
and $\mathit{id}$ is \texttt{goal} (the negation of the assertion),
\texttt{h}$I$ (the $I$-th assumption of $N$, counted from one), or
\texttt{h}$I$\texttt{-rev} (the $I$-th assumption, an equation, read from
right to left).  For an assertion that is an equation, $c$ is a pair of
two such lists, one for each direction.  For example
\texttt{(ineq ((h1 .\ 1) (h3 .\ 2) (goal .\ 1)))} says that the first
assumption, twice the third, and the negated assertion add up to a
false inequality between constants.  The rule of \texttt{sos} is
\texttt{(sos (($c_1$ .\ $\lambda_1$) \ldots))}.
%% pi-ineq!, structure-library/ineq-oracle.scm:218 (tag built at :308);
%% ineq-hyp-constraints :167; sos-oracle.scm:120

The certificate therefore belongs to the deduction graph of the proof,
and lasts as long as that graph.  What \texttt{qed} keeps of a proof is
the statement, the script and the bill; the script records the command
as it was typed, for instance \texttt{(ineq 1 3)}, not the certificate.
Replaying the script, as the page audit does, runs the oracle again, and
the checker verifies the certificate recorded by that run.
%% proof-debt.scm:205 (*proof-oracles* holds oracle names, not certificates)

When it succeeds the oracle prints the certificate:
\begin{VNBcode}
;; ineq: closed by Farkas certificate ((h1 . 1) (h3 . 2) (goal . 1))
\end{VNBcode}
\texttt{sos} prints its decomposition in the same way.  \texttt{proof-tex}
prints, for each step, the command and the sequent it leaves, read from
the command trace recorded while the proof ran; it does not print the
certificate.  The command \texttt{(proof-rules)} reads the graph of the
current proof and lists the rules its inferences recorded.
%% ineq-report, ineq-oracle.scm:215; sos-report, sos-oracle.scm:49;
%% proof-tex--steps-from-live, proof-tex.scm:94; proof-rules, suggest.scm:5795

\subsubsection*{Controls}

Each checker has two controls in the test suite: the proofs of the
library, which it must accept, and an inference that is wrong in a
stated way, which it must refuse.  A checker that accepted everything
would pass the first and fail the second.

\begin{remark}
Checking is on by default.  Setting the environment variable
\texttt{VNB\_NO\_RULE\_CHECK} turns it off.  On the library as loaded on
2026-09-20, 197{,}677 inferences were verified and none refused, and the
load took 30\% longer with checking than without.
\end{remark}

\subsection{The boundary is checked, not merely documented}\label{sec:kernel-checked}

The restriction would be worth little if it were a convention.  The eight
files permitted to record inferences are named in a list in the source
(\texttt{*kernel-caller-files*}).  Every time the system starts, it reads
all of its own source with the Scheme reader and counts, in each file,
the calls to \texttt{dg-apply-rule!}\ and any use of that procedure as a
value --- the second being the way a file could otherwise hand the
capability to code outside the list.  If any file not on the list records
inferences, or if any file passes the procedure on, the system reports
the offending files and \emph{refuses to finish starting}.  A reader is
used rather than a text search so that a mention inside a comment, a
string or quoted data is not mistaken for code.

On a clean tree the check prints one line:

\begin{VNBcode}
;; kernel-callers-audit: ok (69 call site(s) to dg-apply-rule!
;;   in 8 file(s), all allowed; 0 value use(s))
\end{VNBcode}

A companion check, \texttt{kernel-rules-audit}, compares the rules
actually recorded during the load against the documented list in both
directions, so a rule that is trusted but undocumented, or documented but
unreachable, is reported rather than left to be noticed.

Widening the trusted base is possible --- a new decision procedure is a
legitimate thing to add --- but it cannot happen quietly: the new file
must be added to the list by hand, which is a deliberate act that leaves
a record.  The complete inventory (every entry point, the rules it may
record, and which proof commands reach it) is
\texttt{reference/KERNEL-MAP.md}.

%% =========================================================
\section{Starting and inspecting a proof}

\begin{VNBcode}
; Start a new proof of a formula
(define ps (start-proof (make-wff-from-string
  "n in NN and m in NN implies n in NN")))

; Display the current state
(print-proof-state ps)

; Check if done
(proof-done? ps)   ; => #f or #t

; List open goals
(proof-open-goals ps)
\end{VNBcode}

The variable \texttt{ps} holds the \emph{proof state}: a deduction graph
plus the current \emph{focus} sequent node (the subgoal currently being
worked on).  Every proof command returns a new proof state with an
updated focus.

\texttt{start-proof} takes a \texttt{wff} (produced by \texttt{make-wff} or \texttt{make-wff-from-string}).
The root sequent has no assumptions: a proof begins from the bare goal,
and every hypothesis arrives through an inference.  All sequent assumptions and assertions are
\texttt{wff} objects throughout the proof; the library name propagates
automatically to every subformula produced by inference rules.

%% =========================================================
\section{Proof commands}
\label{sec:commands}

Commands are Scheme procedures of the form
\[
  \texttt{(cmd-}\mathit{name}\;\; \mathit{ps}\;\; \mathit{arg}_1 \ldots\texttt{)}
\]
Each command applies one primitive inference to the focus node, then
advances the focus to the first new ungrounded subgoal.

\subsection{Logical connectives}

\begin{itemize}

\item \texttt{(cmd-direct-inference ps)}
  Decompose the goal by its top-level connective:
  \begin{itemize}
  \item \texttt{p and q} $\to$ two subgoals $p$, $q$.
  \item \texttt{p implies q} $\to$ assume $p$, prove $q$.
  \item \texttt{p iff q} $\to$ two subgoals (each direction).
  \item \texttt{not(p)} $\to$ assume $p$, prove \texttt{falsity}.
  \item \texttt{forall([\ldots], p)} $\to$ peel the entire leading
    \texttt{forall} chain in one step, introducing a fresh variable for
    each bound variable.  If the substituted body is
    \texttt{$y$ in $A$ implies \ldots}, absorb \texttt{$y$ in $A$} into
    the assumptions and continue peeling.  Stop at the first
    non-\texttt{forall} head.
    Example: goal \texttt{forall([x in A, y], body)} yields one subgoal
    $x_0 \in A \;\Longrightarrow\; \texttt{body}[x/x_0,\,y/y_0]$.
  \end{itemize}

\item \texttt{(cmd-antecedent-inference ps $f$)}
  Decompose assumption $f$ by its top-level connective:
  \begin{itemize}
  \item \texttt{p and q} $\to$ replace with $p$ and $q$ separately.
  \item \texttt{p or q} $\to$ case split: branch on $p$, branch on $q$.
  \item \texttt{not(p)} $\to$ derive \texttt{falsity} if $p$ also in context.
  \item \texttt{forsome([\ldots], p)} $\to$ introduce a fresh witness for $x$.
  \end{itemize}

\item \texttt{(cmd-assumption ps)}
  Close goal immediately if it matches an assumption or is \texttt{truth}.

\item \texttt{(cmd-arith ps)}
  Close the current goal by ground arithmetic evaluation.  The goal must
  be a closed formula (no free variables) built from numeric literals,
  $+$, $*$, $-$, \texttt{recip}, \texttt{abs}, \texttt{succ},
  \texttt{conjugate}, $\leq$, $=$,
  $\in \mathit{NN}/\mathit{ZZ}/\mathit{QQ}/\mathit{RR}/\mathit{CC}$,
  and the propositional connectives.  Exact (rational) arithmetic only;
  no axiom chain is constructed.
  Related: \texttt{(vnb-do-arith \textit{f})} returns \texttt{\#t}/\texttt{\#f}
  for use outside a proof.

\item \texttt{(cmd-ring-simplify ps)}
  Close a goal by reducing both sides to their canonical form in the free
  associative $\mathbb{Z}$-algebra.  Addition is commutative (both sides
  are brought to the same sorted canonical order); multiplication is
  non-commutative.  Generators are detected automatically and no generator
  list need be supplied: \texttt{(rs)} treats any \emph{term} that
  \texttt{arith-eval} cannot reduce to a number as a generator of the free
  ring --- a bare symbol such as \texttt{x}, or a compound term such as
  \texttt{f(x)} whose head is not one of $+\,-\,*$.  A compound generator is
  a single opaque indeterminate: \texttt{f(x)} is atomic and its argument is
  not examined.

  The goal may be in either of two forms:
  \begin{itemize}
  \item \textbf{Bare equality} \texttt{$e_1$ = $e_2$}: generators are
    certified by \texttt{(IN $v$ $D$)} assumptions already in the sequent
    context.  Typically reached after \texttt{(di)} peels the leading
    \texttt{forall} chain.
  \item \textbf{Universally quantified identity}
    \texttt{forall([$x_1$ in $D_1$, \ldots], $e_1$ = $e_2$)}: the
    quantifier typings themselves certify the generators.  \texttt{(di)}
    is not needed; \texttt{(rs)} closes the goal in one step.
  \end{itemize}

  In both cases every generator must be certified as an element of one of
  the five ring domains $\mathit{NN}$, $\mathit{ZZ}$, $\mathit{QQ}$,
  $\mathit{RR}$, $\mathit{CC}$ --- a symbol \texttt{x} by an
  \texttt{(IN x $D$)} typing, a compound generator \texttt{f(x)} by an
  \texttt{(IN f(x) $D$)} assumption.  This certification is also what keeps
  the rule sound on partial terms.  Since \texttt{=} is strict (it requires
  both sides defined), \texttt{3 + f(x) = 1 + f(x) + 2}
  holds only when \texttt{f(x)} is defined; membership entails definedness,
  so demanding \texttt{(IN f(x) $D$)} supplies exactly that guarantee.  A
  generator that is not certified leaves the goal untouched --- the command
  never makes an unsound step.

  \begin{VNBcode}
;;; Directly on a universally quantified goal -- no (di) needed:
(sp (make-wff-from-string
  "forall([a in ZZ, b in ZZ, c in ZZ], a + b + c = c + b + a)"))
(rs)    ; closes in one step

;;; After (di), a compound term certified in context becomes a generator:
(sp (make-wff-from-string
  "forall([x in ZZ], (f(x) in ZZ) implies 3 + f(x) = 1 + f(x) + 2)"))
(di) (di) (rs)    ; di peels the forall then the hypothesis; rs then closes

;;; Coefficient arithmetic:
(sp (make-wff-from-string "forall([x in ZZ], 2*x + 3*x = 5*x)"))
(rs)
  \end{VNBcode}

  The primitive inference is \texttt{pi-ring-simplify!}, in
  \texttt{structure-library/ring-simplify.scm}; it normalizes each side to a
  canonical sum-of-monomials polynomial (\texttt{vnb->poly}) and compares.  The interactive short form is \texttt{(rs)}.

\item \texttt{(cmd-comm-ring-simplify ps)}
  The commutative counterpart of \texttt{(cmd-ring-simplify ps)}.  Where
  \texttt{(rs)} fixes the order of factors within a monomial,
  \texttt{(crs)} treats multiplication as commutative: a monomial is a
  \emph{multiset} of generators, canonicalised as a sorted generator
  list, so \texttt{x*y} and \texttt{y*x} share one normal form and
  commutativity lives in the representation rather than in any rewrite
  step.  Each side reduces to a sum of such monomials over \texttt{ZZ},
  and the goal closes iff the two normal forms coincide.  Use it for any
  identity whose proof needs reordering of factors, which \texttt{(rs)}
  cannot perform.

  The goal may take the same two shapes as for \texttt{(rs)} (a bare
  equality with generators certified by context typings, or a typed
  universally quantified identity), and additionally a \textbf{generic}
  form over an arbitrary ring: \texttt{forall($R$, IS-COMMUTATIVE-RING($R$)
  implies \ldots)} built from \texttt{(ADD $R$)}, \texttt{(MUL $R$)},
  \texttt{(NEG $R$)}, \texttt{(ZERO $R$)}, \texttt{(ONE $R$)}, with the
  generators ranging over the carrier \texttt{(CARR $R$)}.

  Generators are detected and certified exactly as for \texttt{(rs)}: a
  symbol or an opaque compound term, certified by an \texttt{(IN $v$ $D$)}
  typing in the concrete form or an \texttt{(IN $v$ (CARR $R$))} assumption in
  the generic form.  The same partiality caveat applies --- a compound
  generator fires only once certified, which in VNB guarantees it is
  defined.

  \begin{VNBcode}
;;; Reordering of factors -- (rs) cannot close this, (crs) does:
(sp (make-wff-from-string
  "forall([x in ZZ, y in ZZ], x*y*x = x*x*y)"))
(crs)   ; closes in one step

;;; Commutative expansion:
(sp (make-wff-from-string
  "forall([x in ZZ, y in ZZ], (x+y)*(x+y) = x*x + 2*x*y + y*y)"))
(crs)
  \end{VNBcode}

  The procedure is a trusted oracle (it compares canonical normal forms
  rather than mechanising the ring properties), carrying the warrant
  \texttt{comm-ring-simplify} of kind \texttt{well-known}.  The primitive
  inference is \texttt{pi-comm-ring-simplify!} in
  \texttt{comm-ring-simplify.scm}.  The interactive short form is
  \texttt{(crs)}.

\item \texttt{(cmd-reflexivity ps)}
  Close goal \texttt{a = a} when \texttt{a} is demonstrably defined.
  Fails if \texttt{a} may be undefined.

\item \texttt{(cmd-quasi-reflexivity ps)}
  Close goal \texttt{a == a} unconditionally, for any term \texttt{a}.

\item \texttt{(cmd-cut ps $L$)}
  Introduce lemma $L$: subgoal~1 proves $L$; subgoal~2 has $L$ as assumption.

\item \texttt{(cmd-or-intro-left ps)}
  To prove \texttt{p or q}, prove $p$ instead.

\item \texttt{(cmd-or-intro-right ps)}
  To prove \texttt{p or q}, prove $q$ instead.

\item \texttt{(cmd-proof-by-contradiction ps)}
  To prove $p$, assume \texttt{not(p)} and derive \texttt{falsity}.

\item \texttt{(cmd-instantiate ps $F$\ $t$)}
  From \texttt{forall([$x$], $p$)} $= F$ in context, add $p[x/t]$.

\item \texttt{(cmd-exists-witness ps $t$)}
  To prove \texttt{forsome([$x$], $p$)}, prove $p[x/t]$.

\item \texttt{(cmd-weaken ps $f$)}
  Remove formula $f$ from the assumptions.

\item \texttt{(cmd-backchain ps $F$)}
  From \texttt{p implies goal} $= F$ in context, generate subgoal~$p$.

\item \texttt{(cmd-theorem-assumption ps $\mathit{name}$)}
  Add the named theorem (or axiom) to the assumptions.

\item \texttt{(cmd-apply-macete ps $\mathit{name}$)}
  Apply named macete: rewrite the goal using the associated theorem.

\item \begin{sloppypar}\texttt{(cmd-eq-subst ps $F$)}
  Equality substitution (Leibniz).  From an equality \texttt{$s$ = $t$} $= F$
  in the assumptions (either orientation), rewrite every occurrence of $s$
  in the goal to $t$.  Short form \texttt{(subst $F$)}.  This realises the
  schema \texttt{$s$ = $t$} $\wedge$ \texttt{$P[t]$} $\Rightarrow$
  \texttt{$P[s]$} for an arbitrary context $P$; sound for partial equality
  because \texttt{$s$ = $t$} already certifies both sides
  defined.\end{sloppypar}

\item \texttt{(cmd-qed ps $\mathit{name}$)}
  After a proof completes, install its root assertion as a named theorem.
  Errors if the proof is not complete.  This is the only command that
  records a result as \emph{proven}.

\end{itemize}

\subsection{N-ary constructors}
\label{sec:nary-proof-commands}

\begin{itemize}

\item \begin{sloppypar}\texttt{(cmd-cartesian-intro ps)}
  To prove \texttt{[$a_1$,\ldots,$a_n$] in cartesian($A_1$,\ldots,$A_n$)},
  generate $n$ subgoals \texttt{$a_i$ in $A_i$}.\end{sloppypar}

\item \begin{sloppypar}\texttt{(cmd-cartesian-elim ps $F$\ $k$)}
  From \texttt{[$a_1$,\ldots,$a_n$] in cartesian($A_1$,\ldots,$A_n$)} $= F$
  in context, add \texttt{$a_k$ in $A_k$}.\end{sloppypar}

\item \texttt{(cmd-tuples-intro ps)}
  To prove \texttt{[$a_1$,\ldots,$a_n$] in tuples($A$)},
  generate $n$ subgoals \texttt{$a_i$ in $A$}.
  For the empty sequence \texttt{[]} the goal closes immediately (0 subgoals).

\item \begin{sloppypar}\texttt{(cmd-tuples-elim ps $F$\ $k$)}
  From \texttt{[$a_1$,\ldots,$a_n$] in tuples($A$)} $= F$
  in context, add \texttt{$a_k$ in $A$}.\end{sloppypar}

\item \texttt{(cmd-nth-reduce ps)}
  Rewrite every occurrence of \texttt{nth($k$, [$\ldots$])} in the goal to
  the $k$-th element (1-indexed).

\item \begin{sloppypar}\texttt{(cmd-functoid-beta ps)}
  Rewrite every occurrence of
  \texttt{lambda([$x_1$ in $A_1$,\ldots], $\mathit{body}$)($v_1$,\ldots)}
  (or \texttt{lambdoid}) in the goal to
  $\mathit{body}[x_1 := v_1]\cdots$ (capture-avoiding).\end{sloppypar}

\item \texttt{(cmd-union-intro ps $k$)}
  To prove \texttt{$x$ in union($A_1$,\ldots,$A_n$)},
  prove \texttt{$x$ in $A_k$} instead.

\item \texttt{(cmd-union-elim ps $F$)}
  From \texttt{$x$ in union($A_1$,\ldots,$A_n$)} $= F$ in context,
  case-split: one branch per $A_i$ with \texttt{$x$ in $A_i$} assumed.

\item \begin{sloppypar}\texttt{(cmd-intersection-intro ps)}
  To prove \texttt{$x$ in intersection($A_1$,\ldots,$A_n$)},
  generate $n$ subgoals \texttt{$x$ in $A_i$}.\end{sloppypar}

\item \begin{sloppypar}\texttt{(cmd-intersection-elim ps $F$\ $k$)}
  From \texttt{$x$ in intersection($A_1$,\ldots,$A_n$)} $= F$ in context,
  add \texttt{$x$ in $A_k$}.\end{sloppypar}

\end{itemize}

%% =========================================================
\section{Two surfaces, one script}
\label{sec:two-surfaces}

Every proof step exists in two forms.

The \emph{builder} \texttt{(cmd-\textit{name}\ \textit{ps}\ \textit{arg}
\ldots)} (Section~\ref{sec:commands}) takes a proof state explicitly and
\emph{returns} a new one; it has no side effects and never prints.  It is
the form to call from Scheme code that manipulates proof states as values.

The \emph{interactive short form} \texttt{(\textit{name}\ \textit{arg}
\ldots)} (Section~\ref{sec:interactive}) is a thin wrapper around the
builder: it reads the global proof state \texttt{*ps*}, applies the
builder, stores the result back in \texttt{*ps*}, and calls \texttt{(show)}
so the new state is displayed.  This is the form typed in a live session.

\begin{center}
\begin{tabular}{@{}lll@{}}
\toprule
Builder & Short form & \\
\midrule
\texttt{(cmd-direct-inference ps)}        & \texttt{(di)}    & no argument \\
\texttt{(cmd-antecedent-inference ps $f$)}& \texttt{(ai $f$)}& formula argument \\
\texttt{(cmd-backchain ps $F$)}           & \texttt{(bc $F$)}& formula argument \\
\texttt{(cmd-instantiate ps $F$ $t$)}     & \texttt{(inst $F$ $t$)} & formula + term \\
\texttt{(cmd-apply-macete ps $n$)}        & \texttt{(mac $n$)} & theorem name \\
\bottomrule
\end{tabular}
\end{center}

\noindent
The wrapper does one more thing: it \emph{records itself}.  Each short form
appends an entry \texttt{(\textit{name}\ .\ \textit{args})} to the global
accumulator \texttt{*proof-script*}, which \texttt{(sp \ldots)} clears at
the start of a proof.  So there is no separate scripting language to learn:
\textbf{the script is the transcript} --- the very sequence of short forms
you typed.  \texttt{(qed \textit{name})} and \texttt{(save-proof
\textit{name})} store that sequence; \texttt{(replay-proof \textit{name})}
re-runs it (Section~\ref{sec:replay}).

\subsection{Supplying arguments}
\label{sec:arg-collection}

A short form takes one of a few kinds of argument, and each kind is
collected the same way wherever it appears.

\begin{itemize}
\item \textbf{Nothing.}  The structural steps --- \texttt{di}, \texttt{ass},
  \texttt{pbc}, \texttt{rs}, \texttt{rfl}, \texttt{ci}, \texttt{ni},
  \ldots\ --- read the focus goal and need no argument.

\item \textbf{A formula} (\texttt{ai}, \texttt{bc}, \texttt{wk},
  \texttt{cut}, \texttt{ue}, the assumption cited by \texttt{inst} and
  \texttt{mac-h}).  It may be supplied three interchangeable ways:
  \begin{itemize}
  \item a \textbf{surface string}, parsed for you ---
    \texttt{"n in NN and m in NN"};
  \item a \textbf{raw S-expression} ---
    \texttt{'(and (in n nn) (in m nn))};
  \item for a command that cites an existing \emph{assumption}, a
    \textbf{1-based assumption index}: a positive integer $k$ selects the
    $k$-th assumption of the focus sequent, exactly as numbered in the
    display.  \texttt{(ai 1)} stands in for retyping the first assumption.
  \end{itemize}

\item \textbf{A term} (\texttt{ew}, the witness slot of \texttt{inst},
  \texttt{if-true}/\texttt{if-false}).  A string or an S-expression; there
  is no index form, since a term is not an assumption to point at.

\item \textbf{A name} (\texttt{ta}, \texttt{mac}, \texttt{bc*},
  \texttt{qed}).  A symbol naming an installed theorem or macete.

\item \textbf{A small integer} (\texttt{ui $k$}, \texttt{ce $f$ $k$},
  \texttt{focus $n$}).  Selects a branch, a tuple component, or an open
  goal --- not an assumption index.
\end{itemize}

\begin{VNBcode}
(sp (make-wff-from-string "n in NN and m in NN implies n in NN"))
(di)        ; assume the conjunction (now assumption #1)
(ai 1)      ; crack assumption #1 -- the index, not the retyped formula
(ass)
(qed 'and-elim-demo)

(lookup-proof 'and-elim-demo)
; => ((di) (ai 1) (ass))     -- the index is recorded verbatim
\end{VNBcode}

\begin{remark}
The assumption-index form is accepted by the commands that \emph{cite} an
assumption --- \texttt{ai}, \texttt{bc}, \texttt{wk}, \texttt{inst},
\texttt{ce}, \texttt{te}, \texttt{ie}, \texttt{mac-h}.  It is recorded
verbatim in the script, so a replay re-resolves ``assumption~1'' against
the focus sequent it actually reaches --- robust to the formula's exact
text, but dependent on the assumption \emph{order} being the same on
replay.  When that is in doubt, give the formula itself.
\end{remark}

\subsection{Argument collection in a front end: \texttt{bc*}, \texttt{B}, \texttt{B+}}
\label{sec:bc-star-collection}

\texttt{bc*} (Section~\ref{sec:interactive}) matches a theorem's conclusion
against the goal and instantiates whatever that match determines.  A schema
variable the match leaves \emph{undetermined} must be supplied --- as a
binding list \texttt{((\textit{v}\ \textit{val})\ \ldots)} of Scheme
expressions.  The pure query \texttt{(bc*-undetermined '\textit{name})}
reports exactly which variables those are, without touching the proof.

A front end uses that query to collect the values \emph{by name} instead of
making the user write the quoted binding list: it asks for each undetermined
variable's value as an ordinary \emph{term string} (\texttt{"RAN(f)"},
\texttt{"nn"}) and parses it.  In the Emacs Focus workspace this is the
\texttt{B} key --- the argument \emph{builder}.  Its companion key
\texttt{B+} is a saturating closer: it only commits a citation it
can first \emph{prove} will close the focus goal
(\texttt{bc*-can-close?} matches the conclusion \emph{and} checks every
hypothesis is already an assumption), so it never leaves the proof worse
than it found it.

%% =========================================================
\section{Interactive short-form commands}\label{sec:tactics-ref}
\label{sec:interactive}

\texttt{interactive.scm} provides the short-form commands --- one- or two-letter
wrappers around the \texttt{cmd-*} proof procedures --- that you type in a live
proof session, at the bare REPL or through the Emacs workspace.  Each mutates the
global proof state \texttt{*ps*} and calls \texttt{(show)} automatically; the
exceptions, which return a value rather than mutating state, are \texttt{sp},
\texttt{show}, \texttt{pp}, \texttt{fa}, \texttt{fs}, \texttt{vnb-unwind},
\texttt{vnb-wind}, \texttt{parse-string}, and \texttt{make-wff-from-string}.

The reference below is \emph{generated} from the live \texttt{*tactic-help*}
registry (through \texttt{reference/TACTICS.md}), so it lists the current commands
with their signatures and trust \emph{kind} and cannot drift from the prover.  At
the REPL, \texttt{(tactics)} prints the same menu and \texttt{(tactics 'name)} the
full per-command help.

%% AUTO-GENERATED by docs/gen-tactics-ref.py from reference/TACTICS.md -- do not edit.
%% Regenerate:  python3 docs/gen-tactics-ref.py   (after a prover load refreshes TACTICS.md)

\DeclareRobustCommand{\tactkind}[1]{{\footnotesize\textsc{#1}}}

Every tactic is tagged with a \emph{kind} grounded in the \texttt{dg-apply-rule!} tag it emits, so the trusted base is legible: \texttt{rule} is a single primitive kernel inference (the fixed trusted base); \texttt{oracle} is a trusted decision procedure run as a black box; \texttt{composite} only chains kernel rules, adding no new inference; \texttt{meta} performs no deduction (session, search, navigation).  A proof's trust surface is exactly its \texttt{rule} steps plus whichever \texttt{oracle}s and asserted premises it cites.

The same reference is available at the REPL with \texttt{(tactics 'name)} and in the browser file \texttt{reference/TACTICS.md}; all three are generated from one registry, so they never disagree.

An entry's \emph{Uses} line names every kernel operation the command's own code can cause to be recorded, those recorded by the procedures it calls included. It is the bound on what the command can do, whether or not any proof in the library exercises it, and it is measured by reading the source rather than declared alongside the command. One further set is common to every command and is stated once instead of in every entry: the hook that runs after any command which changed the proof, to close the definedness sequents an instantiation left owed. Section~\ref{sec:command-uses} tabulates the same figures for every command at once.

\subsection*{Starting \& finishing a proof}

\par\smallskip{\raggedright\noindent\texttt{(sp goal)}\quad\tactkind{meta}\index{sp@\texttt{sp} (tactic)}\par}\nobreak
\noindent Start a proof of `goal' (a ``string'', a raw S-expr, or a wff); clears the script.\par
\noindent\textit{Uses:} no kernel operation: it records no inference.\par
\noindent\textit{When useful:} starting a new proof\par

Begin a new proof: state what you want to prove.  The goal may be a ``string'' in surface syntax, a raw S-expr, or an already-built wff (e.g. (wff ``...'') or (make-wff '(...)));  sp coerces a string/S-expr for you.  This becomes your one open goal and clears any previous script.  (Technically: set the proof goal.)

\par\smallskip{\raggedright\noindent\texttt{(qed 'name)}\quad\tactkind{meta}\index{qed@\texttt{qed} (tactic)}\par}\nobreak
\noindent Install the finished proof as theorem `name' and save its replayable script.\par
\noindent\textit{Uses:} no kernel operation: it records no inference.\par
\noindent\textit{When useful:} no open goals remain -- record the result\par

Finish: once there are no open goals, this records the result as a named theorem you can cite later, and saves the replayable script.  (Technically: install-theorem! plus script capture; reports the asserted facts the proof still rests on, `proven modulo \{...\}'.)

\par\smallskip{\raggedright\noindent\texttt{(save-proof 'name)}\quad\tactkind{meta}\index{save-proof@\texttt{save-proof} (tactic)}\par}\nobreak
\noindent Save the current script under a name without finishing.\par
\noindent\textit{Uses:} no kernel operation: it records no inference.\par

Save the current (possibly unfinished) proof script under a name, to resume or replay later.  (Technically: snapshot the script without installing a theorem.)

\par\smallskip{\raggedright\noindent\texttt{(replay-proof 'name [subst])}\quad\tactkind{meta}\index{replay-proof@\texttt{replay-proof} (tactic)}\par}\nobreak
\noindent Re-run a saved script on the current goal, optionally renaming free vars.\par
\noindent\textit{Uses:} every one of the 64 kernel operations.\par

The optional subst is an alist ((old . new) ...) applied to every command argument before replay -- this is how one proof is reused at fresh eigenvariables.


\subsection*{Decomposing the goal}

\par\smallskip{\raggedright\noindent\texttt{(di)}\quad\tactkind{rule}\index{di@\texttt{di} (tactic)}\par}\nobreak
\noindent Direct inference: split an AND goal, move an IMPLIES antecedent into the assumptions, or introduce a fresh eigenvariable for a leading FORALL.\par
\noindent\textit{Uses:} \texttt{and-intro} \texttt{forall-intro} \texttt{iff-intro} \texttt{implies-intro} \texttt{not-intro}.\par
\noindent\textit{When useful:} the goal head is AND / IMPLIES / FORALL -- decompose before deciding\par

Break the goal into its pieces.  A conjunction `A and B' splits into two goals, A and B.  An implication `P implies Q' assumes P (P becomes a hypothesis you may use) and leaves you to prove Q.  A universal `for all x, ...' fixes an arbitrary x (a new constant standing for `any x') and asks for the statement at that x.  Apply it repeatedly until the goal is a single atomic statement.  (Technically: the goal-side introduction rules for AND / IMPLIES / FORALL; the constant introduced for a FORALL is an eigenvariable.  Pairs with mac, which unfolds a definition in the goal.)

\par\smallskip{\raggedright\noindent\texttt{(pbc)}\quad\tactkind{rule}\index{pbc@\texttt{pbc} (tactic)}\par}\nobreak
\noindent Proof by contradiction: assume the goal's negation, prove FALSITY.\par
\noindent\textit{Uses:} \texttt{proof-by-contradiction}.\par
\noindent\textit{When useful:} a direct argument stalls and assuming the negation gives something concrete to work with\par

Argue by contradiction: assume the goal is false and derive an absurdity.  (Technically: classical reductio -- adds the negation of the goal as a hypothesis and changes the goal to FALSITY.)

\par\smallskip{\raggedright\noindent\texttt{(oi-l)}\quad\tactkind{rule}\index{oi-l@\texttt{oi-l} (tactic)}\par}\nobreak
\noindent OR-intro left: reduce an (OR a b) goal to a.\par
\noindent\textit{Uses:} \texttt{or-intro-left}.\par
\noindent\textit{When useful:} the goal is (OR a b) and the LEFT alternative is provable\par

Prove a disjunction `A or B' by proving the LEFT alternative, A.  (Technically: or-introduction, left.)

\par\smallskip{\raggedright\noindent\texttt{(oi-r)}\quad\tactkind{rule}\index{oi-r@\texttt{oi-r} (tactic)}\par}\nobreak
\noindent OR-intro right: reduce an (OR a b) goal to b.\par
\noindent\textit{Uses:} \texttt{or-intro-right}.\par
\noindent\textit{When useful:} the goal is (OR a b) and the RIGHT alternative is provable\par

Prove a disjunction `A or B' by proving the RIGHT alternative, B.  (Technically: or-introduction, right.)

\par\smallskip{\raggedright\noindent\texttt{(ew term)}\quad\tactkind{rule}\index{ew@\texttt{ew} (tactic)}\par}\nobreak
\noindent Existential witness: discharge a FORSOME goal by supplying the witness term.\par
\noindent\textit{Uses:} \texttt{forsome-intro}.\par
\noindent\textit{When useful:} the goal is FORSOME and you have a witness term in mind\par

Prove `there exists an x with property P' by exhibiting a specific witness: you supply the term, and the goal becomes `P holds of that term'.  The everyday `take x = ...' step.  (Technically: existential introduction for a FORSOME goal.)

\par\smallskip{\raggedright\noindent\texttt{(ci)}\quad\tactkind{rule}\index{ci@\texttt{ci} (tactic)}\par}\nobreak
\noindent Cartesian intro: prove a CARTESIAN-product membership component-wise.\par
\noindent\textit{Uses:} \texttt{cartesian-intro}.\par
\noindent\textit{When useful:} the goal is membership in a CARTESIAN product\par

Prove that a pair (or tuple) belongs to a Cartesian product A x B by proving each coordinate lies in its factor.  (Technically: Cartesian-membership introduction, component-wise.)

\par\smallskip{\raggedright\noindent\texttt{(ti)}\quad\tactkind{rule}\index{ti@\texttt{ti} (tactic)}\par}\nobreak
\noindent Tuple intro: prove a tuple/LIST membership component-wise.\par
\noindent\textit{Uses:} \texttt{tuples-intro}.\par
\noindent\textit{When useful:} the goal is membership in a set of tuples / LIST\par

Prove that a list belongs to the set of tuples over A by proving each entry lies in A.  (Technically: tuple-membership introduction, component-wise.)

\par\smallskip{\raggedright\noindent\texttt{(ii)}\quad\tactkind{rule}\index{ii@\texttt{ii} (tactic)}\par}\nobreak
\noindent Intersection intro: prove (IN x (INTERSECTION ...)) for each branch.\par
\noindent\textit{Uses:} \texttt{intersection-intro}.\par
\noindent\textit{When useful:} the goal is membership in an INTERSECTION\par

Prove that x lies in an intersection of sets by proving it lies in each of them.  (Technically: intersection-membership introduction.)

\par\smallskip{\raggedright\noindent\texttt{(ui k)}\quad\tactkind{rule}\index{ui@\texttt{ui} (tactic)}\par}\nobreak
\noindent Union intro: reduce a goal (IN x (UNION ...)) to membership of x in the k-th set (1-based).\par
\noindent\textit{Uses:} \texttt{union-intro}.\par
\noindent\textit{When useful:} the goal is membership in a UNION and you know which branch\par

Prove that x lies in a union of sets by proving it lies in the k-th one (you choose which).  (Technically: union-membership introduction at branch k.)

\par\smallskip{\raggedright\noindent\texttt{(ni)}\quad\tactkind{rule}\index{ni@\texttt{ni} (tactic)}\par}\nobreak
\noindent Natural-number induction on the goal's leading FORALL over NN.\par
\noindent\textit{Uses:} \texttt{nn-induction}.\par
\noindent\textit{When useful:} the goal is `forall n in NN, P(n)'\par

Prove `for all natural numbers n, P(n)' by induction: show P(0), and that P(n) implies P(n+1).  (Technically: natural-number induction on the goal's leading FORALL over NN.)


\subsection*{Using a hypothesis}

\par\smallskip{\raggedright\noindent\texttt{(ai hyp)}\quad\tactkind{rule}\index{ai@\texttt{ai} (tactic)}\par}\nobreak
\noindent Antecedent inference: decompose a cited assumption -- AND-split, OR-into-cases, or FORSOME-elimination to a fresh eigenvariable.\par
\noindent\textit{Uses:} \texttt{and-elim} \texttt{forsome-elim} \texttt{iff-elim} \texttt{not-elim} \texttt{or-elim}.\par
\noindent\textit{When useful:} a hypothesis is an AND / OR / FORSOME to break apart\par

Break a hypothesis into its pieces -- the mirror image of di, but on the assumptions instead of the goal.  From a hypothesis `A and B' you get both A and B; from `A or B' you split into two cases (prove the goal in each); from `there exists x, P(x)' you get a fresh name for such an x together with P(x).  Cite the hypothesis by its number in the display, its formula, or a ``string''.  (Technically: the hypothesis-side elimination rules for AND / OR / FORSOME; the fresh name is an eigenvariable.)

\par\smallskip{\raggedright\noindent\texttt{(ass)}\quad\tactkind{rule}\index{ass@\texttt{ass} (tactic)}\par}\nobreak
\noindent Close the goal by an assumption alpha-equivalent to it.\par
\noindent\textit{Uses:} \texttt{assumption}.\par
\noindent\textit{When useful:} the goal already appears (up to bound-var renaming) among the hypotheses\par

Finish the goal because it is already one of your hypotheses -- the `that is exactly what we assumed' step.  Renaming of dummy/bound variables doesn't matter (`there exists x, P(x)' closes `there exists y, P(y)').  (Technically: the assumption rule, matching up to alpha-equivalence.)

\par\smallskip{\raggedright\noindent\texttt{(inst forall-hyp term)}\quad\tactkind{rule}\index{inst@\texttt{inst} (tactic)}\par}\nobreak
\noindent Instantiate a universally-quantified assumption at `term', adding the instance to context.\par
\noindent\textit{Uses:} \texttt{forall-elim}.\par
\noindent\textit{When useful:} you have a forall-hypothesis and a specific value to use it at\par

Use a `for all x, ...' hypothesis at a particular value: you supply the term, and the statement with x replaced by that term is added to your hypotheses.  (Technically: universal instantiation of an assumption.)

\par\smallskip{\raggedright\noindent\texttt{(inst+ forall-hyp term)}\quad\tactkind{composite}\index{inst+@\texttt{inst+} (tactic)}\par}\nobreak
\noindent Instantiate an in-context universal at `term', then forward-detach any guards whose antecedents are in context.\par
\noindent\textit{Uses:} \texttt{arith-forsome} \texttt{arith-ground} \texttt{arith-simplify} \texttt{cut} \texttt{detach} \texttt{forall-elim}.\par
\noindent\textit{When useful:} a GUARDED forall-hyp whose guards are dischargeable from context (e.g. the metric laws post-grind)\par

inst followed by detach: from `forall x. (x in S) => P(x)' and `t in S' already known, this lands `P(t)' directly (peeling nested guards level by level), instead of leaving the guarded implication for you to detach by hand.  The hypothesis-side analogue of `fact' (which assembles a THEOREM); this assembles an in-context UNIVERSAL.  scout's inst lane emits these to close witness-needing goals like the metric laws.  (Technically: pi-instantiate! then pi-detach! while the consequent stays a guard with an in-context antecedent.)

\par\smallskip{\raggedright\noindent\texttt{(mp)}\quad\tactkind{composite}\index{mp@\texttt{mp} (tactic)}\par}\nobreak
\noindent Close a goal that is an INSTANCE of a universal already in your context.  No arguments: the term is DERIVED by matching, not supplied.\par
\noindent\textit{Uses:} \texttt{arith-forsome} \texttt{arith-ground} \texttt{arith-simplify} \texttt{assumption} \texttt{cut} \texttt{detach} \texttt{forall-elim}.\par
\noindent\textit{When useful:} the goal IS an instance of a universal you already have -- no term to supply, it is derived by matching\par

The syllogism, with nothing to fill in.  From `forall([thing in human], thing in mortal)' together with `socrates in human', it closes `socrates in mortal'.  Where inst+ makes you name the term, mp works it out: it matches the universal's conclusion against the goal, which either determines the bound variable or fails outright, and then checks that the universal's guard at that term is something you already have.  Nothing is searched for and nothing is ranked -- if the goal is an instance there is exactly one term it can be.  It DECLINES, naming what it found, when no universal in the context has the goal as an instance, and when more than one does: a closer that silently picks between two universals is one you cannot predict.  One bound variable in this version; for `forall([a in A, b in B], ...)' use inst+.  (Technically: composite -- inst+ at the derived term, then ass.  It adds no kernel rule and no debt; the bill is what typing the two steps by hand would bill, which is nothing.)

\par\smallskip{\raggedright\noindent\texttt{(grind-and-mp)}\quad\tactkind{composite}\index{grind-and-mp@\texttt{grind-and-mp} (tactic)}\par}\nobreak
\noindent Tidy the goal with grind, then close it by modus ponens -- the one-button form of the syllogism, for a sentence typed exactly as it reads.\par
\noindent\textit{Uses:} \texttt{and-elim} \texttt{and-intro} \texttt{arith-forsome} \texttt{arith-ground} \texttt{arith-simplify} \texttt{assumption} \texttt{cut} \texttt{detach} \texttt{forall-elim} \texttt{forall-intro} \texttt{forsome-elim} \texttt{iff-elim} \texttt{iff-intro} \texttt{implies-intro} \texttt{macete-hyp} \texttt{not-elim} \texttt{not-intro} \texttt{or-elim}.\par
\noindent\textit{When useful:} a sentence typed as it reads -- `all men are mortal and socrates is a man implies socrates is mortal' -- closing in one move\par

Runs the two steps a syllogism actually takes on a sentence you have just typed.  `forall([thing in human], thing in mortal) and socrates in human implies socrates in mortal' is not yet a syllogism: it is an implication whose antecedent is a conjunction.  grind does the BOOKKEEPING -- split the AND, move the antecedents into the assumptions -- which leaves the goal `socrates in mortal' with the universal and `socrates in human' as hypotheses; mp then does the LOGIC (universal instantiation, then modus ponens).  They stay separate tactics because they are different kinds of step, and mp on its own is what to reach for once the goal is already tidy.  The recorded script keeps both steps, so the page reads as the two moves rather than as this wrapper.  (Technically: composite -- grind, quietly since on a tidy goal it legitimately declines, then mp.  Adds no kernel rule and no debt.)

\par\smallskip{\raggedright\noindent\texttt{(detach! impl)}\quad\tactkind{rule}\index{detach"!@\texttt{detach"!} (tactic)}\par}\nobreak
\noindent Forward modus ponens: from an in-context (IMPLIES A B) whose A is also in context, leave B in context.\par
\noindent\textit{Uses:} \texttt{detach}.\par
\noindent\textit{When useful:} you have both P and (P => Q) in context and want Q\par

If you have both `P' and `P implies Q' among your hypotheses, this adds `Q' to them.  It grows what you KNOW (the hypotheses) rather than changing the goal -- the forward counterpart of bc.  (Technically: forward modus ponens, the kernel rule pi-detach!; sound since P and P=>Q give Q.)

\par\smallskip{\raggedright\noindent\texttt{(fact 'thm term ...)}\quad\tactkind{composite}\index{fact@\texttt{fact} (tactic)}\par}\nobreak
\noindent Forward APPLICATION of a theorem: bring it in, instantiate its leading universals with the terms, and auto-detach every antecedent already in context, landing the consequent as a hypothesis.\par
\noindent\textit{Uses:} \texttt{arith-forsome} \texttt{arith-ground} \texttt{arith-simplify} \texttt{cut} \texttt{detach} \texttt{forall-elim} \texttt{theorem-assumption}.\par
\noindent\textit{When useful:} a library law `forall x. H(x) => P(x)' whose P you want landed as a hypothesis\par

Handles interleaved forall/implies (e.g. forall s. IS-X(s) => forall a. a in CARR(s) => P): consumes one term per FORALL, detaches each IMPLIES whose antecedent is in context.  The forward-assembly workhorse -- a law `forall x. H(x) => P(x)' becomes the usable fact P in one call, instead of ta + inst* + cut/backchain.  See theorem-library/module-zero-act.scm.

\par\smallskip{\raggedright\noindent\texttt{(ce hyp k)}\quad\tactkind{rule}\index{ce@\texttt{ce} (tactic)}\par}\nobreak
\noindent Cartesian elim: project the k-th component out of a CARTESIAN-membership assumption.\par
\noindent\textit{Uses:} \texttt{cartesian-elim}.\par
\noindent\textit{When useful:} a hypothesis is CARTESIAN-product membership\par

From a hypothesis that a tuple lies in a Cartesian product, extract that its k-th coordinate lies in the k-th factor.  (Technically: Cartesian-membership elimination, k-th projection.)

\par\smallskip{\raggedright\noindent\texttt{(te hyp k)}\quad\tactkind{rule}\index{te@\texttt{te} (tactic)}\par}\nobreak
\noindent Tuple elim: project the k-th component out of a tuple-membership assumption.\par
\noindent\textit{Uses:} \texttt{tuples-elim}.\par
\noindent\textit{When useful:} a hypothesis is tuple membership\par

From a hypothesis that something is a tuple over A, extract that its k-th entry lies in A.  (Technically: tuple-membership elimination, k-th projection.)

\par\smallskip{\raggedright\noindent\texttt{(ie hyp k)}\quad\tactkind{rule}\index{ie@\texttt{ie} (tactic)}\par}\nobreak
\noindent Intersection elim: extract the k-th branch of an INTERSECTION-membership assumption.\par
\noindent\textit{Uses:} \texttt{intersection-elim}.\par
\noindent\textit{When useful:} a hypothesis is INTERSECTION membership\par

From a hypothesis that x lies in an intersection, extract that x lies in the k-th set.  (Technically: intersection-membership elimination.)

\par\smallskip{\raggedright\noindent\texttt{(ue hyp)}\quad\tactkind{rule}\index{ue@\texttt{ue} (tactic)}\par}\nobreak
\noindent Union elim: split a cited (IN x (UNION ...)) membership assumption into one subgoal per set -- union-side case analysis, the dual of `ui'.\par
\noindent\textit{Uses:} \texttt{union-elim}.\par
\noindent\textit{When useful:} a hypothesis is UNION membership -- case-split on it\par

From a hypothesis that x lies in a union, split into cases -- one for each set x might belong to -- and prove the goal in each.  (Technically: union-membership elimination, case analysis; the dual of ui.)

\par\smallskip{\raggedright\noindent\texttt{(cut formula)}\quad\tactkind{rule}\index{cut@\texttt{cut} (tactic)}\par}\nobreak
\noindent Cut: prove `formula' as a side subgoal, then continue the main goal with `formula' added as an assumption (Gentzen cut).\par
\noindent\textit{Uses:} \texttt{cut}.\par
\noindent\textit{When useful:} you want to prove a lemma on the spot, then use it\par

Introduce a lemma you prove on the spot.  You state a formula; it becomes a side goal to prove, and on the main line you may then use it as a hypothesis.  The standard `we first show that ..., and now using it ...' move.  (Technically: Gentzen's cut -- nothing is left assumed-but-unproved, the side goal discharges it.)

\par\smallskip{\raggedright\noindent\texttt{(wk hyp)}\quad\tactkind{rule}\index{wk@\texttt{wk} (tactic)}\par}\nobreak
\noindent Weaken: drop a cited assumption from the context to tidy the hypothesis list.  `hyp' may be a formula, a ``string'', or a 1-based assumption index.\par
\noindent\textit{Uses:} \texttt{weakening}.\par
\noindent\textit{When useful:} the hypothesis list is cluttered with something no longer needed\par

Discard a hypothesis you no longer need, to keep the assumption list readable.  (Technically: weakening -- removing a hypothesis is always sound.)

\par\smallskip{\raggedright\noindent\texttt{(keep hyp ...)}\index{keep@\texttt{keep} (tactic)}\par}\nobreak
\noindent Keep only the cited assumptions and drop every other one, as ONE recorded step.  Each `hyp' may be a formula, a ``string'', or a 1-based assumption index; they are recorded as formulas.\par
\noindent\textit{When useful:} the context is deep and only a few hypotheses matter for the next step (prop, ineq)\par

Prune the context down to what the next step needs -- what a driver does before `prop' or `ineq' in a deep context.  (Technically: one weakening per dropped hypothesis, on the same leaf, recorded once; the driver kit's dk-only! is this command.)


\subsection*{Rewriting}

\par\smallskip{\raggedright\noindent\texttt{(mac 'name)}\quad\tactkind{rule}\index{mac@\texttt{mac} (tactic)}\par}\nobreak
\noindent Rewrite the GOAL with an equivalence macete (unfold a definition, apply an iff/=/== law).  Fires only where the macete's side-conditions already hold in context.\par
\noindent\textit{Uses:} \texttt{cartesian-decompose} \texttt{macete} \texttt{tuple-equality-decompose}.\par
\noindent\textit{When useful:} the GOAL has a defined predicate to unfold / an identity to apply, side-conditions already met\par

Rewrite the goal using a definition or a known equivalence/equality (a `macete').  For instance, replace a defined predicate by what it stands for, or apply an identity.  It fires only where the law's side-conditions already hold in your hypotheses; where they don't, it leaves that spot untouched and looks deeper inside.  (Technically: goal-side rewriting by an equivalence/equality macete; all-or-nothing -- it does NOT spawn unmet side-conditions as goals.  Its hypothesis-side cousin is mac-h; the minor-premise-spawning variant is macm.)

\par\smallskip{\raggedright\noindent\texttt{(slot 'acc)}\index{slot@\texttt{slot} (tactic)}\par}\nobreak
\noindent Reduce a structure ACCESSOR to its projection: (CARR s) becomes (NTH 1 s).  The one door for accessor reductions -- use it, never mac, on an accessor name.\par
\noindent\textit{Uses:} \texttt{cartesian-decompose} \texttt{macete} \texttt{tuple-equality-decompose}.\par

Replace an accessor by the tuple position it stands for: (CARR s) is slot 1, so it becomes (NTH 1 s).  Pair it with nth-r to compute on a concrete structure -- (MUL ZZ-RING) reduces to bintimes.  It refuses anything that is not an accessor.  (Technically: fires the accessor's macete, which def-structure installs at declaration.  It exists so that ALL accessor reductions go through ONE procedure: the reduction is currently global and unconditional -- index k for every argument -- and making it structure-relative later, guarded by IS-X(s), should change this door and not its callers.  accessor-callsite-audit fails the suite if any file fires an accessor macete by name.)

\par\smallskip{\raggedright\noindent\texttt{(slot-h 'acc hyp)}\index{slot-h@\texttt{slot-h} (tactic)}\par}\nobreak
\noindent Reduce a structure ACCESSOR to its projection inside a cited HYPOTHESIS -- `slot', hypothesis-side.\par
\noindent\textit{Uses:} \texttt{macete-hyp}.\par

The hypothesis-side door of slot: replace (PTS RR-MS) by the tuple position it stands for inside an assumption you cite, instead of in the goal.  Before it existed, a proof needing that reduction in a hypothesis reached for mac-h with the accessor macete's own name (rr-ms@pts) -- firing an accessor macete BY NAME, which is what the accessor pin (accessor-callsite-audit) exists to prevent, and which kept the suite red.  It refuses anything that is not an accessor, and it ERRORS rather than guess when the hypothesis mentions two different instances of the same accessor: after the first rewrite the assumption is a different formula, so a silent half-rewrite would be the failure mode.  (Technically: slot's projection-macete lookup run over the cited assumption, then cmd-apply-macete-to-assumption.)

\par\smallskip{\raggedright\noindent\texttt{(macm 'name)}\quad\tactkind{rule}\index{macm@\texttt{macm} (tactic)}\par}\nobreak
\noindent Like mac, but a CONDITIONAL macete fires even when its side-conditions are not yet in context: each unmet condition is left as a new subgoal.\par
\noindent\textit{Uses:} \texttt{cartesian-decompose} \texttt{macete} \texttt{tuple-equality-decompose}.\par
\noindent\textit{When useful:} the GOAL has a CONDITIONAL identity/definition to apply whose side-conditions you are willing to leave as subgoals\par

The minor-premises variant of mac.  mac only rewrites where the law's side-conditions ALREADY hold; macm rewrites regardless, spawning each unmet condition as a fresh goal to discharge -- the goal-side mirror of what mac-h already does on hypotheses.  This is what makes conditional finite-sum lemmas (finsum-*, sum-*, ...) usable as goal rewrites instead of only forward via fact.  (Technically: goal-side rewriting by a conditional macete; it emits the SAME single `macete' kernel rule as mac -- the unmet conditions become that one step's minor-premise subgoals.  It introduces NO new inference rule.)

\par\smallskip{\raggedright\noindent\texttt{(mac-h 'name hyp)}\quad\tactkind{rule}\index{mac-h@\texttt{mac-h} (tactic)}\par}\nobreak
\noindent Rewrite a cited ASSUMPTION in place with an equivalence macete; any side-condition not already in context is spawned as a new subgoal.\par
\noindent\textit{Uses:} \texttt{macete-hyp}.\par
\noindent\textit{When useful:} a HYPOTHESIS has a definition to unfold / an equivalence to apply\par

Like mac, but rewrites inside one of your HYPOTHESES instead of the goal -- replacing it by an equivalent statement (e.g. unfolding a definition in a hypothesis).  If the rewrite law has a side-condition you haven't established, that side-condition becomes a new goal to prove.  (Technically: hypothesis-side rewriting by Leibniz substitution of equivalents; sound because the macete is a genuine equivalence under its side-conditions, which are spawned as goals.)

\par\smallskip{\raggedright\noindent\texttt{(mac-h*)}\quad\tactkind{composite}\index{mac-h*@\texttt{mac-h*} (tactic)}\par}\nobreak
\noindent Saturating mac-h: repeatedly unfold every defined predicate in the hypotheses and split the conjunctions they expose, until nothing is left folded.  One step in the trace instead of a dozen.\par
\noindent\textit{Uses:} \texttt{and-elim} \texttt{forsome-elim} \texttt{iff-elim} \texttt{macete-hyp} \texttt{not-elim} \texttt{or-elim}.\par
\noindent\textit{When useful:} several hypotheses are folded definitions / conjunctions to open at once\par

The hands-free version of mac-h.  Instead of naming each definition and each conjunction by hand -- (mac-h 'IS-METRIC-SPACE A1), split, (mac-h 'is-metric A2), split, ... -- (mac-h*) keeps unfolding any assumption whose head is a defined predicate and splitting any AND assumption, restarting until a full pass changes nothing.  It records as a SINGLE step, so the proof (and its PDF) is far shorter.  It adds no kernel rule: every unfold is the same kernel-checked mac-h, every split the same ai -- just driven to a fixpoint.  Ask (what-now) to see which hypotheses it would touch first.

\par\smallskip{\raggedright\noindent\texttt{(grind)}\quad\tactkind{composite}\index{grind@\texttt{grind} (tactic)}\par}\nobreak
\noindent Saturate the no-choice moves at the focus: decompose the goal connective (di) and break open the hypotheses (mac-h*), repeatedly, until neither fires.  The whole introduce-everything prefix in one step.\par
\noindent\textit{Uses:} \texttt{and-elim} \texttt{and-intro} \texttt{forall-intro} \texttt{forsome-elim} \texttt{iff-elim} \texttt{iff-intro} \texttt{implies-intro} \texttt{macete-hyp} \texttt{not-elim} \texttt{not-intro} \texttt{or-elim}.\par
\noindent\textit{When useful:} right after sp, or any time the focus has connective/definitional structure -- normalize first\par

The deterministic normalizer.  It keeps applying di -- which strips a leading FORALL/IMPLIES/AND, introducing the bound variables and moving antecedents into your hypotheses -- and mac-h*, which unfolds defined predicates and splits conjunctions in the hypotheses, until the goal head is no longer a connective and no hypothesis is foldable.  None of these moves involves CHOOSING a lemma, so there is never anything to reconsider: grind just puts the focus into normal form (e.g. a metric-law goal becomes `(d s)(x,y) = (d s)(y,x)' with the metric's defining properties sitting unfolded in context).  It collapses the di...di mac-h* prefix that bloats proofs into a single step, and adds no kernel rule.

\par\smallskip{\raggedright\noindent\texttt{(scout [depth [branch [nodes]]])}\quad\tactkind{meta}\index{scout@\texttt{scout} (tactic)}\par}\nobreak
\noindent Speculatively search to a bounded depth across INDEPENDENT scratch branches; RETURNS a nested list (number-of-branches (d b) goal best-partials closing-branches) rather than printing.  Never touches the live proof.\par
\noindent\textit{Uses:} every one of the 64 kernel operations.\par
\noindent\textit{When useful:} you are stuck and want the machine to TRY move-sequences (returns data)\par

The copilot's deep lane.  Where (what-now)/(tt) suggest one move, scout actually TRIES sequences: it clones the focus into a fresh deduction graph per branch (so backtracking is free -- a dead branch is just discarded), and does a BEST-FIRST search (frontier ordered by open subgoals left, ties broken deeper-first so it dives toward a closure) whose alphabet is (grind), the closers (ass/rfl/crs/arith), the top rewrite-index (mac) rules, the parameterless backchain (bc*) lemmas, and the INST LANE -- (inst+ <universal hyp> <context-typed term>), instantiating a `forall' hypothesis at a term the context already types and detaching the guard.  It RETURNS the result as data -- (examined-count (depth branch) goal best-partials closing-branches), where best-partials is ((open-goals-left (form ...)) ...) most-reduced first and closing-branches is ((form ...) ...) shortest first -- so you can pick it apart programmatically; use (scout-show ...) for the readable REPL report.  Every move is the real, kernel-checked tactic, so a closing branch is a genuine proof when you adopt it with (scout-run k).  A CITATION GUARD suppresses circular closures: a branch that discharges the goal by citing a library theorem that is the goal -- DIRECTLY (statement alpha-equal to the goal, `P proved by P') or TRANSITIVELY (the cited theorem was itself proven using the goal, i.e. a name alpha-equal to the goal is in its proof-debt closure -- the `compact=>complete by compact=>bongo + bongo=>complete' trap) -- is dropped from closing-branches, the theorem name recorded in *scout-citations*, and scout-show notes `the goal is already theorem X'.  (Loops built purely from independent ASSERTED facts have no live dependency edge to detect; that is warrant/debt hygiene, not scout's to check.)  Bounds: depth (default 6), per-node fan-out branch (default 3), total nodes (default 600).  El cheapo: it does not reason about which move is wise (only the witness terms are guessed, from what the context already types -- a bad guess just dies on the clone); it brute-forces a small tree and you eyeball the survivors.  The inst lane closes the metric laws past grind; goals that need a witness scout can't type (a fresh existential, a constructed term) still surface under best-partials.

\par\smallskip{\raggedright\noindent\texttt{(scout-show [depth [branch [nodes]]])}\quad\tactkind{meta}\index{scout-show@\texttt{scout-show} (tactic)}\par}\nobreak
\noindent Like (scout), but PRINTS the human-readable report (examined count, goal, closing branches or best partials) to the REPL.  Returns the same nested list (scout) does.\par
\noindent\textit{Uses:} every one of the 64 kernel operations.\par
\noindent\textit{When useful:} same as scout but you want the readable report at the REPL\par

The eyeball version of scout: same search, same return value, but it also prints the numbered CLOSING branches (or, when none close, the best partials with how many goals each leaves open).  Use scout-show interactively, (scout) when you want to consume the result as data.

\par\smallskip{\raggedright\noindent\texttt{(scout-run k)}\quad\tactkind{composite}\index{scout-run@\texttt{scout-run} (tactic)}\par}\nobreak
\noindent Adopt closing branch k from the last (scout)/(scout-show) onto the live proof, running its tactics for real (they record and display normally).\par
\noindent\textit{Uses:} no kernel operation: it records no inference.\par
\noindent\textit{When useful:} scout found a closing branch you want to adopt onto the live proof\par

After scout finds CLOSING branches [1], [2], ..., (scout-run k) replays branch k's tactic forms through the real tactics on your live *ps*, from the same focus scout cloned -- so the proof advances and the steps are recorded for the script / PDF exactly as if you had typed them.  Works after either (scout) or (scout-show); both stash the closing branches.

\par\smallskip{\raggedright\noindent\texttt{(prep 'ineq)}\index{prep@\texttt{prep} (tactic)}\par}\nobreak
\noindent Why will a tactic not fire here, and what must be done first?  Runs the tactic's OWN preconditions one at a time, marks each ok / PREP / STOP, names the library lemma that repairs each unmet one, and prints a PLAN it has checked by running it on a scratch clone.  READ-ONLY: the live proof is never touched.  (prep 'ineq) is the only method so far.\par
\noindent\textit{Uses:} no kernel operation: it records no inference.\par

A tactic reports failure as a boolean, so every unmet precondition comes out as the same \#f and the same warning -- (ineq) says `goal not a linear-RR consequence of the named assumptions' whether your goal merely needs a di, or an atom needs a typing fact the library already proves, or the goal is simply false.  prep runs the same predicates separately and tells you WHICH failed.  For ineq the obligations are: the goal must BE an order relation (repair: di, counted by measuring, since one di consumes a typed FORALL and its guard together); every atom must be certified in RR by a LITERAL (IN t RR) scan -- (IN k NN) does not count, the oracle does no subtype reasoning (repair: a coercion lemma, e.g. nn-in-rr); some assumption must be order-shaped, since a fact like not(k=0) is invisible to Farkas (repair: a bridge lemma, e.g. nn-pos-of-nonzero); and the goal must actually follow, which only Fourier-Motzkin decides -- a STOP there means no prepping will help, which is the useful answer.  The repair search is by SHAPE over the whole library and is verified by SPECULATION: a candidate counts only if ineq then fires, not merely because it landed something order-shaped.  Every repair that works is reported, not just the first -- they come out alphabetical, which is no order of merit, and if the goal is itself a library theorem then citing IT is one of the closers.  Worked case: |- forall k in NN. \textasciitilde{}(k=0) => k < 2k.  prep returns (di) (di) (fact 'nn-in-rr 'k) (fact 'nn-pos-of-nonzero 'k) (ineq 4) -- six steps, where the hand-built cut/crs route through k = k+0 < k+k = 2k took thirteen and billed two extra transitivity lemmas.  A tactic joins the table with (prep-method! 'FUBA proc); crs and ass are the obvious next two.

\par\smallskip{\raggedright\noindent\texttt{(subst '(= s t))  |  (subst '(== s t))}\quad\tactkind{rule}\index{subst@\texttt{subst} (tactic)}\par}\nobreak
\noindent Rewrite s -> t throughout the goal, using an equation s = t -- or a quasi-equation s == t -- that is in context, in EITHER orientation (Leibniz substitution).  Reaches every position, operator (head) position included; a binder whose variable the equation mentions is not entered.\par
\noindent\textit{Uses:} \texttt{eq-subst}.\par
\noindent\textit{When useful:} you have an equation s = t (or a quasi-equation s == t) in context and want to rewrite the goal by it -- wherever the term occurs, operator position included\par

Use an equation among your hypotheses to replace s by t everywhere in the goal.  BOTH equalities license it: `=' is VNB's partial equality, and `s = t' already entails s = s and t = t, so t is defined wherever s was and the rewrite never replaces a defined term by an undefined one; `==' is quasi-equality (same definedness, equal where defined), which is a congruence and so substitutes exactly as `=' does -- which you need, since the partial-op recursion and bridge facts are stated with `=='.  The head you pass need not match the head in context, and neither need the orientation: all four of (= s t), (= t s), (== s t), (== t s) are searched for, and the goal is always rewritten s -> t.  LIMIT: the rewrite walk reaches argument positions only, so it is a silent no-op on a term in OPERATOR position -- (subst '(= (VADD md) ...)) will not touch the goal ((VADD md) x y).  Structure accessors are almost always in operator position; use the equation as a macete there (mac on the goal, mac-h on a hypothesis).  (Technically: Leibniz substitution from an in-context equality or quasi-equality.)

\par\smallskip{\raggedright\noindent\texttt{(beta)}\quad\tactkind{rule}\index{beta@\texttt{beta} (tactic)}\par}\nobreak
\noindent Beta-reduce a functoid application in the goal.\par
\noindent\textit{Uses:} \texttt{functoid-beta}.\par
\noindent\textit{When useful:} the goal has a functoid applied to an argument\par

Simplify a function-expression applied to an argument by substituting the argument into its body.  (Technically: beta-reduction of a functoid application in the goal.)

\par\smallskip{\raggedright\noindent\texttt{(nth-r)}\quad\tactkind{rule}\index{nth-r@\texttt{nth-r} (tactic)}\par}\nobreak
\noindent Reduce an NTH applied to a literal LIST in the goal.\par
\noindent\textit{Uses:} \texttt{nth-reduce}.\par
\noindent\textit{When useful:} the goal has NTH of a literal LIST\par

Simplify `the k-th entry of an explicit list [a1, a2, ...]' to that entry.  (Technically: reduce NTH applied to a literal LIST.)

\par\smallskip{\raggedright\noindent\texttt{(len-r)}\quad\tactkind{rule}\index{len-r@\texttt{len-r} (tactic)}\par}\nobreak
\noindent Reduce a LENGTH applied to a literal LIST in the goal.\par
\noindent\textit{Uses:} \texttt{length-reduce}.\par
\noindent\textit{When useful:} the goal has LENGTH of a literal LIST\par

Simplify `the length of an explicit list [a1, ..., an]' to the count n.  Structural (the dual of nth-r): the spine of a list literal is total, so its length is n regardless of whether the entries are defined.

\par\smallskip{\raggedright\noindent\texttt{(rfl)}\quad\tactkind{rule}\index{rfl@\texttt{rfl} (tactic)}\par}\nobreak
\noindent Close a reflexive equality goal (t = t).\par
\noindent\textit{Uses:} \texttt{reflexivity}.\par
\noindent\textit{When useful:} the goal is `t = t' with t manifestly defined\par

Close a goal `t = t' -- a thing equals itself.  (Technically: reflexivity of equality.)

\par\smallskip{\raggedright\noindent\texttt{(qrfl)}\quad\tactkind{rule}\index{qrfl@\texttt{qrfl} (tactic)}\par}\nobreak
\noindent Quasi-reflexivity: close t = t under the partial-equality definedness reading.\par
\noindent\textit{Uses:} \texttt{quasi-reflexivity}.\par
\noindent\textit{When useful:} the goal is `t = t' but t might be UNDEFINED (partial =)\par

Close `t = t' under the partial-equality reading, where asserting t = t also asserts that t is DEFINED.  Use this rather than rfl when t might be undefined.  (Technically: quasi-reflexivity; see the partial-equality convention, where `t = t' is the definedness predicate.)


\subsection*{Propositional logic}

\par\smallskip{\raggedright\noindent\texttt{(prop)}\index{prop@\texttt{prop} (tactic)}\par}\nobreak
\noindent Decide the goal by PROPOSITIONAL logic from the context, and close it if it follows.  Every non-connective formula -- `x in a', an equation, a whole `forall(...)' -- is one opaque atom; AND/OR/NOT/IMPLIES/IFF are read.  On a goal that does NOT follow it prints the countermodel (which atom must be true, which false) and leaves the proof untouched.  Adds no trust: it decides semantically, then discharges through di/ai/oi/ass/use-em/have!/detach!, so the bill is unchanged and the recorded script is the ordinary step-by-step proof.\par
\noindent\textit{Uses:} \texttt{and-elim} \texttt{and-intro} \texttt{arith-forsome} \texttt{arith-ground} \texttt{arith-simplify} \texttt{assumption} \texttt{cartesian-decompose} \texttt{cut} \texttt{detach} \texttt{eq-subst} \texttt{forall-elim} \texttt{forall-intro} \texttt{forsome-elim} \texttt{iff-elim} \texttt{iff-intro} \texttt{implies-intro} \texttt{macete} \texttt{not-elim} \texttt{not-intro} \texttt{or-elim} \texttt{or-intro-left} \texttt{or-intro-right} \texttt{proof-by-contradiction} \texttt{theorem-assumption} \texttt{tuple-equality-decompose}.\par
\noindent\textit{When useful:} the goal follows from the hypotheses by AND/OR/NOT/IMPLIES/IFF alone -- no quantifier or equality reasoning needed\par

Close a goal that follows from the hypotheses by pure propositional reasoning -- the leaves where you can SEE the answer and still have to pick between (oi-l)(ass), (oi-r) plus a conjunction split, and (ai) on a negation.  It treats each atomic statement as a black box, so it will not instantiate a quantifier or reason about equality: if the goal needs an instance, land the instance first (fact / inst+) and run it again.  When it declines it names a falsifying assignment, which usually tells you exactly which hypothesis is missing.  (Technically: three-valued evaluation over the atoms decides the entailment; the proof is then replayed through the kernel rules, case-splitting with excluded middle where a disjunctive goal needs it.)


\subsection*{Arithmetic \& ring oracles}

\par\smallskip{\raggedright\noindent\texttt{(arith)}\quad\tactkind{oracle}\index{arith@\texttt{arith} (tactic)}\par}\nobreak
\noindent Discharge a ground arithmetic goal by evaluation.\par
\noindent\textit{Uses:} \texttt{arith-forsome} \texttt{arith-ground} \texttt{arith-simplify} \texttt{cartesian-decompose} \texttt{eq-subst} \texttt{macete} \texttt{tuple-equality-decompose}.\par
\noindent\textit{When useful:} the goal is ground arithmetic -- no variables\par

Close a goal that is a concrete numerical fact with no variables -- e.g. 2 + 3 = 5, or 7 in NN -- by just computing it.  (Technically: decision by ground arithmetic evaluation.)

\par\smallskip{\raggedright\noindent\texttt{(rs)}\quad\tactkind{oracle}\index{rs@\texttt{rs} (tactic)}\par}\nobreak
\noindent Ring-simplify the goal (normal form over the ambient ring).\par
\noindent\textit{Uses:} \texttt{ring-simplify}.\par
\noindent\textit{When useful:} the goal is a (possibly non-commutative) ring-word identity\par

Simplify the goal to a normal form in the ambient ring (which need not be commutative).  The non-commutative companion of crs.  (Technically: ring-simplify to a canonical word form.)

\par\smallskip{\raggedright\noindent\texttt{(crs)}\quad\tactkind{oracle}\index{crs@\texttt{crs} (tactic)}\par}\nobreak
\noindent Commutative-ring decision procedure: prove a polynomial identity over ZZ[generators].  Expands literal powers, so (x+y)\textasciicircum{}2 = ... closes directly.\par
\noindent\textit{Uses:} \texttt{comm-ring-simplify}.\par
\noindent\textit{When useful:} the goal is a commutative-ring identity, e.g. (x+y)\textasciicircum{}2 = x\textasciicircum{}2+2xy+y\textasciicircum{}2\par

Prove a polynomial identity that holds in EVERY commutative ring -- e.g. (x+y)\textasciicircum{}2 = x\textasciicircum{}2 + 2xy + y\textasciicircum{}2 -- by reducing both sides to a canonical sum-of-monomials form and checking they agree.  A genuine decision procedure: a true commutative-ring identity closes, a non-identity is refused.  Literal powers are expanded for you.  (Technically: normal form over the free commutative ring ZZ[generators].)

\par\smallskip{\raggedright\noindent\texttt{(simp [target])}\quad\tactkind{oracle}\index{simp@\texttt{simp} (tactic)}\par}\nobreak
\noindent Rewrite a commutative-ring SUBTERM of the goal to canonical form, IN PLACE (e.g. (x+y)\textasciicircum{}2 inside a larger goal becomes x\textasciicircum{}2 + 2*x*y + y\textasciicircum{}2).  Works on BOTH surfaces: concrete number domains (+ * - \textasciicircum{} over NN/ZZ/QQ/RR/CC) and a generic ring s ((ADD s)/(MUL s)/(NEG s), carrier (CARR s)).  No arg = outermost ring subterm; ``term'' targets a specific one.  Sound by cut + crs + eq-subst (no new kernel rule); needs the subterm's generators typed in context (true post-di), else refuses and names them.\par
\noindent\textit{Uses:} \texttt{comm-ring-simplify} \texttt{cut} \texttt{eq-subst}.\par
\noindent\textit{When useful:} a ring SUBTERM of a larger goal should be put in normal form in place\par
\par\smallskip{\raggedright\noindent\texttt{(supply)}\index{supply@\texttt{supply} (tactic)}\par}\nobreak
\noindent Close an inequality goal the way a hand proof would: beta-reduce any applied lambda, land the typing certificates and the standard bounds the linear oracle cannot see -- including the TRIANGLE inequality at the summands of an abs of a sum -- then call ineq over every premise.  Rehearses the whole sequence on a throwaway copy and does NOTHING unless it closes; on a miss it prints the sequence it tried.  Adds no trust: it drives lam-b, fact and ineq, each of which records itself.\par
\noindent\textit{Uses:} every one of the 64 kernel operations.\par

The committing form of what-now's SUPPLY THE ORACLE lane.  `ineq' reads abs(...), max(...) and (dist(s))(x,y) as opaque atoms and refuses any premise whose atoms are not certified real, so an inequality that is obviously true can fail for want of a typing nobody mentioned.  This lands them.  The one real idea is subadditivity: where the goal bounds abs(u + v), the useful intermediate term is abs(u) + abs(v), and the decomposition is in the term itself -- nothing is searched for.  Because a lambda reduction cannot be undone and can owe an unprovable leaf, the whole sequence is rehearsed before any of it is run.  (Technically: certificate closure to depth 2 over the forward-citation lane, plus the curated bound table, then Fourier-Motzkin.)

\par\smallskip{\raggedright\noindent\texttt{(ineq i1 i2 ...)}\quad\tactkind{oracle}\index{ineq@\texttt{ineq} (tactic)}\par}\nobreak
\noindent Close a linear-inequality goal over RR (<= < > >= = between RR terms) as a consequence of the named assumptions (1-based indices), via the Fourier-Motzkin/Farkas oracle.  Linearizes over + - * and the binplus/binneg/bintimes aliases; every MAXIMAL non-arithmetic subterm is an atom that must be certified in RR.  (Does NOT see through a generic ring's (ADD s)/(MUL s) -- those become opaque atoms.)\par
\noindent\textit{Uses:} \texttt{ineq}.\par
\noindent\textit{When useful:} a LINEAR <= / < over RR that follows from inequalities you can cite\par

Close a LINEAR inequality over the reals that follows from inequalities you cite (by their hypothesis numbers) -- by chaining them, adding them, and scaling by positive constants.  Anything that is not built from + - and multiplication-by-constants is treated as an opaque quantity, so it handles e.g. the triangle inequality where the distances are unknowns.  For NONLINEAR (polynomial) inequalities use sos instead.  (Technically: Fourier-Motzkin / Farkas over the ordered field RR; every atom must be certified real.)

\par\smallskip{\raggedright\noindent\texttt{(sos ``c1'' ``c2'' ...)}\quad\tactkind{oracle}\index{sos@\texttt{sos} (tactic)}\par}\nobreak
\noindent Sum-of-squares closer for a nonstrict polynomial inequality a <= b over RR (the nonlinear companion of (ineq)).  You supply the terms to be SQUARED; it finds the nonnegative coefficients.  (tactics 'sos) for the worked example.\par
\noindent\textit{Uses:} \texttt{sos}.\par
\noindent\textit{When useful:} a NONSTRICT POLYNOMIAL <= over RR (the nonlinear cousin of ineq)\par

A sum-of-squares closer.  To prove  a <= b  it is enough to exhibit b - a  as a sum of squares (each obviously >= 0), reducing the inequality to an algebraic identity.  You hand sos the terms to be SQUARED -- the c\_i, NOT the squares -- and it solves for nonnegative rationals lambda\_i with

\begin{VNBsnip}
      b - a  =  lambda_1 c_1^2 + ... + lambda_n c_n^2
\end{VNBsnip}

(matched coefficient-by-coefficient in crs's commutative-ring normal form), then PRINTS the lambda\_i it found.  You do not supply them.

Worked example.  Goal  forall([x in rr, y in rr], x*y <= x\textasciicircum{}2 + y\textasciicircum{}2).  Here b - a = x\textasciicircum{}2 - x*y + y\textasciicircum{}2 = 1/2(x-y)\textasciicircum{}2 + 1/2 x\textasciicircum{}2 + 1/2 y\textasciicircum{}2,  so the things squared are  x-y, x, y:

\begin{VNBsnip}
      (sos "x - y" "x" "y")
          ;; closes, printing  1/2(x-y)^2 + 1/2 x^2 + 1/2 y^2
\end{VNBsnip}

2*x*y <= x\textasciicircum{}2 + y\textasciicircum{}2 needs only one square:  (sos ``x - y'').  The three-variable a*b+b*c+c*a <= a\textasciicircum{}2+b\textasciicircum{}2+c\textasciicircum{}2  wants  (sos ``a - b'' ``b - c'' ``c - a'').

Notes. - Write the c\_i with the goal's own variable names.  di preserves them, so sos works before OR after di; skip di and sos peels the typed forall([... in rr]) wrapper itself. - Every generator (variable) must be certified in RR -- a goal binder x in rr  or an  (IN g RR)  assumption. - A wrong or insufficient certificate refuses cleanly: nothing is closed, you simply try other squares. - Strict (<) goals are refused -- a square can be 0, so squares alone never force a strict inequality.


\subsection*{Chained reasoning}

\par\smallskip{\raggedright\noindent\texttt{(calc L0 (rel1 L1 [just1]) (rel2 L2 [just2]) ...)}\quad\tactkind{composite}\index{calc@\texttt{calc} (tactic)}\par}\nobreak
\noindent Ground the focus goal (REL L0 Ln) by a CHAIN of intermediaries L0 rel1 L1 rel2 L2 ... reln Ln: proves each link and composes them into the endpoint.  A link no lane can close is LEFT OPEN as a leaf -- the refinement point, and the PSS candidate.  Relations: = == (folded by transitivity), < <= (folded through the co-*-trans lemmas, with = steps riding along), IFF (chained per direction).  RETURNS the links, the open ones, and their PSS candidates as an alist.\par
\noindent\textit{Uses:} every one of the 64 kernel operations.\par
\noindent\textit{When useful:} the goal is (REL L0 Ln) and you can WRITE the chain of intermediaries that gets there -- or you want the one link that will not close isolated as a PSS candidate\par

Write the argument the way you would on paper -- as a chain through intermediate terms -- and let the machine prove each link:

\begin{VNBsnip}
      Goal  k < 2*k   (with k in NN, ~(k=0))
      (calc 'k '(= (+ k 0)) '(< (+ k k)) '(= (* 2 k)))
          ;; i.e.  k = k+0 < k+k = 2*k
\end{VNBsnip}

Each step names a RELATION and the next LINE; calc cuts the link (rel L(i-1) Li), dispatches a lane at it, and on success composes the whole chain into the goal.  It is pure bookkeeping over the trusted tactics -- no kernel rule, no axiom of its own -- so a calc proof bills exactly the lemmas its lanes and composers cite.

The COMPOSER is forced by the relations you used, and there are three families: `cong' for = and ==, folded by eq-trans; `order' for < and <=, folded through the co-lt-trans / co-le-lt-trans / co-eq-lt-trans ... lemmas (a = step rewrites an endpoint of the running relation and rides along, which is why the chain above needs no separate arithmetic); and `iff', which di's the goal into its two directions and chains each with ai/detach!, because VNB cannot quantify over propositions and so has NO first-order iff-trans lemma to fold with.

The JUSTIFICATION of a link is optional and defaults to 'auto -- try the relation's lanes: crs then arith for = / ==; for < / <= the NN->RR bridge then ineq over the order-shaped premises; grind for IFF.  Otherwise pass 'scout (a small scout search), a macete/theorem NAME (mac it, then close), an explicit tactic form like '(fact 'nn-pos-of-nonzero 'k) eval'd at the link's focus, or 'open / \#f to leave the link open ON PURPOSE.

Sanity first: calc ERRORS before touching the proof if the goal's head is not the composed relation, or its LHS is not L0, or its RHS is not Ln.  A link already in the main context is taken as a GIVEN and skipped -- cutting it would be an alpha self-loop that opens no leaf.

What it is FOR is the open links.  A chain whose links are all closed is a proof; a chain with one open link has isolated your obstacle to a single formula, printed with its free variables and their context typing -- which is the PSS candidate to state as a lemma.  Refine by inserting more intermediaries until each link is discoverable.  (See also (prep 'ineq), which answers the other question: why a lane will not fire.)

\par\smallskip{\raggedright\noindent\texttt{(have! CLAIM [THUNK])}\quad\tactkind{composite}\index{have"!@\texttt{have"!} (tactic)}\par}\nobreak
\noindent Assert CLAIM as an intermediate step and carry on with it as a hypothesis: cut CLAIM, discharge the side goal from context (or by THUNK), and stay on the main branch.\par
\noindent\textit{Uses:} \texttt{and-intro} \texttt{arith-forsome} \texttt{arith-ground} \texttt{arith-simplify} \texttt{assumption} \texttt{cartesian-decompose} \texttt{cut} \texttt{detach} \texttt{eq-subst} \texttt{forall-elim} \texttt{forall-intro} \texttt{iff-intro} \texttt{implies-intro} \texttt{macete} \texttt{not-intro} \texttt{theorem-assumption} \texttt{tuple-equality-decompose}.\par
\noindent\textit{When useful:} you want to state an intermediate fact that follows immediately from context and continue with it -- the `we have X' step\par

The `we have X' step.  You state an intermediate fact CLAIM that follows immediately from what is already known; have! cuts it, proves the resulting side goal automatically, and returns you to the main branch with CLAIM now available as a hypothesis.  The automatic discharge (`from-context!') closes the everyday cases -- a conjunction, splitwise; an (IN (a*b) NN) by nn-mul-closed on the factors; a numeral membership by arith; otherwise the goal is already a context assumption up to alpha.  When the side goal needs more than that, pass a THUNK -- any tactic sequence -- as the second argument to discharge it your way.  It ERRORS, never silently no-ops, if CLAIM is already in context up to alpha: the cut would self-loop, opening one child and no main branch.  So a script reads as the argument does -- `we have q0*q0 = 3*(k*k); we have k =/= 0; ...'.  (Technically: composite -- cut then from-context! or THUNK; adds no kernel rule.)


\subsection*{Backchaining with a theorem}

\par\smallskip{\raggedright\noindent\texttt{(ta 'name)}\quad\tactkind{rule}\index{ta@\texttt{ta} (tactic)}\par}\nobreak
\noindent Theorem-assumption: bring the named installed theorem into context as an assumption.\par
\noindent\textit{Uses:} \texttt{theorem-assumption}.\par
\noindent\textit{When useful:} you want a library theorem available as a hypothesis\par

Bring an already-proved theorem into your current hypotheses so you can use it (instantiate it, detach from it, ...).  (Technically: adds the named installed theorem as an assumption.)

\par\smallskip{\raggedright\noindent\texttt{(wbc ['name])}\quad\tactkind{composite}\index{wbc@\texttt{wbc} (tactic)}\par}\nobreak
\noindent Witness-shape backchain: on an `exists v. ...' goal, cite a PSS lemma that MANUFACTURES a witness of that shape, so you can finish with inst+/grind/ew.\par
\noindent\textit{Uses:} \texttt{and-elim} \texttt{arith-forsome} \texttt{arith-ground} \texttt{arith-simplify} \texttt{cut} \texttt{detach} \texttt{forall-elim} \texttt{forsome-elim} \texttt{iff-elim} \texttt{not-elim} \texttt{or-elim} \texttt{theorem-assumption}.\par
\noindent\textit{When useful:} an EXISTENCE goal whose witness a PSS lemma manufactures (diagonalization / block-family)\par

The discovery move for existence proofs.  Faced with `there exists phi with ...', it looks up the produces-witness index for lemmas whose conclusion builds the RIGHT KIND of object -- the same existential-witness shape -- even when the rest of the statement differs (which is why plain bc* misses them: bc* demands the WHOLE conclusion unify).  E.g. on `exists phi. STRICTLY-MONO-NN(phi) and IS-CAUCHY-SEQ(SUBSEQ f phi)' it reaches for `diagonalization' (which makes a STRICTLY-MONO-NN phi from a nested family); on a nested-family goal it reaches for `block-family'.  With no argument it cites the top retrieved producer; (wbc 'name) cites the one you name.  It brings the lemma in (via ta); you then discharge its hypotheses (inst+), skolemize its existential (grind), and supply its witness to your goal (ew).  Circular `headline' lemmas that already conclude your exact goal are filtered out.  (See (witness-producers goal) for the raw retrieval; what-now lists the (wbc 'name) candidates on any existential goal.)

\par\smallskip{\raggedright\noindent\texttt{(bc impl)}\quad\tactkind{rule}\index{bc@\texttt{bc} (tactic)}\par}\nobreak
\noindent Backchain the goal through an (IMPLIES A B) already in context: if the goal matches B, the new goal is A.\par
\noindent\textit{Uses:} \texttt{backchain}.\par
\noindent\textit{When useful:} you have (P => Q) in context and the goal is Q\par

Work backwards through an implication you already have.  If `P implies Q' is among your hypotheses and your goal is Q, this reduces the goal to proving P.  The everyday `to get Q it's enough to show P'.  (Technically: backchaining the goal through an in-context implication whose conclusion matches.)

\par\smallskip{\raggedright\noindent\texttt{(bc* 'thm [((v val)...)] h1 h2 ...)}\quad\tactkind{composite}\index{bc*@\texttt{bc*} (tactic)}\par}\nobreak
\noindent Matching backchain: unify the theorem's conclusion against the goal, then leave its (instantiated) antecedents as subgoals -- optional handlers hk run on the k-th subgoal.\par
\noindent\textit{Uses:} \texttt{assumption} \texttt{backchain} \texttt{cut} \texttt{forall-elim} \texttt{theorem-assumption}.\par
\noindent\textit{When useful:} a library theorem's CONCLUSION matches your goal\par

Apply a known theorem to your goal.  If you have a theorem `if A and B then C' and your goal is its conclusion C -- matching the theorem's variables to yours, and the names of any dummy/bound variables don't matter (`there exists a net N' applies to a goal `there exists a net F') -- this replaces `prove the goal' with `prove A' and `prove B', the theorem's hypotheses.  The everyday `by Theorem X it suffices to show A and B'.  You may attach a handler to each hypothesis to dispatch it; values the match can't determine you supply as ((v val) ...).  (Technically: peels the theorem's leading universals and implications, matches the conclusion -- alpha-aware on bound variables -- and replays ta/inst/cut/bc automatically.)


\subsection*{Lambda, comprehension \& description}

\par\smallskip{\raggedright\noindent\texttt{(lam-t)}\quad\tactkind{rule}\index{lam-t@\texttt{lam-t} (tactic)}\par}\nobreak
\noindent VNB-LAMBDA typing: reduce (IN (VNB-LAMBDA ...) (FUN A B)) to TWO obligations -- the body, and that the domain is a set.\par
\noindent\textit{Uses:} \texttt{lambda-type}.\par
\noindent\textit{When useful:} the goal types a VNB-LAMBDA into FUN(A,B)\par

Show that a function defined by a formula (`x |-> ...') maps A into B -- reduces to showing that, for an arbitrary input in A, the value lies in B, AND that A is a set.  The second is not a formality: carrying the domain in the term stops one lambda from typing into FUN(A,B) for every A, but says nothing about A being a set, and a lambda over a proper class is not a function.  The lambda's declared domain must also be the FUN's, or the rule does not apply at all.  (Technically: VNB-LAMBDA typing into FUN A B.)

\par\smallskip{\raggedright\noindent\texttt{(lam-b)}\quad\tactkind{rule}\index{lam-b@\texttt{lam-b} (tactic)}\par}\nobreak
\noindent VNB-LAMBDA beta: reduce an applied lambda to its substituted body.\par
\noindent\textit{Uses:} \texttt{lambda-beta}.\par
\noindent\textit{When useful:} the goal has a VNB-LAMBDA applied to an argument\par

Simplify a function `x |-> e(x)' applied to an argument a to e(a) -- the body with the argument substituted in.  A lambda carries its DOMAIN, and it is only defined on that domain, so the reduction is licensed only where a is known to be in it: if the membership is neither in the context nor supplied by an enclosing binder, the step still fires but leaves (IN a A) as an extra subgoal.  Land the typing fact BEFORE the lam-b and no subgoal appears.  (Technically: VNB-LAMBDA beta-reduction.)

\par\smallskip{\raggedright\noindent\texttt{(lam-b-h hyp)}\quad\tactkind{rule}\index{lam-b-h@\texttt{lam-b-h} (tactic)}\par}\nobreak
\noindent VNB-LAMBDA beta in a cited ASSUMPTION -- what mac-h is to mac.\par
\noindent\textit{Uses:} \texttt{lambda-beta-hyp}.\par

Beta-reduce an applied lambda inside a hypothesis, in place.  Needed because a `fact' that instantiates a theorem's function variable at a lambda lands the APPLIED lambda in the CONTEXT, where the goal-side `lam-b' cannot reach it: union-of-opens-open at the identity family g := x |-> x lands is-open(md, big-union(i, fam, (x |-> x)(i))).  Without this the proof must detour through a cut beta-equation and a subst.  Cites nothing, so it adds no debt.

\par\smallskip{\raggedright\noindent\texttt{(sep-set)}\quad\tactkind{rule}\index{sep-set@\texttt{sep-set} (tactic)}\par}\nobreak
\noindent Separation sethood: the separation set \{x in A | p\} is a set.\par
\noindent\textit{Uses:} \texttt{sep-sethood}.\par
\noindent\textit{When useful:} the goal asserts a separation set \{ x in A | p \} is a set\par

Show that a set-builder set \{x in A | p(x)\} is genuinely a set.  (Technically: separation sethood -- a subclass of a set is a set.)

\par\smallskip{\raggedright\noindent\texttt{(sep-mi)}\quad\tactkind{rule}\index{sep-mi@\texttt{sep-mi} (tactic)}\par}\nobreak
\noindent Separation membership intro: prove (IN t \{x in A | p\}).\par
\noindent\textit{Uses:} \texttt{sep-mem-intro}.\par
\noindent\textit{When useful:} the goal is membership in a separation set \{ x in A | p \}\par

Prove that a term t belongs to \{x in A | p(x)\} by showing t lies in A and satisfies the condition p.  (Technically: separation-membership introduction.)

\par\smallskip{\raggedright\noindent\texttt{(sep-me hyp)}\quad\tactkind{rule}\index{sep-me@\texttt{sep-me} (tactic)}\par}\nobreak
\noindent Separation membership elim: split a separation-membership assumption into A-membership and the predicate.\par
\noindent\textit{Uses:} \texttt{sep-mem-elim}.\par
\noindent\textit{When useful:} a hypothesis is separation-set membership -- unpack it\par

From a hypothesis that t belongs to \{x in A | p(x)\}, extract the two facts it packs: t lies in A, and p(t) holds.  (Technically: separation-membership elimination.)

\par\smallskip{\raggedright\noindent\texttt{(comp-mi)}\quad\tactkind{rule}\index{comp-mi@\texttt{comp-mi} (tactic)}\par}\nobreak
\noindent Comprehension membership intro.\par
\noindent\textit{Uses:} \texttt{comp-mem-intro}.\par
\noindent\textit{When useful:} the goal is comprehension-set membership\par

Prove that something belongs to a comprehension set by establishing the set's defining condition for it.  (Technically: comprehension-membership introduction.)

\par\smallskip{\raggedright\noindent\texttt{(comp-me hyp)}\quad\tactkind{rule}\index{comp-me@\texttt{comp-me} (tactic)}\par}\nobreak
\noindent Comprehension membership elim.\par
\noindent\textit{Uses:} \texttt{comp-mem-elim}.\par
\noindent\textit{When useful:} a hypothesis is comprehension-set membership\par

From a comprehension-membership hypothesis, extract its defining condition.  (Technically: comprehension-membership elimination.)

\par\smallskip{\raggedright\noindent\texttt{(iota-d term)}\quad\tactkind{rule}\index{iota-d@\texttt{iota-d} (tactic)}\par}\nobreak
\noindent Definite-description: discharge the IOTA uniqueness obligation for `term'.\par
\noindent\textit{Uses:} \texttt{iota-def}.\par
\noindent\textit{When useful:} the goal involves IOTA and you must discharge its uniqueness obligation\par

Justify `the unique x such that p(x)' (a definite description) by proving that exactly one such x exists.  (Technically: IOTA -- discharge the uniqueness obligation for the described term.)

\par\smallskip{\raggedright\noindent\texttt{(bu-set)}\quad\tactkind{rule}\index{bu-set@\texttt{bu-set} (tactic)}\par}\nobreak
\noindent Big-union sethood.\par
\noindent\textit{Uses:} \texttt{big-union-sethood}.\par
\noindent\textit{When useful:} the goal asserts a BIG-UNION is a set\par

Show that a big union (the union of an indexed family of sets) is itself a set.  (Technically: big-union sethood.)

\par\smallskip{\raggedright\noindent\texttt{(bu-mi w)}\quad\tactkind{rule}\index{bu-mi@\texttt{bu-mi} (tactic)}\par}\nobreak
\noindent Big-union membership intro via the index witness w.\par
\noindent\textit{Uses:} \texttt{big-union-mem-intro}.\par
\noindent\textit{When useful:} the goal is BIG-UNION membership -- you have the index witness\par

Prove that an element lies in a big union by exhibiting one index w whose set already contains it.  (Technically: big-union-membership introduction via the index witness.)

\par\smallskip{\raggedright\noindent\texttt{(bu-me hyp)}\quad\tactkind{rule}\index{bu-me@\texttt{bu-me} (tactic)}\par}\nobreak
\noindent Big-union membership elim.\par
\noindent\textit{Uses:} \texttt{big-union-mem-elim}.\par
\noindent\textit{When useful:} a hypothesis is BIG-UNION membership -- obtain its index\par

From a hypothesis that an element lies in a big union, obtain an index whose set contains it (a fresh name for that index).  (Technically: big-union-membership elimination.)


\subsection*{Conditional terms}

\par\smallskip{\raggedright\noindent\texttt{(if-true t)}\quad\tactkind{rule}\index{if-true@\texttt{if-true} (tactic)}\par}\nobreak
\noindent Reduce an (IF p a b) term on the p branch: spawns p as a subgoal; the continuation gains (= (IF p a b) a).\par
\noindent\textit{Uses:} \texttt{if-true}.\par
\noindent\textit{When useful:} the goal has an (IF p a b) term to evaluate on the p branch\par

Evaluate a conditional term `if p then a else b' on the assumption that p holds: you take on p as a side goal, and may then use that the conditional equals a.  (Technically: if-true reduction, spawning p.)

\par\smallskip{\raggedright\noindent\texttt{(if-false t)}\quad\tactkind{rule}\index{if-false@\texttt{if-false} (tactic)}\par}\nobreak
\noindent Reduce an (IF p a b) term on the not-p branch: spawns (NOT p); the continuation gains (= (IF p a b) b).\par
\noindent\textit{Uses:} \texttt{if-false}.\par
\noindent\textit{When useful:} the goal has an (IF p a b) term to evaluate on the not-p branch\par

Evaluate a conditional term `if p then a else b' on the assumption that p fails: you take on `not p' as a side goal, and may then use that the conditional equals b.  (Technically: if-false reduction, spawning not p.)


\subsection*{Navigation, display \& tacticals}

\par\smallskip{\raggedright\noindent\texttt{(focus n)}\index{focus@\texttt{focus} (tactic)}\par}\nobreak
\noindent Switch the focus to the n-th open goal (1-based).\par
\noindent\textit{Uses:} no kernel operation: it records no inference.\par
\par\smallskip{\raggedright\noindent\texttt{(backup-one)}\quad\tactkind{meta}\index{backup-one@\texttt{backup-one} (tactic)}\par}\nobreak
\noindent Take back the last recorded command: restore the goal, the deduction graph and the proof script to the state before it.  Repeat to walk further back; (sp) clears the stack.  `undo' is an alias.\par
\noindent\textit{Uses:} no kernel operation: it records no inference.\par
\noindent\textit{When useful:} the last command was a mistake -- a greedy `di' that ate the induction, a `cut' you did not mean, a branch you want back\par

The one undo in VNB.  It restores the deduction graph -- not merely the focus -- by rolling back the journalled arrow and grounding writes and DROPPING every sequent and inference node posted since, so no node from the abandoned branch survives to be counted as an open leaf and `qed' cannot be handed a phantom obligation.  *proof-script* and the live trace are rewound with it, so the saved script and the printed proof are the proof you actually kept.  It backs up over the LAST RECORDED command: a command that changed nothing was never recorded and is not a step to back up over.  One thing is deliberately not restored -- *fresh-counter*, the eigenvariable source -- so backing up over an `ai' or `ew' that minted `u\_4' and re-running it mints `u\_5'; the proof is the same, the witness has a different name.  Depth is capped at *vnb-undo-depth* (64).

\par\smallskip{\raggedright\noindent\texttt{(undo)}\quad\tactkind{meta}\index{undo@\texttt{undo} (tactic)}\par}\nobreak
\noindent Alias for (backup-one).\par
\noindent\textit{Uses:} no kernel operation: it records no inference.\par
\noindent\textit{When useful:} the last command was a mistake -- alias for backup-one\par
\par\smallskip{\raggedright\noindent\texttt{(show)}\index{show@\texttt{show} (tactic)}\par}\nobreak
\noindent Redisplay the current proof state.\par
\noindent\textit{Uses:} no kernel operation: it records no inference.\par
\par\smallskip{\raggedright\noindent\texttt{(dial notch)}\index{dial@\texttt{dial} (tactic)}\par}\nobreak
\noindent The PRESENTATION RESOLUTION DIAL: how much of the sequent to show.  r0 kernel S-expression, r1 surface syntax (the default, and the last INVERTIBLE notch -- what is printed re-parses to what was printed), r2 guards and typing hypotheses elided, r3 also structures destructured to their declared slot names and deep subterms elided.  (dial) reports the notch, (dial n) sets it, (dial 'auto) reads the current goal and suggests one, (dial 'why) shows the score behind the suggestion.  For ONE FORMULA rather than the proof state, (dial-wff x [n]) -- x a wff, a ``surface string'' or an S-expression, n defaulting to the dial's current notch: (dial-wff (lookup-theorem 'matmul-assoc) 3).\par
\noindent\textit{Uses:} no kernel operation: it records no inference.\par

It is a FILTER, not a set of modes, and two constraints make it one.  The rungs NEST: turning it up only ever SUBTRACTS, so everything readable at a higher notch is readable below it -- checked in the suite, atom by atom.  And the APERTURE IS ON SCREEN: every notch above r1 names itself on the line and counts what it hid, because a lossy render that does not announce its own lossiness is the defect class of `(f)' printing as `f'.  The auto-notch SUGGESTS a starting notch and then STAYS PUT: it is consulted when you ask, never per sequent, so the lens cannot move while you are reading a proof.

\par\smallskip{\raggedright\noindent\texttt{(dial-wff x notch)}\index{dial-wff@\texttt{dial-wff} (tactic)}\par}\nobreak
\noindent Present ONE FORMULA at a resolution notch, instead of the current proof state: x is a wff, a ``surface string'' or an S-expression, and notch defaults to the dial's current setting.  (dial-wff (lookup-theorem 'matmul-assoc) 3).  Same rungs, same aperture line, no proof required.\par
\noindent\textit{Uses:} no kernel operation: it records no inference.\par

`sequent->string' is the display path for a proof STATE, so the dial reaches a formula only while a proof is holding it.  A theorem looked up, a statement being drafted, a hypothesis pasted from somewhere wants the same dial, and this is it.

\par\smallskip{\raggedright\noindent\texttt{(goal-status)}\index{goal-status@\texttt{goal-status} (tactic)}\par}\nobreak
\noindent One-line summary: done / N open goals.\par
\noindent\textit{Uses:} no kernel operation: it records no inference.\par
\par\smallskip{\raggedright\noindent\texttt{(repeat thunk [cap])}\index{repeat@\texttt{repeat} (tactic)}\par}\nobreak
\noindent Run thunk until it stops changing the proof state (LCF REPEAT).\par
\noindent\textit{Uses:} no kernel operation: it records no inference.\par
\par\smallskip{\raggedright\noindent\texttt{(orelse t1 t2 ...)}\index{orelse@\texttt{orelse} (tactic)}\par}\nobreak
\noindent Run thunks in order, stop at the first that makes progress (LCF ORELSE).\par
\noindent\textit{Uses:} no kernel operation: it records no inference.\par
\par\smallskip{\raggedright\noindent\texttt{(quietly thunk)}\index{quietly@\texttt{quietly} (tactic)}\par}\nobreak
\noindent Run thunk with state-dump output suppressed; returns its value.\par
\noindent\textit{Uses:} no kernel operation: it records no inference.\par

\subsection*{Choosing and naming witnesses}

\par\smallskip{\raggedright\noindent\texttt{(choose ex witness [prover])}\index{choose@\texttt{choose} (tactic)}\par}\nobreak
\noindent Choose eps such that FUBA(eps): cut the existential ex = `forsome([eps], FUBA(eps))', discharge it by exhibiting WITNESS (proving FUBA(WITNESS) by PROVER, or from the context), eliminate it, and RETURN the fresh eigenvariable with FUBA of it landed.\par
\noindent\textit{Uses:} \texttt{and-elim} \texttt{and-intro} \texttt{arith-forsome} \texttt{arith-ground} \texttt{arith-simplify} \texttt{assumption} \texttt{cartesian-decompose} \texttt{cut} \texttt{detach} \texttt{eq-subst} \texttt{forall-elim} \texttt{forall-intro} \texttt{forsome-elim} \texttt{forsome-intro} \texttt{iff-elim} \texttt{iff-intro} \texttt{implies-intro} \texttt{macete} \texttt{not-elim} \texttt{not-intro} \texttt{or-elim} \texttt{theorem-assumption} \texttt{tuple-equality-decompose}.\par

The `we may assume without loss of generality that FUBA(eps)' step: a hypothesis `forall([eps in rr], GUBA(eps))' is about to be used at some eps satisfying FUBA, so choose one first and then instantiate the universal at the name this returns (use-at / obtain-at).  You pay for the existence on the spot -- WITNESS is the value and PROVER (a thunk) shows FUBA holds of it; with no PROVER the side goal is closed from the context.  The rest of the proof then depends on FUBA(eps) alone and never on which witness paid for it, which is what makes it read as `without loss of generality'.  If the existential is already in context the cut is skipped and this is plain elimination; a formula that is not an existential is an error.  choose-pos is this tactic with FUBA frozen at pos-rr and witness 1: `let eps > 0 be given' when nothing gives it.  (Technically: composite -- cut / ew / ai and the AND-splitting of what lands; no kernel rule, no choice principle, and the recorded script is the expansion.)

\par\smallskip{\raggedright\noindent\texttt{(minimize! '(v1 ... vk) GUARD MEASURE)}\quad\tactkind{composite}\index{minimize"!@\texttt{minimize"!} (tactic)}\par}\nobreak
\noindent Choose v1..vk satisfying GUARD with the NN-valued MEASURE as small as possible: lands the witnesses and their minimality, and opens the two obligations any minimisation owes (MEASURE lands in NN; GUARD is satisfiable).\par
\noindent\textit{Uses:} \texttt{and-elim} \texttt{and-intro} \texttt{arith-forsome} \texttt{arith-ground} \texttt{arith-simplify} \texttt{assumption} \texttt{cartesian-decompose} \texttt{cut} \texttt{detach} \texttt{eq-subst} \texttt{forall-elim} \texttt{forall-intro} \texttt{forsome-elim} \texttt{forsome-intro} \texttt{iff-elim} \texttt{iff-intro} \texttt{implies-intro} \texttt{macete} \texttt{not-elim} \texttt{not-intro} \texttt{or-elim} \texttt{reflexivity} \texttt{sep-mem-elim} \texttt{sep-mem-intro} \texttt{theorem-assumption} \texttt{tuple-equality-decompose}.\par
\noindent\textit{When useful:} the goal falls to a `least such' / minimal-counterexample argument -- descent proofs (sqrt 2, sqrt 3 irrational), least-degree or least-pivot witnesses\par

The `least such' / minimal-counterexample step, mechanised.  You give three things: the variables to choose, a GUARD formula they must satisfy, and a natural-number-valued MEASURE term to make small.  On the main branch it hands you fresh eigenconstants w1..wk with two hypotheses -- GUARD holds of them, and nothing satisfying GUARD has a strictly smaller MEASURE -- and it opens exactly the two side goals a minimisation genuinely owes: that MEASURE lands in NN wherever GUARD holds (the TYPE goal), and that GUARD holds somewhere (the NONEMPTY goal).  Everything in between -- forming the value set, well-ordering it, unpacking the least witness, and restating minimality in terms of your variables rather than the set -- is done for you.  Every vi must occur in MEASURE (a variable the measure ignores is not one you are minimising over).  Because it quantifies over the vk at the META level, no lambda and no FUN(U,NN) typing obligation ever enters the logic -- it takes a formula and a term, which is what a driver has in hand, and never forms the set U at all.  Returns (list (w1 ... wk) TYPE-node NONEMPTY-node); either node is \#f when that obligation was already in context up to alpha, so test before refocusing.  (Technically: composite -- it drives cut / di / ai / ew / sep-mi / sep-me / fact / inst / detach! / subst / ass / rfl and adds no kernel rule; its one mathematical appeal is `nn-least-element', the well-ordering of NN, proven from ord-well-ordered.  NO choice principle is used: well-ordering returns a MEMBER of a separation set, and sep-me recovers the witness.)

\par\smallskip{\raggedright\noindent\texttt{(obtain LANE)}\quad\tactkind{composite}\index{obtain@\texttt{obtain} (tactic)}\par}\nobreak
\noindent Run LANE -- a thunk whose forward step lands a `there exists' into context -- then eliminate that existential and RETURN the fresh witness's name, read off the proof state rather than guessed.\par
\noindent\textit{Uses:} every one of the 64 kernel operations.\par
\noindent\textit{When useful:} a forward step yields `there exists ...' and you want to name the witness for later use, without guessing the engine's eigenvariable\par

The `let w be such an x' step, with the eigenvariable named for you.  LANE is a zero-argument procedure whose effect is to land some `there exists v. P(v)' among your hypotheses -- e.g. (lambda () (fact 'nn-3-div-square w)).  obtain runs it, spots the existential that newly appeared, eliminates it (introducing a fresh eigenvariable and landing P at that witness), unpacks any conjunction in the body, and hands back the eigenvariable's NAME -- identified by free-variable set-difference (the symbol now in context that was not there before), so the driver never guesses what the engine called it.  This is the robust way to name a descended witness: (define k (obtain (lambda () ...))), then use k.  Returns \#f (with a note) if the lane landed no existential.  (Technically: composite -- a guarded forward discharge followed by forsome-elim (ai) and AND-splitting.  Plain existential elimination, so it owes nothing and uses no choice.)

\par\smallskip{\raggedright\noindent\texttt{(vlet (n1 ...) FORMER)}\quad\tactkind{composite}\index{vlet@\texttt{vlet} (tactic)}\par}\nobreak
\noindent Bind proof-object names from the current state.  FORMER is (match PATTERN) -- bind each name to its slot in a context formula matching PATTERN -- or (choice [v body]) -- eliminate an in-context existential and bind its witness.\par
\noindent\textit{Uses:} \texttt{and-elim} \texttt{cut} \texttt{forsome-elim} \texttt{iff-elim} \texttt{not-elim} \texttt{or-elim}.\par
\noindent\textit{When useful:} you need to NAME a witness or a matched subterm from the proof state, rather than navigate to it by shape\par

A binding form for the pieces of a proof, so a driver NAMES what it needs instead of navigating to it by shape.  Two FORMERs.  (match PATTERN): the listed names ARE the holes (wildcards) in PATTERN; everything else is literal.  vlet searches the context for a formula or subterm matching PATTERN and binds each name to what filled its slot -- pure selection, no proof step, no obligation; a hard error if nothing matches (you named a piece of something absent).  (choice): eliminate the sole existential in context and bind its witness.  (choice v body): present-else-debt -- if `there exists v. body' is already in context up to alpha, eliminate it; otherwise cut it, LEAVE the existence side-goal open as a debt leaf, and eliminate on the main branch -- either way the witness is bound and body[witness] is landed.  vlet's (choice) and `obtain' are the same underlying mechanism -- eliminate an existential, name the witness by free-variable difference -- differing only in what they take: obtain runs a lane and names what it just produced, vlet names an existential already (or, with a body, about to be) in context.  No choice AXIOM is used; a single witness is plain existential-elimination.  (Technically: a define-syntax expanding to (define n ...) over vlet--match / vlet--choice!; composite, no kernel rule beyond forsome-elim and, for present-else-debt, cut.)




%% =========================================================
\section{Focus management}

After each command the focus advances automatically to the first
ungrounded subgoal just introduced.  If there are no new subgoals, it
moves to any remaining open goal.

\begin{remark}
Every proof command operates exclusively on the focus node; the other
open goals are untouched.  There can be several open goals simultaneously
(for example, after \texttt{(ci)} on a conjunction goal, or after
\texttt{(cut f)} which always creates two subgoals).  \texttt{(show)}
displays them all, numbering the non-focus goals so you can use those
numbers with \texttt{(focus $n$)}.  Until you call \texttt{(focus $n$)},
the system always works on whichever goal it chose automatically.
\end{remark}

In the interactive session use:
\begin{VNBcode}
(show)       ; list all open goals (numbered) and the current focus
(focus n)    ; switch focus to the n-th open goal (1-based)
\end{VNBcode}

The underlying API (for use outside the interactive session) is:
\begin{VNBcode}
(focus-on ps sqn)            ; focus on a specific sequent node
(focus-on-first-open ps)     ; focus on the first open goal
(proof-open-goals ps)        ; list all open goals
\end{VNBcode}

%% =========================================================
\section{Macetes}
\label{sec:macetes}

A \emph{macete} is a named, composable rewrite strategy.  When a
theorem is installed, its macete is automatically derived by extracting a
source pattern and replacement pattern from the theorem's logical form:

\begin{center}
\begin{tabular}{ll}
\toprule
Theorem form & Rewrite \\
\midrule
$\forall\bar{x}.\;(s = r)$ & $s \mapsto r$ \\
$\forall\bar{x}.\;(s \Leftrightarrow r)$ & $s \mapsto r$ \\
$\forall\bar{x}.\;(\text{conds} \Rightarrow (s = r))$ & $s \mapsto r$ when conds hold \\
$\forall\bar{x}.\;p$ & $p \mapsto \top$ \\
\bottomrule
\end{tabular}
\end{center}

The quantifiers need not all be leading.  A \texttt{forall} occurring in
positive position --- under an \texttt{implies} consequent or inside an
\texttt{and} --- is pulled to the front before the source and replacement
patterns are extracted, so a theorem such as
\texttt{forall([a,b], $H$ implies forall(x, $s$ iff $r$))} still yields the
rewrite $s \mapsto r$ (conditional on $H$).

Macetes can be composed:
\begin{VNBcode}
(macete-series m1 m2 m3)   ; apply in sequence
(macete-repeat m)           ; apply until no change
(macete-try m)              ; apply m, succeed even if m fails
(macete-choice m1 m2 m3)   ; first one that applies
\end{VNBcode}

To install a custom macete:
\begin{VNBcode}
(install-theorem! 'eq-symm
  (make-wff-from-string "forall([x, y], x = y iff y = x)"))

; Use it:
(cmd-apply-macete ps 'eq-symm)
\end{VNBcode}

\subsection{Applying macetes to wffs outside a proof}

\begin{sloppypar}The function \texttt{apply-macete} applies a named macete directly to a
\texttt{wff} without creating a proof state:\end{sloppypar}

\begin{itemize}
\item \texttt{(apply-macete \textit{name} \textit{w})} --- rewrites \textit{w};
  returns a new \texttt{wff} if any rewrite fires, or \texttt{\#f} if the
  pattern matches nowhere in \textit{w}.
\end{itemize}

This is purely syntactic: no deduction graph is involved and no
theorem is proved.  It is useful for inspecting rewrite behavior,
preprocessing formulas, or building higher-level tools outside of proofs.

\paragraph{Unconditional equation macete.}
\begin{VNBcode}
(install-theorem! 'nn-zero-succ
  (make-wff-from-string "forall([n in NN], succ(n) = n + 1)"))

(apply-macete 'nn-zero-succ (make-wff-from-string "succ(3) = x"))
; => wff for "3 + 1 = x"

(apply-macete 'nn-zero-succ (make-wff-from-string "3 in nn"))
; => #f               -- macete doesn't fire
\end{VNBcode}

\paragraph{Biconditional (\texttt{iff}) macete.}
\begin{VNBcode}
(install-theorem! 'eq-symm
  (make-wff-from-string "forall([x, y], x = y iff y = x)"))

(apply-macete 'eq-symm (make-wff-from-string "n + m = 0"))
; => wff for "0 = n + m"

(apply-macete 'eq-symm (make-wff-from-string "n = m implies truth"))
; => wff for "m = n implies truth"

(apply-macete 'eq-symm (make-wff-from-string "n in nn"))
; => #f
\end{VNBcode}

\paragraph{Closure axiom macete (conditional).}
\begin{VNBcode}
;; nn-add-closed: forall([n in NN, m in NN], n + m in NN)
;; Condition (n in NN and m in NN) must be present in local context.

(apply-macete 'nn-add-closed
  (make-wff-from-string "n in nn and m in nn implies n + m in nn"))
; => wff for "n in nn and m in nn implies truth"

(apply-macete 'nn-add-closed (make-wff-from-string "n + m in nn"))
; => #f  -- conditions not in context; rewrite correctly suppressed
\end{VNBcode}

\begin{remark}
\texttt{apply-macete} applies the rewrite depth-first throughout the
formula (outermost first within each subexpression).  Conditions of a
conditional macete are checked against the \emph{local context}
accumulated as the rewriter descends: when descending into the
consequent of \texttt{implies}, the antecedent is added to the local
context; when descending into the right conjunct of \texttt{and}, the
left conjunct is added; and so on (Monk~\cite{Monk1988}).  A condition that
matches a formula already present in the local context is considered
discharged; if any condition is unmet at a candidate site, the rewrite
is suppressed at that site.
\end{remark}

%% =========================================================
\section{Example proofs}

\subsection{Identity}

\begin{VNBcode}
(sp (make-wff-from-string "x in NN implies x in NN"))
(di)    ; assume (x in nn), prove (x in nn)
(ass)   ; it is in context
; => Proof complete.
\end{VNBcode}

\subsection{AND elimination}

\begin{VNBcode}
(sp (make-wff-from-string "n in NN and m in NN implies n in NN"))
(di)                          ; assume (n in nn) and (m in nn)
(ai "n in NN and m in NN")   ; split into (n in nn) and (m in nn)
(ass)                         ; prove (n in nn)
\end{VNBcode}

\subsection{VNB axiom instance}

Proving \texttt{"a in b implies a in SET"} from the base axiom:

\begin{VNBcode}
(sp (make-wff-from-string "a in b implies a in SET"))
(di)
(ta 'membership-implies-sethood)
(inst "forall([a, b], a in b implies a in set)" 'a)
(inst "forall([b], a in b implies a in set)" 'b)
(bc "a in b implies a in set")
(ass)
\end{VNBcode}

\subsection{Double negation (classical)}

\begin{VNBcode}
(sp (make-wff-from-string "not(not(x in NN)) implies x in NN"))
(di)                              ; assume not(not((x in nn))), prove (x in nn)
(pbc)                             ; assume not((x in nn)), prove falsity
(ai "not(not(x in NN))")         ; not(not(P)) and not(P) yield falsity
\end{VNBcode}

\subsection{Ordered tuples and Cartesian products}

\begin{VNBcode}
; From x in A, y in B, z in C prove [x,y,z] in cartesian(A,B,C)
(sp (make-wff-from-string
      "x in A and y in B and z in C implies [x, y, z] in cartesian(A, B, C)"))
(di)
(ai "x in A and y in B and z in C")
(ai "y in B and z in C")
(ci)              ; cartesian-intro generates 3 subgoals
(ass) (ass) (ass)

; NTH projection: nth(2, [x, y, z]) = y
(sp (make-wff-from-string "nth(2, [x, y, z]) = y"))
(nth-r)           ; reduces goal to (y = y)
(rfl)             ; closes by reflexivity
\end{VNBcode}

\subsection{Mathematical structures}

Algebraic structures are declared with \texttt{declare-structure}.
No quoting is needed: carrier names and operation signatures are written
directly.

\begin{VNBcode}
(declare-structure GROUP
  (carriers CARR)
  (op OPR (CARTESIAN CARR CARR) CARR)  ; operation: CARR x CARR -> CARR
  (constant IDEN CARR)                 ; identity: IDEN in CARR
  (op INV CARR CARR))                  ; inverse: CARR -> CARR
; Accessor indices: CARR -> 1, OPR -> 2, IDEN -> 3, INV -> 4
; (the group laws are folded in by additional (property ...) clauses)
\end{VNBcode}

\begin{sloppypar}The \texttt{(op name domain range)} clause declares
$\texttt{name}(s) \in \texttt{FUN}(\texttt{domain}, \texttt{range})$.
Higher-arity operations write the Cartesian product explicitly:
\texttt{(op OPR (CARTESIAN CARR CARR) CARR)} for a binary operation.
A \texttt{(constant name S)} clause declares an element of set $S$.
Accessor indices are assigned sequentially: carriers first (starting
at~1), then operations in declaration order.\end{sloppypar}

A \texttt{declare-structure} call installs two things automatically:
\begin{itemize}
\item \emph{Accessor macetes}: \texttt{CARR}$(s)$ rewrites to
  \texttt{nth(1,$s$)}, \texttt{OPR}$(s)$ to \texttt{nth(2,$s$)}, etc.
\item \emph{Definitional biconditional} \texttt{IS-NAME}: the axiom
  $\forall s.\;\texttt{IS-GROUP}(s) \Leftrightarrow
   (\texttt{CARR}(s)\in\mathrm{SET} \wedge \texttt{OPR}(s)\in\cdots)$.
\end{itemize}

Given \texttt{IS-GROUP(g)} in context, the \texttt{IS-GROUP} macete
unfolds to:
\begin{VNBsnip}
CARR(g) in SET
and OPR(g) in fun(cartesian(CARR(g), CARR(g)), CARR(g))
and IDEN(g) in CARR(g)
and INV(g) in fun(CARR(g), CARR(g))
\end{VNBsnip}

\paragraph{Pre-loaded structures.}
The following structures are defined in \texttt{structure-library/} (one
file per structure) and loaded automatically.

\begin{center}\scriptsize
\begin{tabular}{@{}p{2.3cm}p{4.3cm}p{5.3cm}@{}}
\toprule
\textbf{Structure} & \textbf{Accessor slots (index)} & \textbf{Characteristic axioms} \\
\midrule
\texttt{SEMIGROUP} &
  \texttt{CARR}\,(1), \texttt{OPR}\,(2) &
  \texttt{semigroup-assoc} \\[4pt]
\texttt{MONOID} &
  \texttt{CARR}\,(1), \texttt{OPR}\,(2), \texttt{IDEN}\,(3) &
  \texttt{monoid-assoc},
  \texttt{monoid-left-id},
  \texttt{monoid-right-id} \\[4pt]
\texttt{GROUP} &
  \texttt{CARR}\,(1), \texttt{OPR}\,(2), \texttt{IDEN}\,(3), \texttt{INV}\,(4) &
  \texttt{group-assoc},
  \texttt{group-left-id},
  \texttt{group-left-inv} \\[4pt]
\texttt{RING} &
  \texttt{CARR}\,(1), \texttt{ADD}\,(2), \texttt{MUL}\,(3),
  \texttt{NEG}\,(4), \texttt{ZERO}\,(5), \texttt{ONE}\,(6) &
  \texttt{ring-add-assoc}, \texttt{ring-add-comm},
  \texttt{ring-add-left-id}, \texttt{ring-add-left-inv},
  \texttt{ring-mul-assoc}, \texttt{ring-mul-left-id},
  \texttt{ring-mul-right-id},
  \texttt{ring-left-dist}, \texttt{ring-right-dist} \\[4pt]
\texttt{METRIC-SPACE} &
  \texttt{PTS}\,(1), \texttt{DIST}\,(2) &
  \texttt{metric-pos},
  \texttt{metric-self-zero},
  \texttt{metric-zero-eq},
  \texttt{metric-sym},
  \texttt{metric-triangle} \\
\bottomrule
\end{tabular}
\end{center}

\begin{remark}
Each entry above is citable by name with \texttt{(ta 'name)}, which transports
it into the current proof.  They do \emph{not} all have the same standing,
however, and the column heading is a convenience rather than a claim.

A structure declaration generates an \texttt{IS-X} defining biconditional, and
each operation-property conjunct of it, unfolded, \emph{is} one of the
structure's laws.  Stating such a law separately therefore adds no mathematical
content, and the question is only how the system records that fact.  Three
answers are in use, and they cost a citing \texttt{qed} the same --- nothing ---
by three different routes:

\begin{itemize}
\item \emph{Proved.}  The \textsc{Metric-Space} laws (\texttt{metric-pos},
  \texttt{metric-self-zero}, \texttt{metric-zero-eq}, \texttt{metric-sym},
  \texttt{metric-triangle}) are established by \texttt{prove-metric-law!} in
  \texttt{structure-library/metric-laws.scm}; the \textsc{Group} laws
  (\texttt{group-assoc}, \texttt{group-left-id}, \texttt{group-left-inv},
  \texttt{group-identity-in}) and \texttt{abelian-group-opr-comm} by
  \texttt{stl--project!} in \texttt{structure-library/subtype-laws.scm}.  The
  unfold is \emph{checked}, not asserted to exist.
\item \emph{Installed and stamped \texttt{definitional}.}  The \textsc{Ring}
  laws are installed by \texttt{add-axiom!} in
  \texttt{structure-library/ring.scm} and then swept with
  \texttt{register-provenance!\ \ldots\ 'definitional} in the same file.  A
  \texttt{definitional} fact contributes the empty set to every bill, exactly as
  a \texttt{primitive} one does.
\item \emph{Wrapped at installation.}  \texttt{structure-library/module.scm}
  reaches the same place by the other door, installing its projections inside
  \texttt{(fluid-let ((*current-provenance* 'definitional)) \ldots)}.
\end{itemize}

Two laws in the table are genuinely \emph{derived} rather than projected, and
they do carry a warrant: \texttt{ring-mul-zero-left} and
\texttt{ring-mul-zero-right} ($0\cdot a = 0$, which needs distributivity and
cancellation, not merely the definition of a ring).
\texttt{reference/PROOF-DEBT.md} is the authority on which is which; this table
is not.  For example, to apply
\texttt{group-left-inv} (which asserts
$\forall s,a.\;\texttt{IS-GROUP}(s)\Rightarrow a\in\texttt{CARR}(s)
 \Rightarrow \texttt{OPR}(s)(\texttt{INV}(s)(a),a)=\texttt{IDEN}(s)$),
supply the group and the element as instantiation witnesses.
\end{remark}

If you prefer to work with raw tuples directly, the n-ary constructors
are still available:
\begin{VNBcode}
; Prove a tuple is in a Cartesian product
(sp (make-wff-from-string
  "G in power(U) and op in fun(cartesian(G,G),G) and e in G and inv in fun(G,G)
   implies [G, op, e, inv]
   in cartesian(power(U), fun(cartesian(G,G),G), G, fun(G,G))"))
(di)
(ai "G in power(U) and op in fun(cartesian(G,G),G) and e in G and inv in fun(G,G)")
(ai "op in fun(cartesian(G,G),G) and e in G and inv in fun(G,G)")
(ai "e in G and inv in fun(G,G)")
(ci)                              ; cartesian-intro: 4 subgoals
(ass) (ass) (ass) (ass)
\end{VNBcode}

\subsection{Applying a macete inside a proof}

The short-form command \texttt{(mac 'name)} applies the named macete to
the \emph{current goal}.  It unifies the macete's source pattern against
the goal formula (outermost-first, depth-first) and replaces the first
matched subterm with the replacement pattern.

\paragraph{Biconditional macete.}
\texttt{pairing-membership} states
$\forall a\,b\,x.\;x \in \mathrm{pair}(a,b) \Leftrightarrow (x=a \vee x=b)$.
Its source pattern is \verb|x in pair(a, b)| and its replacement is
\verb|x = a or x = b|.

\begin{VNBcode}
(sp (make-wff-from-string "2 in pair(2, 3)"))
; proof state:  []  |-  2 in pair(2, 3)
(mac 'pairing-membership)
; source: x in pair(a, b)   replacement: x = a or x = b
; unifies: x=2, a=2, b=3
; goal rewritten:  2 in pair(2, 3)  ->  2 = 2 or 2 = 3
; proof state:  []  |-  2 = 2 or 2 = 3
(oi-l)      ; or-intro-left: new goal  2 = 2
(rfl)       ; reflexivity
; => Proof complete.
\end{VNBcode}

\paragraph{Rewrite fires inside a larger formula.}
The rewrite finds its source pattern at the first matching position,
which may be nested inside a connective.

\begin{VNBcode}
(sp (make-wff-from-string "x in pair(a, b) implies x = a or x = b"))
; proof state:  []  |-  (x in pair(a,b)) implies (x = a or x = b)
(mac 'pairing-membership)
; pattern "x in pair(a, b)" matched at the left of implies
; goal rewritten:  (x = a or x = b) implies (x = a or x = b)
(di)        ; discharge implies: push antecedent to context
; context: ["x = a or x = b"]   goal: "x = a or x = b"
(ass)
; => Proof complete.
\end{VNBcode}

\begin{remark}
Macetes rewrite the goal; they do not rewrite the context.  If the
matching pattern appears only in an assumption, \texttt{(mac 'name)}
will not fire.  To use an assumption as a rewrite source, retrieve it
with \texttt{(ai ...)} so that it appears in the goal.
\end{remark}

\paragraph{Conditional macete.}
For a conditional theorem
$\forall\bar{x}.\;(\mathit{cond} \Rightarrow s = r)$, the macete still
fires unconditionally: it rewrites $s \mapsto r$ without generating
side-goals for $\mathit{cond}$.  The goal $s$ is replaced by~$r$
(or by $\top$ when the conclusion is not an equation or biconditional),
and the resulting $\top$ goal is closed with \texttt{(ass)}.

\begin{VNBcode}
; nn-add-closed:
;   forall([a in NN, b in NN], a + b in NN)
; As a macete: source = "a + b in nn"
;              replacement = truth
(sp (make-wff-from-string "2 + 3 in NN"))
; proof state:  []  |-  2 + 3 in NN
(mac 'nn-add-closed)
; unifies a=2, b=3
; conditions (a in NN, b in NN) not discharged
; goal rewritten:  2 + 3 in NN  ->  truth
(ass)       ; truth is trivially grounded
; => Proof complete.
\end{VNBcode}

\begin{remark}
The conditions of a conditional macete are silently dropped.  This
makes \texttt{(mac 'name)} fast for goals whose conditions are obviously
satisfied, but it is the user's responsibility to verify them; the
proof sketch above is technically unsound without separate proofs that
$2 \in \mathit{NN}$ and $3 \in \mathit{NN}$.  For fully verified
condition discharge, use \texttt{(ta 'name)} followed by explicit
instantiation and back-chaining.
\end{remark}

\subsection{Ground arithmetic with \texttt{(arith)}}

The \texttt{(arith)} command (primitive \texttt{cmd-arith}) closes any
goal whose formula reduces to true under exact arithmetic evaluation:

\begin{VNBcode}
(sp (make-wff-from-string "1/3 + 1/6 = 1/2"))
(arith)                              ; closes by exact rational arithmetic

(sp (make-wff-from-string "2 ^ 10 = 1024"))
(arith)                              ; 2^10 evaluates to 1024

(sp (make-wff-from-string "3 + 4 <= 2 * 5"))
(arith)                              ; 7 <= 10
\end{VNBcode}

\subsection{The calculator pattern}

\texttt{(arith)} also recognises a goal of the form
\begin{VNBsnip}
forsome([n in CLASS], EXPR = n)
\end{VNBsnip}
where \texttt{EXPR} is a ground arithmetic expression and \texttt{CLASS}
is one of \texttt{NN}, \texttt{ZZ}, \texttt{QQ}, \texttt{RR},
\texttt{CC}.  It evaluates \texttt{EXPR}, checks that the result belongs
to \texttt{CLASS}, and closes the goal with that value as the witness.
This lets the prover double as a calculator:

\begin{VNBcode}
(sp (make-wff-from-string "forsome([n in NN], 2^10 = n)"))
(arith)   ; closes with witness 1024

(sp (make-wff-from-string "forsome([q in QQ], 1/3 + 1/4 = q)"))
(arith)   ; closes with witness 7/12
\end{VNBcode}

\begin{remark}
The symmetric form \texttt{n = EXPR} is also accepted.
The pattern is intentionally narrow: \texttt{EXPR} must be fully ground
(no free variables), and the existential variable must appear alone on
one side of the equality.
\end{remark}

\subsection{\texorpdfstring{$\beta$}{beta}-reduction with \texttt{(beta)}}

Applied \texttt{lambda} (or \texttt{lambdoid}) terms are reduced
in-place, then the underlying numeric or logical content is closed
by another rule:

\begin{VNBcode}
(sp (make-wff-from-string "lambda([x in nn], 2 * x)(5) = 10"))
(beta)   ; goal becomes (2 * 5 = 10)
(arith)  ; closes
\end{VNBcode}

\subsection{Auto-installing arithmetic theorems}

\texttt{vnb-create-num-theorem} verifies a closed arithmetic wff by
ground evaluation and installs it as a named theorem.  The theorem is
then available via \texttt{(ta \textit{name})}:

\begin{VNBcode}
(vnb-create-num-theorem 'pythag-3-4-5 "3 ^ 2 + 4 ^ 2 = 5 ^ 2")
; => pythag-3-4-5

(sp (make-wff-from-string "truth implies 3 ^ 2 + 4 ^ 2 = 5 ^ 2"))
(di) (ta 'pythag-3-4-5) (ass)
; => Proof complete.
\end{VNBcode}

False or undecidable formulas are rejected:

\begin{VNBcode}
(vnb-create-num-theorem 'bogus "1 + 1 = 3")
; ERROR: formula evaluates to FALSE

(vnb-create-num-theorem 'bogus2 "x = 5")
; ERROR: not a decidable closed arithmetic sentence
\end{VNBcode}

\subsection{Installing a proven theorem}
\label{sec:qed-and-subst}

A completed proof does not by itself add anything to the library.  Use
\texttt{(qed \textit{name})} to install the root assertion as a named
theorem:

\begin{VNBcode}
(sp (make-wff-from-string "n in NN and m in NN implies m in NN and n in NN"))
(di) (ai "n in NN and m in NN") (di) (ass) (ass)
(qed 'and-comm)
; => and-comm
; (lookup-theorem 'and-comm) returns the formula S-expression.
\end{VNBcode}

\paragraph{What \texttt{subst-wff} is and is not.}
\texttt{subst-wff} applies capture-avoiding substitution to a wff's
underlying S-expression and returns a new \texttt{wff} (re-validated
through \texttt{make-wff}).  It is \emph{not} a logical inference;
after calling it you must reprove the result.  It is useful for
generating related goals from a known formula:

\begin{VNBcode}
(define w
  (subst-wff '((n . (in x nn)) (m . (in y rr)))
             (lookup-theorem 'and-comm)))
; w is the wff for (x in nn) and (y in rr) implies (y in rr) and (x in nn)
\end{VNBcode}

\subsection{Saving and replaying proof scripts}
\label{sec:replay}

Every interactive command you issue between \texttt{(sp\ \ldots)} and
\texttt{(qed\ \ldots)} is appended to a global accumulator
\texttt{*proof-script*}.  \texttt{(qed \textit{name})} installs the
proven theorem and saves the script under the same name; you can also
call \texttt{(save-proof \textit{name})} on demand.

\begin{VNBcode}
(sp (make-wff-from-string "n in NN and m in NN implies m in NN and n in NN"))
(di) (ai "n in NN and m in NN") (di) (ass) (ass)
(qed 'and-comm)
; theorem AND script installed under the same name

(lookup-proof 'and-comm)
; => ((di) (ai (and (in n nn) (in m nn))) (di) (ass) (ass))
\end{VNBcode}

\paragraph{Replay verbatim.}
\texttt{(replay-proof \textit{name})} re-runs the saved commands
sequentially against the current proof state.  Useful for redoing a
proof in a fresh session, or applying an identical script to a sequent
of the same shape.

%% =========================================================
\subsection{NN induction for type properties}
\label{sec:ex-ni-type}

The \texttt{(ni)} command reduces a goal
$\forall n \in \mathit{NN}.\; P(n)$ to two subgoals:
a base case $P(0)$ (focused first) and a step case
$\forall n \in \mathit{NN}.\; P(n) \Rightarrow P(\mathrm{succ}(n))$.
The following proof shows that $\mathit{PROD\text{-}ORD}(m, f, n)$
stays in the carrier $\mathit{CARR}(m)$ for every $n \in \mathit{NN}$,
deriving \texttt{prod-ord-type} from first principles.

\begin{VNBcode}
(sp (make-wff '(IMPLIES (IS-MONOID m)
               (IMPLIES (IN f (FUN NN (CARR m)))
                 (FORALL n (IMPLIES (IN n NN)
                   (IN (PROD-ORD m f n) (CARR m))))))))
(di) (di)   ; IS-MONOID(m) and f in FUN(NN,CARR(m)) enter context
(ni)        ; NN-induction; BASE gets focus first

;; --- BASE: PROD-ORD(m,f,0) in CARR(m) ---
;; Strategy: PROD-ORD(m,f,0) = IDEN(m) (prod-ord-zero),
;;           IDEN(m) in CARR(m)            (monoid-identity-in),
;;           so by eq-subst-membership the goal closes.
(ta 'eq-subst-membership)
(inst ... '(PROD-ORD m f 0)) (inst ... '(IDEN m)) (inst ... '(CARR m))
(bc '(IMPLIES (AND (= (PROD-ORD m f 0) (IDEN m)) (IN (IDEN m) (CARR m)))
              (IN (PROD-ORD m f 0) (CARR m))))
(di)    ; AND-split: focus = equality goal
(ta 'prod-ord-zero) (inst ... 'm) (inst ... 'f) (ass)
        ; focus jumps: manually refocus on membership goal
(ta 'monoid-identity-in) (inst ... 'm)
(bc '(IMPLIES (IS-MONOID m) (IN (IDEN m) (CARR m)))) (ass)
        ; BASE complete; focus moves to STEP

;; --- STEP: forall([n in NN], IH => PROD-ORD(m,f,succ n) in CARR(m)) ---
(di)    ; peel FORALL n -> n_k; (IN n_k NN) enters context
;; n_k is the bound name introduced by (di); fetch it via:
;;   (cadr (wff-formula
;;           (car (sequent-node-assumptions
;;                  (proof-state-focus *ps*)))))
(di)    ; peel IH: PROD-ORD(m,f,n_k) in CARR(m) enters context
;; Strategy:
;;   PROD-ORD(m,f,succ n_k)
;;     = OPR(m)(PROD-ORD(m,f,n_k), f(n_k))   (prod-ord-succ)
;;   then close via monoid-carrier-closed-opr.
(ta 'eq-subst-membership) (inst ...) (inst ...) (inst ...)
(bc ...)
(di)    ; AND-split: focus = equality goal
(ta 'prod-ord-succ)
(inst ...) (inst ...) (inst ...)
(bc ...) (ass)   ; n_k in NN closes it
;; focus moves to membership goal: OPR(m)(...) in CARR(m)
(ta 'monoid-carrier-closed-opr)
(inst ...) (inst ...) (inst ...)
(bc ...)
(di) (ass)   ; IS-MONOID(m)
(di) (ass)   ; PROD-ORD(m,f,n_k) in CARR(m) [IH]
;; f(n_k) in CARR(m) via fun-apply-type:
(ta 'fun-apply-type)
(inst ...) (inst ...) (inst ...) (inst ...)
(bc ...) (di) (ass) (ass)
;; closes f in FUN(NN,CARR(m)) and n_k in NN
\end{VNBcode}

\begin{remark}
The \texttt{inst} chains are elided above; each peels one leading
\texttt{forall} from a theorem schema.  Two subtleties arise in practice.
First, \texttt{(di)} on an AND-goal always focuses on the left (equality)
conjunct first.  After \texttt{(ass)} closes that conjunct, automatic
focus advancement jumps to the earliest ungrounded goal in the deduction
graph --- which may be the STEP node, not the right (membership) conjunct.
The fix is to capture the membership node
immediately after the AND-split:
\begin{VNBcode}
(di)   ; AND-split
(let* ((dg       (proof-state-dg *ps*))
       (mem-node (car (reverse (dg-sequent-nodes dg)))))
  ... prove equality ... (ass)   ; focus jumps
  (set-proof-state-focus! *ps* mem-node)
  ... prove membership ...)
\end{VNBcode}

Second, \texttt{alpha-equiv?} --- used by \texttt{(inst)} and
\texttt{(bc)} when searching assumptions --- relies on structural
equality of formula heads.  Compound heads such as
\texttt{((OPR m) a b)} (a curried two-argument application; see
Section~\ref{sec:apply-tupling}) are compared correctly only when
the head pair is the same \emph{object}.  If \texttt{alpha-equiv?}
fails to locate a formula you know is in context, the likely cause
is a compound-head mismatch; the fix is in
\texttt{expressions.scm}'s \texttt{alpha-equiv-under?}.
\end{remark}

\begin{remark}[The curried/tupled apply convention]
\label{sec:apply-tupling}
VNB's S-expression form \texttt{(f a$_1$ \ldots a$_n$)} is what the
string-syntax parser emits from \texttt{f(a$_1$,\,\ldots,\,a$_n$)} ---
verbatim, no implicit tupling.  An operation typed
\texttt{f $\in$ fun(cartesian(A$_1$,\,\ldots,\,A$_n$),~B)}, however,
takes one argument from the Cartesian product.  The two views are
reconciled by the library axioms
\begin{equation*}
\texttt{apply-tupling-$n$:} \qquad
(f\ a_1\ \ldots\ a_n)\ \mathrel{==}\ (f\ (\texttt{LIST}\ a_1\ \ldots\ a_n))
\quad\text{for each arity }n
\end{equation*}
(installed for $n=1,2,3$ in \texttt{theorem-library/axioms.scm};
add more arities by declaration as needed).  Quasi-equality
$\mathrel{==}$, not partial $=$: the convention holds vacuously when
either side is undefined.  Combined with \texttt{fun-apply-type} and
\texttt{quasi-eq-def}, these give the $n$-ary apply closure
\((f\ a_1\ \ldots\ a_n)\in B\) from the operation's signature and
the membership of each $a_i$ in its declared component.
\end{remark}

\subsection{Infinite descent: \texorpdfstring{$\sqrt3$}{sqrt3} is irrational}
\label{sec:ex-sqrt3}

Stated without reals or division --- the form the checker proves, entirely
inside $\mathit{NN}$ --- $\sqrt3$ irrational is
\[
  \forall p,q \in \mathit{NN}.\quad
  q \neq 0 \;\Longrightarrow\; p\cdot p \neq 3\cdot(q\cdot q) ,
\]
that is, the goal
\begin{VNBsnip}
(FORALL p (IMPLIES (IN p NN) (FORALL q (IMPLIES (IN q NN)
  (IMPLIES (NOT (= q 0)) (NOT (= (* p p) (* 3 (* q q)))))))))
\end{VNBsnip}

\paragraph{The argument.}
Infinite descent.  Suppose the equation has a solution with $q \neq 0$, and
among all such solutions choose one whose numerator $w$ is \emph{least}.  Since
$3 \mid w\cdot w$ we have $3 \mid w$, say $w = 3k$; substituting and cancelling
the factor $3$ gives $q_0\cdot q_0 = 3\cdot(k\cdot k)$, whence $3 \mid q_0$, say
$q_0 = 3m_0$; cancelling once more gives $k\cdot k = 3\cdot(m_0\cdot m_0)$.  Thus
$(k,m_0)$ is again a solution, with $k \neq 0$ and $k < 3k = w$ --- a numerator
strictly below the least one.  Contradiction.

The entire mathematical content is the single implication
\begin{equation}\label{eq:3divsq}
  3 \mid p\cdot p \;\Longrightarrow\; 3 \mid p ;
\end{equation}
everything else --- clearing denominators, the two cancellations, the order
step $k < 3k$ --- is bookkeeping.

\paragraph{Why parity is not special.}
The $\sqrt2$ proof is the same skeleton with \emph{even} for \emph{divisible by
$3$}: there \eqref{eq:3divsq} reads ``$p\cdot p$ even $\Rightarrow p$ even'' and
follows from parity, every natural being $2k$ or $2k+1$.  For $\sqrt3$ it follows
from the trichotomy --- every natural is $3k$, $3k+1$, or $3k+2$ --- together
with the fact that the three residues square to $0,1,1$ modulo $3$, so a square
is divisible by $3$ only when its root is.  Parity is the $r=2$ instance of a
uniform statement about residues modulo $r$; the $\sqrt3$ development is the
$\sqrt2$ development with the residue system widened by one, and nothing changes
in kind.

\paragraph{The lemma ladder.}
The mod-$3$ arithmetic (\texttt{theorem-library/nn-mod3-proof.scm}) is built by
induction from the two Peano recursion equations for $+$ and $\cdot$, in exact
parallel with the parity ladder (\texttt{nn-parity-proof.scm}).  Each lemma is
\emph{proven}; none is asserted.

\begin{center}
\begin{tabular}{lll}
\toprule
parity ($r=2$) & this proof ($r=3$) & statement \\
\midrule
\texttt{nn-two-mul-succ}     & \texttt{nn-three-mul-succ} & $r\cdot\mathrm{succ}(k) = \mathrm{succ}^{\,r}(r\cdot k)$ \\
\texttt{nn-parity}           & \texttt{nn-trichotomy-3}   & every $n$ is $r\cdot k + i$, residue $i<r$ \\
\texttt{nn-parity-exclusive} & \texttt{nn-3-not-succ}     & $r\cdot x$ is never $r\cdot y + 1$ \\
\texttt{nn-even-square}      & \texttt{nn-3-div-square}   & $r \mid p\cdot p \Rightarrow r \mid p$ \\
\bottomrule
\end{tabular}
\end{center}

\noindent For $r=3$ a single residue-inequality suffices: both nonzero residues
square to $r\cdot j + 1$, so only ``not $\equiv 1$'' is ever needed.

\paragraph{The descent, mechanized.}
The driver (\texttt{theorem-library/sqrt3-proof.scm}) states the goal and hands
the descent to three tools; the intermediate steps are elided.

\begin{VNBcode}
(sp (make-wff '(FORALL p (IMPLIES (IN p NN) (FORALL q (IMPLIES (IN q NN)
     (IMPLIES (NOT (= q 0)) (NOT (= (* p p) (* 3 (* q q)))))))))))
(di)(di)(di)(di)(di)(di)                 ; intro p,q, hyps; assume p*p = 3*(q*q)

(define R (minimize! '(pv) (guard 'pv) 'pv))    ; choose the LEAST numerator w
;; ... discharge minimize!'s two obligations: measure in NN; guard satisfiable

(define k (obtain (lambda () (fact 'nn-3-div-square w))))   ; 3|w  =>  w = 3k
;; ... q0*q0 = 3*(k*k) by have! ; 3|q0 => q0 = 3*m0 by obtain ; k*k = 3*(m0*m0)

(inst+ minf k)              ; minimality gives  w <= k
(fact 'nn-lt-triple k)      ; but k < 3k = w -- contradiction
\end{VNBcode}

\begin{itemize}
\item \texttt{minimize!} supplies the least-numerator witness together with its
  two honest obligations --- the measure lands in $\mathit{NN}$, and the guard
  is satisfiable --- and appeals to no choice principle.
\item \texttt{obtain} names each descended witness --- the $k$ with $w = 3k$,
  the $m_0$ with $q_0 = 3m_0$ --- by \emph{reading it off the proof state} (the
  variable left free once the existential is eliminated), never by guessing a
  machine-generated name.  This is what makes the descent survive a re-run.
\item \texttt{have!} discharges each intermediate equality and typing from
  context, so the script reads as the argument does: ``we have
  $q_0\cdot q_0 = 3\cdot(k\cdot k)$; we have $k \neq 0$; \ldots''.
\end{itemize}

\begin{remark}[The \texttt{minimize!} step]
\label{remark:minimize}
The enigmatic-looking \texttt{(minimize! '(pv) (guard 'pv) 'pv)} decodes as:
choose the single number \texttt{pv} satisfying the descent guard, with
\texttt{pv} itself as the quantity to minimise --- that is, take the
\emph{least} numerator meeting the guard.

The second argument to \texttt{minimize!} is the \textbf{guard formula} itself:
an ordinary VNB formula, free in the variables being chosen (here \texttt{pv}),
stating the property those values must satisfy.  \texttt{guard} is \emph{not} a
built-in and is nothing special to \texttt{minimize!} --- it is a helper the
proof author writes (named \texttt{s3-guard} in \texttt{sqrt3-proof.scm}), a
one-argument function that takes a term and returns the guard formula about it.
So \texttt{(guard 'pv)} yields the formula with \texttt{pv} free --- exactly what
\texttt{minimize!} wants --- and the \emph{same} helper, called later at the
chosen witness $w$ as \texttt{(guard $w$)}, restates that property of $w$.  Here
the property is ``\texttt{pv} is a natural number that is the numerator of a
fraction whose square is~$3$''; written out, the helper is just a formula
constructor:

\begin{VNBcode}
(define (s3-guard pv)
  (list 'AND (list 'IN pv 'NN)
        (list 'FORSOME 'q_
          (list 'AND (list 'IN 'q_ 'NN)
                (list 'AND (list 'NOT (list '= 'q_ 0))
                      (list '= (list '* pv pv) (list '* 3 (list '* 'q_ 'q_))))))))
\end{VNBcode}

Passing the formula literally would work just as well; the helper only spares the
author from writing it twice.  In general,
\[
  \texttt{(minimize! '($v_1$ \ldots\ $v_k$) GUARD MEASURE)}
\]
means ``choose $v_1,\ldots,v_k$ satisfying the formula \texttt{GUARD} with the
$\mathit{NN}$-valued term \texttt{MEASURE} as small as possible.''  On the main
branch it lands two hypotheses, for fresh eigenconstants $w_1,\ldots,w_k$: that
\texttt{GUARD} holds of them, and that nothing satisfying \texttt{GUARD} has a
strictly smaller \texttt{MEASURE}.  It leaves exactly the two obligations any
minimisation owes --- that \texttt{MEASURE} lands in $\mathit{NN}$ wherever
\texttt{GUARD} holds, and that \texttt{GUARD} holds somewhere.  Everything
between --- forming the value set, well-ordering it, unpacking the least
witness, restating minimality over the $v_i$ rather than the set --- it does for
you.  Every $v_i$ must occur in \texttt{MEASURE} (a variable the measure ignores
is not one you are minimising over).

It adds no kernel rule --- it chains \texttt{cut}, \texttt{di}, \texttt{ai},
\texttt{ew}, \ldots --- and uses \emph{no choice principle}: its one appeal is
\texttt{nn-least-element}, the well-ordering of $\mathit{NN}$, and well-ordering
returns a genuine \emph{member} of the value set, from which the witness is read
back.  Stated instead as a first-order lemma, the same fact would force every
call site to tuple its variables into a Cartesian product, build a lambda, and
prove that lambda inhabits $\mathit{FUN}(U,\mathit{NN})$ --- a typing obligation
larger than the minimisation it buys.  As a tactic it quantifies over the $v_i$
at the meta level and never forms $U$ at all.
\end{remark}

\begin{remark}[\texttt{obtain} versus \texttt{vlet}]
\texttt{obtain} and the \texttt{choice} former of \texttt{vlet} are the same
mechanism --- eliminate an existential and read its witness off by
free-variable difference --- but they differ in what they take.  \texttt{obtain}
takes a \emph{lane}: it runs a forward step and names the witness that step
\emph{just produced}.  That is the right fit inside a descent, where the relevant
existential is freshly landed and typically sits beside others: at the
``$3 \mid w$'' step the conclusion $\exists k.\,w = 3k$ coexists in context with
the evidence $\exists m.\,w\cdot w = 3m$ cut in to derive it, and \texttt{obtain}
selects the new one automatically.  Concretely, that one step reads
\begin{VNBsnip}
;; as written in the proof:
(define k (obtain (lambda () (fact 'nn-3-div-square w))))
;; the vlet alternative -- land the existential, then name it by its body:
(fact 'nn-3-div-square w)
(define k (vlet--choice! 'k '(AND (IN k NN) (= w (* 3 k)))))
\end{VNBsnip}
The \texttt{vlet} body names the eigenconstant $w$, a value produced at run
time, so this is the procedure \texttt{vlet--choice!}, not the
\texttt{(vlet \ldots)} macro (whose body is a fixed literal pattern).
\texttt{vlet} instead destructures a
\emph{named, stable} formula: its \texttt{match} former binds the holes of a
pattern (``the $a$ and $b$ in $x\cdot b = a$''), and its \texttt{choice} former
eliminates a \emph{designated} existential.  The rule of thumb: \texttt{obtain}
to name whatever a forward step produced, \texttt{vlet} to take apart a formula
you can already name.
\end{remark}

\begin{remark}
The proof adds no asserted step of its own: the mod-$3$ ladder is proven and the
descent cites it, exactly as \texttt{sqrt2-proof.scm} cites the parity ladder.
A new irrationality --- $\sqrt r$ for any non-square $r$ --- is then a matter of
widening the residue ladder, not of inventing an argument.  That division of
labour is the lesson: the author states the intermediate claims, and the machine
supplies the least witness (\texttt{minimize!}), the fresh names
(\texttt{obtain}), and the bookkeeping (\texttt{have!}, \texttt{from-context!}).
\end{remark}

\paragraph{The same proof at the interactive surface.}
By the transcript principle (Section~\ref{sec:two-surfaces}), most of this
script is \emph{already} a sequence of interactive short forms: \texttt{di},
\texttt{ai}, \texttt{ew}, \texttt{fact}, \texttt{cut}, \texttt{subst},
\texttt{crs}, \texttt{ass}, and \texttt{minimize!} are exactly what a user
types.  What the driver adds is the orchestration a human otherwise does by eye,
and it is of two kinds.  \emph{Focus}: after a branching command one selects the
next open leaf from the frontier, where the script calls \texttt{dk-focus!}.
\emph{Names}: each \texttt{(ai\ (FORSOME \ldots))} skolemizes to a machine-chosen
eigenconstant, printed in the resulting sequent, which the user reads and types
into the following commands.  The $3\mid w$ step, done by hand, is
\begin{VNBsnip}
(fact 'nn-3-div-square w)                          ; lands  exists k. w = 3k
(ai '(FORSOME k (AND (IN k NN) (= w (* 3 k)))))    ; skolemize; sequent now shows k_899
;; ... read k_899 off the sequent; type it into the subsequent (subst ...) / (fact ...)
\end{VNBsnip}
\texttt{obtain} is exactly this pair automated --- run the forward step, read the
fresh name back by free-variable difference --- so that a \emph{script}, unlike a
transcript with \texttt{k\_899} hard-coded, replays unchanged when the
eigenvariable counter lands on a different number.

%% =========================================================


%% =========================================================
\chapter{Mathematical content}

\section{Number systems} \label{sec:nums}

The file \texttt{number-systems.scm} declares five number systems as
class constants and installs their arithmetic axioms in the
library.  All declarations take effect at load time.

\subsection{The standard number domains}

\begin{itemize}
\item \texttt{nn} ($\mathbb{N}$) Non-negative integers $\{0, 1, 2, \ldots\}$
\item \texttt{zz} ($\mathbb{Z}$) Integers
\item \texttt{qq} ($\mathbb{Q}$) Rationals
\item \texttt{rr} ($\mathbb{R}$) Reals
\item \texttt{cc} ($\mathbb{C}$) Complex numbers
\end{itemize}

Sethood axioms (\texttt{nn-is-set} through \texttt{cc-is-set}) and the inclusion chain
\[
  \mathbb{N} \subseteq \mathbb{Z} \subseteq \mathbb{Q} \subseteq \mathbb{R} \subseteq \mathbb{C}
\]
\begin{sloppypar}are provided as axioms \texttt{nn-subset-zz}, \texttt{zz-subset-qq},
\texttt{qq-subset-rr}, \texttt{rr-subset-cc}.\end{sloppypar}

\subsection{Arithmetic operations}

Arithmetic operations are shared across all number systems: there is one
\texttt{+}, one \texttt{*}, one \texttt{-}, and so on.  Membership
determines which system an argument belongs to; closure axioms guarantee
the result stays in the appropriate system.  For example,
\texttt{forall([x in nn, y in nn], x + y in nn)} is the
closure axiom for \texttt{+} on \texttt{nn}.

\begin{center}
\begin{tabular}{lll}
\toprule
System & Operations & Distinguished constants \\
\midrule
\texttt{nn} &
  \texttt{+}, \texttt{*}, \texttt{succ}, \texttt{<=} &
  \texttt{0}, \texttt{1} \\[3pt]
\texttt{zz} &
  \texttt{+}, \texttt{*}, \texttt{-} (unary), \texttt{<=} &
  \texttt{0}, \texttt{1} \\[3pt]
\texttt{qq} &
  \texttt{+}, \texttt{*}, \texttt{-}, \texttt{recip}, \texttt{<=} &
  \texttt{0}, \texttt{1} \\[3pt]
\texttt{rr} &
  \texttt{+}, \texttt{*}, \texttt{-}, \texttt{recip}, \texttt{abs}, \texttt{<=} &
  \texttt{0}, \texttt{1} \\[3pt]
\texttt{cc} &
  \texttt{+}, \texttt{*}, \texttt{-}, \texttt{recip}, \texttt{conjugate} &
  \texttt{0}, \texttt{1} \\
\bottomrule
\end{tabular}
\end{center}

\begin{remark}\label{remark:27135-34942-842417}
%
The infix notation \texttt{x \^{} n} is the preferred way to write
exponentiation; it parses to the same internal form as
\texttt{power(x, n)}.  The typing axiom is the natural closure form,
matching the two-argument application syntax:
\[
  \forall x\,\forall n.\;
    (x \in \mathtt{cc} \land n \in \mathtt{nn})
    \;\Rightarrow\; \texttt{power}(x,n) \in \mathtt{cc}
  \qquad\text{(axiom \texttt{power-typing-nonneg})}
\]
The defining axioms are \texttt{power-zero} ($x^0 = 1$) and
\texttt{power-succ} ($x^{\mathrm{succ}(n)} = x \cdot x^n$ for
$n \in \mathtt{nn}$).  A conditional equation \texttt{power-neg}
covers negative integer exponents: $x^{-n} = \mathrm{recip}(x^n)$
for $n \in \mathtt{nn}$, $x \neq 0$; but this does not extend the
\texttt{fun} typing to $\mathtt{zz}$, since $0^{-1}$ is undefined.
Complex exponentiation $x^z$ for $z \in \mathtt{cc}$ will be defined
later via the power series for \texttt{exp}.
%
\end{remark}
%
\begin{remark}\label{remark:27135-34995-146430}
%
In particular, $0^0 = 1$.  This is the mathematically correct value
for combinatorial and power-series purposes, as Knuth argues, and it is
what MIT~Scheme's exact arithmetic returns.  GNU/Maxima raises an error
here.  IMPS treated $0^0$ as undefined, which is logically sound but
forces theorems such as the binomial theorem to carry an explicit
$x \neq 0$ or $y \neq 0$ hypothesis to exclude the special case;
this system avoids that overhead by defining $0^0 = 1$.
%
\end{remark}

The arithmetic decision procedure \texttt{vnb-do-arith} (and the
\texttt{(arith)} command) evaluates \texttt{x \^{} n} on ground
numeric arguments.

\noindent
Integer literals are constants directly: \texttt{0}, \texttt{1},
\texttt{7}, \texttt{-3}, etc.\ are elements of \texttt{zz} (and of
\texttt{qq}, \texttt{rr}, \texttt{cc} by inclusion); non-negative integer
literals are additionally elements of \texttt{nn}.  No subscripted
forms (\texttt{0\_ZZ}, \texttt{1\_RR}, etc.) are used.  In MIT~Scheme
numeric literals are first-class objects, so this representation is
natural.

Integer literals must be in canonical decimal form: no leading zeros
(other than \texttt{0} itself), so \texttt{007} and \texttt{000075} are
rejected by the parser.  \texttt{-0} normalises to \texttt{0}.

The ordinal operators \texttt{ord-le}, \texttt{ord-lt}, and
\texttt{succ\_ord} are named apart from the numeric \texttt{<=} and
\texttt{succ}.  The first two were written \texttt{<=\_ORD} and
\texttt{<\_ORD} until 2026-08-24; those spellings could not be tokenized
(\texttt{read-op} consumes the \texttt{<} and stops, and \texttt{\_}
cannot begin a token), so no formula mentioning them could be retyped
from its own printed form.  \texttt{succ\_ord} keeps its subscript: it
begins with a letter and reads back.

\subsection{Axioms installed}

Each number system comes with: a sethood axiom, inclusion from the
previous system, closure of each operation, and algebraic identities
appropriate to its structure.

\begin{itemize}
\item \texttt{nn}: closed under \texttt{+} and \texttt{*}; constants
  \texttt{0} (additive unit) and \texttt{1} (multiplicative unit);
  commutativity, associativity, and distributivity for \texttt{+} and
  \texttt{*}.  Unary \texttt{-} and \texttt{recip} are not total here.
\item \texttt{zz}: $\mathtt{nn}$ axioms plus closure under unary
  \texttt{-} (additive inverse); $x + (-x) = 0$.
\item \texttt{qq}: $\mathtt{zz}$ axioms plus \texttt{recip} on nonzero
  elements; ordering \texttt{<=} compatible with \texttt{+} and \texttt{*}.
\item \texttt{rr}: $\mathtt{qq}$ axioms plus \texttt{abs} satisfying the
  triangle inequality, with \texttt{abs($x$) = 0} iff \texttt{$x$ = 0}.
\item \texttt{cc}: $\mathtt{rr}$ axioms plus \texttt{conjugate}, an
  involution compatible with $+$ and $\cdot$, with
  $a \cdot \overline{a} \in \mathbb{R}$ and $a \cdot \overline{a} \geq 0$.
  No ordering on $\mathtt{cc}$.
\end{itemize}

\noindent The abstract structural names---commutative monoid,
commutative ring, ordered field---and the corresponding officialised
ring instances are deferred to Section~\ref{sec:basic-rings}.  At this
point we are simply listing the operations and identities each
domain carries.

Because there is a single \texttt{+}, \texttt{*}, etc.\ shared across all
systems, no compatibility axioms between adjacent systems are needed:
there is nothing to be compatible with.

\begin{remark}
Each of \texttt{nn}, \texttt{zz}, and \texttt{qq} carries an infinite
family of membership axioms, one per numeral appearing in the library.
When a positive-integer numeral $n$ is introduced, it is implicitly
axiomatized as a iterated sum $1 + 1 + \cdots + 1$ ($n$ times), so
$n \in \mathtt{nn}$.  For \texttt{zz}, negative numerals are similarly
axiomatized as the additive inverse of such sums, giving $-n \in
\mathtt{zz}$.  These axioms are generated on demand: entering a numeral
literal in a formula automatically extends the library with the
corresponding membership fact.
\end{remark}

\subsection{Example use in a proof}

\begin{VNBcode}
; Prove: 0 is in RR
(sp (make-wff-from-string "0 in RR"))
(ta 'rr-zero-in)                ; introduce theorem (0 in rr)
(ass)                           ; goal is now in assumptions

; Prove: for reals a, b, (a + b) is in RR
(sp (make-wff-from-string "a in RR and b in RR implies a + b in RR"))
(di)                              ; assume (a in rr) and (b in rr)
(ai "a in RR and b in RR")       ; split into (a in rr) and (b in rr)
(ta 'rr-add-closed)               ; closure axiom: forall([a, b], ...)
(inst "forall([a, b], a in rr and b in rr implies a + b in rr)" 'a)
(inst "forall([b], a in rr and b in rr implies a + b in rr)" 'b)
(bc "a in rr and b in rr implies a + b in rr")
(ass)
; => Proof complete.
\end{VNBcode}

%% =========================================================
\section{Ordinals}
\label{sec:ordinals}

\subsection{The class \texttt{ord}}

\texttt{ord} is a primitive class constant, declared in the same style as
\texttt{nn}, \texttt{rr}, and \texttt{cc}: no explicit construction of
ordinals is needed.  The class is characterised entirely by its axioms.

\begin{itemize}
\item \texttt{ord} the class of all ordinals (a proper class)
\item \texttt{$\alpha$ in ord} $\alpha$ is an ordinal
\end{itemize}

\subsection{The ordering predicate}

Ordinals are equipped with a primitive ordering predicate \texttt{ord-le}
(and its strict companion \texttt{ord-lt}) that is independent of set
membership.  The ordering has no required connection to $\in$.

\begin{itemize}
\item \texttt{ord-le($\alpha$, $\beta$)} \quad $\alpha \leq \beta$ in the ordinal ordering
\item \texttt{ord-lt($\alpha$, $\beta$)} \quad $\alpha < \beta$ (strict)
\item \texttt{succ\_ord($\alpha$)} \quad successor ordinal of $\alpha$
\end{itemize}

\subsection{Axioms}

The following axioms are installed in \texttt{ordinals.scm}.

\begin{itemize}
\item \texttt{burali-forti}
  $\mathtt{ord} \notin \mathtt{set}$: \texttt{ord} is a proper class.
  Every \emph{set} of ordinals has a supremum in \texttt{ord}
  (axiom \texttt{sup-ord-in}); the supremum need not lie in the set itself
  (e.g.\ $\sup \mathtt{nn} = \omega \notin \mathtt{nn}$).
  If \texttt{ord} were a set, its supremum $\alpha$ would be an ordinal,
  hence $\mathit{succ\_ord}(\alpha) \in \mathtt{ord}$ would strictly
  exceed $\alpha$---contradicting that $\alpha$ bounds \texttt{ord}.
\item \texttt{ord-le-closure}
  $\alpha \leq_O \beta \;\Rightarrow\; \alpha \in \mathtt{ord} \land \beta \in \mathtt{ord}$
\item \texttt{ord-le-refl}
  $\alpha \in \mathtt{ord} \;\Rightarrow\; \alpha \leq_O \alpha$
\item \texttt{ord-le-antisymm}
  $\alpha \leq_O \beta \land \beta \leq_O \alpha \;\Rightarrow\; \alpha = \beta$
\item \texttt{ord-le-trans}
  $\alpha \leq_O \beta \land \beta \leq_O \gamma \;\Rightarrow\; \alpha \leq_O \gamma$
\item \texttt{ord-le-total}
  $\alpha,\beta \in \mathtt{ord} \;\Rightarrow\; \alpha \leq_O \beta \lor \beta \leq_O \alpha$
\item \texttt{ord-zero-least}
  $\alpha \in \mathtt{ord} \;\Rightarrow\; 0 \leq_O \alpha$
\item \texttt{ord-lt-iff}
  $\alpha <_O \beta \;\Leftrightarrow\; \alpha \leq_O \beta \land \alpha \neq \beta$
\item \texttt{ord-succ-in}
  $\alpha \in \mathtt{ord} \;\Rightarrow\; \mathit{succ\_ord}(\alpha) \in \mathtt{ord}$
\item \texttt{ord-succ-above}
  $\alpha \in \mathtt{ord} \;\Rightarrow\; \alpha <_O \mathit{succ\_ord}(\alpha)$
\item \texttt{ord-succ-immediate}
  $\alpha,\beta \in \mathtt{ord},\; \alpha <_O \beta \;\Rightarrow\; \mathit{succ\_ord}(\alpha) \leq_O \beta$
\item \texttt{ord-succ-injective}
  $\mathit{succ\_ord}(\alpha) = \mathit{succ\_ord}(\beta) \;\Rightarrow\; \alpha = \beta$
\item \begin{sloppypar}\texttt{limit-ord-iff}
  $\mathit{limit\text{-}ord}(\lambda) \;\Leftrightarrow\; \lambda \in \mathtt{ord} \land \lambda \neq 0 \land \lnot\exists\alpha.\,\lambda = \mathit{succ\_ord}(\alpha)$\end{sloppypar}
\item \texttt{ord-segment-is-set}
  $\alpha \in \mathtt{ord} \;\Rightarrow\; \mathit{ord\text{-}segment}(\alpha) \in \mathtt{set}$
\item \texttt{ord-segment-membership}
  $x \in \mathit{ord\text{-}segment}(\alpha) \;\Leftrightarrow\; x <_O \alpha$
\item \texttt{ord-segment-zero}
  $\mathit{ord\text{-}segment}(0) = \emptyset$
\item \texttt{ord-segment-succ}
  \[x \in \mathit{ord\text{-}segment}(\mathit{succ\_ord}(\alpha))
  \;\Leftrightarrow\; x \in \mathit{ord\text{-}segment}(\alpha) \lor x = \alpha\]
\item \texttt{sup-ord-in}
  $A \in \mathtt{set},\; A \subseteq \mathtt{ord} \;\Rightarrow\; \mathit{sup\text{-}ord}(A) \in \mathtt{ord}$
\item \texttt{sup-ord-upper}
  $\mathit{sup\text{-}ord}(A)$ is an upper bound of $A$
\item \texttt{sup-ord-least}
  $\mathit{sup\text{-}ord}(A)$ is the least upper bound
\item \texttt{sup-ord-empty}
  $\mathit{sup\text{-}ord}(\emptyset) = 0$
\item \texttt{transfinite-induction}
  see \S\ref{sec:ord-induction}
\end{itemize}

Here $\leq_O$ and $<_O$ abbreviate \texttt{ord-le} and \texttt{ord-lt}.

\subsection{Relationship to number systems}

\texttt{nn} is a subset of \texttt{ord} (axiom \texttt{nn-subset-ord}):
\[
  \forall n.\; n \in \mathtt{nn} \;\Rightarrow\; n \in \mathtt{ord}
\]
Consequently $0 \in \mathtt{ord}$ (from \texttt{nn-zero-in} and
\texttt{nn-subset-ord}).  The ordinal successor agrees with the numeric
successor on $\mathtt{nn}$ (axiom \texttt{ord-succ-nn}), and
\texttt{ord-le} restricted to $\mathtt{nn}$ coincides with the numeric
\texttt{<=} (axiom \texttt{ord-le-nn-compat}).  The number systems
\texttt{zz}, \texttt{qq}, \texttt{rr}, \texttt{cc} are not subsets of
\texttt{ord}.

\subsection{Transfinite induction}
\label{sec:ord-induction}

\paragraph{The induction axiom.}

\begin{sloppypar}The axiom \texttt{transfinite-induction} is the \emph{complete} (strong)
form: a class $C$ that contains every ordinal whose strict predecessors are
all in $C$ must contain every ordinal.\end{sloppypar}
\[
  \forall C.\;
  \Bigl(\forall\alpha \in \mathtt{ord}.\;
        (\forall\beta <_O \alpha.\; \beta \in C) \Rightarrow \alpha \in C\Bigr)
  \;\Rightarrow\;
  \forall\alpha \in \mathtt{ord}.\; \alpha \in C
\]

\paragraph{The \texttt{(tfi)} proof command.}

\texttt{(tfi)} applies transfinite induction to the current focus goal.
The goal must have the form
\[
  \forall \alpha.\; \alpha \in \mathtt{ord} \;\Rightarrow\; P(\alpha)
\]
and \texttt{(tfi)} replaces it with the single inductive-step subgoal
\[
  \forall \alpha.\;
    \bigl(\alpha \in \mathtt{ord} \;\land\;
          \forall\beta <_O \alpha.\; P(\beta)\bigr)
    \;\Rightarrow\; P(\alpha)
\]

\paragraph{The \texttt{(tfi3)} proof command.}

\texttt{(tfi3)} splits the same goal into three subgoals matching the
standard Peano-style cases:
\begin{enumerate}
\item \textbf{Base:} $P(0)$.
\item \textbf{Successor:} $\forall\alpha \in \mathtt{ord}.\; P(\alpha) \Rightarrow P(\mathit{succ\_ord}(\alpha))$.
\item \textbf{Limit:} $\forall\lambda.\;
  \mathit{limit\text{-}ord}(\lambda) \land
  (\forall\beta <_O \lambda.\; P(\beta)) \;\Rightarrow\; P(\lambda)$.
\end{enumerate}

\subsection{Transfinite recursion}
\label{sec:ord-recursion}

The theoretical basis for transfinite recursive definitions is the
following pair of propositions.  Both are proved by the standard
least-counterexample argument using well-foundedness of $\leq_O$.

\begin{sloppypar}
\begin{prop}[Transfinite recursion, sets]\label{prop:tfrec-set}
Let $U = \mathit{ord\text{-}segment}(\alpha)$ for some $\alpha \in \mathtt{ord}$,
let $A \in \mathtt{set}$, and let
\[
  \mathcal{F}(U, A) \;=\; \bigcup_{x \in U}\,\mathtt{fun}(\mathit{ord\text{-}segment}(x),\, A)
\]
be the set of all functions defined on an initial segment within $U$ with
values in $A$.  For any $\mathcal{E} \in \mathtt{fun}(\mathcal{F}(U,A),\, A)$ there
is a unique $f \in \mathtt{fun}(U, A)$ such that
\[
  f(x) \;=\; \mathcal{E}\bigl(f \!\restriction\! \mathit{ord\text{-}segment}(x)\bigr)
  \quad \text{for every } x \in U.
\]
\end{prop}
\end{sloppypar}

\begin{prop}[Transfinite recursion, class version]\label{prop:tfrec-class}
Let $A$ be a class, and let $\mathcal{E}$ be a functoid from the class of
all functions defined on initial ordinal segments with values in $A$ to $A$
(a proper-class domain, so $\mathcal{E}$ cannot be an element of
$\mathtt{fun}$).
Then there is a unique functoid $f$ on $\mathtt{ord}$ with values in $A$ such that
\[
  f(\xi) \;=\; \mathcal{E}\bigl(f \!\restriction\! \mathit{ord\text{-}segment}(\xi)\bigr)
  \quad \text{for every } \xi \in \mathtt{ord}.
\]
\end{prop}

\begin{remark}
Both $\mathcal{E}$ and $f$ in Proposition~\ref{prop:tfrec-class} are functoids, not
elements of $\mathtt{fun}(A,B)$: their domains are proper classes ($\mathtt{ord}$
for $f$; the class of all initial-segment functions for $\mathcal{E}$), and
$\mathtt{fun}$ requires set-valued domains.  Each is nonetheless a well-defined
class of ordered pairs with the function property.
\end{remark}

The Scheme helper \texttt{def-by-ord-recursion} implements this pattern;
see Section~\ref{sec:recursive-defs}.

%% =========================================================
\section{The base library}
\label{sec:axioms}

The base library \texttt{VNB-LIBRARY} is loaded automatically and
provides the following named axioms (all usable via
\texttt{cmd-theorem-assumption}):

\begin{itemize}
\item \texttt{membership-implies-sethood}
  $\forall a\,\forall b.\;(a \in b) \Rightarrow a \in \mathtt{set}$
\item \texttt{extensionality}
  $\forall a\,\forall b.\;(a \in \mathtt{set} \land b \in \mathtt{set}) \Rightarrow
  (a = b \Leftrightarrow \forall x.\,(x \in a \Leftrightarrow x \in b))$
\item \texttt{empty-set-is-set}
  $\texttt{empty-set} \in \mathtt{set}$
\item \texttt{empty-set-has-no-members}
  $\forall x.\;\lnot(x \in \texttt{empty-set})$
\item \texttt{pairing}
  $\forall a\,\forall b.\;(a \in \mathtt{set} \land b \in \mathtt{set}) \Rightarrow \{a,b\} \in \mathtt{set}$
\item \texttt{pairing-membership}
  $\forall a\,\forall b.\;(a \in \mathtt{set} \land b \in \mathtt{set}) \Rightarrow
  \forall x.\;(x \in \{a,b\}) \Leftrightarrow (x=a \lor x=b)$
\item \texttt{power-set}
  $\forall a.\;a \in \mathtt{set} \Rightarrow \texttt{power}(a) \in \mathtt{set}$
\item \texttt{power-set-membership}
  $\forall a\,\forall x.\;(x \in \texttt{power}(a)) \Leftrightarrow
  (x \in \mathtt{set} \land \forall z.\,(z \in x \Rightarrow z \in a))$
\end{itemize}

A theorem is added by proving it.  State it with \texttt{sp}, prove it,
and install it with \texttt{qed}, which also installs its macete:
\begin{VNBcode}
(sp "forall([x in nn], x in nn)")
(di) (ass)
(qed 'my-thm)
\end{VNBcode}
\texttt{qed} refuses unless the proof is complete, and it is the only
command that records a result as proven.  A statement that is to be
assumed for the time being is added with \texttt{support} and given a
stated justification with \texttt{warrant!}
(Chapter~\ref{ch:classification}).  It is then recorded as asserted, and
every theorem that depends on it says so when it is installed.


%% Definitions moved to ch-defs.tex


%% =========================================================
\section{Pending axioms}
\label{sec:pending-axioms}

The axioms in this section are needed for full set-theoretic reasoning
but are not yet installed.  They fall into two groups.

\subsection*{Installed in \texttt{axioms.scm}}

The following were previously pending but are now installed.
They are straightforward first-order statements requiring no new
expression primitives.

\begin{itemize}

\item \texttt{equality-symmetry}
  $a = b \Rightarrow b = a$.

\item \texttt{equality-transitivity}
  $a = b \land b = c \Rightarrow a = c$.

\item \texttt{fun-apply-type}
  $(f \in \mathtt{fun}(A,B)) \land (x \in A) \Rightarrow f(x) \in B$.
  Derivable from \texttt{fun-codomain-iff} but named for direct use.

\item \texttt{list-sethood}
  $A \in \mathtt{set} \land L \in \mathtt{tuples}(A) \Rightarrow L \in \mathtt{set}$.
  Derivable from \texttt{subset-set} (elements of any class are sets)
  but named for direct use in structure instance proofs.

\item \texttt{eq-subst-membership}
  $a = b \land b \in S \Rightarrow a \in S$.
  A first-order instance of \texttt{equality-substitution} restricted to
  membership; derivable from \texttt{equality-transitivity} and
  extensionality, but named for direct use in type proofs.

\item \texttt{quasi-eq-reflexivity}\ $(a \mathbin{==} a)$.

\item \texttt{quasi-eq-symmetry}\ $(a \mathbin{==} b) \Rightarrow (b \mathbin{==} a)$.

\item \texttt{quasi-eq-transitivity}\
  $(a \mathbin{==} b) \land (b \mathbin{==} c) \Rightarrow (a \mathbin{==} c)$.

\item \texttt{quasi-eq-def}\
  $(a = a) \Rightarrow ((a \mathbin{==} b) \Leftrightarrow (a = b))$;
  when $a$ is defined, quasi-equality coincides with strict equality.

\end{itemize}

The following were already installed (noted here for completeness):
\texttt{fun-domain-apply-def}, \texttt{fun-domain-extensionality},
\texttt{fun-codomain-iff} and \texttt{choice-axiom} in \texttt{library.scm};
\texttt{nn-subset-ord} and all \texttt{ord-*} axioms in
\texttt{structure-library/ordinals.scm}, not in \texttt{library.scm}
(see Section~\ref{sec:ordinals}).  The ordinal axioms carry
\texttt{primitive} provenance and so contribute nothing to any proof's bill.

\subsection*{List recursion and list induction}
\label{sec:list-recursion}

Both were listed as pending in earlier editions of this manual, and both are
now installed.  \texttt{structure-library/list-recursion.scm} adds the
constructor \texttt{cons}$(x, L)$ --- the name is \texttt{cons}, not
\texttt{prepend} --- with four defining equations, stamped
\texttt{definitional} because they constrain a new head and so add no
consequence about the old vocabulary:
\[
  \texttt{length}(\texttt{cons}(x,L)) = \texttt{succ}(\texttt{length}(L)),
  \qquad
  \texttt{nth}(1, \texttt{cons}(x,L)) = x,
\]
with \texttt{nth-cons-succ} shifting the remaining indices and
\texttt{makeset-cons} relating \texttt{make-set} to a union with a singleton.
Alongside them are two \emph{generation} axioms, which are claims about
\texttt{tuples} rather than about \texttt{cons} and are therefore asserted and
warranted: \texttt{tuple-length-zero} (a tuple of length $0$ is the empty
tuple) and \texttt{tuple-cons-decompose} (a tuple of length $\texttt{succ}(n)$
is $x$ prepended to a tuple of length $n$).

The induction principle is then a \emph{theorem}, not a schema-level primitive
inference: \texttt{theorem-library/tuples-induction.scm} proves
\[
  \forall A, C.\;
  [\,]\in C \;\wedge\;
  (\forall x, M.\; x \in A \wedge M \in \mathtt{tuples}(A) \wedge M \in C
     \Rightarrow \texttt{cons}(x,M) \in C)
  \;\Rightarrow\;
  \forall L \in \mathtt{tuples}(A).\; L \in C
\]
by induction on the \emph{length}, using \texttt{nn-induction} and the two
generation axioms.  It is stated in class form, quantifying over a class $C$
with membership in it as the conclusion, exactly as \texttt{nn-induction} and
\texttt{finite-set-induction} are; VNB quantifies over classes, so one theorem
does the work a schema would do in a first-order presentation.  Its bill reads
\texttt{modulo \{tuple-length-zero, tuple-cons-decompose\}}, which is the honest
report: the induction is derived, and what it rests on is the statement that
tuples are finite.

\subsection*{Still pending}

The remaining items require schema-level infrastructure --- a primitive
inference rather than a first-order axiom.  (An earlier edition also listed
\texttt{separation}, \texttt{union-set}, \texttt{comp} and \texttt{prepend}
as missing expression primitives.  They are not missing: \texttt{sep},
\texttt{comp} and \texttt{big-union} are registered term-forming heads
(\texttt{wff.scm}) with intro/elim rules in \texttt{primitive-inferences.scm},
and they are documented in \texttt{reference/KERNEL-RULES.md}.)

\begin{itemize}

\item \texttt{equality-substitution}
  $a = b \land P(a) \Rightarrow P(b)$ (for any formula $P$).
  Requires a substitution schema; cannot be stated as a single
  first-order sentence.

\item \texttt{separation}
  $A \in \mathtt{set} \Rightarrow \{x \in A \mid p\} \in \mathtt{set}$.
  Requires the \texttt{separation} expression primitive.

\item \texttt{union-set} and \texttt{union-set-membership}
  $A \in \mathtt{set} \Rightarrow \bigcup A \in \mathtt{set}$, and
  $x \in \bigcup A \Leftrightarrow \exists a \in A.\; x \in a$.
  Requires the \texttt{union-set} (big union) expression primitive,
  distinct from the binary \texttt{union} constructor.

\item \texttt{cartesian-n-ary-sethood}
  All $A_i \in \mathtt{set} \Rightarrow A_1 \times \cdots \times A_n \in \mathtt{set}$.
  The binary case is installed (\texttt{cartesian-set-iff}); the $n$-ary
  case reduces to binary by associativity but requires explicit axioms.

\item \texttt{comp-membership}
  $x \in \{y \mid p\} \Leftrightarrow p[y \mapsto x]$.
  Requires the \texttt{comp} (class comprehension) expression primitive.

\item \texttt{lambda-type}
  $(\forall x_i \in A_i.\;\mathit{body} \in B) \Rightarrow
   \lambda([x_1,\ldots], \mathit{body}) \in \mathtt{fun}(A_1{\times}\cdots{\times}A_n, B)$.
  Requires a schema-level primitive inference (\texttt{pi-lambda-type}).

\item \texttt{lambda-beta}
  $\lambda([x_1,\ldots], \mathit{body})(v_1,\ldots,v_n)
   = \mathit{body}[x_i \mapsto v_i]$.
  Requires a beta-reduction primitive inference.

\item \texttt{iota-def}
  $(\exists!\, x.\;p(x)) \Rightarrow p(\texttt{iota}(x,p))$.
  Requires a schema-level primitive inference.

\item \texttt{infinity}
  There exists an inductive set:
  $\exists A \in \mathtt{set}.\;\emptyset \in A \land
   \forall x \in A.\; x \cup \{x\} \in A$.

\end{itemize}

%% =========================================================


%% =========================================================
\chapter{Definitions}

The forms described in this chapter are \emph{Scheme procedures} called
at the meta-level --- in a \texttt{.scm} source file or at the Scheme
REPL --- not expressions in the VNB object language.  Each call installs
axioms or macetes into the current theory as a side effect; nothing is
proved.  The distinction between meta-level definition and object-level
proof bears repeating: \texttt{declare-structure}, \texttt{def-constant},
and their companions manipulate the theory data structure directly, while
proofs are constructed separately using the proof commands of
Chapter~\ref{ch:proofs}.

%% =========================================================
\section{Constants and predicates}
\label{sec:def-constant}

A \emph{defined constant} is a new constant symbol introduced by a
collection of characterizing axioms.  It differs from a plain
\texttt{theory-add-axiom!} call only in bookkeeping: the axioms are
recorded under a separate \texttt{definitions} slot in the theory record,
so the system can distinguish what was \emph{postulated} from what was
\emph{introduced by definition}.

\begin{VNBcode}
(def-constant 'NAME
  '(axiom-name-1 formula-1)
  '(axiom-name-2 formula-2)
  ...)
\end{VNBcode}

\begin{sloppypar}Each extra argument is a two-element list \texttt{(axiom-name formula)}.
The formulas are installed in the global theorem and macete tables exactly
as \texttt{theory-add-axiom!} would install them, so they are immediately
available to proof commands and macetes.\end{sloppypar}

\textbf{Existence and uniqueness are the user's responsibility.}  The
system accepts the definition without checking that the described object
exists or is unique.  For constants where this matters (e.g., a least
element), you should prove existence and uniqueness before invoking
\texttt{def-constant}, or at least note the obligation explicitly.

Example---introducing the first uncountable ordinal $\omega_1$:
\begin{VNBcode}
(def-constant 'OMEGA1
  (list 'omega1-is-ordinal
        (parse-string "OMEGA1 in ORD"))
  (list 'omega1-is-uncountable
        (parse-string "not(countable(OMEGA1))"))
  (list 'omega1-is-least-unc
        (parse-string "forall([alpha in ORD],
          not(countable(alpha)) implies OMEGA1 <= alpha)")))
\end{VNBcode}

To display all defined constants and their characterizing axioms:
\begin{VNBcode}
(display-definitions)
\end{VNBcode}

%% =========================================================
\section{Predicates}
\label{sec:def-predicate}

An $n$-ary \emph{defined predicate} is introduced by \texttt{def-predicate}.
It installs a single characterising axiom --- an \texttt{iff} --- that
unfolds the predicate in terms of its arguments.

\begin{VNBcode}
(def-predicate 'PRED-NAME '(p1 p2 ...)  body)
\end{VNBcode}

Installs one axiom named \texttt{PRED-NAME}:
\begin{VNBsnip}
forall([p1, p2, ...], PRED-NAME(p1, p2, ...) iff body)
\end{VNBsnip}

The axiom is immediately available via \texttt{(ta 'PRED-NAME)} in a proof.

\paragraph{Example: Cauchy sequence in a metric space.}
\begin{VNBcode}
(def-predicate 'IS-CAUCHY-SEQ '(X f)
  '(AND (IS-METRIC-SPACE X)
        (IN f (FUN NN (CARRIER X)))
        (FORALL eps
          (IMPLIES (AND (IN eps RR) (< 0 eps))
            (FORSOME N (IN N NN)
              (FORALL m (IMPLIES (IN m NN)
                (FORALL n (IMPLIES (IN n NN)
                  (IMPLIES (AND (<= N m) (<= N n))
                    (< ((D X) (f m) (f n)) eps)))))))))))
\end{VNBcode}

This installs:
\begin{VNBsnip}
forall([X, f],
  IS-CAUCHY-SEQ(X, f) iff
    IS-METRIC-SPACE(X)
    and f in fun(NN, CARRIER(X))
    and forall([eps in RR], 0 < eps implies
          forsome([N in NN],
            forall([m in NN, n in NN], N <= m and N <= n implies
              D(X)(f(m), f(n)) < eps)))
\end{VNBsnip}

%% =========================================================
\section{Recursive definitions}
\label{sec:recursive-defs}

Two Scheme helpers introduce function constants by recursion and install
the defining equations as axioms.  Both rely on the same
\texttt{theory-add-definition!} mechanism as \texttt{def-constant}
(Section~\ref{sec:def-constant}), so every axiom is immediately
available in proofs.

\begin{prop}[Primitive recursion on $\mathit{NN}$]\label{prop:nn-rec}
Let $A \in \mathit{SET}$ and let
\[
  \mathcal{F}(A) \;=\; \bigcup_{n \in \mathit{NN}}\,\mathit{FUN}(\mathit{ORD\text{-}SEGMENT}(n),\, A)
\]
be the set of all functions defined on a finite initial segment of $\mathit{NN}$
with values in $A$ (including the empty function at $n = 0$, since
$\mathit{ORD\text{-}SEGMENT}(0) = \emptyset$).
For any $\mathcal{E} \in \mathit{FUN}(\mathcal{F}(A),\, A)$ there is a unique
$f \in \mathit{FUN}(\mathit{NN},\, A)$ such that
\[
  f(n) \;=\; \mathcal{E}\!\bigl(f \!\restriction\! \mathit{ORD\text{-}SEGMENT}(n)\bigr)
  \qquad \text{for every } n \in \mathit{NN}.
\]
\end{prop}

\begin{remark}
This is the special case of Proposition~\ref{prop:tfrec-set} with
$U = \mathit{NN}$; it follows from the well-ordering of $\mathit{NN}$ by the
standard least-counterexample argument.
Since $\mathit{ORD\text{-}SEGMENT}(n)$ is the integer interval $\{0, 1, \ldots, n-1\}$,
the extension function $\mathcal{E}$ sees the entire history
$f(0), f(1), \ldots, f(n-1)$ when computing $f(n)$.
The \emph{predecessor form} --- $f(0) = b$ and
$f(\mathrm{succ}(n)) = s(n,\, f(n))$ for $n \in \mathit{NN}$ --- is the special case
where $\mathcal{E}$ only inspects the last value; this covers sums, products,
and most concrete recursive definitions.
\end{remark}

\subsection{\texttt{def-by-nn-recursion}: primitive recursion on \texttt{NN}}

Defines a function $F(p_1, \ldots, p_k, n)$ by two-clause recursion on the
last argument $n \in \mathit{NN}$, with arbitrary extra parameters
$p_1, \ldots, p_k$:
\[
  F(p_1,\ldots,p_k,0) = b, \qquad
  F(p_1,\ldots,p_k,\mathrm{succ}(n)) = s(p_1,\ldots,p_k,n,F(p_1,\ldots,p_k,n)).
\]

\begin{VNBcode}
(def-by-nn-recursion 'F
  '(p1 p2 ...)      ; extra parameters ('() for a plain NN->A function)
  base-val          ; F(p1,p2,..., 0) = base-val
  '(n val)          ; step variables
  succ-expr)        ; F(p1,p2,..., succ(n)) = succ-expr
                    ;   (val stands for F(p1,p2,...,n))
\end{VNBcode}

Two axioms are installed:
\begin{VNBsnip}
F-zero: forall([p1,p2,...], F(p1,p2,...,0) = base-val)
F-succ: forall([p1,p2,..., n in NN],
          F(p1,p2,...,succ(n)) = succ-expr[val:=F(p1,p2,...,n)])
\end{VNBsnip}

\paragraph{Example: factorial (no parameters).}
\begin{VNBcode}
(def-by-nn-recursion 'FACT '()
  1                           ; FACT(0) = 1
  '(n val) '(* (succ n) val)) ; FACT(succ(n)) = succ(n) * FACT(n)
\end{VNBcode}

\paragraph{Example: power of 2 (no parameters).}
\begin{VNBcode}
(def-by-nn-recursion 'POW2 '()
  1                       ; POW2(0) = 1
  '(n val) '(* 2 val))   ; POW2(succ(n)) = 2 * POW2(n)
\end{VNBcode}

\paragraph{Example: partial sum (one parameter).}
\begin{sloppypar}For $f \in \mathit{FUN}(\mathit{NN}, A)$, define
$\mathrm{SUM}(f, n) = \sum_{k=0}^{n} f(k)$:\end{sloppypar}
\begin{VNBcode}
(def-by-nn-recursion 'SUM '(f)
  '(f 0)                          ; SUM(f, 0) = f(0)
  '(n val) '(+ val (f (succ n)))) ; SUM(f, succ(n)) = SUM(f,n) + f(succ(n))

; Installs SUM-zero: forall([f], SUM(f,0) = f(0))
; Installs SUM-succ: forall([f, n in NN], SUM(f,succ(n)) = SUM(f,n) + f(succ(n)))
\end{VNBcode}

\subsection{\texttt{def-by-ord-recursion}: transfinite recursion on \texttt{ORD}}

For functions on all ordinals, use the three-clause schema
(Proposition~\ref{prop:tfrec-class}, Section~\ref{sec:ord-recursion}):
\begin{itemize}[label=$\circ$, leftmargin=0.9cm]
\item \textbf{Base} ($\xi = 0$):\quad $F(0) = b$
\item \textbf{Successor} ($\xi = \mathrm{succ\_ORD}(\alpha)$):\quad
      $F(\mathrm{succ\_ORD}(\alpha)) = s(\alpha,\, F(\alpha))$
\item \textbf{Limit} ($\mathrm{LIMIT\text{-}ORD}(\xi)$):\quad
      $F(\xi) = \ell\!\left(\xi,\, F\!\restriction\!\mathrm{ORD\text{-}SEGMENT}(\xi)\right)$
\end{itemize}

\begin{VNBcode}
(def-by-ord-recursion 'F
  base-val            ; F(0) = base-val
  '(alpha val)        ; successor: F(succ_ORD(alpha)) = succ-expr
  succ-expr           ;   where val stands for F(alpha)
  '(lambda)           ; limit: F(lambda) = lim-expr
  lim-expr)           ;   where lambda is the limit ordinal
\end{VNBcode}

Three axioms are installed:
\begin{VNBsnip}
F-zero:  F(0) = base-val
F-succ:  forall([alpha in ORD],
           F(succ_ORD(alpha)) = succ-expr[val:=F(alpha)])
F-limit: forall([lambda], LIMIT-ORD(lambda) implies F(lambda) = lim-expr)
\end{VNBsnip}

\paragraph{Example: identity function on ordinals.}
\begin{VNBcode}
(def-by-ord-recursion 'COPY-ORD
  0                                         ; COPY-ORD(0) = 0
  '(alpha val) '(succ_ORD val)              ; COPY-ORD(succ(a)) = succ(COPY-ORD(a))
  '(lam) '(SUP-ORD (ORD-SEGMENT lam)))      ; COPY-ORD(lam) = sup of predecessors
\end{VNBcode}

\paragraph{Example: ordinal addition $\alpha + \beta$ (fixing $\alpha$).}

Defining $A_\alpha(\beta) = \alpha + \beta$ as a function of $\beta$:
\begin{VNBcode}
(def-by-ord-recursion 'ORD-ADD
  'alpha-fixed                              ; ORD-ADD(0) = alpha-fixed
  '(beta val) '(succ_ORD val)              ; ORD-ADD(succ(b)) = succ(ORD-ADD(b))
  '(lam) '(SUP-ORD (ORD-SEGMENT lam)))     ; ORD-ADD(lam) = sup of ORD-ADD on lam
\end{VNBcode}

(Here \texttt{alpha-fixed} is the constant playing the role of $\alpha$.
A full definition would use nested ordinal recursion.)

%% =========================================================
\section{Mathematical structures}
\label{sec:declare-structure}

\subsection{The \texttt{declare-structure} form}

\texttt{declare-structure} declares a named mathematical structure at the
Scheme level.  It registers the structure, installs NTH-based accessor
macetes for every carrier set and every operation, and posts a
definitional axiom for the associated \texttt{IS-\textit{NAME}} predicate
in the current theory.  No quoting is needed.

\begin{VNBcode}
(declare-structure NAME
  (carriers C1 C2 ...)
  (op OPNAME (D1 D2 ...) RANGE)
  (constant CNAME SET))
\end{VNBcode}

\begin{itemize}[label = $\circ$, leftmargin=1cm]

\item \texttt{(carriers C1 C2 ...)}
  Names the carrier sets.  A structure with $m$ carriers has its $i$-th
  carrier at \texttt{nth($i$, s)} in any instance~\texttt{s}.

\item \texttt{(op NAME (D1 D2 ...) RANGE)}
  Declares a function-valued slot.  If more than one domain component is
  listed, the domain is \texttt{CARTESIAN(D1, D2, ...)};
  a single-element list gives a unary domain.  All carrier and op names
  in domain and range positions are automatically expanded to accessor
  applications.

\item \texttt{(constant NAME SET)}
  Declares a constant-valued slot: $\mathit{NAME}(s) \in \mathit{SET}(s)$.

\end{itemize}

Accessor indices are assigned sequentially: carriers occupy indices
$1, \ldots, m$ in declaration order; operations and constants follow
at indices $m+1, m+2, \ldots$.

\begin{remark}
A \texttt{declare-structure} declaration involves three distinct things,
which it is worth keeping clearly separated.
\begin{enumerate}
\item \textbf{The name} --- a symbol such as \texttt{RING} or
  \texttt{METRIC-SPACE}, used as a label in declarations and prose.
\item \textbf{An instance} --- a specific VNB list, which is an
  \emph{element of \textsc{set}} (a concrete set-theoretic object).
  For example,
  \[
    \mathtt{CC\text{-}MS} \;=\;
      [\,\mathtt{CC},\;\lambda([\mathtt{x},\mathtt{y}],\,
      \mathtt{magnitude}(\mathtt{x}-\mathtt{y}))\,]
  \]
  is a 2-element VNB list.  Internally it is \texttt{(LIST CC (VNB-LAMBDA
  (LIST x y) (magnitude (- x y))))}.  The accessors \texttt{X(s)} and
  \texttt{D(s)} are \texttt{nth(1,~s)} and \texttt{nth(2,~s)}.
\item \textbf{The associated class} --- the collection
  $\{\,s : \mathit{IS\text{-}NAME}(s)\,\}$ of \emph{all} VNB lists
  satisfying the structure's typing conditions.  \emph{This class
  shares the structure's name}: \texttt{declare-structure NAME}
  installs both the predicate \texttt{IS-NAME} and the class
  \texttt{NAME}, related by the axiom \texttt{NAME-class}:
  \[
    \forall s.\; s \in \mathit{NAME} \;\Leftrightarrow\; \mathit{IS\text{-}NAME}(s).
  \]
  The class is a proper class (defined by a predicate), not an
  element of \textsc{set}.  Naming the class is what makes the
  natural quantification idiom
  $\texttt{forall}([s \in \mathit{NAME}], \ldots)$
  work --- without it, one would have to write the unbounded form
  $\texttt{forall}([s], \mathit{IS\text{-}NAME}(s) \Rightarrow \ldots)$.
\end{enumerate}
This matches the standard textbook presentation in which a ring, say,
is defined as a 6-tuple $(R,+,\times,-,0,1)$ satisfying certain axioms:
the tuple is the instance (point~2) and the collection of all such tuples
is the class (point~3), now an addressable name.
\end{remark}

\subsection{What \texttt{declare-structure} installs}

\paragraph{Accessor macetes.}
For each carrier $C_i$ (at NTH index $i$) and each operation
$\mathit{op}_j$ (at NTH index $m+j$), a macete is installed that
rewrites the accessor term to the corresponding NTH projection:
\[
  C_i(s) \;\longmapsto\; \texttt{nth}(i,\; s)
  \qquad
  \mathit{op}_j(s) \;\longmapsto\; \texttt{nth}(m{+}j,\; s)
\]
The schema variable matches any expression, so the macete fires anywhere
in a goal.

\paragraph{IS-\textit{NAME} axiom.}
The predicate \texttt{IS-\textit{NAME}} is axiomatized as a conjunction
whose first conjunct asserts that $s$ is a VNB list of exactly the right
length, followed by typing conditions on each accessor:
\begin{align*}
  \forall s.\;\mathit{IS\text{-}NAME}(s) \;\Leftrightarrow\; &
  \mathtt{length}(s) = n \\
  &\land\; C_1(s) \in \mathit{SET} \land \cdots \land C_m(s) \in \mathit{SET}\\
  &\land\; \mathit{op}_1(s) \in \mathit{FUN}(\mathit{dom}_1, \mathit{rng}_1)
  \land \cdots
\end{align*}
where $n = m + k$ is the total number of components ($m$ carriers and $k$
operations/constants).  The length condition ties the predicate to the
instance side of the three-part picture: only a VNB list of exactly $n$
elements can satisfy \texttt{IS-\textit{NAME}}.

\begin{sloppypar}The axiom is named \texttt{IS-\textit{NAME}} and is immediately available
via \texttt{cmd-theorem-assumption} and \texttt{cmd-apply-macete}.\end{sloppypar}

\paragraph{\textit{NAME}-class axiom.}
A second axiom named \texttt{\textit{NAME}-class} ties the predicate to
the class:
\[
  \forall s.\; s \in \mathit{NAME} \;\Leftrightarrow\; \mathit{IS\text{-}NAME}(s).
\]
The symbol $\mathit{NAME}$ thereby denotes the proper class
$\{s \mid \mathit{IS\text{-}NAME}(s)\}$, ready to use as the
quantification range:
\begin{VNBsnip}
forall([r in RING], blah(r))             ; bounded over the ring class
forall([E in NORMEDSPACE], blah(E))      ; bounded over normed spaces
\end{VNBsnip}
which the parser expands to
$\forall s.\; s \in \mathit{NAME} \Rightarrow \ldots$; the
\texttt{\textit{NAME}-class} axiom then converts the antecedent to
the predicate form $\mathit{IS\text{-}NAME}(s)$ exactly when needed.

Characteristic axioms (associativity, distributivity, etc.) are
\emph{not} generated automatically; they are installed by separate
\texttt{theory-add-axiom!} calls after the declaration.

\subsection{Pre-loaded algebraic structures}
\label{sec:algebraic}

The file \texttt{algebraic.scm} is loaded automatically and declares the
following structures, each paired with its characteristic axioms.

\subsubsection{Semigroup}

\begin{VNBcode}
(declare-structure SEMIGROUP
  (carriers A)
  (op MUL (A A) A))
; A -> (NTH 1 s),  MUL -> (NTH 2 s)
\end{VNBcode}

One axiom:
\begin{VNBsnip}
semigroup-assoc:
  forall([s], IS-SEMIGROUP(s) implies
    forall([x in A(s), y in A(s), z in A(s)],
      MUL(s)(MUL(s)(x, y), z) = MUL(s)(x, MUL(s)(y, z))))
\end{VNBsnip}

\subsubsection{Monoid}

\begin{VNBcode}
(declare-structure MONOID
  (carriers A)
  (op MUL (A A) A)
  (constant E A))
; A -> (NTH 1 s),  MUL -> (NTH 2 s),  E -> (NTH 3 s)
\end{VNBcode}

Three axioms: \texttt{monoid-assoc}, \texttt{monoid-left-id}, \texttt{monoid-right-id}.
\begin{VNBsnip}
monoid-left-id:
  forall([s], IS-MONOID(s) implies
    forall([a in A(s)], MUL(s)(E(s), a) = a))

monoid-right-id:
  forall([s], IS-MONOID(s) implies
    forall([a in A(s)], MUL(s)(a, E(s)) = a))
\end{VNBsnip}

Two convenience lemmas are installed in \texttt{axioms.scm}:
\begin{description}
\item[\texttt{monoid-identity-in}]
  $\mathit{IS\text{-}MONOID}(m) \Rightarrow \mathit{E}(m) \in \mathit{A}(m)$.
\item[\texttt{monoid-carrier-closed-mul}]
  $\mathit{IS\text{-}MONOID}(m) \land a \in \mathit{A}(m) \land b \in \mathit{A}(m)
  \Rightarrow \mathit{MUL}(m)(a,b) \in \mathit{A}(m)$.
\end{description}

\subsubsection{Group}

\begin{VNBcode}
(declare-structure GROUP
  (carriers A)
  (op MUL (A A) A)
  (constant E A)
  (op INV (A) A))
; A -> (NTH 1 s),  MUL -> (NTH 2 s),  E -> (NTH 3 s),  INV -> (NTH 4 s)
\end{VNBcode}

\begin{sloppypar}Three axioms suffice (left-only): \texttt{group-assoc}, \texttt{group-left-id},
\texttt{group-left-inv}.\end{sloppypar}
\begin{VNBsnip}
group-left-inv:
  forall([s], IS-GROUP(s) implies
    forall([a in A(s)], MUL(s)(INV(s)(a), a) = E(s)))
\end{VNBsnip}

\subsubsection{Ring}

\begin{VNBcode}
(declare-structure RING
  (carriers A)
  (op ADD  (A A) A)
  (op MUL  (A A) A)
  (op NEG  (A)   A)
  (constant ZERO A)
  (constant ONE  A))
; A->1, ADD->2, MUL->3, NEG->4, ZERO->5, ONE->6
\end{VNBcode}

\begin{sloppypar}
Nine axioms: \texttt{ring-add-assoc}, \texttt{ring-add-comm},
\texttt{ring-add-left-id}, \texttt{ring-add-left-inv},
\texttt{ring-mul-assoc}, \texttt{ring-mul-left-id}, \texttt{ring-mul-right-id},
\texttt{ring-left-dist}, \texttt{ring-right-dist}.
\end{sloppypar}

$(A, \mathit{ADD}, \mathit{ZERO}, \mathit{NEG})$ is an abelian group;
$(A, \mathit{MUL}, \mathit{ONE})$ is a monoid; $\mathit{MUL}$ distributes
over $\mathit{ADD}$.

Two convenience lemmas are installed in \texttt{axioms.scm}:
\begin{description}
\item[\texttt{ring-zero-in}]
  $\mathit{IS\text{-}RING}(r) \Rightarrow \mathit{ZERO}(r) \in \mathit{A}(r)$.
\item[\texttt{ring-carrier-closed-add}]
  $\mathit{IS\text{-}RING}(r) \land a \in \mathit{A}(r) \land b \in \mathit{A}(r)
  \Rightarrow \mathit{ADD}(r)(a,b) \in \mathit{A}(r)$.
\end{description}

\subsubsection{Metric space}

\begin{VNBcode}
(declare-structure METRIC-SPACE
  (carriers X)
  (op D (X X) RR))
; X -> (NTH 1 s),  D -> (NTH 2 s)
\end{VNBcode}

Five axioms: \texttt{metric-pos}, \texttt{metric-self-zero},
\texttt{metric-zero-eq}, \texttt{metric-sym}, \texttt{metric-triangle}.
\begin{VNBsnip}
metric-triangle:
  forall([s], IS-METRIC-SPACE(s) implies
    forall([u in X(s), v in X(s), w in X(s)],
      D(s)(u, w) <= D(s)(u, v) + D(s)(v, w)))
\end{VNBsnip}

\subsection{Further structure definitions}

The same mechanism handles any arity of carriers and operations.  For
instance, a Banach space over $\mathbb{C}$:

\begin{VNBcode}
(declare-structure NORMEDSPACE
  (carriers VECTORS)
  (op SCALAR-TIMES (CC VECTORS) VECTORS)
  (op PLUS         (VECTORS VECTORS) VECTORS)
  (op MINUS        (VECTORS VECTORS) VECTORS)
  (op NORM         (VECTORS) RR))
; VECTORS->1, SCALAR-TIMES->2, PLUS->3, MINUS->4, NORM->5
\end{VNBcode}

This installs five accessor macetes and adds the following axiom:
\begin{VNBsnip}
forall([s],
  IS-NORMEDSPACE(s) iff
    VECTORS(s) in SET
    and SCALAR-TIMES(s) in fun(cartesian(CC, VECTORS(s)), VECTORS(s))
    and PLUS(s) in fun(cartesian(VECTORS(s), VECTORS(s)), VECTORS(s))
    and MINUS(s) in fun(cartesian(VECTORS(s), VECTORS(s)), VECTORS(s))
    and NORM(s) in fun(VECTORS(s), RR))
\end{VNBsnip}

A characteristic axiom relating \texttt{MINUS} to \texttt{PLUS} and
\texttt{SCALAR-TIMES} is installed separately.  Here the axiom is written
in S-expression form rather than string notation, because the string
parser reads \texttt{-1} as unary minus applied to \texttt{1} rather than
as a negative integer literal; the S-expression \texttt{-1} is a genuine
Scheme integer and avoids that ambiguity.

\begin{VNBcode}
(theory-add-axiom! *current-theory* 'normedspace-minus
  '(FORALL s
     (IMPLIES (IS-NORMEDSPACE s)
       (FORALL x (IMPLIES (IN x (VECTORS s))
         (FORALL y (IMPLIES (IN y (VECTORS s))
           (= ((MINUS s) x y)
              ((PLUS s) x ((SCALAR-TIMES s) -1 y))))))))))
\end{VNBcode}

A vector space over an arbitrary field:
\begin{VNBcode}
(declare-structure VECTORSPACE
  (carriers FIELD VECTORS)
  (op SCALAR-TIMES (FIELD VECTORS) VECTORS)
  (op PLUS         (VECTORS VECTORS) VECTORS)
  (constant ZERO VECTORS)
  (op NEG          (VECTORS) VECTORS))
; FIELD->1, VECTORS->2, SCALAR-TIMES->3, PLUS->4, ZERO->5, NEG->6
\end{VNBcode}

%%----------------------------------------------------------
\subsection{Structure specialization: \texttt{specialize-structure}}
\label{sec:specialize-structure}

Every characteristic axiom and every theorem proved about a structure class
has the generic shape
\[
  \forall s.\;\mathit{IS\text{-}NAME}(s) \;\Rightarrow\; P[s].
\]
Once a specific instance $h$ has been \emph{certified}---meaning a theorem
of the form $\mathit{IS\text{-}RING}(h)$ has been proved and installed---every
such generic statement can be specialized to $h$ by universal instantiation
followed by modus ponens, yielding the concrete statement $P[h]$.

\texttt{specialize-structure} automates this transport in bulk:

\begin{VNBcode}
(specialize-structure 'HARP 'RING 'harp-is-ring)
; => 17 theorems specialized for HARP.
\end{VNBcode}

The interactive short form is \texttt{(spec instance struct is-thm)}.

\paragraph{Arguments.}
\begin{itemize}[label=$\circ$, leftmargin=1cm]
\item \textit{instance-name} --- symbol naming the specific structure, e.g.\
  \texttt{HARP}.
\item \textit{struct-name} --- the structure class, e.g.\ \texttt{RING}.
  The predicate \texttt{IS-RING} is derived automatically.
\item \textit{is-thm-name} --- name of a theorem in the current theory
  whose formula is exactly $\mathit{IS\text{-}RING}(\mathit{HARP})$.
\end{itemize}

\paragraph{What gets specialized.}
Every entry in the global theorem table (axioms and proved theorems alike)
whose formula matches
\[
  \forall s.\;\mathit{IS\text{-}RING}(s) \;\Rightarrow\; P[s]
\]
--- with $s$ the bound variable appearing as the sole argument to
$\mathit{IS\text{-}RING}$ --- is specialized.  The result $P[\mathit{HARP}]$
is installed under the name
$\langle\textit{original\text{-}name}\rangle\texttt{-harp}$.
Axioms such as \texttt{ring-add-assoc} are included, so
\texttt{ring-add-assoc-harp} is immediately available after the call.

\paragraph{Soundness.}
No new axioms are introduced.  Each specialized theorem $P[h]$ follows
from the certified $\mathit{IS\text{-}RING}(h)$ and the generic
$\forall s.\;\mathit{IS\text{-}RING}(s)\Rightarrow P[s]$ by two standard
inferences: universal instantiation and modus ponens.
\texttt{specialize-structure} performs these in bulk without requiring a
separate interactive proof for each one.  The certification check at the
start of the call verifies that \textit{is-thm-name} has the expected form
before any specialization occurs; an incorrect theorem name raises an error.

\begin{remark}
Theorems about sub-structures are \emph{not} automatically transported.
For instance, theorems of the form
$\forall s.\;\mathit{IS\text{-}GROUP}(s)\Rightarrow\ldots$
are not touched by \texttt{(spec 'HARP 'RING ...)}.  To obtain them you
would need to certify $\mathit{IS\text{-}GROUP}(\mathit{HARP})$ separately
and call \texttt{spec} again for the \texttt{GROUP} class.
\end{remark}

\paragraph{Example: the complex numbers as a ring.}
The constant \texttt{CC-RING} is defined in \texttt{complex.scm} as the
6-tuple $[\mathit{CC},\,+,\,\times,\,-,\,0,\,1]$.  The axiom
\texttt{cc-is-ring} asserts $\mathit{IS\text{-}RING}(\mathit{CC\text{-}RING})$.

\begin{VNBcode}
(spec 'CC-RING 'RING 'cc-is-ring)
; => n theorems specialized for CC-RING.
; Installs: ring-add-assoc-cc-ring, ring-add-comm-cc-ring,
;           ring-mul-assoc-cc-ring, ring-left-dist-cc-ring, ...
\end{VNBcode}

A proof about $\mathbb{C}$ can then retrieve any of these directly:
\begin{VNBcode}
(ta 'ring-add-comm-cc-ring)
; adds  ADD(CC-RING)(x, y) = ADD(CC-RING)(y, x)  to the context
\end{VNBcode}

%% =========================================================
\section{\texttt{def-functoid}}
\label{sec:def-functoid}

\texttt{def-functoid} introduces a named functoid at the proof-checker
level.  It is a definitional device, not an object of the theory: the
name it installs can be applied and reasoned about, but there is no
formal quantifier ranging over functoids.

The general form is:
\begin{VNBcode}
(def-functoid name (x) body)
\end{VNBcode}
where \texttt{x} is a formal parameter ranging over classes and
\texttt{body} is a class-valued expression, typically built from
accessor macetes installed by \texttt{declare-structure}.

\paragraph{Example: forgetful functor NORMEDSPACE $\to$ METRIC-SPACE.}
A normed space carries an induced metric $d(x,y) = \|x - y\|$.
Since the domain is the proper class of all normed spaces, \texttt{fun}
cannot express this mapping and a functoid is required:

\begin{VNBcode}
(def-functoid NormedToMetric (S)
  [(VECTORS S),
   (lambda [[x, y] in cartesian((VECTORS S), (VECTORS S))],
     ((NORM S) ((MINUS S) x y)))])
\end{VNBcode}

The result is a tuple \texttt{[VECTORS(S), D]} matching the
\texttt{METRIC-SPACE} signature (carrier \texttt{X}, distance \texttt{D}).
That it satisfies \texttt{IS-METRIC-SPACE} is a theorem, not automatic:
the metric axioms follow from the norm axioms and the definition of
\texttt{MINUS}.

\paragraph{Example: forgetful functor NORMEDSPACE $\to$ VECTORSPACE.}
Dropping the norm recovers the underlying vector space:

\begin{VNBcode}
(def-functoid NormedToVector (S)
  [(VECTORS S), (SCALAR-TIMES S), (PLUS S), (MINUS S)])
\end{VNBcode}

\paragraph{Example: the carrier functor.}
\begin{VNBcode}
(def-functoid GroupCarrier (G) (CARRIER G))
\end{VNBcode}
This maps every group to its underlying carrier set.  The result is
always a set (\texttt{IS-GROUP} entails \texttt{CARRIER(G) in SET}),
even though the domain is a proper class.

\paragraph{Transfinite recursion and functoids.}
\begin{sloppypar}The transfinite recursion scheme
(Section~\ref{sec:ord-recursion}) produces a functoid, not a
set-theoretic function: its domain is the proper class $\mathit{ORD}$
of all ordinals, so the result cannot live in $\mathit{FUN}(\mathit{ORD}, B)$
for any set-valued $B$.  A typical pattern is to define the step
functor $\mathcal{E}$ as a functoid via \texttt{def-functoid} and
then invoke \texttt{def-by-ord-recursion}:\end{sloppypar}
The actual signature has six arguments: the new functoid name, the
base value, plus a binding spec and body for both the successor and
limit cases.  The successor body is evaluated with the bound name
referring to the recursive value at the predecessor; the limit body
is evaluated with the bound name referring to the partial functoid
restricted to all earlier ordinals.
\begin{VNBcode}
(def-by-ord-recursion 'MySeq
  'a0                                 ; base value at 0
  '(prev) '(succ_ORD prev)            ; successor: MySeq(succ alpha) = succ_ORD(prev)
  '(part) '(SUP-ORD (image part)))    ; limit:    MySeq(lambda) = sup of partial seq
; => installs axioms
;      MySeq(0) = a0
;      IN alpha ORD => MySeq(succ_ORD alpha) = succ_ORD(MySeq alpha)
;      LIMIT-ORD lambda => MySeq(lambda) = SUP-ORD(image-of MySeq below lambda)
\end{VNBcode}

%%% -----------------------------------------------------------------------
\section{Cardinality: \texttt{CARD}}
\label{sec:cardinality}

\texttt{CARD} is a functoid mapping every set $A$ to its cardinality
$\mathit{CARD}(A) \in \mathit{ORD}$.  Under the axiom of choice,
$\mathit{CARD}(A)$ is the least ordinal $\alpha$ that bijects with $A$
(the Hartogs--Scott construction).  We axiomatise the key properties
directly in \texttt{cardinality.scm}.

\begin{remark}
Because the domain of $\mathit{CARD}$ is the proper class $\mathit{SET}$,
it cannot be expressed as an element of $\mathit{FUN}(\mathit{SET}, \mathit{ORD})$
(that class is not a set).  It is therefore a functoid, on the same
footing as $\mathit{ORD\text{-}SEGMENT}$.
\end{remark}

\begin{remark}
$\mathit{CARD}$ is total over all classes, but characterised only when
its argument is a set.  The kernel does not define $\mathit{CARD}(A)$
for a proper class $A$: the term $\mathit{CARD}(\mathit{ORD})$, for
instance, is syntactically well-formed but no axiom constrains its
value.  This is consistent with the totality framing (every functoid
returns \emph{some} value), but means that any equational use of
$\mathit{CARD}(A)$ outside the established sethood-conditioned axioms
will not discharge.  Convention: only invoke $\mathit{CARD}$ on a
class you have proven to be a set.
\end{remark}

\subsection*{Axioms}

\begin{description}
\item[\texttt{card-in-ord}]
  $A \in \mathit{SET} \Rightarrow \mathit{CARD}(A) \in \mathit{ORD}$.

\item[\texttt{card-empty}]
  $\mathit{CARD}(\emptyset) = 0$.

\item[\texttt{card-insert}]
  $A \in \mathit{SET} \wedge x \notin A \;\Rightarrow\;
   \mathit{CARD}(A \cup \{x\}) = \mathit{succ_{ORD}}(\mathit{CARD}(A))$.\\
  Here $\{x\} = \mathit{PAIR}(x,x)$ (by the pairing axiom with $a = b = x$).

\item[\texttt{card-segment}]
  $n \in \mathit{NN} \;\Rightarrow\; \mathit{CARD}(\mathit{ORD\text{-}SEGMENT}(n)) = n$.\\
  Bridge between cardinality and the $\mathit{NN}$-indexed model of finite sets.

\item[\texttt{card-finite-bij}]
  $A \in \mathit{SET} \wedge \mathit{CARD}(A) \in \mathit{NN} \;\Rightarrow\;$
  there exists a bijection
  $\varphi : \mathit{ORD\text{-}SEGMENT}(\mathit{CARD}(A)) \to A$.\\
  (Existence follows from AC; stated as an axiom here.)

\item[\texttt{card-union-disjoint}]
  $A, B \in \mathit{SET},\;
   \mathit{CARD}(A), \mathit{CARD}(B) \in \mathit{NN},\;
   A \cap B = \emptyset \;\Rightarrow\;
   \mathit{CARD}(A \cup B) = \mathit{CARD}(A) + \mathit{CARD}(B)$.\\
  Finite additivity; the $+$ on the right is $\mathit{NN}$ addition.
\end{description}

\subsection*{Finiteness criterion}

A set $A$ is \emph{finite} iff $\mathit{CARD}(A) \in \mathit{NN}$.
For finite sets, \texttt{card-segment} says
$\mathit{CARD}(\mathit{ORD\text{-}SEGMENT}(n)) = n$, so the initial
segment $\{0, 1, \ldots, n-1\}$ has cardinality exactly $n$.  The
axiom \texttt{card-finite-bij} supplies an enumeration bijection for
any finite set, which is the key ingredient for defining unordered
finite sums in Section~\ref{sec:sequences}.

%%% -----------------------------------------------------------------------
\section{Monoid products and ring sums}
\label{sec:sequences}

\texttt{sequences.scm} defines two constants by primitive recursion on
$\mathit{NN}$:

\begin{itemize}
\item $\mathit{PROD\text{-}ORD}(m, f, n)$ --- product of $f(0) * f(1) * \cdots * f(n{-}1)$
  in a monoid $m$.
\item $\mathit{SUM}(r, f, n)$ --- sum $f(0) + f(1) + \cdots + f(n{-}1)$
  in a ring $r$.
\end{itemize}

Both are $\mathit{NN}$-indexed: the third argument $n$ counts how many
terms are included.  The empty product ($n = 0$) is the identity;
$\mathit{SUM}(r, f, 0) = \mathit{ZERO}(r)$.

\subsection{\texttt{PROD-ORD}: monoid product}

\paragraph{Signature.}
$\mathit{PROD\text{-}ORD}(m, f, n)$ where $m$ is a \textsc{Monoid},
$f \in \mathit{FUN}(\mathit{NN}, \mathit{A}(m))$, $n \in \mathit{NN}$.

\paragraph{Defining recursion.}
Installed by \texttt{def-by-nn-recursion} as two named axioms:

\begin{align*}
\mathtt{prod\text{-}ord\text{-}zero}: &\quad
  \forall m\, f.\;
  \mathit{PROD\text{-}ORD}(m, f, 0) = \mathit{E}(m) \\
\mathtt{prod\text{-}ord\text{-}succ}: &\quad
  \forall m\, f\, n \in \mathit{NN}.\;
  \mathit{PROD\text{-}ORD}(m, f, \mathit{succ}(n))
  = \mathit{PROD\text{-}ORD}(m, f, n) * f(n)
\end{align*}

where $*$ denotes $\mathit{MUL}(m)$.

\paragraph{Additional axioms.}
\begin{description}
\item[\texttt{prod-ord-type}]
  $\mathit{IS\text{-}MONOID}(m) \wedge f \in \mathit{FUN}(\mathit{NN}, \mathit{A}(m))
  \wedge n \in \mathit{NN} \;\Rightarrow\;
  \mathit{PROD\text{-}ORD}(m, f, n) \in \mathit{A}(m)$.
  Installed as an axiom for convenience; formally derivable from
  \texttt{prod-ord-zero}, \texttt{prod-ord-succ}, \texttt{monoid-identity-in},
  and \texttt{monoid-carrier-closed-mul} by $\mathit{NN}$-induction
  (see Section~\ref{sec:ex-ni-type}).

\item[\texttt{prod-ord-singleton}]
  $\mathit{IS\text{-}MONOID}(m) \wedge f \in \mathit{FUN}(\mathit{NN}, \mathit{A}(m))
  \;\Rightarrow\; \mathit{PROD\text{-}ORD}(m, f, 1) = f(0)$.\\
  (Derivable from \texttt{prod-ord-succ} at $n=0$, \texttt{prod-ord-zero},
  and \texttt{monoid-left-id}.)
\end{description}

\subsection{\texttt{SUM}: ring summation}

\texttt{SUM} is the special case of \texttt{PROD-ORD} for the additive
monoid of a ring---\textit{i.e.}, $\mathit{MUL} \to \mathit{ADD}$,
$\mathit{E} \to \mathit{ZERO}$---defined independently for direct use.

\paragraph{Signature.}
$\mathit{SUM}(r, f, n)$ where $r$ is a \textsc{Ring},
$f \in \mathit{FUN}(\mathit{NN}, \mathit{A}(r))$, $n \in \mathit{NN}$.

\paragraph{Defining recursion.}

\begin{align*}
\mathtt{sum\text{-}zero}: &\quad
  \forall r\, f.\; \mathit{SUM}(r, f, 0) = \mathit{ZERO}(r) \\
\mathtt{sum\text{-}succ}: &\quad
  \forall r\, f\, n \in \mathit{NN}.\;
  \mathit{SUM}(r, f, \mathit{succ}(n))
  = \mathit{SUM}(r, f, n) + f(n)
\end{align*}

\paragraph{Additional axioms.}
\begin{description}
\item[\texttt{sum-type}]
  $\mathit{IS\text{-}RING}(r) \wedge f \in \mathit{FUN}(\mathit{NN}, \mathit{A}(r))
  \wedge n \in \mathit{NN} \;\Rightarrow\; \mathit{SUM}(r, f, n) \in \mathit{A}(r)$.
  Installed as an axiom for convenience; formally derivable from
  \texttt{sum-zero}, \texttt{sum-succ}, \texttt{ring-zero-in},
  and \texttt{ring-carrier-closed-add} by $\mathit{NN}$-induction
  (Section~\ref{sec:ex-ni-type}).

\item[\texttt{sum-singleton}]
  $\mathit{IS\text{-}RING}(r) \wedge f \in \mathit{FUN}(\mathit{NN}, \mathit{A}(r))
  \;\Rightarrow\; \mathit{SUM}(r, f, 1) = f(0)$.

\item[\texttt{sum-left-scalar}]
  $a \in \mathit{A}(r) \wedge f \in \mathit{FUN}(\mathit{NN}, \mathit{A}(r))
  \wedge n \in \mathit{NN} \;\Rightarrow\;$\\
  $\mathit{SUM}(r,\, \lambda i.\, a \cdot f(i),\, n)
  = a \cdot \mathit{SUM}(r, f, n)$.\\
  (Derivable by $\mathit{NN}$-induction using \texttt{ring-left-dist}.)
\end{description}

\begin{remark}
The unordered finite sum $\sum_{x \in S} f(x)$ over an arbitrary finite
set $S$ (not indexed by $\mathit{NN}$) requires commutativity of
addition and a well-definedness proof (independence of the enumeration
bijection from \texttt{card-finite-bij}).  This more general
$\mathit{SUM\text{-}SET}$ operation is planned but not yet implemented.
\end{remark}

\begin{remark}
The connection $\mathit{SUM}(r, f, n) = \mathit{PROD\text{-}ORD}(m, f, n)$
where $m = [\mathit{A}(r),\, \mathit{ADD}(r),\, \mathit{ZERO}(r)]$ is the
additive monoid of $r$ is a theorem (by structural induction on $n$),
not an axiom.
\end{remark}

\paragraph{Typical proof pattern.}
To use \texttt{SUM} in a proof, unfold the recursion manually with
\texttt{(ta 'sum-succ)} and \texttt{(ta 'sum-zero)}, then close
numeric identities with \texttt{(rs)} or \texttt{(arith)}.

\begin{VNBcode}
;; Prove SUM(r, f, 2) = f(0) + f(1)
(sp (make-wff-from-string
  "forall([r, f], IS-RING(r) and f in FUN(NN, A(r)) implies
     SUM(r, f, 2) = ADD(r)(SUM(r, f, 1), f(1))"))
(ta 'sum-succ)    ; unfold outer step
(ta 'sum-singleton)
(rfl)
\end{VNBcode}

%% =========================================================
\section{Ring products}
\label{sec:ring-prod}

\texttt{algebraic.scm} defines the binary ring product, and
\texttt{sequences.scm} extends it to finite indexed products.
Both follow Path~2: the operations are \emph{total} functoids, defined
for any arguments; typing theorems fire conditionally when the inputs
are rings.

\subsection{Zero ring: \texttt{ZERO-RING}}

\texttt{ZERO-RING} is the terminal ring: carrier $\{0\}$, all operations
returning $0$, with $\mathit{ZERO} = \mathit{ONE} = 0$.

\begin{VNBsnip}
ZERO-RING = [{0},
             lambda([p, q], 0),   ; ADD
             lambda([p, q], 0),   ; MUL
             lambda([p], 0),      ; NEG
             0,                   ; ZERO
             0]                   ; ONE
\end{VNBsnip}

\begin{description}
\item[\texttt{zero-ring-def}]
  $\mathit{ZERO\text{-}RING} = [\{0\},\,
   \lambda(p,q).\,0,\; \lambda(p,q).\,0,\; \lambda(p).\,0,\; 0,\; 0]$.
\item[\texttt{zero-ring-is-ring}]
  $\mathit{IS\text{-}RING}(\mathit{ZERO\text{-}RING})$.
  (Taken as an axiom; the proof would verify the nine ring axioms for
  the trivial one-element ring.)
\end{description}

\begin{remark}
The zero ring is the terminal object in the category of rings: there is
exactly one ring homomorphism from \emph{any} ring $R$ to
$\mathit{ZERO\text{-}RING}$ (mapping everything to $0$).  It is also
the identity for \texttt{RING-PROD} up to isomorphism:
$\mathit{RING\text{-}PROD}(\mathit{ZERO\text{-}RING}, R) \cong R$.
\end{remark}

\subsection{Binary product: \texttt{RING-PROD}}

\texttt{RING-PROD(X,~Y)} is the product ring with carrier
$\mathit{CARTESIAN}(\mathit{A}(X),\,\mathit{A}(Y))$ and
componentwise operations.

\paragraph{Definition.}
Installed as a \texttt{def-functoid} rewrite:

\begin{VNBsnip}
RING-PROD(X, Y) =
  [CARTESIAN(A(X), A(Y)),
   lambda([p, q], [ADD(X)(nth(1,p), nth(1,q)),
                   ADD(Y)(nth(2,p), nth(2,q))]),
   lambda([p, q], [MUL(X)(nth(1,p), nth(1,q)),
                   MUL(Y)(nth(2,p), nth(2,q))]),
   lambda([p],    [NEG(X)(nth(1,p)),
                   NEG(Y)(nth(2,p))]),
   [ZERO(X), ZERO(Y)],
   [ONE(X),  ONE(Y)]]
\end{VNBsnip}

Elements of $\mathit{A}(\mathit{RING\text{-}PROD}(X,Y))$ are pairs $[a,b]$
with $a \in \mathit{A}(X)$ and $b \in \mathit{A}(Y)$;
all six operations act componentwise.

\paragraph{Accessor expansion.}
The \texttt{RING-PROD} macete chain is:
\[
  \mathit{A}(\mathit{RING\text{-}PROD}(X,Y))
  \;\longrightarrow\; \mathit{CARTESIAN}(\mathit{A}(X),\,\mathit{A}(Y))
\]
and similarly for \texttt{ADD}, \texttt{MUL}, \texttt{NEG}, \texttt{ZERO},
\texttt{ONE}: each expands by first applying the \texttt{RING-PROD}
rewrite and then the accessor macete.

\paragraph{Typing axiom.}
\begin{description}
\item[\texttt{ring-prod-is-ring}]
  $\mathit{IS\text{-}RING}(X) \wedge \mathit{IS\text{-}RING}(Y)
  \;\Rightarrow\; \mathit{IS\text{-}RING}(\mathit{RING\text{-}PROD}(X,Y))$.
  (Taken as an axiom.  The proof would verify that the componentwise
  operations satisfy each ring axiom, given that $X$ and $Y$ do.)
\end{description}

\begin{remark}[Totality]
\texttt{RING-PROD} is total: $\mathit{RING\text{-}PROD}(X,Y)$ is a
well-defined tuple for \emph{any} $X$ and $Y$.  When neither is a ring
--- when bongos are passed in --- the result is some tuple that satisfies
no ring axioms.  \texttt{ring-prod-is-ring} is silent on this case (its
antecedent is false), so no contradiction arises.  This is the VNB
tradition: totality and typing are separate concerns.
\end{remark}

\subsection{Finite indexed product: \texttt{RING-PROD-N}}

$\mathit{RING\text{-}PROD\text{-}N}(f, n)$ is the $n$-fold product ring
$f(0) \times f(1) \times \cdots \times f(n-1)$, defined by primitive
recursion on~$n$.

\paragraph{Signature.}
$f$ is any function on $\mathit{NN}$ (not required to map to rings;
the typing axiom states when the result is a ring).
$n \in \mathit{NN}$.

\paragraph{Defining recursion.}
Installed by \texttt{def-by-nn-recursion}:

\begin{align*}
\mathtt{ring\text{-}prod\text{-}n\text{-}zero}: &\quad
  \forall f.\; \mathit{RING\text{-}PROD\text{-}N}(f, 0) = \mathit{ZERO\text{-}RING} \\
\mathtt{ring\text{-}prod\text{-}n\text{-}succ}: &\quad
  \forall f,\, n \in \mathit{NN}.\;
  \mathit{RING\text{-}PROD\text{-}N}(f, \mathit{succ}(n))
  = \mathit{RING\text{-}PROD}(\mathit{RING\text{-}PROD\text{-}N}(f, n),\; f(n))
\end{align*}

\paragraph{Typing axiom.}
\begin{description}
\item[\texttt{ring-prod-n-is-ring}]
  $(\forall i \in \mathit{NN}.\;\mathit{IS\text{-}RING}(f(i)))
  \;\wedge\; n \in \mathit{NN}
  \;\Rightarrow\; \mathit{IS\text{-}RING}(\mathit{RING\text{-}PROD\text{-}N}(f,n))$.\\
  (Derivable by $\mathit{NN}$-induction using \texttt{zero-ring-is-ring}
  and \texttt{ring-prod-is-ring}; installed as an axiom for direct use.)
\end{description}

\begin{remark}
The indexed-family product $\prod_{i \in I} f(i)$ for an arbitrary
index set $I$ (not necessarily $\mathit{NN}$) requires the big-union
constructor $\bigcup_{i \in I} \mathit{A}(f(i))$ to serve as the
function codomain.  This depends on \texttt{UNION-SET}, which is
currently pending.  The family case will be added once
\texttt{UNION-SET} is in place.
\end{remark}


%% =========================================================
\appendix
\chapter{Source file overview}
\label{app:source}

The prover sources live in \texttt{\textasciitilde/prover/} and are
organised into three groups:

\begin{itemize}[label=\textendash, leftmargin=0.85cm]
\item \textbf{Kernel} (root of \texttt{\textasciitilde/prover/}) ---
  the proof checker engine, parser, theory record, and interactive
  command layer.
\item \textbf{\texttt{structure-library/}} --- algebraic structures
  (semigroup, monoid, group, ring, metric-space), the numeric-system
  instances (NN/ZZ/QQ/RR/CC as rings), ordinals, cardinality,
  sequences, the complex-number library, and the polynomial
  ring-identity decision procedure.
\item \textbf{\texttt{theorem-library/}} --- axioms not yet derivable
  from the primitive inference kernel.  Today this is just
  \texttt{axioms.scm}; future lemma collections will accumulate here.
\end{itemize}

\medskip
\noindent
File-by-file:

\begin{itemize}[label = $\circ$, leftmargin=0.85cm]

\item \begin{sloppypar}\texttt{expressions.scm}
  Expression representation, free variables, substitution
  (capture-avoiding), alpha-equivalence.\end{sloppypar}

\item \texttt{sequents.scm}
  Sequents (list of assumptions + assertion), context operations.

\item \texttt{deduction-graphs.scm}
  Sequent nodes, inference nodes, the deduction graph record type,
  grounding propagation.

\item \texttt{primitive-inferences.scm}
  The kernel: all proof rules as Scheme procedures.  The only code
  allowed to extend the deduction graph.

\item \texttt{macetes.scm}
  Pattern matching, elementary macete construction from theorems,
  compound macete combinators.

\item \texttt{theory.scm}
  Theory record, axiom/theorem registration, the VNB base theory.

\item \begin{sloppypar}\texttt{structures.scm}
  \texttt{declare-structure} (user-facing macro) and \texttt{def-structure}
  (lower-level Scheme procedure): register structure definitions, install
  NTH-based accessor macetes, auto-generate \texttt{IS-\textit{NAME}}
  axioms.\end{sloppypar}

\item \texttt{number-systems.scm}
  Declares \texttt{NN}, \texttt{ZZ}, \texttt{QQ}, \texttt{RR},
  \texttt{CC} as constants and installs their arithmetic axioms
  (sethood, inclusion chain, ring/field laws, absolute value,
  conjugation).

\item \begin{sloppypar}\texttt{wff.scm}
  \texttt{<wff>} record type (\texttt{formula}, \texttt{theory},
  \texttt{contexts} fields); \texttt{wff?}, \texttt{wff-formula},
  \texttt{wff-theory}; \texttt{wff-child} (inherit theory from parent),
  \texttt{wff-in-theory} (bare constructor), \texttt{wff-equiv?}.\end{sloppypar}

\item \begin{sloppypar}\texttt{contexts.scm}
  \texttt{make-wff} (snapshot current theory and context stack),
  \texttt{wff-free-vars}, \texttt{declare-local-context},
  \texttt{undeclare-local-context}, \texttt{get-context},
  \texttt{unravel}, \texttt{print-wff}, \texttt{display-contents}
  (browse a wff's formula, theory, and kind).\end{sloppypar}

\item \texttt{arith-eval.scm}
  Ground arithmetic decision procedure.  Evaluates closed formulas over
  $\mathit{NN}/\mathit{ZZ}/\mathit{QQ}/\mathit{RR}/\mathit{CC}$ using
  Scheme's exact arithmetic.  Handles \texttt{+}, \texttt{*},
  \texttt{-}, \texttt{recip}, \texttt{abs}, \texttt{conjugate},
  \texttt{succ}, \texttt{<=}, \texttt{=}, \texttt{IN}, and
  \texttt{power}.  Provides \texttt{vnb-do-arith} (standalone
  predicate), \texttt{vnb-create-num-theorem} (verify and install a
  numeric theorem), and \texttt{pi-arith!} / \texttt{cmd-arith} (closes
  proof goals).

\item \texttt{proof-commands.scm}
  High-level command wrappers with automatic focus management.  Also
  \texttt{cmd-qed} (install a completed proof as a theorem) and
  \texttt{subst-wff} (meta-level substitution into a wff).

\item \texttt{interactive.scm}
  Interactive short-form commands (\texttt{sp}, \texttt{di},
  \texttt{ai}, \ldots), \texttt{fa}/\texttt{fs} quantification helpers,
  and the \texttt{show} display function.  Loaded at session startup.

\item \texttt{vnb.el}
  Emacs interface.  REPL window (\texttt{*VNB*}),
  proof-state display window (\texttt{*VNB State*}), command scratch
  window (\texttt{*VNB Commands*}), command alist and completion,
  \texttt{vnb-eval-string} for programmatic evaluation.

\item \texttt{vnb-launch.el}
  Workspace launcher.  Loaded by the \texttt{VNB} shell script; sets
  up the styled Home, Focus, and Proof Overview workspaces, the VNB
  menu bar and toolbar, the Display preferences panel, and the error
  panel that catches \texttt{;; VNB error:} lines streaming back from
  the prover.

\item \texttt{structure-library/} --- one file per declared structure
  (\texttt{semigroup.scm}, \texttt{monoid.scm} (includes \texttt{COMM-MONOID}),
  \texttt{group.scm}, \texttt{ring.scm} (includes \texttt{RING-PROD}
  and \texttt{ZERO-RING}), \texttt{metric-space.scm}), plus
  \texttt{ring-simplify.scm} (polynomial-normal-form decision procedure),
  \texttt{numeric-instances.scm} (NN/ZZ/QQ/RR/CC as algebraic
  structures), \texttt{complex.scm}, \texttt{extended-reals.scm},
  \texttt{ordinals.scm}, \texttt{cardinality.scm}, and
  \texttt{sequences.scm}.  New structures land here automatically when
  saved via the \textbf{Build Structure} button on the Home Workspace.

\item \texttt{theorem-library/axioms.scm} --- axioms still taken as
  primitive: equality substitution and other facts not yet derivable
  from the kernel.  The accompanying source header lists what is
  installed and what is still pending.

\end{itemize}

%% =========================================================



%% =========================================================
\chapter{Regression and demonstration test suite}
\label{app:tests}

The file \texttt{test-suite.scm} is the canonical regression suite for
the VNB proof checker.  It exercises every major subsystem in order:
ground arithmetic, axiom installation, the parser, wff validation,
macetes with local context, proof commands, the number-system chain,
and arithmetic over~$\mathbb{Z}$, $\mathbb{Q}$, $\mathbb{R}$,
$\mathbb{C}$.  Running it requires only \texttt{load.scm}; each test
prints \texttt{PASS} or \texttt{FAIL} and the final line gives the
aggregate count.

\begin{VNBcode}
;;; Load the system, then load this file:
;;;   mit-scheme --quiet --load load.scm < test-suite.scm
(load "load.scm")
\end{VNBcode}

\section{Test harness}

\begin{sloppypar}The harness uses three predicates (\texttt{check}, \texttt{check-true},
\texttt{check-false}), one error-expectation helper
(\texttt{check-error}), and one proof-completion helper
(\texttt{check-proof}).  Two convenience wrappers, \texttt{wff-str} and
\texttt{macete-str}, keep individual test lines short.\end{sloppypar}

\begin{VNBcode}
(define *pass-count* 0)
(define *fail-count* 0)

(define (check label thunk expected)
  (let ((got (thunk)))
    (if (equal? got expected)
        (begin (set! *pass-count* (+ *pass-count* 1))
               (display "  PASS  ") (display label) (newline))
        (begin (set! *fail-count* (+ *fail-count* 1))
               (display "  FAIL  ") (display label)
               (display " -- expected ") (write expected)
               (display " got ") (write got) (newline)))))

(define (check-true  label thunk) (check label thunk #t))
(define (check-false label thunk) (check label thunk #f))

(define (check-error label thunk)
  ;; Pass when thunk returns a <vnb-error> (i.e. input was rejected).
  (let ((result (thunk)))
    (if (vnb-error? result)
        (begin (set! *pass-count* (+ *pass-count* 1))
               (display "  PASS  ") (display label) (newline))
        (begin (set! *fail-count* (+ *fail-count* 1))
               (display "  FAIL  ") (display label)
               (display " -- expected error, got ") (write result) (newline)))))

(define (check-proof label thunk)
  ;; Pass when thunk completes without raising.
  (call-with-current-continuation
    (lambda (k)
      (with-exception-handler
        (lambda (e)
          (set! *fail-count* (+ *fail-count* 1))
          (display "  FAIL  ") (display label)
          (display " -- exception: ")
          (display (condition/report-string e)) (newline)
          (k #f))
        thunk)
      (set! *pass-count* (+ *pass-count* 1))
      (display "  PASS  ") (display label) (newline))))

;;; wff-str: extract display string from a wff (strips surrounding quotes).
(define (wff-str w)
  (let ((s (wff->string w)))
    (substring s 1 (- (string-length s) 1))))

;;; macete-str: apply a named macete to a formula given as a string.
;;; Returns the result display string, or #f if the macete doesn't fire.
(define (macete-str name str)
  (let ((r (apply-macete name (make-wff-from-string str))))
    (and r (wff-str r))))

;;; roundtrip: parse a formula string and print it back.
(define (roundtrip str)
  (let ((w (make-wff-from-string str)))
    (if (vnb-error? w) #f (wff-str w))))
\end{VNBcode}

\section{Ground arithmetic evaluation}

The built-in evaluator decides membership and equalities over exact
numbers without invoking the proof kernel.

\begin{VNBcode}
(check-true  "0 in NN"        (lambda () (arith-membership-check 0   'NN)))
(check-true  "100 in NN"      (lambda () (arith-membership-check 100 'NN)))
(check-false "-1 not in NN"   (lambda () (arith-membership-check -1  'NN)))
(check-false "1/2 not in NN"  (lambda () (arith-membership-check 1/2 'NN)))

(check-true  "0 in ZZ"        (lambda () (arith-membership-check 0   'ZZ)))
(check-true  "-3 in ZZ"       (lambda () (arith-membership-check -3  'ZZ)))
(check-false "1/2 not in ZZ"  (lambda () (arith-membership-check 1/2 'ZZ)))

(check-true  "1/3 in QQ"      (lambda () (arith-membership-check 1/3 'QQ)))
(check-false "1/3 not in ZZ"  (lambda () (arith-membership-check 1/3 'ZZ)))

(check "2 + 3 = 5"       (lambda () (arith-eval-term '(+ 2 3)))     5)
(check "3 * 4 = 12"      (lambda () (arith-eval-term '(* 3 4)))     12)
(check "10 + (-3) = 7"   (lambda () (arith-eval-term '(+ 10 (- 3)))) 7)
(check "succ(0) = 1"     (lambda () (arith-eval-term '(succ 0)))    1)
(check "power(2,3) = 8"  (lambda () (arith-eval-term '(power 2 3))) 8)

(check-true  "= 2+3 5"   (lambda () (arith-eval-formula '(= (+ 2 3) 5))))
(check-false "= 2+3 6"   (lambda () (arith-eval-formula '(= (+ 2 3) 6))))
(check-true  "<= 0 5"    (lambda () (arith-eval-formula '(<= 0 5))))
(check-false "<= 5 0"    (lambda () (arith-eval-formula '(<= 5 0))))
\end{VNBcode}

\section{Axioms installed}

Every axiom loaded by \texttt{load.scm} is simultaneously a theorem
(available for \texttt{ta} and backchaining) and a macete entry.
The \texttt{axiom-present?} helper checks the theorem table.

\begin{VNBcode}
(define (axiom-present? name)
  (if (hash-table-ref/default *theorem-table* name #f) #t #f))

(for-each (lambda (name)
            (check-true (symbol->string name)
                        (lambda () (axiom-present? name))))
  '(membership-implies-sethood extensionality
    empty-set-is-set empty-set-has-no-members
    pairing pairing-membership
    power-set power-set-membership
    union-set-closure intersection-set-closure
    fun-set-iff fun-elements-are-sets
    fun-domain-apply-def fun-codomain-iff
    cartesian-set-iff tuples-sethood
    make-set-membership make-set-sethood make-set-empty
    length-of-empty length-in-nn nth-in-range
    nn-is-set nn-zero-in nn-succ-closed
    nn-add-closed nn-mul-closed nn-add-comm nn-mul-comm
    nn-add-assoc nn-mul-assoc nn-distributive nn-one-mul nn-add-zero
    power-typing-nonneg power-zero power-succ power-neg
    zz-is-set zz-zero-in zz-one-in
    qq-is-set rr-is-set cc-is-set
    nn-subset-zz zz-subset-qq qq-subset-rr rr-subset-cc
    burali-forti nn-subset-ord ord-succ-in transfinite-induction))

;;; ZZplus must be completely absent (superseded by ZZ axioms)
(check-false "zzplus-is-set absent"
             (lambda () (axiom-present? 'zzplus-is-set)))
\end{VNBcode}

\section{Parser round-trips}

The parser produces an internal S-expression; \texttt{wff->string} prints
it back.  Round-tripping through \texttt{make-wff-from-string} and
\texttt{wff-str} is the primary correctness check for the surface syntax.

\begin{VNBcode}
(check "n in nn"         (lambda () (roundtrip "n in nn"))         "n in nn")
(check "n = m"           (lambda () (roundtrip "n = m"))           "n = m")
(check "truth"           (lambda () (roundtrip "truth"))           "truth")
(check "falsity"         (lambda () (roundtrip "falsity"))         "falsity")
(check "n + m = m + n"   (lambda () (roundtrip "n + m = m + n"))   "n + m = m + n")

(check "A and B"         (lambda () (roundtrip "n in nn and m in nn"))
                         "n in nn and m in nn")
(check "A implies B"     (lambda () (roundtrip "n in nn implies n in zz"))
                         "n in nn implies n in zz")
(check "not A"           (lambda () (roundtrip "not(n in nn)"))    "not(n in nn)")

(check-true "forall bare"
  (lambda () (wff? (make-wff-from-string "forall([x], x in nn implies x in zz)"))))
(check-true "forall restricted"
  (lambda () (wff? (make-wff-from-string "forall([x in nn], x in zz)"))))
(check-true "forsome"
  (lambda () (wff? (make-wff-from-string "forsome([x in nn], x = 0)"))))
(check-true "sep {x in A | p}"
  (lambda () (wff? (make-wff-from-string "z in {x in a | x in nn}"))))
\end{VNBcode}

\section{Wff validation}

\begin{sloppypar}\texttt{make-wff} and \texttt{make-wff-from-string} validate the formula
and return a \texttt{<vnb-error>} (rather than raising) on
ill-formed input.\end{sloppypar}

\begin{VNBcode}
(check-true  "IN accepted"      (lambda () (wff? (make-wff-from-string
                                                   "n in nn"))))
(check-true  "IMPLIES accepted" (lambda () (wff? (make-wff-from-string
                                                   "n in nn implies n in zz"))))
(check-true  "NOT accepted"     (lambda () (wff? (make-wff-from-string
                                                   "not(n in nn)"))))
(check-true  "FORALL accepted"  (lambda () (wff? (make-wff-from-string
                                                   "forall([x in nn], x in zz)"))))

;;; Invalid: bare number in wff position
(check-error "bare number rejected"
             (lambda () (make-wff-from-string "1")))

;;; Invalid: term-forming head (UNION) in wff position
(check-error "union in wff position rejected"
             (lambda () (make-wff '(UNION x y))))
\end{VNBcode}

\section[\texttt{apply-macete} and local context (Monk 1988)]{\texttt{apply-macete} and local context (Monk~\cite{Monk1988})}

Conditions of a conditional macete are checked against the local
context accumulated as the rewriter descends.  Entering the
consequent of \texttt{implies} adds the antecedent; entering the
right conjunct of \texttt{and} adds the left conjunct; entering the
right disjunct of \texttt{or} adds the negation of the left disjunct.

\subsection*{Unconditional macetes}

\begin{VNBcode}
(install-theorem! 'eq-symm
  (make-wff-from-string "forall([x, y], x = y iff y = x)"))

(check "eq-symm: 3+5=8 -> 8=3+5"
       (lambda () (macete-str 'eq-symm "3 + 5 = 8"))
       "8 = 3 + 5")
(check "eq-symm: rewrites inside implies"
       (lambda () (macete-str 'eq-symm "n = m implies truth"))
       "m = n implies truth")
(check "eq-symm: no match on n in nn"
       (lambda () (macete-str 'eq-symm "n in nn"))
       #f)
(check "eq-symm: both sides rewritten"
       (lambda () (macete-str 'eq-symm "n = m and m = k"))
       "m = n and k = m")

(install-theorem! 'and-comm-iff
  (make-wff-from-string "forall([P, Q], P and Q iff Q and P)"))

(check "and-comm: (A and B) -> (B and A)"
       (lambda () (macete-str 'and-comm-iff "n in nn and m in nn"))
       "m in nn and n in nn")
(check "and-comm: fires inside implies antecedent"
       (lambda () (macete-str 'and-comm-iff "n in nn and m in nn implies truth"))
       "m in nn and n in nn implies truth")
\end{VNBcode}

\subsection*{Conditional macetes}

\begin{VNBcode}
;;; nn-add-comm: forall([a,b in NN], a+b = b+a)
;;; Without context the rewrite is suppressed.
(check "nn-add-comm: no fire on bare sum"
       (lambda () (macete-str 'nn-add-comm "n + m = 0"))   #f)

;;; Antecedent (AND (IN n NN) (IN m NN)) is flattened into local
;;; context when entering the consequent.
(check "nn-add-comm: fires inside implies consequent"
       (lambda () (macete-str 'nn-add-comm
                              "n in nn and m in nn implies n + m = m + n"))
       "n in nn and m in nn implies m + n = n + m")

;;; nn-add-zero: forall([n in NN], n + 0 = n)
(check "nn-add-zero: no fire without context"
       (lambda () (macete-str 'nn-add-zero "n + 0 = 5"))   #f)
(check "nn-add-zero: fires inside implies consequent"
       (lambda () (macete-str 'nn-add-zero "n in nn implies n + 0 = 5"))
       "n in nn implies n = 5")
;;; Entering right conjunct of AND adds left conjunct to context.
(check "nn-add-zero: fires in right conjunct of AND"
       (lambda () (macete-str 'nn-add-zero "k in nn and k + 0 = 5"))
       "k in nn and k = 5")
;;; Wrong order: condition is right conjunct — no context available yet.
(check "nn-add-zero: no fire when condition is right conjunct"
       (lambda () (macete-str 'nn-add-zero "k + 0 = 5 and k in nn"))  #f)

;;; nn-add-closed: rewrites n+m in NN to truth when conditions hold.
(check "nn-add-closed: fires inside implies, rewrites to truth"
       (lambda () (macete-str 'nn-add-closed
                              "n in nn and m in nn implies n + m in nn"))
       "n in nn and m in nn implies truth")
(check "nn-add-closed: no fire without context"
       (lambda () (macete-str 'nn-add-closed "n + m in nn"))  #f)

;;; power-succ: forall([x in CC, n in NN], power(x,succ(n)) = x*power(x,n))
(check "power-succ: fires with two conditions in context"
       (lambda () (macete-str 'power-succ
                              "z in cc and k in nn implies power(z, succ(k)) = w"))
       "z in cc and k in nn implies z * power(z, k) = w")
\end{VNBcode}

\section{Proof commands}

Each sub-test proves a simple theorem using one or two rules.
\texttt{check-proof} verifies that the proof completes without exception.

\begin{VNBcode}
;;; Arithmetic (arith) — ground facts
(check-proof "arith: 0 in NN"
  (lambda () (sp (make-wff-from-string "0 in nn")) (arith)))
(check-proof "arith: 2 + 3 = 5"
  (lambda () (sp (make-wff-from-string "2 + 3 = 5")) (arith)))
(check-proof "arith: -5 in ZZ"
  (lambda () (sp (make-wff-from-string "-5 in zz")) (arith)))
(check-proof "arith: 1/3 in QQ"
  (lambda () (sp (make-wff-from-string "1/3 in qq")) (arith)))

;;; Reflexivity (rfl) and quasi-reflexivity (qrfl)
(check-proof "rfl: n = n"
  (lambda () (sp (make-wff-from-string "n = n")) (rfl)))
(check-proof "qrfl: n == n"
  (lambda () (sp (make-wff-from-string "n == n")) (qrfl)))

;;; Direct inference (di) — introduce implication, split conjunction/AND
(check-proof "di + ass: A implies A"
  (lambda ()
    (sp (make-wff-from-string "n in nn implies n in nn"))
    (di) (ass)))

;;; Antecedent inference (ai) — decompose an assumption
(check-proof "di + ai: prove A and B from A, B"
  (lambda ()
    (sp (make-wff-from-string "n in nn and m in nn implies m in nn and n in nn"))
    (di)
    (ai (make-wff-from-string "n in nn and m in nn"))
    (di) (ass) (ass)))

;;; OR introduction (oi-l, oi-r)
(check-proof "oi-l: A implies (A or B)"
  (lambda ()
    (sp (make-wff-from-string "n in nn implies n in nn or m in nn"))
    (di) (oi-l) (ass)))
(check-proof "oi-r: A implies (B or A)"
  (lambda ()
    (sp (make-wff-from-string "n in nn implies m in nn or n in nn"))
    (di) (oi-r) (ass)))

;;; Theorem assumption (ta) + universal instantiation (ui)
(check-proof "ta + ui: n+0 = n"
  (lambda ()
    (sp (make-wff-from-string "n in nn implies n + 0 = n"))
    (di) (ta 'nn-add-zero) (ui "n") (ass)))
(check-proof "ta + ui: n+m = m+n"
  (lambda ()
    (sp (make-wff-from-string "n in nn and m in nn implies n + m = m + n"))
    (di) (ta 'nn-add-comm) (ui "n") (ui "m") (ass)))
(check-proof "ta + ui: n in NN => n in ZZ"
  (lambda ()
    (sp (make-wff-from-string "n in nn implies n in zz"))
    (di) (ta 'nn-subset-zz) (ui "n") (ass)))
(check-proof "ta + ui: power(x,succ n)"
  (lambda ()
    (sp (make-wff-from-string
          "x in cc and n in nn implies power(x, succ(n)) = x * power(x, n)"))
    (di) (ta 'power-succ) (ui "x") (ui "n") (ass)))

;;; Backchain (bc) — use an implication theorem as an introduction rule
(check-proof "bc: n+m in NN"
  (lambda ()
    (sp (make-wff-from-string "n in nn and m in nn implies n + m in nn"))
    (di) (bc 'nn-add-closed) (ass) (ass)))

;;; Proof by contradiction (pbc)
(check-proof "pbc: double negation"
  (lambda ()
    (sp (make-wff-from-string "not(not(n in nn)) implies n in nn"))
    (di) (pbc) (ass)))

;;; Macete inside a proof (mac)
(check-proof "mac: eq-symm rewrites goal"
  (lambda ()
    (sp (make-wff-from-string "n = m implies m = n"))
    (di) (mac 'eq-symm) (ass)))

;;; Exists-witness (ew)
(check-proof "ew: exists x in NN. x = 0"
  (lambda ()
    (sp (make-wff-from-string "forsome([x in nn], x = 0)"))
    (ew "0") (arith) (arith)))

;;; Cut — introduce an intermediate lemma
(check-proof "cut: introduce intermediate fact"
  (lambda ()
    (sp (make-wff-from-string "n in nn implies n + 0 = n"))
    (di)
    (cut (make-wff-from-string "n in nn implies n + 0 = n"))
    (di) (ta 'nn-add-zero) (ui "n") (ass)
    (di) (ass)))

;;; Functoid (lambda) beta reduction (beta)
(check-proof "beta: lambda([x in nn], x+1)(3) = 3+1"
  (lambda ()
    (sp (make-wff-from-string "lambda([x in nn], x + 1)(3) = 3 + 1"))
    (beta) (rfl)))
\end{VNBcode}

\section{Number-system inclusion chain}

\begin{VNBcode}
(for-each
  (lambda (triple)
    (let ((val (car triple)) (sys (cadr triple)) (expected (caddr triple)))
      (check (string-append (number->string val) " in " (symbol->string sys))
             (lambda () (arith-membership-check val sys))
             expected)))
  '((0    NN #t)  (1    NN #t)  (-1   NN #f)  (1/2  NN #f)
    (0    ZZ #t)  (-5   ZZ #t)  (1/2  ZZ #f)
    (1/3  QQ #t)  (-7   QQ #t)  (1/3  ZZ #f)
    (0    RR #t)  (-7   RR #t)
    (0    CC #t)  (-5   CC #t)))
\end{VNBcode}

\section{Arithmetic in \texorpdfstring{$\mathbb{Z}$, $\mathbb{Q}$, $\mathbb{R}$, $\mathbb{C}$}{Z, Q, R, C}}

\begin{VNBcode}
(check-proof "arith: 0 in ZZ"
  (lambda () (sp (make-wff-from-string "0 in zz")) (arith)))
(check-proof "arith: -3 in ZZ"
  (lambda () (sp (make-wff-from-string "-3 in zz")) (arith)))
(check-proof "arith: 1/3 in QQ"
  (lambda () (sp (make-wff-from-string "1/3 in qq")) (arith)))
(check-proof "arith: 1/3 not in ZZ"
  (lambda () (sp (make-wff-from-string "not(1/3 in zz)")) (arith)))

(check "1/3 in qq" (lambda () (arith-eval-formula '(IN 1/3 QQ)))  #t)
(check "1/3 in zz" (lambda () (arith-eval-formula '(IN 1/3 ZZ)))  #f)
(check "-3 in zz"  (lambda () (arith-eval-formula '(IN -3 ZZ)))   #t)
(check "-3 in qq"  (lambda () (arith-eval-formula '(IN -3 QQ)))   #t)
\end{VNBcode}

\section{Summary line}

\begin{VNBcode}
(newline)
(display "=== SUMMARY: ")
(display *pass-count*) (display " passed, ")
(display *fail-count*) (display " failed ===\n")
\end{VNBcode}

\noindent A clean run prints \verb|165 passed, 0 failed|.

%% =========================================================


%% =========================================================
% app-glossary-generated.tex -- GENERATED by docs/gen-glossary.py from
% reference/GLOSSARY.md.  Do not edit; re-run `make gen' in docs/.

\chapter{Glossary: every name, A to Z}
\label{app:glossary}

Every other listing in this manual and in the browser reference is grouped by
something: operators by class, structures by structure, tactics by task.  This
one has the ordering nobody has to learn.  It carries the whole
\emph{vocabulary} -- 516 names: 16 structure, 7 refinement, 8 instance, 17 view, 163 predicate, 112 functoid, 25 accessor, 18 defined-fn, 60 operator, 8 primitive, 82 tactic -- and no results: the theorems, axioms and
supports are catalogued in \texttt{THEOREMS.md}, \texttt{BY-OPERATOR.md} and
\texttt{PROOF-DEBT.md}, and each entry here says how many of them mention it.

At the REPL, \texttt{(glossary 'NAME)} prints one entry and
\texttt{(glossary "pat")} every name containing \texttt{pat}.

\begingroup
\small
\setlength{\parindent}{0pt}

\section*{Symbols}
\addcontentsline{toc}{section}{Symbols}

\noindent\textbf{\texttt{*}}\quad\textit{operator}\\
a kernel term-former \\
mentioned by 238 result(s)

\smallskip

\noindent\textbf{\texttt{+}}\quad\textit{operator}\\
a kernel term-former \\
mentioned by 190 result(s)

\smallskip

\noindent\textbf{\texttt{-}}\quad\textit{operator}\\
a kernel term-former \\
mentioned by 131 result(s)

\smallskip

\noindent\textbf{\texttt{/}}\quad\textit{operator}\\
a kernel term-former \\
mentioned by 16 result(s)

\smallskip

\noindent\textbf{\texttt{<(x, y)}}\quad\textit{predicate}\\
reads: \$1 is less than \$2 \\
defined by \texttt{<} \\
mentioned by 129 result(s)

\smallskip

\noindent\textbf{\texttt{<=}}\quad\textit{primitive}\\
reads: \$1 is at most \$2 \\
a kernel relation \\
mentioned by 300 result(s)

\smallskip

\noindent\textbf{\texttt{<=\_ord}}\quad\textit{predicate}\\
a predicate \\
mentioned by 16 result(s)

\smallskip

\noindent\textbf{\texttt{<\_ord}}\quad\textit{predicate}\\
a predicate \\
mentioned by 13 result(s)

\smallskip

\noindent\textbf{\texttt{=}}\quad\textit{primitive}\\
reads: \$1 equals \$2 \\
a kernel relation

\smallskip

\noindent\textbf{\texttt{==}}\quad\textit{primitive}\\
reads: \$1 is identical to \$2 \\
a kernel relation

\smallskip

\noindent\textbf{\texttt{>}}\quad\textit{primitive}\\
reads: \$1 is greater than \$2 \\
a kernel relation

\smallskip

\noindent\textbf{\texttt{>=}}\quad\textit{primitive}\\
reads: \$1 is at least \$2 \\
a kernel relation

\smallskip


\section*{A}
\addcontentsline{toc}{section}{A}

\noindent\textbf{\texttt{abelian-group}}\quad\textit{structure}\\
predicate \texttt{is-abelian-group}; slots: carr opr iden inv

\smallskip

\noindent\textbf{\texttt{abelian-group-as-monoid(s)}}\quad\textit{view}\\
an abelian-group viewed as a monoid \\
mentioned by 49 result(s)

\smallskip

\noindent\textbf{\texttt{abs}}\quad\textit{operator}\\
a kernel term-former \\
mentioned by 60 result(s)

\smallskip

\noindent\textbf{\texttt{act(s)}}\quad\textit{accessor}\\
a structure slot \\
mentioned by 95 result(s)

\smallskip

\noindent\textbf{\texttt{add(s)}}\quad\textit{accessor}\\
a structure slot \\
mentioned by 420 result(s)

\smallskip

\noindent\textbf{\texttt{ai}}\quad\textit{tactic}\\
\texttt{(ai hyp)} -- Antecedent inference: decompose a cited assumption -- AND-split, OR-into-cases, or FORSOME-elimination to a fresh eigenvariable. \\
declared in \texttt{tactics-help.scm}

\smallskip

\noindent\textbf{\texttt{apply-functoid}}\quad\textit{operator}\\
a kernel term-former

\smallskip

\noindent\textbf{\texttt{arith}}\quad\textit{tactic}\\
\texttt{(arith)} -- Discharge a ground arithmetic goal by evaluation. \\
declared in \texttt{tactics-help.scm}

\smallskip

\noindent\textbf{\texttt{ass}}\quad\textit{tactic}\\
\texttt{(ass)} -- Close the goal by an assumption alpha-equivalent to it. \\
declared in \texttt{tactics-help.scm}

\smallskip

\noindent\textbf{\texttt{ass-all-frontier!}}\quad\textit{tactic}\\
\texttt{(ass-all-frontier!)} -- Close every frontier leaf whose goal is already among its assumptions.  [proof-local] \\
declared in \texttt{tactics-help.scm}

\smallskip


\section*{B}
\addcontentsline{toc}{section}{B}

\noindent\textbf{\texttt{ball(s, c, r)}}\quad\textit{functoid}\\
= \{y in pts(s): (dist(s))(c, y) <= r and not((dist(s))(c, y) = r)\} \\
mentioned by 20 result(s)

\smallskip

\noindent\textbf{\texttt{ball-cover(s, r)}}\quad\textit{functoid}\\
= image(vnb-lambda(c, pts(s), ball(s, c, r)), pts(s)) \\
mentioned by 4 result(s)

\smallskip

\noindent\textbf{\texttt{bc}}\quad\textit{tactic}\\
\texttt{(bc impl)} -- Backchain the goal through an (IMPLIES A B) already in context: if the goal matches B, the new goal is A. \\
declared in \texttt{tactics-help.scm}

\smallskip

\noindent\textbf{\texttt{bc*}}\quad\textit{tactic}\\
\texttt{(bc* 'thm [((v val)...)] h1 h2 ...)} -- Matching backchain: unify the theorem's conclusion against the goal, then leave its (instantiated) antecedents as subgoals -- optional handlers hk run on the k-th subgoal. \\
declared in \texttt{tactics-help.scm}

\smallskip

\noindent\textbf{\texttt{bdd-metric(s)}}\quad\textit{functoid}\\
= [pts(s), vnb-lambda([u, v], cartesian(pts(s), pts(s)), /((dist(s))(u, v), 1 + (dist(s))(u, v)))] \\
mentioned by 11 result(s)

\smallskip

\noindent\textbf{\texttt{beta}}\quad\textit{tactic}\\
\texttt{(beta)} -- Beta-reduce a functoid application in the goal. \\
declared in \texttt{tactics-help.scm}

\smallskip

\noindent\textbf{\texttt{big-union}}\quad\textit{operator}\\
a kernel term-former \\
mentioned by 14 result(s)

\smallskip

\noindent\textbf{\texttt{bijection}}\quad\textit{operator}\\
a kernel term-former \\
mentioned by 36 result(s)

\smallskip

\noindent\textbf{\texttt{binneg}}\quad\textit{operator}\\
a kernel term-former \\
mentioned by 6 result(s)

\smallskip

\noindent\textbf{\texttt{binplus}}\quad\textit{operator}\\
a kernel term-former \\
mentioned by 14 result(s)

\smallskip

\noindent\textbf{\texttt{bintimes}}\quad\textit{operator}\\
a kernel term-former \\
mentioned by 10 result(s)

\smallskip

\noindent\textbf{\texttt{block(p, k, l)}}\quad\textit{functoid}\\
= matof(k, l, vnb-lambda([i, j], cartesian(interval(1, k), interval(1, l)), entry(p, i, j))) \\
mentioned by 7 result(s)

\smallskip

\noindent\textbf{\texttt{border(a, b, m, p, q)}}\quad\textit{functoid}\\
= matof(succ(p), succ(q), vnb-lambda([i, j], cartesian(interval(1, succ(p)), interval(1, succ(q))), if(i = 1, if ... \\
mentioned by 22 result(s)

\smallskip

\noindent\textbf{\texttt{bu-me}}\quad\textit{tactic}\\
\texttt{(bu-me hyp)} -- Big-union membership elim. \\
declared in \texttt{tactics-help.scm}

\smallskip

\noindent\textbf{\texttt{bu-mi}}\quad\textit{tactic}\\
\texttt{(bu-mi w)} -- Big-union membership intro via the index witness w. \\
declared in \texttt{tactics-help.scm}

\smallskip

\noindent\textbf{\texttt{bu-set}}\quad\textit{tactic}\\
\texttt{(bu-set)} -- Big-union sethood. \\
declared in \texttt{tactics-help.scm}

\smallskip


\section*{C}
\addcontentsline{toc}{section}{C}

\noindent\textbf{\texttt{calc}}\quad\textit{tactic}\\
\texttt{(calc L0 (rel1 L1 [just1]) (rel2 L2 [just2]) ...)} -- Ground the focus goal (REL L0 Ln) by a CHAIN of intermediaries L0 rel1 L1 rel2 L2 ... reln Ln: proves each link and composes them into the endpoint.  A link no lane can close is LEFT OPEN as a leaf -- the refinement point, and the PSS candidate.  Relations: = == (folded by transitivity), < <= (folded through the co-*-trans lemmas, with = steps riding along), IFF (chained per direction).  RETURNS the links, the open ones, and their PSS candidates as an alist. \\
declared in \texttt{tactics-help.scm}

\smallskip

\noindent\textbf{\texttt{card}}\quad\textit{operator}\\
a kernel term-former \\
mentioned by 153 result(s)

\smallskip

\noindent\textbf{\texttt{carr(s)}}\quad\textit{accessor}\\
a structure slot \\
mentioned by 1072 result(s)

\smallskip

\noindent\textbf{\texttt{cartesian}}\quad\textit{operator}\\
a kernel term-former \\
mentioned by 95 result(s)

\smallskip

\noindent\textbf{\texttt{cauchy-setoid(m)}}\quad\textit{functoid}\\
= [cseq(m), crel(m)] \\
mentioned by 1 result(s)

\smallskip

\noindent\textbf{\texttt{cc-ms}}\quad\textit{instance}\\
an instance of metric-space; predicate \texttt{is-cc-ms}; same shape as metric-space

\smallskip

\noindent\textbf{\texttt{cc-normed-field}}\quad\textit{instance}\\
an instance of normed-field; predicate \texttt{is-cc-normed-field}; same shape as normed-field

\smallskip

\noindent\textbf{\texttt{ccint(a, b)}}\quad\textit{functoid}\\
= \{x in rr: a <= x and x <= b\} \\
mentioned by 27 result(s)

\smallskip

\noindent\textbf{\texttt{ce}}\quad\textit{tactic}\\
\texttt{(ce hyp k)} -- Cartesian elim: project the k-th component out of a CARTESIAN-membership assumption. \\
declared in \texttt{tactics-help.scm}

\smallskip

\noindent\textbf{\texttt{centre-set(s, r, f)}}\quad\textit{functoid}\\
= image(vnb-lambda(b, f, choice(centres(s, b, r))), f) \\
mentioned by 2 result(s)

\smallskip

\noindent\textbf{\texttt{centres(s, b, r)}}\quad\textit{functoid}\\
= \{c in pts(s): ball(s, c, r) = b\} \\
mentioned by 6 result(s)

\smallskip

\noindent\textbf{\texttt{choice}}\quad\textit{operator}\\
a kernel term-former \\
mentioned by 9 result(s)

\smallskip

\noindent\textbf{\texttt{choose(n, m)}}\quad\textit{functoid}\\
= card(choose-set(n, m)) \\
mentioned by 9 result(s)

\smallskip

\noindent\textbf{\texttt{choose-set(n, m)}}\quad\textit{functoid}\\
= \{a in power(ord-segment(n)): card(a) = m\}

\smallskip

\noindent\textbf{\texttt{ci}}\quad\textit{tactic}\\
\texttt{(ci)} -- Cartesian intro: prove a CARTESIAN-product membership component-wise. \\
declared in \texttt{tactics-help.scm}

\smallskip

\noindent\textbf{\texttt{class(s, a)}}\quad\textit{functoid}\\
= \{b in pts(s): related(s, a, b)\} \\
mentioned by 25 result(s)

\smallskip

\noindent\textbf{\texttt{closure(s, a)}}\quad\textit{functoid}\\
= \{x in pts(s): forall([u], is-open(s, u) and x in u implies forsome([y in u], y in a))\} \\
mentioned by 2 result(s)

\smallskip

\noindent\textbf{\texttt{cluster-point(s, f, x)}}\quad\textit{predicate}\\
reads: \$3 is a cluster point of the sequence \$2 in \$1 \\
defined by \texttt{cluster-point} \\
mentioned by 6 result(s)

\smallskip

\noindent\textbf{\texttt{comb-kk}}\quad\textit{defined-fn}\\
defined by \texttt{comb-kk-zero} \texttt{comb-kk-succ} \\
mentioned by 13 result(s)

\smallskip

\noindent\textbf{\texttt{comm-monoid}}\quad\textit{structure}\\
predicate \texttt{is-comm-monoid}; slots: carr opr iden

\smallskip

\noindent\textbf{\texttt{commutative-ring}}\quad\textit{refinement}\\
predicate \texttt{is-commutative-ring}; same shape as ring

\smallskip

\noindent\textbf{\texttt{commutative-ring-additive-ag(s)}}\quad\textit{view}\\
a commutative-ring viewed as an abelian-group \\
mentioned by 20 result(s)

\smallskip

\noindent\textbf{\texttt{commutative-ring-multiplicative-cm(s)}}\quad\textit{view}\\
a commutative-ring viewed as a comm-monoid \\
mentioned by 5 result(s)

\smallskip

\noindent\textbf{\texttt{comp}}\quad\textit{operator}\\
a kernel term-former

\smallskip

\noindent\textbf{\texttt{comp-me}}\quad\textit{tactic}\\
\texttt{(comp-me hyp)} -- Comprehension membership elim. \\
declared in \texttt{tactics-help.scm}

\smallskip

\noindent\textbf{\texttt{comp-mi}}\quad\textit{tactic}\\
\texttt{(comp-mi)} -- Comprehension membership intro. \\
declared in \texttt{tactics-help.scm}

\smallskip

\noindent\textbf{\texttt{complement-in}}\quad\textit{operator}\\
a kernel term-former \\
mentioned by 11 result(s)

\smallskip

\noindent\textbf{\texttt{completion(m)}}\quad\textit{functoid}\\
= [quotient(cauchy-setoid(m)), completion-dist(m)] \\
mentioned by 5 result(s)

\smallskip

\noindent\textbf{\texttt{completion-dist(m)}}\quad\textit{functoid}\\
= vnb-lambda(p, cartesian(quotient(cauchy-setoid(m)), quotient(cauchy-setoid(m))), iota(dval, forsome([f, g], f  ...

\smallskip

\noindent\textbf{\texttt{compose(f, g)}}\quad\textit{functoid}\\
= vnb-lambda(z\_, dom(g), f(g(z\_))) \\
mentioned by 11 result(s)

\smallskip

\noindent\textbf{\texttt{conjugate}}\quad\textit{operator}\\
a kernel term-former \\
mentioned by 13 result(s)

\smallskip

\noindent\textbf{\texttt{converges(s, f)}}\quad\textit{predicate}\\
reads: \$2 converges in \$1 \\
defined by \texttt{converges} \\
mentioned by 14 result(s)

\smallskip

\noindent\textbf{\texttt{converges-along(s, g, b, p)}}\quad\textit{predicate}\\
reads: \$2 converges to \$4 along \$3 in \$1 \\
defined by \texttt{converges-along} \\
mentioned by 4 result(s)

\smallskip

\noindent\textbf{\texttt{converges-on(s, fam, dseq)}}\quad\textit{predicate}\\
reads: \$2 converges at every point of the sequence \$3 in \$1 \\
defined by \texttt{converges-on} \\
mentioned by 4 result(s)

\smallskip

\noindent\textbf{\texttt{converges-to(s, f, l)}}\quad\textit{predicate}\\
reads: \$2 converges to \$3 in \$1 \\
defined by \texttt{converges-to} \\
mentioned by 19 result(s)

\smallskip

\noindent\textbf{\texttt{converges-uniformly(s, seq, g)}}\quad\textit{predicate}\\
reads: \$2 converges uniformly to \$3 on \$1 \\
defined by \texttt{converges-uniformly} \\
mentioned by 5 result(s)

\smallskip

\noindent\textbf{\texttt{cos}}\quad\textit{operator}\\
a kernel term-former

\smallskip

\noindent\textbf{\texttt{crel(m)}}\quad\textit{functoid}\\
= \{p in cartesian(cseq(m), cseq(m)): cseq-equiv(m, nth(1, p), nth(2, p))\}

\smallskip

\noindent\textbf{\texttt{crs}}\quad\textit{tactic}\\
\texttt{(crs)} -- Commutative-ring decision procedure: prove a polynomial identity over ZZ[generators].  Expands literal powers, so (x+y)\textasciicircum{}2 = ... closes directly. \\
declared in \texttt{tactics-help.scm}

\smallskip

\noindent\textbf{\texttt{cseq(m)}}\quad\textit{functoid}\\
= \{f in fun(nn, pts(m)): is-cauchy-seq(m, f)\}

\smallskip

\noindent\textbf{\texttt{cseq-equiv(m, f, g)}}\quad\textit{predicate}\\
reads: \$2 and \$3 are equivalent Cauchy sequences in \$1 \\
defined by \texttt{cseq-equiv} \\
mentioned by 2 result(s)

\smallskip

\noindent\textbf{\texttt{cut}}\quad\textit{tactic}\\
\texttt{(cut formula)} -- Cut: prove \texttt{formula' as a side subgoal, then continue the main goal with }formula' added as an assumption (Gentzen cut). \\
declared in \texttt{tactics-help.scm}

\smallskip

\noindent\textbf{\texttt{cut-mem!}}\quad\textit{tactic}\\
\texttt{(cut-mem! mem A)} -- Prove a membership (IN (f x) B) by fun-apply-type with domain A, leaving it in context.  [proof-local] \\
declared in \texttt{tactics-help.scm}

\smallskip


\section*{D}
\addcontentsline{toc}{section}{D}

\noindent\textbf{\texttt{delete-at}}\quad\textit{operator}\\
a kernel term-former \\
mentioned by 6 result(s)

\smallskip

\noindent\textbf{\texttt{deriv(f, a)}}\quad\textit{functoid}\\
= iota(l, is-diff-at(f, a, l)) \\
mentioned by 4 result(s)

\smallskip

\noindent\textbf{\texttt{deriv-v(m, f, a)}}\quad\textit{functoid}\\
= iota(l, is-diff-at-v(m, f, a, l)) \\
mentioned by 2 result(s)

\smallskip

\noindent\textbf{\texttt{descend(s, f)}}\quad\textit{functoid}\\
= vnb-lambda(c, quotient(s), iota(z, forsome([a in c], z = f(a)))) \\
mentioned by 5 result(s)

\smallskip

\noindent\textbf{\texttt{descend2(s, f)}}\quad\textit{functoid}\\
= vnb-lambda([c, d], cartesian(quotient(s), quotient(s)), iota(z, forsome([a in c, b in d], z = f(a, b)))) \\
mentioned by 7 result(s)

\smallskip

\noindent\textbf{\texttt{det}}\quad\textit{primitive}\\
reads: det(\$3) \\
defined by \texttt{det-zero} \texttt{det-cofactor} \\
mentioned by 15 result(s)

\smallskip

\noindent\textbf{\texttt{detach!}}\quad\textit{tactic}\\
\texttt{(detach! impl)} -- Forward modus ponens: from an in-context (IMPLIES A B) whose A is also in context, leave B in context. \\
declared in \texttt{tactics-help.scm}

\smallskip

\noindent\textbf{\texttt{di}}\quad\textit{tactic}\\
\texttt{(di)} -- Direct inference: split an AND goal, move an IMPLIES antecedent into the assumptions, or introduce a fresh eigenvariable for a leading FORALL. \\
declared in \texttt{tactics-help.scm}

\smallskip

\noindent\textbf{\texttt{difference}}\quad\textit{operator}\\
a kernel term-former \\
mentioned by 13 result(s)

\smallskip

\noindent\textbf{\texttt{dist(s)}}\quad\textit{accessor}\\
a structure slot \\
mentioned by 76 result(s)

\smallskip

\noindent\textbf{\texttt{dist-seq(m, f, g)}}\quad\textit{functoid}\\
= vnb-lambda(n, nn, (dist(m))(f(n), g(n))) \\
mentioned by 2 result(s)

\smallskip

\noindent\textbf{\texttt{dom}}\quad\textit{operator}\\
a kernel term-former \\
mentioned by 18 result(s)

\smallskip

\noindent\textbf{\texttt{dual-norm(m, f)}}\quad\textit{functoid}\\
= iota(c\_, c\_ in rr and 0 <= c\_ and forall([x\_ in vec(m)], abs(f(x\_)) <= c\_ * (vnrm(m))(x\_)) and forall([d\_], d\_ ... \\
mentioned by 6 result(s)

\smallskip

\noindent\textbf{\texttt{dual-norm-on(m, s, f)}}\quad\textit{functoid}\\
= iota(c\_, c\_ in rr and 0 <= c\_ and forall([x\_ in s], abs(f(x\_)) <= c\_ * (vnrm(m))(x\_)) and forall([d\_], d\_ in r ... \\
mentioned by 8 result(s)

\smallskip


\section*{E}
\addcontentsline{toc}{section}{E}

\noindent\textbf{\texttt{elem-f(a, n, k, l)}}\quad\textit{functoid}\\
= matof(n, n, vnb-lambda([i, j], cartesian(interval(1, n), interval(1, n)), if(i = k and j = l or i = l and j =  ... \\
mentioned by 42 result(s)

\smallskip

\noindent\textbf{\texttt{elem-g(a, n, r, k, l)}}\quad\textit{functoid}\\
= matof(n, n, vnb-lambda([i, j], cartesian(interval(1, n), interval(1, n)), if(i = j, one(a), if(i = k and j = l ... \\
mentioned by 40 result(s)

\smallskip

\noindent\textbf{\texttt{elem-h(a, n, r, k)}}\quad\textit{functoid}\\
= matof(n, n, vnb-lambda([i, j], cartesian(interval(1, n), interval(1, n)), if(i = j, if(i = k, r, one(a)), zero ... \\
mentioned by 17 result(s)

\smallskip

\noindent\textbf{\texttt{embed(m)}}\quad\textit{functoid}\\
= vnb-lambda(u, pts(m), class(cauchy-setoid(m), embed-seq(m, u))) \\
mentioned by 3 result(s)

\smallskip

\noindent\textbf{\texttt{embed-seq(m, u)}}\quad\textit{functoid}\\
= vnb-lambda(n, nn, u)

\smallskip

\noindent\textbf{\texttt{entry(m, i, j)}}\quad\textit{functoid}\\
= nth(j, nth(i, m)) \\
mentioned by 230 result(s)

\smallskip

\noindent\textbf{\texttt{enum-fam(ag, f, phi, n)}}\quad\textit{functoid}\\
= vnb-lambda(i, nn, if(i in ord-segment(n), f(phi(i)), iden(ag))) \\
mentioned by 5 result(s)

\smallskip

\noindent\textbf{\texttt{eplus}}\quad\textit{operator}\\
a kernel term-former \\
mentioned by 6 result(s)

\smallskip

\noindent\textbf{\texttt{esum}}\quad\textit{operator}\\
a kernel term-former \\
mentioned by 7 result(s)

\smallskip

\noindent\textbf{\texttt{esup}}\quad\textit{operator}\\
a kernel term-former \\
mentioned by 5 result(s)

\smallskip

\noindent\textbf{\texttt{euclidean-gauges(s)}}\quad\textit{functoid}\\
= \{dg in fun(carr(s), nn): has-div-remainder(s, dg)\} \\
mentioned by 3 result(s)

\smallskip

\noindent\textbf{\texttt{euclidean-ring}}\quad\textit{refinement}\\
predicate \texttt{is-euclidean-ring}; same shape as integral-domain

\smallskip

\noindent\textbf{\texttt{ew}}\quad\textit{tactic}\\
\texttt{(ew term)} -- Existential witness: discharge a FORSOME goal by supplying the witness term. \\
declared in \texttt{tactics-help.scm}

\smallskip

\noindent\textbf{\texttt{exp}}\quad\textit{operator}\\
a kernel term-former

\smallskip

\noindent\textbf{\texttt{extends-on(s, g, f)}}\quad\textit{predicate}\\
reads: \$2 agrees with \$3 on \$1 \\
defined by \texttt{extends-on} \\
mentioned by 9 result(s)

\smallskip


\section*{F}
\addcontentsline{toc}{section}{F}

\noindent\textbf{\texttt{fact}}\quad\textit{tactic}\\
\texttt{(fact 'thm term ...)} -- Forward APPLICATION of a theorem: bring it in, instantiate its leading universals with the terms, and auto-detach every antecedent already in context, landing the consequent as a hypothesis. \\
declared in \texttt{tactics-help.scm}

\smallskip

\noindent\textbf{\texttt{factorial}}\quad\textit{defined-fn}\\
defined by \texttt{factorial-zero} \texttt{factorial-succ} \\
mentioned by 20 result(s)

\smallskip

\noindent\textbf{\texttt{falling}}\quad\textit{defined-fn}\\
defined by \texttt{falling-zero} \texttt{falling-succ} \\
mentioned by 9 result(s)

\smallskip

\noindent\textbf{\texttt{fam-of-list}}\quad\textit{defined-fn}\\
defined by \texttt{fam-of-list-apply} \\
mentioned by 18 result(s)

\smallskip

\noindent\textbf{\texttt{field}}\quad\textit{structure}\\
predicate \texttt{is-field}; slots: carr add mul neg zero one non-zero recip

\smallskip

\noindent\textbf{\texttt{field-additive-ag(s)}}\quad\textit{view}\\
a field viewed as an abelian-group \\
mentioned by 3 result(s)

\smallskip

\noindent\textbf{\texttt{field-as-euclidean-ring(s)}}\quad\textit{view}\\
a field viewed as an euclidean-ring \\
mentioned by 4 result(s)

\smallskip

\noindent\textbf{\texttt{field-as-integral-domain(s)}}\quad\textit{view}\\
a field viewed as an integral-domain \\
mentioned by 2 result(s)

\smallskip

\noindent\textbf{\texttt{field-multiplicative-group(s)}}\quad\textit{view}\\
a field viewed as a group \\
mentioned by 1 result(s)

\smallskip

\noindent\textbf{\texttt{fin-enum(s)}}\quad\textit{functoid}\\
= choice(bijection(ord-segment(card(s)), s)) \\
mentioned by 1 result(s)

\smallskip

\noindent\textbf{\texttt{finprod(m, f, s)}}\quad\textit{functoid}\\
= finsum(m, f, s)

\smallskip

\noindent\textbf{\texttt{finsum(ag, f, s)}}\quad\textit{functoid}\\
= sum-ag(ag, enum-fam(ag, f, fin-enum(s), card(s)), card(s)) \\
mentioned by 145 result(s)

\smallskip

\noindent\textbf{\texttt{finsupp(a, m)}}\quad\textit{functoid}\\
reads: the finitely-supported functions \$2 -> \$1 \\
= \{f in fun(carr(m), carr(a)): card(supp(a, m, f)) in nn\} \\
mentioned by 8 result(s)

\smallskip

\noindent\textbf{\texttt{fnrm(s)}}\quad\textit{accessor}\\
a structure slot \\
mentioned by 13 result(s)

\smallskip

\noindent\textbf{\texttt{focus}}\quad\textit{tactic}\\
\texttt{(focus n)} -- Switch the focus to the n-th open goal (1-based). \\
declared in \texttt{tactics-help.scm}

\smallskip

\noindent\textbf{\texttt{focus-leaf!}}\quad\textit{tactic}\\
\texttt{(focus-leaf! substr)} -- Focus the frontier leaf whose goal contains substr (never trust auto-advance).  [proof-local] \\
declared in \texttt{tactics-help.scm}

\smallskip

\noindent\textbf{\texttt{fr-ball(m, fam, n, eps)}}\quad\textit{functoid}\\
= \{x in vec(m): forall([k], k in nn and k <= n implies (fam(k))(x) < eps)\} \\
mentioned by 2 result(s)

\smallskip

\noindent\textbf{\texttt{fr-cauchy(m, fam, seq)}}\quad\textit{predicate}\\
reads: \$3 is Cauchy in (\$1, \$2) \\
defined by \texttt{fr-cauchy} \\
mentioned by 4 result(s)

\smallskip

\noindent\textbf{\texttt{fr-conv(m, fam, seq, lim)}}\quad\textit{predicate}\\
reads: \$3 converges to \$4 in (\$1, \$2) \\
defined by \texttt{fr-conv} \\
mentioned by 8 result(s)

\smallskip

\noindent\textbf{\texttt{fun}}\quad\textit{operator}\\
a kernel term-former \\
mentioned by 506 result(s)

\smallskip


\section*{G}
\addcontentsline{toc}{section}{G}

\noindent\textbf{\texttt{gauge(s)}}\quad\textit{functoid}\\
= choice(euclidean-gauges(s)) \\
mentioned by 15 result(s)

\smallskip

\noindent\textbf{\texttt{generates(md, n, u)}}\quad\textit{predicate}\\
reads: the \$2 vectors \$3 generate \$1 \\
defined by \texttt{generates} \\
mentioned by 8 result(s)

\smallskip

\noindent\textbf{\texttt{goal-status}}\quad\textit{tactic}\\
\texttt{(goal-status)} -- One-line summary: done / N open goals. \\
declared in \texttt{tactics-help.scm}

\smallskip

\noindent\textbf{\texttt{good-sub(m, s, f, t)}}\quad\textit{predicate}\\
reads: \$3 has a norm-preserving extension from \$2 to \$4, in \$1 \\
defined by \texttt{good-sub} \\
mentioned by 7 result(s)

\smallskip

\noindent\textbf{\texttt{greedy-chain(phi, grd, porel)}}\quad\textit{functoid}\\
= zkept(phi, grd, porel, card(grd)) \\
mentioned by 1 result(s)

\smallskip

\noindent\textbf{\texttt{grind}}\quad\textit{tactic}\\
\texttt{(grind)} -- Saturate the no-choice moves at the focus: decompose the goal connective (di) and break open the hypotheses (mac-h*), repeatedly, until neither fires.  The whole introduce-everything prefix in one step. \\
declared in \texttt{tactics-help.scm}

\smallskip

\noindent\textbf{\texttt{group}}\quad\textit{structure}\\
predicate \texttt{is-group}; slots: carr opr iden inv

\smallskip


\section*{H}
\addcontentsline{toc}{section}{H}

\noindent\textbf{\texttt{has-closed-graph(m1, fam1, m2, fam2, tt)}}\quad\textit{predicate}\\
reads: \$5 has a closed graph from (\$1, \$2) to (\$3, \$4) \\
defined by \texttt{has-closed-graph} \\
mentioned by 3 result(s)

\smallskip

\noindent\textbf{\texttt{has-div-remainder(s, deg)}}\quad\textit{predicate}\\
reads: \$1 has division with remainder for the degree function \$2 \\
defined by \texttt{has-div-remainder} \\
mentioned by 8 result(s)

\smallskip

\noindent\textbf{\texttt{has-fip(s, c)}}\quad\textit{predicate}\\
reads: \$2 has the finite intersection property in \$1 \\
defined by \texttt{has-fip} \\
mentioned by 4 result(s)

\smallskip

\noindent\textbf{\texttt{has-inverses(op, unit, invop, crr)}}\quad\textit{predicate}\\
reads: every element of \$4 has an inverse under \$1, given by \$3, with unit \$2 \\
defined by \texttt{has-inverses} \\
mentioned by 20 result(s)

\smallskip

\noindent\textbf{\texttt{have!}}\quad\textit{tactic}\\
\texttt{(have! CLAIM [THUNK])} -- Assert CLAIM as an intermediate step and carry on with it as a hypothesis: cut CLAIM, discharge the side goal from context (or by THUNK), and stay on the main branch. \\
declared in \texttt{tactics-help.scm}

\smallskip


\section*{I}
\addcontentsline{toc}{section}{I}

\noindent\textbf{\texttt{iden(s)}}\quad\textit{accessor}\\
a structure slot \\
mentioned by 102 result(s)

\smallskip

\noindent\textbf{\texttt{identmat(a, n)}}\quad\textit{functoid}\\
= matof(n, n, vnb-lambda([i, j], cartesian(interval(1, n), interval(1, n)), if(i = j, one(a), zero(a)))) \\
mentioned by 27 result(s)

\smallskip

\noindent\textbf{\texttt{idl(s)}}\quad\textit{accessor}\\
a structure slot \\
mentioned by 12 result(s)

\smallskip

\noindent\textbf{\texttt{ie}}\quad\textit{tactic}\\
\texttt{(ie hyp k)} -- Intersection elim: extract the k-th branch of an INTERSECTION-membership assumption. \\
declared in \texttt{tactics-help.scm}

\smallskip

\noindent\textbf{\texttt{if}}\quad\textit{operator}\\
a kernel term-former \\
mentioned by 46 result(s)

\smallskip

\noindent\textbf{\texttt{if-false}}\quad\textit{tactic}\\
\texttt{(if-false t)} -- Reduce an (IF p a b) term on the not-p branch: spawns (NOT p); the continuation gains (= (IF p a b) b). \\
declared in \texttt{tactics-help.scm}

\smallskip

\noindent\textbf{\texttt{if-true}}\quad\textit{tactic}\\
\texttt{(if-true t)} -- Reduce an (IF p a b) term on the p branch: spawns p as a subgoal; the continuation gains (= (IF p a b) a). \\
declared in \texttt{tactics-help.scm}

\smallskip

\noindent\textbf{\texttt{ii}}\quad\textit{tactic}\\
\texttt{(ii)} -- Intersection intro: prove (IN x (INTERSECTION ...)) for each branch. \\
declared in \texttt{tactics-help.scm}

\smallskip

\noindent\textbf{\texttt{imag-part}}\quad\textit{operator}\\
a kernel term-former

\smallskip

\noindent\textbf{\texttt{image}}\quad\textit{operator}\\
a kernel term-former \\
mentioned by 13 result(s)

\smallskip

\noindent\textbf{\texttt{in}}\quad\textit{primitive}\\
reads: \$1 is in \$2 \\
a kernel relation \\
mentioned by 2576 result(s)

\smallskip

\noindent\textbf{\texttt{ineq}}\quad\textit{tactic}\\
\texttt{(ineq i1 i2 ...)} -- Close a linear-inequality goal over RR (<= < > >= = between RR terms) as a consequence of the named assumptions (1-based indices), via the Fourier-Motzkin/Farkas oracle.  Linearizes over + - * and the binplus/binneg/bintimes aliases; every MAXIMAL non-arithmetic subterm is an atom that must be certified in RR.  (Does NOT see through a generic ring's (ADD s)/(MUL s) -- those become opaque atoms.) \\
declared in \texttt{tactics-help.scm}

\smallskip

\noindent\textbf{\texttt{inf-subsets(a)}}\quad\textit{functoid}\\
= \{s in power(a): not(card(s) in nn)\} \\
mentioned by 13 result(s)

\smallskip

\noindent\textbf{\texttt{injection}}\quad\textit{operator}\\
a kernel term-former \\
mentioned by 16 result(s)

\smallskip

\noindent\textbf{\texttt{injective*(f)}}\quad\textit{predicate}\\
reads: \$1 is injective \\
defined by \texttt{injective*} \\
mentioned by 3 result(s)

\smallskip

\noindent\textbf{\texttt{insert-last(phi, x, n)}}\quad\textit{functoid}\\
= vnb-lambda(i, nn, if(i in ord-segment(n), phi(i), x))

\smallskip

\noindent\textbf{\texttt{inst}}\quad\textit{tactic}\\
\texttt{(inst forall-hyp term)} -- Instantiate a universally-quantified assumption at `term', adding the instance to context. \\
declared in \texttt{tactics-help.scm}

\smallskip

\noindent\textbf{\texttt{inst+}}\quad\textit{tactic}\\
\texttt{(inst+ forall-hyp term)} -- Instantiate an in-context universal at `term', then forward-detach any guards whose antecedents are in context. \\
declared in \texttt{tactics-help.scm}

\smallskip

\noindent\textbf{\texttt{integral-domain}}\quad\textit{refinement}\\
predicate \texttt{is-integral-domain}; same shape as commutative-ring

\smallskip

\noindent\textbf{\texttt{interior(s, a)}}\quad\textit{functoid}\\
= \{x in pts(s): forsome([u], is-open(s, u) and x in u and u subset a)\} \\
mentioned by 2 result(s)

\smallskip

\noindent\textbf{\texttt{intersection}}\quad\textit{operator}\\
a kernel term-former \\
mentioned by 22 result(s)

\smallskip

\noindent\textbf{\texttt{interval(a, b)}}\quad\textit{functoid}\\
= \{i in nn: a <= i and i <= b\} \\
mentioned by 218 result(s)

\smallskip

\noindent\textbf{\texttt{inv(s)}}\quad\textit{accessor}\\
a structure slot \\
mentioned by 29 result(s)

\smallskip

\noindent\textbf{\texttt{inverse-bij(phi, x, y)}}\quad\textit{functoid}\\
= vnb-lambda(y\_, y, choice(\{x\_ in x: phi(x\_) = y\_\})) \\
mentioned by 6 result(s)

\smallskip

\noindent\textbf{\texttt{iota}}\quad\textit{operator}\\
a kernel term-former \\
mentioned by 2 result(s)

\smallskip

\noindent\textbf{\texttt{iota-d}}\quad\textit{tactic}\\
\texttt{(iota-d term)} -- Definite-description: discharge the IOTA uniqueness obligation for `term'. \\
declared in \texttt{tactics-help.scm}

\smallskip

\noindent\textbf{\texttt{is-abelian-group(s)}}\quad\textit{predicate}\\
reads: an abelian group \\
a predicate \\
mentioned by 106 result(s)

\smallskip

\noindent\textbf{\texttt{is-absolutely-summable(grp, f)}}\quad\textit{predicate}\\
reads: \$2 is absolutely summable in \$1 \\
defined by \texttt{is-absolutely-summable} \\
mentioned by 3 result(s)

\smallskip

\noindent\textbf{\texttt{is-algebra-of-sets(omega, ca)}}\quad\textit{predicate}\\
reads: \$2 is an algebra of subsets of \$1 \\
defined by \texttt{is-algebra-of-sets} \\
mentioned by 2 result(s)

\smallskip

\noindent\textbf{\texttt{is-associative(op, crr)}}\quad\textit{predicate}\\
reads: \$1 is associative on \$2 \\
defined by \texttt{is-associative} \\
mentioned by 26 result(s)

\smallskip

\noindent\textbf{\texttt{is-bounded-linear-functional(m, f)}}\quad\textit{predicate}\\
reads: \$2 is a bounded linear functional on \$1 \\
defined by \texttt{is-bounded-linear-functional} \\
mentioned by 15 result(s)

\smallskip

\noindent\textbf{\texttt{is-bounded-linear-functional-on(m, s, f)}}\quad\textit{predicate}\\
reads: \$3 is a bounded linear functional on the subspace \$2 of \$1 \\
defined by \texttt{is-bounded-linear-functional-on} \\
mentioned by 10 result(s)

\smallskip

\noindent\textbf{\texttt{is-bounded-metric-space(s)}}\quad\textit{predicate}\\
reads: a bounded metric space \\
defined by \texttt{is-bounded-metric-space} \\
mentioned by 7 result(s)

\smallskip

\noindent\textbf{\texttt{is-cauchy-seq(s, f)}}\quad\textit{predicate}\\
reads: Cauchy \\
defined by \texttt{is-cauchy-seq} \\
mentioned by 11 result(s)

\smallskip

\noindent\textbf{\texttt{is-chain(grd, porel, ch)}}\quad\textit{predicate}\\
reads: \$3 is a chain in \$1 under \$2 \\
defined by \texttt{is-chain} \\
mentioned by 11 result(s)

\smallskip

\noindent\textbf{\texttt{is-closed(s, a)}}\quad\textit{predicate}\\
reads: \$2 is closed in \$1 \\
defined by \texttt{is-closed} \\
mentioned by 6 result(s)

\smallskip

\noindent\textbf{\texttt{is-comm-monoid(s)}}\quad\textit{predicate}\\
reads: a comm monoid \\
a predicate \\
mentioned by 32 result(s)

\smallskip

\noindent\textbf{\texttt{is-commutative(op, crr)}}\quad\textit{predicate}\\
reads: \$1 is commutative on \$2 \\
defined by \texttt{is-commutative} \\
mentioned by 20 result(s)

\smallskip

\noindent\textbf{\texttt{is-commutative-ring(s)}}\quad\textit{predicate}\\
reads: a commutative ring \\
a predicate \\
mentioned by 108 result(s)

\smallskip

\noindent\textbf{\texttt{is-compact(s)}}\quad\textit{predicate}\\
reads: \$1 is compact \\
defined by \texttt{is-compact} \\
mentioned by 18 result(s)

\smallskip

\noindent\textbf{\texttt{is-complete(s)}}\quad\textit{predicate}\\
reads: complete \\
defined by \texttt{is-complete} \\
mentioned by 12 result(s)

\smallskip

\noindent\textbf{\texttt{is-cont-lin(m1, fam1, m2, fam2, tt)}}\quad\textit{predicate}\\
reads: \$5 is a continuous linear map from (\$1, \$2) to (\$3, \$4) \\
defined by \texttt{is-cont-lin} \\
mentioned by 5 result(s)

\smallskip

\noindent\textbf{\texttt{is-continuous(s, t, f)}}\quad\textit{predicate}\\
reads: \$3 is continuous from \$1 to \$2 \\
defined by \texttt{is-continuous} \\
mentioned by 16 result(s)

\smallskip

\noindent\textbf{\texttt{is-continuous-at(s, t, f, a)}}\quad\textit{predicate}\\
reads: \$3 is continuous at \$4 \\
defined by \texttt{is-continuous-at} \\
mentioned by 45 result(s)

\smallskip

\noindent\textbf{\texttt{is-countable-pseudometric-family(fam, ground)}}\quad\textit{predicate}\\
reads: \$1 is a countable family of pseudometrics on \$2 \\
defined by \texttt{is-countable-pseudometric-family} \\
mentioned by 4 result(s)

\smallskip

\noindent\textbf{\texttt{is-dense-seq(s, dseq)}}\quad\textit{predicate}\\
reads: \$2 is a dense sequence in \$1 \\
defined by \texttt{is-dense-seq} \\
mentioned by 4 result(s)

\smallskip

\noindent\textbf{\texttt{is-diagonal(a, m, n, d)}}\quad\textit{predicate}\\
reads: \$4 is a diagonal \$2-by-\$3 matrix over \$1 \\
defined by \texttt{is-diagonal} \\
mentioned by 7 result(s)

\smallskip

\noindent\textbf{\texttt{is-diff-at(f, a, l)}}\quad\textit{predicate}\\
reads: \$1 is differentiable at \$2, with derivative \$3 \\
defined by \texttt{is-diff-at} \\
mentioned by 36 result(s)

\smallskip

\noindent\textbf{\texttt{is-diff-at-v(m, f, a, l)}}\quad\textit{predicate}\\
reads: \$2 is differentiable at \$3 with derivative \$4, as a curve in \$1 \\
defined by \texttt{is-diff-at-v} \\
mentioned by 4 result(s)

\smallskip

\noindent\textbf{\texttt{is-distributive(addop, mulop, crr)}}\quad\textit{predicate}\\
reads: \$2 distributes over \$1 on \$3 \\
defined by \texttt{is-distributive} \\
mentioned by 10 result(s)

\smallskip

\noindent\textbf{\texttt{is-eps-cauchy-seq(s, eps, y)}}\quad\textit{predicate}\\
reads: \$3 is \$2-Cauchy in \$1 \\
defined by \texttt{is-eps-cauchy-seq} \\
mentioned by 3 result(s)

\smallskip

\noindent\textbf{\texttt{is-equicontinuous(s, t, fam)}}\quad\textit{predicate}\\
reads: \$3 is equicontinuous from \$1 to \$2 \\
defined by \texttt{is-equicontinuous} \\
mentioned by 5 result(s)

\smallskip

\noindent\textbf{\texttt{is-equivalence(rho, crr)}}\quad\textit{predicate}\\
reads: \$1 is an equivalence relation on \$2 \\
defined by \texttt{is-equivalence} \\
mentioned by 5 result(s)

\smallskip

\noindent\textbf{\texttt{is-euclidean-ring(s)}}\quad\textit{predicate}\\
reads: a Euclidean ring \\
a predicate \\
mentioned by 33 result(s)

\smallskip

\noindent\textbf{\texttt{is-field(s)}}\quad\textit{predicate}\\
reads: a field \\
a predicate \\
mentioned by 80 result(s)

\smallskip

\noindent\textbf{\texttt{is-finite-cover(c, a)}}\quad\textit{predicate}\\
reads: \$1 is a finite cover of \$2 \\
defined by \texttt{is-finite-cover} \\
mentioned by 5 result(s)

\smallskip

\noindent\textbf{\texttt{is-finite-dimensional(m)}}\quad\textit{predicate}\\
reads: \$1 is finite dimensional \\
defined by \texttt{is-finite-dimensional} \\
mentioned by 8 result(s)

\smallskip

\noindent\textbf{\texttt{is-frechet-structure(m, fam)}}\quad\textit{predicate}\\
reads: (\$1, \$2) is a Frechet space \\
defined by \texttt{is-frechet-structure} \\
mentioned by 5 result(s)

\smallskip

\noindent\textbf{\texttt{is-fun}}\quad\textit{predicate}\\
a predicate \\
mentioned by 2 result(s)

\smallskip

\noindent\textbf{\texttt{is-gauge-countable(s)}}\quad\textit{predicate}\\
reads: \$1 has a countable pseudometric gauge \\
defined by \texttt{is-gauge-countable} \\
mentioned by 7 result(s)

\smallskip

\noindent\textbf{\texttt{is-group(s)}}\quad\textit{predicate}\\
reads: a group \\
a predicate \\
mentioned by 20 result(s)

\smallskip

\noindent\textbf{\texttt{is-group-norm(nm, op, invop, unit, crr)}}\quad\textit{predicate}\\
reads: \$1 is a group norm for \$2 on \$5 \\
defined by \texttt{is-group-norm} \\
mentioned by 4 result(s)

\smallskip

\noindent\textbf{\texttt{is-hausdorff(s)}}\quad\textit{predicate}\\
reads: \$1 is Hausdorff \\
defined by \texttt{is-hausdorff} \\
mentioned by 6 result(s)

\smallskip

\noindent\textbf{\texttt{is-hom-abelian-group(a, b, f)}}\quad\textit{predicate}\\
reads: \$3 is a homomorphism from \$1 to \$2 \\
a predicate \\
mentioned by 9 result(s)

\smallskip

\noindent\textbf{\texttt{is-hom-comm-monoid(a, b, f)}}\quad\textit{predicate}\\
reads: \$3 is a homomorphism from \$1 to \$2 \\
a predicate \\
mentioned by 3 result(s)

\smallskip

\noindent\textbf{\texttt{is-hom-commutative-ring(a, b, f)}}\quad\textit{predicate}\\
reads: \$3 is a homomorphism from \$1 to \$2 \\
a predicate \\
mentioned by 7 result(s)

\smallskip

\noindent\textbf{\texttt{is-hom-euclidean-ring(a, b, f)}}\quad\textit{predicate}\\
reads: \$3 is a homomorphism from \$1 to \$2 \\
a predicate \\
mentioned by 3 result(s)

\smallskip

\noindent\textbf{\texttt{is-hom-field(a, b, f)}}\quad\textit{predicate}\\
reads: \$3 is a homomorphism from \$1 to \$2 \\
a predicate \\
mentioned by 5 result(s)

\smallskip

\noindent\textbf{\texttt{is-hom-group(a, b, f)}}\quad\textit{predicate}\\
reads: \$3 is a homomorphism from \$1 to \$2 \\
a predicate \\
mentioned by 2 result(s)

\smallskip

\noindent\textbf{\texttt{is-hom-integral-domain(a, b, f)}}\quad\textit{predicate}\\
reads: \$3 is a homomorphism from \$1 to \$2 \\
a predicate \\
mentioned by 8 result(s)

\smallskip

\noindent\textbf{\texttt{is-hom-metric-space(a, b, f)}}\quad\textit{predicate}\\
reads: \$3 is an isometry from \$1 to \$2 \\
a predicate \\
mentioned by 4 result(s)

\smallskip

\noindent\textbf{\texttt{is-hom-metrizable-top-space(s, t, f)}}\quad\textit{predicate}\\
reads: \$3 is continuous from \$1 to \$2 \\
a predicate \\
mentioned by 3 result(s)

\smallskip

\noindent\textbf{\texttt{is-hom-module(a, b, f)}}\quad\textit{predicate}\\
reads: \$3 is a homomorphism from \$1 to \$2 \\
a predicate \\
mentioned by 6 result(s)

\smallskip

\noindent\textbf{\texttt{is-hom-monoid(a, b, f)}}\quad\textit{predicate}\\
reads: \$3 is a homomorphism from \$1 to \$2 \\
a predicate \\
mentioned by 4 result(s)

\smallskip

\noindent\textbf{\texttt{is-hom-normed-ag(a, b, f)}}\quad\textit{predicate}\\
reads: \$3 is a homomorphism from \$1 to \$2 \\
a predicate \\
mentioned by 4 result(s)

\smallskip

\noindent\textbf{\texttt{is-hom-normed-field(a, b, f)}}\quad\textit{predicate}\\
reads: \$3 is a homomorphism from \$1 to \$2 \\
a predicate \\
mentioned by 5 result(s)

\smallskip

\noindent\textbf{\texttt{is-hom-normed-vector-space(a, b, f)}}\quad\textit{predicate}\\
reads: \$3 is a homomorphism from \$1 to \$2 \\
a predicate \\
mentioned by 4 result(s)

\smallskip

\noindent\textbf{\texttt{is-hom-pid(a, b, f)}}\quad\textit{predicate}\\
reads: \$3 is a homomorphism from \$1 to \$2 \\
a predicate \\
mentioned by 2 result(s)

\smallskip

\noindent\textbf{\texttt{is-hom-pseudometric-space(a, b, f)}}\quad\textit{predicate}\\
reads: \$3 is a homomorphism from \$1 to \$2 \\
a predicate \\
mentioned by 2 result(s)

\smallskip

\noindent\textbf{\texttt{is-hom-ring(a, b, f)}}\quad\textit{predicate}\\
reads: \$3 is a homomorphism from \$1 to \$2 \\
a predicate \\
mentioned by 8 result(s)

\smallskip

\noindent\textbf{\texttt{is-hom-ringoid(a, b, f)}}\quad\textit{predicate}\\
reads: \$3 is a ringoid homomorphism from \$1 to \$2 \\
a predicate \\
mentioned by 2 result(s)

\smallskip

\noindent\textbf{\texttt{is-hom-semigroup(a, b, f)}}\quad\textit{predicate}\\
reads: \$3 is a homomorphism from \$1 to \$2 \\
a predicate \\
mentioned by 2 result(s)

\smallskip

\noindent\textbf{\texttt{is-hom-setoid(a, b, f1, f2)}}\quad\textit{predicate}\\
reads: \$3 is a setoid homomorphism from \$1 to \$2 \\
a predicate \\
mentioned by 2 result(s)

\smallskip

\noindent\textbf{\texttt{is-hom-top-space(s, t, f)}}\quad\textit{predicate}\\
reads: \$3 is continuous from \$1 to \$2 \\
a predicate \\
mentioned by 2 result(s)

\smallskip

\noindent\textbf{\texttt{is-hom-vector-space(a, b, f)}}\quad\textit{predicate}\\
reads: \$3 is a homomorphism from \$1 to \$2 \\
a predicate \\
mentioned by 2 result(s)

\smallskip

\noindent\textbf{\texttt{is-ideal(s, i)}}\quad\textit{predicate}\\
reads: \$2 is an ideal of \$1 \\
defined by \texttt{is-ideal} \\
mentioned by 9 result(s)

\smallskip

\noindent\textbf{\texttt{is-ideal-in(i\_, carr\_, add\_, mul\_, neg\_, zero\_)}}\quad\textit{predicate}\\
reads: \$1 is a two-sided ideal of the ring with carrier \$2 \\
defined by \texttt{is-ideal-in} \\
mentioned by 4 result(s)

\smallskip

\noindent\textbf{\texttt{is-identity(op, unit, crr)}}\quad\textit{predicate}\\
reads: \$2 is an identity for \$1 on \$3 \\
defined by \texttt{is-identity} \\
mentioned by 24 result(s)

\smallskip

\noindent\textbf{\texttt{is-integral-domain(s)}}\quad\textit{predicate}\\
reads: an integral domain \\
a predicate \\
mentioned by 19 result(s)

\smallskip

\noindent\textbf{\texttt{is-invertible-mat(a, n, u)}}\quad\textit{predicate}\\
reads: \$3 is an invertible \$2-by-\$2 matrix over \$1 \\
defined by \texttt{is-invertible-mat} \\
mentioned by 17 result(s)

\smallskip

\noindent\textbf{\texttt{is-isometry(s, t, f)}}\quad\textit{predicate}\\
reads: \$3 is an isometry from \$1 to \$2 \\
defined by \texttt{is-isometry} \\
mentioned by 2 result(s)

\smallskip

\noindent\textbf{\texttt{is-k-linear(m, f)}}\quad\textit{predicate}\\
reads: \$2 is a K-linear functional on \$1 \\
defined by \texttt{is-k-linear} \\
mentioned by 3 result(s)

\smallskip

\noindent\textbf{\texttt{is-k-linear-on(m, s, f)}}\quad\textit{predicate}\\
reads: \$3 is a K-linear functional on the submodule \$2 of \$1 \\
defined by \texttt{is-k-linear-on} \\
mentioned by 3 result(s)

\smallskip

\noindent\textbf{\texttt{is-linear-functional(m, f)}}\quad\textit{predicate}\\
reads: \$2 is a linear functional on \$1 \\
defined by \texttt{is-linear-functional} \\
mentioned by 6 result(s)

\smallskip

\noindent\textbf{\texttt{is-linear-functional-on(m, s, f)}}\quad\textit{predicate}\\
reads: \$3 is a linear functional on the subspace \$2 of \$1 \\
defined by \texttt{is-linear-functional-on} \\
mentioned by 12 result(s)

\smallskip

\noindent\textbf{\texttt{is-linear-map(m1, m2, tt)}}\quad\textit{predicate}\\
reads: \$3 is a linear map from \$1 to \$2 \\
defined by \texttt{is-linear-map} \\
mentioned by 9 result(s)

\smallskip

\noindent\textbf{\texttt{is-maximal(grd, porel, mx)}}\quad\textit{predicate}\\
reads: \$3 is maximal in \$1 under \$2 \\
defined by \texttt{is-maximal} \\
mentioned by 4 result(s)

\smallskip

\noindent\textbf{\texttt{is-meager(s, a)}}\quad\textit{predicate}\\
reads: \$2 is meager in \$1 \\
defined by \texttt{is-meager} \\
mentioned by 4 result(s)

\smallskip

\noindent\textbf{\texttt{is-metric(dst, crr)}}\quad\textit{predicate}\\
reads: \$1 is a metric on \$2 \\
defined by \texttt{is-metric} \\
mentioned by 4 result(s)

\smallskip

\noindent\textbf{\texttt{is-metric-space(s)}}\quad\textit{predicate}\\
reads: a metric space \\
a predicate \\
mentioned by 125 result(s)

\smallskip

\noindent\textbf{\texttt{is-metrizable-top-space(s)}}\quad\textit{predicate}\\
reads: \$1 is a metrizable topological space \\
a predicate \\
mentioned by 13 result(s)

\smallskip

\noindent\textbf{\texttt{is-module(s)}}\quad\textit{predicate}\\
reads: a module \\
a predicate \\
mentioned by 146 result(s)

\smallskip

\noindent\textbf{\texttt{is-monoid(s)}}\quad\textit{predicate}\\
reads: a monoid \\
a predicate \\
mentioned by 45 result(s)

\smallskip

\noindent\textbf{\texttt{is-ms-sequence(ms)}}\quad\textit{predicate}\\
reads: \$1 is a sequence of metric spaces \\
defined by \texttt{is-ms-sequence} \\
mentioned by 12 result(s)

\smallskip

\noindent\textbf{\texttt{is-noetherian(m)}}\quad\textit{predicate}\\
reads: \$1 is Noetherian \\
defined by \texttt{is-noetherian} \\
mentioned by 5 result(s)

\smallskip

\noindent\textbf{\texttt{is-nonmeager(s, a)}}\quad\textit{predicate}\\
reads: \$2 is nonmeager in \$1 \\
defined by \texttt{is-nonmeager} \\
mentioned by 3 result(s)

\smallskip

\noindent\textbf{\texttt{is-norm(nm, addop, mulop, zr, crr)}}\quad\textit{predicate}\\
reads: \$1 is a norm on \$5 \\
defined by \texttt{is-norm} \\
mentioned by 4 result(s)

\smallskip

\noindent\textbf{\texttt{is-normed-ag(s)}}\quad\textit{predicate}\\
reads: a normed ag \\
a predicate \\
mentioned by 61 result(s)

\smallskip

\noindent\textbf{\texttt{is-normed-field(s)}}\quad\textit{predicate}\\
reads: a normed field \\
a predicate \\
mentioned by 161 result(s)

\smallskip

\noindent\textbf{\texttt{is-normed-vector-space(s)}}\quad\textit{predicate}\\
reads: a normed vector space \\
a predicate \\
mentioned by 231 result(s)

\smallskip

\noindent\textbf{\texttt{is-nowhere-dense(s, a)}}\quad\textit{predicate}\\
reads: \$2 is nowhere dense in \$1 \\
defined by \texttt{is-nowhere-dense} \\
mentioned by 4 result(s)

\smallskip

\noindent\textbf{\texttt{is-open(s, u)}}\quad\textit{predicate}\\
reads: \$2 is open in \$1 \\
defined by \texttt{is-open} \\
mentioned by 15 result(s)

\smallskip

\noindent\textbf{\texttt{is-open-cover(s, c)}}\quad\textit{predicate}\\
reads: \$2 is an open cover of \$1 \\
defined by \texttt{is-open-cover} \\
mentioned by 7 result(s)

\smallskip

\noindent\textbf{\texttt{is-open-gauge(fam, ground, u)}}\quad\textit{predicate}\\
reads: \$3 is open in the gauge topology of the family \$1 on \$2 \\
defined by \texttt{is-open-gauge} \\
mentioned by 2 result(s)

\smallskip

\noindent\textbf{\texttt{is-open-lin-map(m1, fam1, m2, fam2, tt)}}\quad\textit{predicate}\\
reads: \$5 is an open linear map from (\$1, \$2) to (\$3, \$4) \\
defined by \texttt{is-open-lin-map} \\
mentioned by 3 result(s)

\smallskip

\noindent\textbf{\texttt{is-ord(x)}}\quad\textit{predicate}\\
reads: \$1 is an ordinal \\
a predicate \\
mentioned by 2 result(s)

\smallskip

\noindent\textbf{\texttt{is-partial-order(grd, porel)}}\quad\textit{predicate}\\
reads: \$2 partially orders \$1 \\
defined by \texttt{is-partial-order} \\
mentioned by 13 result(s)

\smallskip

\noindent\textbf{\texttt{is-pid(s)}}\quad\textit{predicate}\\
reads: a pid \\
a predicate \\
mentioned by 6 result(s)

\smallskip

\noindent\textbf{\texttt{is-pseudometric(dst, crr)}}\quad\textit{predicate}\\
reads: \$1 is a pseudometric on \$2 \\
defined by \texttt{is-pseudometric} \\
mentioned by 4 result(s)

\smallskip

\noindent\textbf{\texttt{is-pseudometric-space(s)}}\quad\textit{predicate}\\
reads: a pseudometric space \\
a predicate \\
mentioned by 8 result(s)

\smallskip

\noindent\textbf{\texttt{is-r-net(s, f, a, r)}}\quad\textit{predicate}\\
reads: \$2 is an \$4-net for \$3 in \$1 \\
defined by \texttt{is-r-net} \\
mentioned by 5 result(s)

\smallskip

\noindent\textbf{\texttt{is-ring(s)}}\quad\textit{predicate}\\
reads: a ring \\
a predicate \\
mentioned by 387 result(s)

\smallskip

\noindent\textbf{\texttt{is-ringoid(s)}}\quad\textit{predicate}\\
reads: a ringoid \\
a predicate \\
mentioned by 78 result(s)

\smallskip

\noindent\textbf{\texttt{is-semigroup(s)}}\quad\textit{predicate}\\
reads: a semigroup \\
a predicate \\
mentioned by 8 result(s)

\smallskip

\noindent\textbf{\texttt{is-seminorm(m, p)}}\quad\textit{predicate}\\
reads: \$2 is a seminorm on \$1 \\
defined by \texttt{is-seminorm} \\
mentioned by 5 result(s)

\smallskip

\noindent\textbf{\texttt{is-seminorm-family(m, fam)}}\quad\textit{predicate}\\
reads: \$2 is a separating countable family of seminorms on \$1 \\
defined by \texttt{is-seminorm-family} \\
mentioned by 4 result(s)

\smallskip

\noindent\textbf{\texttt{is-separable(s)}}\quad\textit{predicate}\\
reads: \$1 is separable \\
defined by \texttt{is-separable} \\
mentioned by 3 result(s)

\smallskip

\noindent\textbf{\texttt{is-set(x)}}\quad\textit{predicate}\\
reads: \$1 is a set \\
a predicate \\
mentioned by 2 result(s)

\smallskip

\noindent\textbf{\texttt{is-setoid(s)}}\quad\textit{predicate}\\
reads: a setoid \\
a predicate \\
mentioned by 25 result(s)

\smallskip

\noindent\textbf{\texttt{is-sigma-algebra(omega, ca)}}\quad\textit{predicate}\\
reads: \$2 is a sigma-algebra of subsets of \$1 \\
defined by \texttt{is-sigma-algebra} \\
mentioned by 2 result(s)

\smallskip

\noindent\textbf{\texttt{is-strictly-below(porel, x, y)}}\quad\textit{predicate}\\
reads: \$2 is strictly below \$3 under \$1 \\
defined by \texttt{is-strictly-below} \\
mentioned by 11 result(s)

\smallskip

\noindent\textbf{\texttt{is-submodule(m, s)}}\quad\textit{predicate}\\
reads: \$2 is a submodule of \$1 \\
defined by \texttt{is-submodule} \\
mentioned by 39 result(s)

\smallskip

\noindent\textbf{\texttt{is-subsequence(s, y, f)}}\quad\textit{predicate}\\
reads: \$2 is a subsequence of \$3 in \$1 \\
defined by \texttt{is-subsequence} \\
mentioned by 2 result(s)

\smallskip

\noindent\textbf{\texttt{is-subspace(m, s)}}\quad\textit{predicate}\\
reads: \$2 is a subspace of \$1 \\
defined by \texttt{is-subspace} \\
mentioned by 2 result(s)

\smallskip

\noindent\textbf{\texttt{is-summable(grp, f)}}\quad\textit{predicate}\\
reads: \$2 is summable in \$1 \\
defined by \texttt{is-summable} \\
mentioned by 3 result(s)

\smallskip

\noindent\textbf{\texttt{is-top-space(s)}}\quad\textit{predicate}\\
reads: \$1 is a topological space \\
a predicate \\
mentioned by 19 result(s)

\smallskip

\noindent\textbf{\texttt{is-unif-cauchy(s, fam)}}\quad\textit{predicate}\\
reads: \$2 is uniformly Cauchy on \$1 \\
defined by \texttt{is-unif-cauchy} \\
mentioned by 4 result(s)

\smallskip

\noindent\textbf{\texttt{is-uniformly-continuous(s, t, f)}}\quad\textit{predicate}\\
reads: \$3 is uniformly continuous from \$1 to \$2 \\
defined by \texttt{is-uniformly-continuous} \\
mentioned by 3 result(s)

\smallskip

\noindent\textbf{\texttt{is-upper-bound(grd, porel, ch, b)}}\quad\textit{predicate}\\
reads: \$4 is an upper bound of \$3 in \$1 under \$2 \\
defined by \texttt{is-upper-bound} \\
mentioned by 6 result(s)

\smallskip

\noindent\textbf{\texttt{is-vector-space(s)}}\quad\textit{predicate}\\
reads: a vector space \\
a predicate \\
mentioned by 9 result(s)

\smallskip


\section*{K}
\addcontentsline{toc}{section}{K}

\noindent\textbf{\texttt{keep-set(phi, grd, porel, kset, alpha)}}\quad\textit{functoid}\\
= \{y in grd: y = phi(alpha) and forall([z in kset], [z, y] in porel)\} \\
mentioned by 5 result(s)

\smallskip


\section*{L}
\addcontentsline{toc}{section}{L}

\noindent\textbf{\texttt{lam-b}}\quad\textit{tactic}\\
\texttt{(lam-b)} -- VNB-LAMBDA beta: reduce an applied lambda to its substituted body. \\
declared in \texttt{tactics-help.scm}

\smallskip

\noindent\textbf{\texttt{lam-b-h}}\quad\textit{tactic}\\
\texttt{(lam-b-h hyp)} -- VNB-LAMBDA beta in a cited ASSUMPTION -- what mac-h is to mac. \\
declared in \texttt{tactics-help.scm}

\smallskip

\noindent\textbf{\texttt{lam-t}}\quad\textit{tactic}\\
\texttt{(lam-t)} -- VNB-LAMBDA typing: reduce (IN (VNB-LAMBDA ...) (FUN A B)) to TWO obligations -- the body, and that the domain is a set. \\
declared in \texttt{tactics-help.scm}

\smallskip

\noindent\textbf{\texttt{lastcoeff-set(md, p, u, sm)}}\quad\textit{functoid}\\
= \{r\_ in carr(scal(md)): forsome([c\_ in mat(1, succ(p), carr(scal(md)))], entry(c\_, 1, succ(p)) = r\_ and entry(m ... \\
mentioned by 3 result(s)

\smallskip

\noindent\textbf{\texttt{len-r}}\quad\textit{tactic}\\
\texttt{(len-r)} -- Reduce a LENGTH applied to a literal LIST in the goal. \\
declared in \texttt{tactics-help.scm}

\smallskip

\noindent\textbf{\texttt{length}}\quad\textit{operator}\\
a kernel term-former \\
mentioned by 42 result(s)

\smallskip

\noindent\textbf{\texttt{limit-ord}}\quad\textit{operator}\\
a kernel term-former \\
mentioned by 11 result(s)

\smallskip

\noindent\textbf{\texttt{line(m, v)}}\quad\textit{functoid}\\
= \{y\_ in vec(m): forsome([r\_ in rr], y\_ = (act(m))(r\_, v))\} \\
mentioned by 3 result(s)

\smallskip

\noindent\textbf{\texttt{list}}\quad\textit{operator}\\
a kernel term-former \\
mentioned by 111 result(s)

\smallskip

\noindent\textbf{\texttt{little-o-at(g, a)}}\quad\textit{predicate}\\
reads: \$1 is little-o at \$2 \\
defined by \texttt{little-o-at} \\
mentioned by 6 result(s)

\smallskip


\section*{M}
\addcontentsline{toc}{section}{M}

\noindent\textbf{\texttt{mac}}\quad\textit{tactic}\\
\texttt{(mac 'name)} -- Rewrite the GOAL with an equivalence macete (unfold a definition, apply an iff/=/== law).  Fires only where the macete's side-conditions already hold in context. \\
declared in \texttt{tactics-help.scm}

\smallskip

\noindent\textbf{\texttt{mac-h}}\quad\textit{tactic}\\
\texttt{(mac-h 'name hyp)} -- Rewrite a cited ASSUMPTION in place with an equivalence macete; any side-condition not already in context is spawned as a new subgoal. \\
declared in \texttt{tactics-help.scm}

\smallskip

\noindent\textbf{\texttt{mac-h*}}\quad\textit{tactic}\\
\texttt{(mac-h*)} -- Saturating mac-h: repeatedly unfold every defined predicate in the hypotheses and split the conjunctions they expose, until nothing is left folded.  One step in the trace instead of a dozen. \\
declared in \texttt{tactics-help.scm}

\smallskip

\noindent\textbf{\texttt{macm}}\quad\textit{tactic}\\
\texttt{(macm 'name)} -- Like mac, but a CONDITIONAL macete fires even when its side-conditions are not yet in context: each unmet condition is left as a new subgoal. \\
declared in \texttt{tactics-help.scm}

\smallskip

\noindent\textbf{\texttt{magnitude}}\quad\textit{operator}\\
a kernel term-former \\
mentioned by 15 result(s)

\smallskip

\noindent\textbf{\texttt{make-set}}\quad\textit{operator}\\
a kernel term-former \\
mentioned by 7 result(s)

\smallskip

\noindent\textbf{\texttt{mat(m, n, x)}}\quad\textit{functoid}\\
= \{p in matrix(x): size(p) = [m, n]\} \\
mentioned by 204 result(s)

\smallskip

\noindent\textbf{\texttt{mat-equiv(a, m, n, c, d)}}\quad\textit{predicate}\\
reads: \$4 and \$5 are equivalent \$2-by-\$3 matrices over \$1 \\
defined by \texttt{mat-equiv} \\
mentioned by 24 result(s)

\smallskip

\noindent\textbf{\texttt{mat-ring(a, n)}}\quad\textit{functoid}\\
= [mat(n, n, carr(a)), vnb-lambda([p, q], cartesian(mat(n, n, carr(a)), mat(n, n, carr(a))), matadd(a, p, q)), v ... \\
mentioned by 15 result(s)

\smallskip

\noindent\textbf{\texttt{matact(md, p, u)}}\quad\textit{functoid}\\
= matof(nth(1, size(p)), nth(2, size(u)), vnb-lambda([i, c], cartesian(interval(1, nth(1, size(p))), interval(1, ... \\
mentioned by 42 result(s)

\smallskip

\noindent\textbf{\texttt{matadd(a, p, q)}}\quad\textit{functoid}\\
= matof(nth(1, size(p)), nth(2, size(p)), vnb-lambda([i, j], cartesian(interval(1, nth(1, size(p))), interval(1, ... \\
mentioned by 24 result(s)

\smallskip

\noindent\textbf{\texttt{matmul(a, p, q)}}\quad\textit{functoid}\\
= matof(nth(1, size(p)), nth(2, size(q)), vnb-lambda([i, k], cartesian(interval(1, nth(1, size(p))), interval(1, ... \\
mentioned by 62 result(s)

\smallskip

\noindent\textbf{\texttt{matneg(a, p)}}\quad\textit{functoid}\\
= matof(nth(1, size(p)), nth(2, size(p)), vnb-lambda([i, j], cartesian(interval(1, nth(1, size(p))), interval(1, ... \\
mentioned by 7 result(s)

\smallskip

\noindent\textbf{\texttt{matof(m, n, g)}}\quad\textit{functoid}\\
= iota(p, p in mat(m, n, image(g, cartesian(interval(1, m), interval(1, n)))) and forall([i in interval(1, m), j ... \\
mentioned by 3 result(s)

\smallskip

\noindent\textbf{\texttt{matrix}}\quad\textit{operator}\\
a kernel term-former \\
mentioned by 3 result(s)

\smallskip

\noindent\textbf{\texttt{matscale(a, r, p)}}\quad\textit{functoid}\\
= matof(nth(1, size(p)), nth(2, size(p)), vnb-lambda([i, j], cartesian(interval(1, nth(1, size(p))), interval(1, ... \\
mentioned by 5 result(s)

\smallskip

\noindent\textbf{\texttt{matunit(a, n, k, l)}}\quad\textit{functoid}\\
= matof(n, n, vnb-lambda([i, j], cartesian(interval(1, n), interval(1, n)), if(i = k and j = l, one(a), zero(a)) ... \\
mentioned by 10 result(s)

\smallskip

\noindent\textbf{\texttt{metric-space}}\quad\textit{structure}\\
predicate \texttt{is-metric-space}; slots: pts dist

\smallskip

\noindent\textbf{\texttt{metric-sym-eq!}}\quad\textit{tactic}\\
\texttt{(metric-sym-eq! S P Q)} -- Add (= ((DIST S) P Q) ((DIST S) Q P)) to context via the metric-sym axiom.  [proof-local] \\
declared in \texttt{tactics-help.scm}

\smallskip

\noindent\textbf{\texttt{metric-top(md)}}\quad\textit{functoid}\\
reads: the metric topology of \$1 \\
= [pts(md), \{u in power(pts(md)): is-open(md, u)\}] \\
mentioned by 12 result(s)

\smallskip

\noindent\textbf{\texttt{metrizable-top-space}}\quad\textit{refinement}\\
predicate \texttt{is-metrizable-top-space}; same shape as top-space

\smallskip

\noindent\textbf{\texttt{minimize!}}\quad\textit{tactic}\\
\texttt{(minimize! '(v1 ... vk) GUARD MEASURE)} -- Choose v1..vk satisfying GUARD with the NN-valued MEASURE as small as possible: lands the witnesses and their minimality, and opens the two obligations any minimisation owes (MEASURE lands in NN; GUARD is satisfiable). \\
declared in \texttt{tactics-help.scm}

\smallskip

\noindent\textbf{\texttt{minor(s, r, c, n)}}\quad\textit{functoid}\\
= matof(n, n, vnb-lambda([i, j], cartesian(interval(1, n), interval(1, n)), entry(s, if(i < r, i, succ(i)), if(j ... \\
mentioned by 3 result(s)

\smallskip

\noindent\textbf{\texttt{module}}\quad\textit{structure}\\
predicate \texttt{is-module}; slots: scal vec vadd vzero vneg act

\smallskip

\noindent\textbf{\texttt{module-vector-ag(s)}}\quad\textit{view}\\
a module viewed as an abelian-group \\
mentioned by 36 result(s)

\smallskip

\noindent\textbf{\texttt{monalg(a, m)}}\quad\textit{functoid}\\
reads: the monoid algebra \$1[\$2] \\
= [finsupp(a, m), vnb-lambda([f, g], cartesian(finsupp(a, m), finsupp(a, m)), monalg-add(a, m, f, g)), vnb-lambd ... \\
mentioned by 2 result(s)

\smallskip

\noindent\textbf{\texttt{monalg-add(a, m, f, g)}}\quad\textit{functoid}\\
= vnb-lambda(x\_, carr(m), (add(a))(f(x\_), g(x\_))) \\
mentioned by 3 result(s)

\smallskip

\noindent\textbf{\texttt{monalg-mul(a, m, f, g)}}\quad\textit{functoid}\\
= vnb-lambda(x\_, carr(m), finsum(ring-additive-ag(a), vnb-lambda(p, \{p in cartesian(supp(a, m, f), supp(a, m, g) ... \\
mentioned by 7 result(s)

\smallskip

\noindent\textbf{\texttt{monalg-neg(a, m, f)}}\quad\textit{functoid}\\
= vnb-lambda(x\_, carr(m), (neg(a))(f(x\_)))

\smallskip

\noindent\textbf{\texttt{monalg-one(a, m)}}\quad\textit{functoid}\\
= vnb-lambda(x\_, carr(m), if(x\_ = iden(m), one(a), zero(a))) \\
mentioned by 2 result(s)

\smallskip

\noindent\textbf{\texttt{monalg-zero(a, m)}}\quad\textit{functoid}\\
= vnb-lambda(x\_, carr(m), zero(a))

\smallskip

\noindent\textbf{\texttt{monoid}}\quad\textit{structure}\\
predicate \texttt{is-monoid}; slots: carr opr iden

\smallskip

\noindent\textbf{\texttt{mpow}}\quad\textit{defined-fn}\\
defined by \texttt{mpow-zero} \texttt{mpow-succ} \\
mentioned by 51 result(s)

\smallskip

\noindent\textbf{\texttt{mul(s)}}\quad\textit{accessor}\\
a structure slot \\
mentioned by 239 result(s)

\smallskip


\section*{N}
\addcontentsline{toc}{section}{N}

\noindent\textbf{\texttt{nag-metric-space(nag)}}\quad\textit{functoid}\\
= [carr(nag), vnb-lambda([u, v], cartesian(carr(nag), carr(nag)), (nrm(nag))((opr(nag))(u, (inv(nag))(v))))] \\
mentioned by 11 result(s)

\smallskip

\noindent\textbf{\texttt{neg(s)}}\quad\textit{accessor}\\
a structure slot \\
mentioned by 101 result(s)

\smallskip

\noindent\textbf{\texttt{neg-rr}}\quad\textit{predicate}\\
reads: a negative real \\
a predicate

\smallskip

\noindent\textbf{\texttt{nf-metric-space(nf)}}\quad\textit{functoid}\\
= [carr(nf), vnb-lambda([x, y], cartesian(carr(nf), carr(nf)), (fnrm(nf))((add(nf))(x, (neg(nf))(y))))] \\
mentioned by 5 result(s)

\smallskip

\noindent\textbf{\texttt{ni}}\quad\textit{tactic}\\
\texttt{(ni)} -- Natural-number induction on the goal's leading FORALL over NN. \\
declared in \texttt{tactics-help.scm}

\smallskip

\noindent\textbf{\texttt{nn-add-monoid}}\quad\textit{instance}\\
an instance of comm-monoid; predicate \texttt{is-nn-add-monoid}; same shape as comm-monoid

\smallskip

\noindent\textbf{\texttt{nn-enum(s)}}\quad\textit{functoid}\\
= choice(\{f in fun(nn, s): forall([m in nn, n in nn], m < n implies f(m) < f(n))\}) \\
mentioned by 1 result(s)

\smallskip

\noindent\textbf{\texttt{nn-minus}}\quad\textit{defined-fn}\\
defined by \texttt{nn-minus-def} \\
mentioned by 11 result(s)

\smallskip

\noindent\textbf{\texttt{nnfst(n\_)}}\quad\textit{functoid}\\
reads: the first component of \$1 \\
= iota(a\_, a\_ in nn and forsome([b\_ in nn], nnpair(a\_, b\_) = n\_)) \\
mentioned by 3 result(s)

\smallskip

\noindent\textbf{\texttt{nnpair(i\_, j\_)}}\quad\textit{functoid}\\
reads: the Cantor code of (\$1, \$2) \\
= trinum(i\_ + j\_) + j\_ \\
mentioned by 7 result(s)

\smallskip

\noindent\textbf{\texttt{nnsnd(n\_)}}\quad\textit{functoid}\\
reads: the second component of \$1 \\
= iota(b\_, b\_ in nn and forsome([a\_ in nn], nnpair(a\_, b\_) = n\_)) \\
mentioned by 3 result(s)

\smallskip

\noindent\textbf{\texttt{non-zero(s)}}\quad\textit{accessor}\\
a structure slot \\
mentioned by 23 result(s)

\smallskip

\noindent\textbf{\texttt{nonneg-rr}}\quad\textit{predicate}\\
reads: a nonnegative real \\
a predicate

\smallskip

\noindent\textbf{\texttt{normed-ag}}\quad\textit{structure}\\
predicate \texttt{is-normed-ag}; slots: carr opr iden inv nrm

\smallskip

\noindent\textbf{\texttt{normed-ag-as-abelian-group(s)}}\quad\textit{view}\\
a normed-ag viewed as an abelian-group \\
mentioned by 7 result(s)

\smallskip

\noindent\textbf{\texttt{normed-field}}\quad\textit{structure}\\
predicate \texttt{is-normed-field}; slots: carr add mul neg zero one fnrm

\smallskip

\noindent\textbf{\texttt{normed-field-additive-ag(s)}}\quad\textit{view}\\
a normed-field viewed as an abelian-group \\
mentioned by 3 result(s)

\smallskip

\noindent\textbf{\texttt{normed-field-as-commutative-ring(s)}}\quad\textit{view}\\
a normed-field viewed as a commutative-ring \\
mentioned by 6 result(s)

\smallskip

\noindent\textbf{\texttt{normed-field-as-integral-domain(s)}}\quad\textit{view}\\
a normed-field viewed as an integral-domain \\
mentioned by 2 result(s)

\smallskip

\noindent\textbf{\texttt{normed-vector-space}}\quad\textit{structure}\\
predicate \texttt{is-normed-vector-space}; slots: scal vec vadd vzero vneg act vnrm

\smallskip

\noindent\textbf{\texttt{normed-vector-space-as-module(s)}}\quad\textit{view}\\
a normed-vector-space viewed as a module \\
mentioned by 4 result(s)

\smallskip

\noindent\textbf{\texttt{normed-vector-space-as-normed-ag(s)}}\quad\textit{view}\\
a normed-vector-space viewed as a normed-ag \\
mentioned by 9 result(s)

\smallskip

\noindent\textbf{\texttt{npe(m, s, f, t, g)}}\quad\textit{predicate}\\
reads: \$5 is a norm-preserving extension of \$3 from \$2 to \$4, in \$1 \\
defined by \texttt{npe} \\
mentioned by 4 result(s)

\smallskip

\noindent\textbf{\texttt{nrm(s)}}\quad\textit{accessor}\\
a structure slot \\
mentioned by 15 result(s)

\smallskip

\noindent\textbf{\texttt{nth}}\quad\textit{operator}\\
a kernel term-former \\
mentioned by 21 result(s)

\smallskip

\noindent\textbf{\texttt{nth-deriv}}\quad\textit{defined-fn}\\
defined by \texttt{nth-deriv-zero} \texttt{nth-deriv-succ} \\
mentioned by 15 result(s)

\smallskip

\noindent\textbf{\texttt{nth-deriv-v}}\quad\textit{defined-fn}\\
defined by \texttt{nth-deriv-v-zero} \texttt{nth-deriv-v-succ} \\
mentioned by 15 result(s)

\smallskip

\noindent\textbf{\texttt{nth-r}}\quad\textit{tactic}\\
\texttt{(nth-r)} -- Reduce an NTH applied to a literal LIST in the goal. \\
declared in \texttt{tactics-help.scm}

\smallskip

\noindent\textbf{\texttt{null-rr-seq(rad)}}\quad\textit{predicate}\\
reads: \$1 is a sequence of positive reals tending to zero \\
defined by \texttt{null-rr-seq} \\
mentioned by 5 result(s)

\smallskip

\noindent\textbf{\texttt{nvs-metric-space(m)}}\quad\textit{functoid}\\
= [vec(m), vnb-lambda([x, y], cartesian(vec(m), vec(m)), (vnrm(m))((vadd(m))(x, (vneg(m))(y))))] \\
mentioned by 5 result(s)

\smallskip


\section*{O}
\addcontentsline{toc}{section}{O}

\noindent\textbf{\texttt{obtain}}\quad\textit{tactic}\\
\texttt{(obtain LANE)} -- Run LANE -- a thunk whose forward step lands a `there exists' into context -- then eliminate that existential and RETURN the fresh witness's name, read off the proof state rather than guessed. \\
declared in \texttt{tactics-help.scm}

\smallskip

\noindent\textbf{\texttt{oi-l}}\quad\textit{tactic}\\
\texttt{(oi-l)} -- OR-intro left: reduce an (OR a b) goal to a. \\
declared in \texttt{tactics-help.scm}

\smallskip

\noindent\textbf{\texttt{oi-r}}\quad\textit{tactic}\\
\texttt{(oi-r)} -- OR-intro right: reduce an (OR a b) goal to b. \\
declared in \texttt{tactics-help.scm}

\smallskip

\noindent\textbf{\texttt{one(s)}}\quad\textit{accessor}\\
a structure slot \\
mentioned by 153 result(s)

\smallskip

\noindent\textbf{\texttt{opens(s)}}\quad\textit{accessor}\\
reads: the topology of \$1 \\
a structure slot \\
mentioned by 8 result(s)

\smallskip

\noindent\textbf{\texttt{opr(s)}}\quad\textit{accessor}\\
a structure slot \\
mentioned by 157 result(s)

\smallskip

\noindent\textbf{\texttt{ord-segment}}\quad\textit{operator}\\
a kernel term-former \\
mentioned by 49 result(s)

\smallskip

\noindent\textbf{\texttt{orelse}}\quad\textit{tactic}\\
\texttt{(orelse t1 t2 ...)} -- Run thunks in order, stop at the first that makes progress (LCF ORELSE). \\
declared in \texttt{tactics-help.scm}

\smallskip


\section*{P}
\addcontentsline{toc}{section}{P}

\noindent\textbf{\texttt{pair}}\quad\textit{operator}\\
a kernel term-former \\
mentioned by 38 result(s)

\smallskip

\noindent\textbf{\texttt{partial-fun}}\quad\textit{operator}\\
a kernel term-former \\
mentioned by 5 result(s)

\smallskip

\noindent\textbf{\texttt{pbc}}\quad\textit{tactic}\\
\texttt{(pbc)} -- Proof by contradiction: assume the goal's negation, prove FALSITY. \\
declared in \texttt{tactics-help.scm}

\smallskip

\noindent\textbf{\texttt{permutations(n)}}\quad\textit{functoid}\\
= injection(ord-segment(n), ord-segment(n)) \\
mentioned by 4 result(s)

\smallskip

\noindent\textbf{\texttt{pid}}\quad\textit{refinement}\\
predicate \texttt{is-pid}; same shape as integral-domain

\smallskip

\noindent\textbf{\texttt{pointwise-bounded(s, fam)}}\quad\textit{predicate}\\
reads: \$2 is pointwise bounded on \$1 \\
defined by \texttt{pointwise-bounded} \\
mentioned by 3 result(s)

\smallskip

\noindent\textbf{\texttt{poly(a)}}\quad\textit{functoid}\\
reads: \$1[x] \\
= monalg(a, nn-add-monoid) \\
mentioned by 1 result(s)

\smallskip

\noindent\textbf{\texttt{pos-rr(r)}}\quad\textit{predicate}\\
reads: a positive real \\
defined by \texttt{pos-rr} \\
mentioned by 50 result(s)

\smallskip

\noindent\textbf{\texttt{power}}\quad\textit{operator}\\
a kernel term-former \\
mentioned by 75 result(s)

\smallskip

\noindent\textbf{\texttt{preimage(s, f, v)}}\quad\textit{functoid}\\
= \{a in pts(s): f(a) in v\} \\
mentioned by 14 result(s)

\smallskip

\noindent\textbf{\texttt{prep}}\quad\textit{tactic}\\
\texttt{(prep 'ineq)} -- Why will a tactic not fire here, and what must be done first?  Runs the tactic's OWN preconditions one at a time, marks each ok / PREP / STOP, names the library lemma that repairs each unmet one, and prints a PLAN it has checked by running it on a scratch clone.  READ-ONLY: the live proof is never touched.  (prep 'ineq) is the only method so far. \\
declared in \texttt{tactics-help.scm}

\smallskip

\noindent\textbf{\texttt{principal-ideal(s, a)}}\quad\textit{functoid}\\
= \{x in carr(s): forsome([r in carr(s)], x = (mul(s))(r, a))\} \\
mentioned by 6 result(s)

\smallskip

\noindent\textbf{\texttt{prod-ord}}\quad\textit{defined-fn}\\
defined by \texttt{prod-ord-zero} \texttt{prod-ord-succ} \\
mentioned by 7 result(s)

\smallskip

\noindent\textbf{\texttt{prod-ring(r, f, s)}}\quad\textit{functoid}\\
= finprod(commutative-ring-multiplicative-cm(r), f, s) \\
mentioned by 9 result(s)

\smallskip

\noindent\textbf{\texttt{prod-set}}\quad\textit{operator}\\
a kernel term-former \\
mentioned by 7 result(s)

\smallskip

\noindent\textbf{\texttt{product-carrier(ms)}}\quad\textit{functoid}\\
= \{x in fun(nn, big-union(n, nn, pts(ms(n)))): forall([n in nn], x(n) in pts(ms(n)))\} \\
mentioned by 7 result(s)

\smallskip

\noindent\textbf{\texttt{product-metric(ms)}}\quad\textit{functoid}\\
= product-metric-w(ms, vnb-lambda(n, nn, /(1, 2 \textasciicircum{} (n + 1)))) \\
mentioned by 2 result(s)

\smallskip

\noindent\textbf{\texttt{product-metric-w(ms, w)}}\quad\textit{functoid}\\
= [product-carrier(ms), vnb-lambda([x, y], cartesian(product-carrier(ms), product-carrier(ms)), iota(l, series-c ... \\
mentioned by 7 result(s)

\smallskip

\noindent\textbf{\texttt{product-proj(ms, n)}}\quad\textit{functoid}\\
= vnb-lambda(x, product-carrier(ms), x(n)) \\
mentioned by 1 result(s)

\smallskip

\noindent\textbf{\texttt{proj(s)}}\quad\textit{functoid}\\
= vnb-lambda(a, pts(s), class(s, a)) \\
mentioned by 1 result(s)

\smallskip

\noindent\textbf{\texttt{ps-absolutely-converges-at(coef, x)}}\quad\textit{predicate}\\
reads: the power series with coefficients \$1 converges absolutely at \$2 \\
defined by \texttt{ps-absolutely-converges-at} \\
mentioned by 3 result(s)

\smallskip

\noindent\textbf{\texttt{ps-converges-at(coef, x)}}\quad\textit{predicate}\\
reads: the power series with coefficients \$1 converges at \$2 \\
defined by \texttt{ps-converges-at} \\
mentioned by 8 result(s)

\smallskip

\noindent\textbf{\texttt{ps-converges-to-at(coef, x, l)}}\quad\textit{predicate}\\
reads: the power series with coefficients \$1 converges to \$3 at \$2 \\
defined by \texttt{ps-converges-to-at} \\
mentioned by 2 result(s)

\smallskip

\noindent\textbf{\texttt{ps-partial-sum(coef, x, k)}}\quad\textit{functoid}\\
= sum-ag(normed-field-additive-ag(rr-normed-field), vnb-lambda(n, nn, coef(n) * x \textasciicircum{} n), k) \\
mentioned by 6 result(s)

\smallskip

\noindent\textbf{\texttt{ps-ratio-limit(coef, l)}}\quad\textit{predicate}\\
reads: the coefficients \$1 have ratio limit \$2 \\
defined by \texttt{ps-ratio-limit} \\
mentioned by 3 result(s)

\smallskip

\noindent\textbf{\texttt{pseudo-gauge-top(fam, ground)}}\quad\textit{functoid}\\
reads: the topology generated by the pseudometric family \$1 on \$2 \\
= [ground, \{u in power(ground): is-open-gauge(fam, ground, u)\}] \\
mentioned by 2 result(s)

\smallskip

\noindent\textbf{\texttt{pseudometric-space}}\quad\textit{structure}\\
predicate \texttt{is-pseudometric-space}; slots: pts dist

\smallskip

\noindent\textbf{\texttt{pts(s)}}\quad\textit{accessor}\\
a structure slot \\
mentioned by 165 result(s)

\smallskip


\section*{Q}
\addcontentsline{toc}{section}{Q}

\noindent\textbf{\texttt{qed}}\quad\textit{tactic}\\
\texttt{(qed 'name)} -- Install the finished proof as theorem `name' and save its replayable script. \\
declared in \texttt{tactics-help.scm}

\smallskip

\noindent\textbf{\texttt{qq-field}}\quad\textit{instance}\\
an instance of field; predicate \texttt{is-qq-field}; same shape as field

\smallskip

\noindent\textbf{\texttt{qq-ring}}\quad\textit{instance}\\
an instance of ring; predicate \texttt{is-qq-ring}; same shape as integral-domain

\smallskip

\noindent\textbf{\texttt{qrfl}}\quad\textit{tactic}\\
\texttt{(qrfl)} -- Quasi-reflexivity: close t = t under the partial-equality definedness reading. \\
declared in \texttt{tactics-help.scm}

\smallskip

\noindent\textbf{\texttt{quietly}}\quad\textit{tactic}\\
\texttt{(quietly thunk)} -- Run thunk with state-dump output suppressed; returns its value. \\
declared in \texttt{tactics-help.scm}

\smallskip

\noindent\textbf{\texttt{quotient(s)}}\quad\textit{functoid}\\
= image(proj(s), pts(s)) \\
mentioned by 9 result(s)

\smallskip


\section*{R}
\addcontentsline{toc}{section}{R}

\noindent\textbf{\texttt{ran(f)}}\quad\textit{functoid}\\
= image(f, dom(f)) \\
mentioned by 7 result(s)

\smallskip

\noindent\textbf{\texttt{real-part}}\quad\textit{operator}\\
a kernel term-former

\smallskip

\noindent\textbf{\texttt{recip(s)}}\quad\textit{accessor}\\
a structure slot \\
mentioned by 45 result(s)

\smallskip

\noindent\textbf{\texttt{reduce}}\quad\textit{defined-fn}\\
defined by \texttt{reduce-one} \texttt{reduce-succ} \\
mentioned by 26 result(s)

\smallskip

\noindent\textbf{\texttt{rel(s)}}\quad\textit{accessor}\\
a structure slot \\
mentioned by 4 result(s)

\smallskip

\noindent\textbf{\texttt{rel-free(md, n, u)}}\quad\textit{predicate}\\
reads: the \$2 vectors \$3 are linearly independent in \$1 \\
defined by \texttt{rel-free} \\
mentioned by 5 result(s)

\smallskip

\noindent\textbf{\texttt{related(s, a, b)}}\quad\textit{functoid}\\
= [a, b] in rel(s) \\
mentioned by 2 result(s)

\smallskip

\noindent\textbf{\texttt{repeat}}\quad\textit{tactic}\\
\texttt{(repeat thunk [cap])} -- Run thunk until it stops changing the proof state (LCF REPEAT). \\
declared in \texttt{tactics-help.scm}

\smallskip

\noindent\textbf{\texttt{replay-proof}}\quad\textit{tactic}\\
\texttt{(replay-proof 'name [subst])} -- Re-run a saved script on the current goal, optionally renaming free vars. \\
declared in \texttt{tactics-help.scm}

\smallskip

\noindent\textbf{\texttt{res}}\quad\textit{operator}\\
a kernel term-former \\
mentioned by 4 result(s)

\smallskip

\noindent\textbf{\texttt{respects(s, f)}}\quad\textit{functoid}\\
= forall([a in pts(s), b in pts(s)], related(s, a, b) implies f(a) = f(b)) \\
mentioned by 4 result(s)

\smallskip

\noindent\textbf{\texttt{respects2(s, f)}}\quad\textit{functoid}\\
= forall([a, b, a\_, b\_], a in pts(s) implies b in pts(s) implies a\_ in pts(s) implies b\_ in pts(s) implies relat ... \\
mentioned by 3 result(s)

\smallskip

\noindent\textbf{\texttt{restvar}}\quad\textit{operator}\\
a kernel term-former \\
mentioned by 4 result(s)

\smallskip

\noindent\textbf{\texttt{rfl}}\quad\textit{tactic}\\
\texttt{(rfl)} -- Close a reflexive equality goal (t = t). \\
declared in \texttt{tactics-help.scm}

\smallskip

\noindent\textbf{\texttt{ring}}\quad\textit{structure}\\
predicate \texttt{is-ring}; slots: carr add mul neg zero one

\smallskip

\noindent\textbf{\texttt{ring-additive-ag(s)}}\quad\textit{view}\\
a ring viewed as an abelian-group \\
mentioned by 151 result(s)

\smallskip

\noindent\textbf{\texttt{ring-multiplicative-monoid(s)}}\quad\textit{view}\\
a ring viewed as a monoid \\
mentioned by 4 result(s)

\smallskip

\noindent\textbf{\texttt{ring-power(r, x, n)}}\quad\textit{functoid}\\
= mpow(commutative-ring-multiplicative-cm(r), x, n) \\
mentioned by 13 result(s)

\smallskip

\noindent\textbf{\texttt{ring-prod(x, y)}}\quad\textit{functoid}\\
= [cartesian(carr(x), carr(y)), vnb-lambda([p, q], cartesian(cartesian(carr(x), carr(y)), cartesian(carr(x), car ... \\
mentioned by 4 result(s)

\smallskip

\noindent\textbf{\texttt{ring-prod-n}}\quad\textit{defined-fn}\\
defined by \texttt{ring-prod-n-zero} \texttt{ring-prod-n-succ} \\
mentioned by 5 result(s)

\smallskip

\noindent\textbf{\texttt{ringoid}}\quad\textit{structure}\\
predicate \texttt{is-ringoid}; slots: carr add mul neg zero one idl

\smallskip

\noindent\textbf{\texttt{ringoid-as-ring(s)}}\quad\textit{view}\\
a ringoid viewed as a ring \\
mentioned by 4 result(s)

\smallskip

\noindent\textbf{\texttt{ringoid-quotient(r)}}\quad\textit{functoid}\\
= quotient(ringoid-setoid(r))

\smallskip

\noindent\textbf{\texttt{ringoid-quotient-ring(r)}}\quad\textit{functoid}\\
= [quotient(ringoid-setoid(r)), descend2(ringoid-setoid(r), vnb-lambda([a, b], cartesian(carr(r), carr(r)), clas ... \\
mentioned by 15 result(s)

\smallskip

\noindent\textbf{\texttt{ringoid-rel(r)}}\quad\textit{functoid}\\
= \{p in cartesian(carr(r), carr(r)): (add(r))(nth(1, p), (neg(r))(nth(2, p))) in idl(r)\} \\
mentioned by 3 result(s)

\smallskip

\noindent\textbf{\texttt{ringoid-setoid(r)}}\quad\textit{functoid}\\
= [carr(r), ringoid-rel(r)] \\
mentioned by 17 result(s)

\smallskip

\noindent\textbf{\texttt{rpow}}\quad\textit{operator}\\
a kernel term-former \\
mentioned by 27 result(s)

\smallskip

\noindent\textbf{\texttt{rr+*-add-monoid}}\quad\textit{refinement}\\
predicate \texttt{is-rr+*-add-monoid}; same shape as comm-monoid

\smallskip

\noindent\textbf{\texttt{rr-bounded-above(s)}}\quad\textit{predicate}\\
reads: \$1 is bounded above \\
defined by \texttt{rr-bounded-above} \\
mentioned by 5 result(s)

\smallskip

\noindent\textbf{\texttt{rr-bounded-ms(())}}\quad\textit{functoid}\\
= bdd-metric(rr-ms)

\smallskip

\noindent\textbf{\texttt{rr-ms}}\quad\textit{instance}\\
an instance of metric-space; predicate \texttt{is-rr-ms}; same shape as metric-space

\smallskip

\noindent\textbf{\texttt{rr-normed-field}}\quad\textit{instance}\\
an instance of normed-field; predicate \texttt{is-rr-normed-field}; same shape as normed-field

\smallskip

\noindent\textbf{\texttt{rr-upper-bound(s, b)}}\quad\textit{predicate}\\
reads: \$2 is an upper bound of \$1 \\
defined by \texttt{rr-upper-bound} \\
mentioned by 6 result(s)

\smallskip

\noindent\textbf{\texttt{rs}}\quad\textit{tactic}\\
\texttt{(rs)} -- Ring-simplify the goal (normal form over the ambient ring). \\
declared in \texttt{tactics-help.scm}

\smallskip


\section*{S}
\addcontentsline{toc}{section}{S}

\noindent\textbf{\texttt{save-proof}}\quad\textit{tactic}\\
\texttt{(save-proof 'name)} -- Save the current script under a name without finishing. \\
declared in \texttt{tactics-help.scm}

\smallskip

\noindent\textbf{\texttt{scal(s)}}\quad\textit{accessor}\\
a structure slot \\
mentioned by 125 result(s)

\smallskip

\noindent\textbf{\texttt{scout}}\quad\textit{tactic}\\
\texttt{(scout [depth [branch [nodes]]])} -- Speculatively search to a bounded depth across INDEPENDENT scratch branches; RETURNS a nested list (number-of-branches (d b) goal best-partials closing-branches) rather than printing.  Never touches the live proof. \\
declared in \texttt{tactics-help.scm}

\smallskip

\noindent\textbf{\texttt{scout-run}}\quad\textit{tactic}\\
\texttt{(scout-run k)} -- Adopt closing branch k from the last (scout)/(scout-show) onto the live proof, running its tactics for real (they record and display normally). \\
declared in \texttt{tactics-help.scm}

\smallskip

\noindent\textbf{\texttt{scout-show}}\quad\textit{tactic}\\
\texttt{(scout-show [depth [branch [nodes]]])} -- Like (scout), but PRINTS the human-readable report (examined count, goal, closing branches or best partials) to the REPL.  Returns the same nested list (scout) does. \\
declared in \texttt{tactics-help.scm}

\smallskip

\noindent\textbf{\texttt{semigroup}}\quad\textit{structure}\\
predicate \texttt{is-semigroup}; slots: carr opr

\smallskip

\noindent\textbf{\texttt{sep}}\quad\textit{operator}\\
a kernel term-former \\
mentioned by 9 result(s)

\smallskip

\noindent\textbf{\texttt{sep-me}}\quad\textit{tactic}\\
\texttt{(sep-me hyp)} -- Separation membership elim: split a separation-membership assumption into A-membership and the predicate. \\
declared in \texttt{tactics-help.scm}

\smallskip

\noindent\textbf{\texttt{sep-mi}}\quad\textit{tactic}\\
\texttt{(sep-mi)} -- Separation membership intro: prove (IN t \{x in A | p\}). \\
declared in \texttt{tactics-help.scm}

\smallskip

\noindent\textbf{\texttt{sep-set}}\quad\textit{tactic}\\
\texttt{(sep-set)} -- Separation sethood: the separation set \{x in A | p\} is a set. \\
declared in \texttt{tactics-help.scm}

\smallskip

\noindent\textbf{\texttt{seq-compact(s)}}\quad\textit{predicate}\\
reads: \$1 is sequentially compact \\
defined by \texttt{seq-compact} \\
mentioned by 7 result(s)

\smallskip

\noindent\textbf{\texttt{series-converges(f)}}\quad\textit{predicate}\\
reads: the series \$1 converges \\
defined by \texttt{series-converges} \\
mentioned by 11 result(s)

\smallskip

\noindent\textbf{\texttt{series-converges-to(f, l)}}\quad\textit{predicate}\\
reads: the series \$1 converges to \$2 \\
defined by \texttt{series-converges-to} \\
mentioned by 3 result(s)

\smallskip

\noindent\textbf{\texttt{series-partial-sum(f, k)}}\quad\textit{functoid}\\
= sum-ag(normed-field-additive-ag(rr-normed-field), f, k) \\
mentioned by 12 result(s)

\smallskip

\noindent\textbf{\texttt{setoid}}\quad\textit{structure}\\
predicate \texttt{is-setoid}; slots: pts rel

\smallskip

\noindent\textbf{\texttt{show}}\quad\textit{tactic}\\
\texttt{(show)} -- Redisplay the current proof state. \\
declared in \texttt{tactics-help.scm}

\smallskip

\noindent\textbf{\texttt{simp}}\quad\textit{tactic}\\
\texttt{(simp [target])} -- Rewrite a commutative-ring SUBTERM of the goal to canonical form, IN PLACE (e.g. (x+y)\textasciicircum{}2 inside a larger goal becomes x\textasciicircum{}2 + 2*x*y + y\textasciicircum{}2).  Works on BOTH surfaces: concrete number domains (+ * - \textasciicircum{} over NN/ZZ/QQ/RR/CC) and a generic ring s ((ADD s)/(MUL s)/(NEG s), carrier (CARR s)).  No arg = outermost ring subterm; "term" targets a specific one.  Sound by cut + crs + eq-subst (no new kernel rule); needs the subterm's generators typed in context (true post-di), else refuses and names them. \\
declared in \texttt{tactics-help.scm}

\smallskip

\noindent\textbf{\texttt{sin}}\quad\textit{operator}\\
a kernel term-former

\smallskip

\noindent\textbf{\texttt{singleton}}\quad\textit{operator}\\
a kernel term-former \\
mentioned by 8 result(s)

\smallskip

\noindent\textbf{\texttt{size(m)}}\quad\textit{functoid}\\
= [length(m), length(nth(1, m))]

\smallskip

\noindent\textbf{\texttt{slot}}\quad\textit{tactic}\\
\texttt{(slot 'acc)} -- Reduce a structure ACCESSOR to its projection: (CARR s) becomes (NTH 1 s).  The one door for accessor reductions -- use it, never mac, on an accessor name. \\
declared in \texttt{tactics-help.scm}

\smallskip

\noindent\textbf{\texttt{smith-staircase(a, m, n, d, k)}}\quad\textit{predicate}\\
reads: \$4 is in Smith staircase form of rank \$5, as an \$2-by-\$3 matrix over \$1 \\
defined by \texttt{smith-staircase} \\
mentioned by 3 result(s)

\smallskip

\noindent\textbf{\texttt{snoc-col(w, n, x)}}\quad\textit{functoid}\\
= matof(succ(n), 1, vnb-lambda([i\_, j\_], cartesian(interval(1, succ(n)), interval(1, 1)), if(i\_ = succ(n), x, en ... \\
mentioned by 7 result(s)

\smallskip

\noindent\textbf{\texttt{snoc-row(c, n, r)}}\quad\textit{functoid}\\
= matof(1, succ(n), vnb-lambda([i\_, j\_], cartesian(interval(1, 1), interval(1, succ(n))), if(j\_ = succ(n), r, en ... \\
mentioned by 7 result(s)

\smallskip

\noindent\textbf{\texttt{sos}}\quad\textit{tactic}\\
\texttt{(sos "c1" "c2" ...)} -- Sum-of-squares closer for a nonstrict polynomial inequality a <= b over RR (the nonlinear companion of (ineq)).  You supply the terms to be SQUARED; it finds the nonnegative coefficients.  (tactics 'sos) for the worked example. \\
declared in \texttt{tactics-help.scm}

\smallskip

\noindent\textbf{\texttt{sp}}\quad\textit{tactic}\\
\texttt{(sp goal)} -- Start a proof of `goal' (a "string", a raw S-expr, or a wff); clears the script. \\
declared in \texttt{tactics-help.scm}

\smallskip

\noindent\textbf{\texttt{span(md, n, u)}}\quad\textit{functoid}\\
= \{x\_ in vec(md): forsome([c\_ in mat(1, n, carr(scal(md)))], x\_ = entry(matact(md, c\_, u), 1, 1))\} \\
mentioned by 6 result(s)

\smallskip

\noindent\textbf{\texttt{span-add-one(m, s, v)}}\quad\textit{functoid}\\
= \{y\_ in vec(m): forsome([x\_ in s, r\_ in rr], y\_ = (vadd(m))(x\_, (act(m))(r\_, v)))\} \\
mentioned by 8 result(s)

\smallskip

\noindent\textbf{\texttt{spans(md, n, u, sm)}}\quad\textit{predicate}\\
reads: the \$2 vectors \$3 span \$4 in \$1 \\
defined by \texttt{spans} \\
mentioned by 10 result(s)

\smallskip

\noindent\textbf{\texttt{splice}}\quad\textit{operator}\\
a kernel term-former \\
mentioned by 4 result(s)

\smallskip

\noindent\textbf{\texttt{split-ands!}}\quad\textit{tactic}\\
\texttt{(split-ands!)} -- Flatten every AND assumption of the focus into separate assumptions.  [proof-local] \\
declared in \texttt{tactics-help.scm}

\smallskip

\noindent\textbf{\texttt{sqrt}}\quad\textit{operator}\\
a kernel term-former \\
mentioned by 13 result(s)

\smallskip

\noindent\textbf{\texttt{strictly-mono-nn(phi)}}\quad\textit{predicate}\\
reads: \$1 is a strictly increasing sequence of naturals \\
defined by \texttt{strictly-mono-nn} \\
mentioned by 17 result(s)

\smallskip

\noindent\textbf{\texttt{submat(s, p, q)}}\quad\textit{functoid}\\
= matof(p, q, vnb-lambda([i, j], cartesian(interval(1, p), interval(1, q)), entry(s, succ(i), succ(j)))) \\
mentioned by 6 result(s)

\smallskip

\noindent\textbf{\texttt{subseq(f, phi)}}\quad\textit{functoid}\\
= vnb-lambda(k, nn, f(phi(k))) \\
mentioned by 12 result(s)

\smallskip

\noindent\textbf{\texttt{subset}}\quad\textit{primitive}\\
reads: \$1 is a subset of \$2 \\
a kernel relation \\
mentioned by 97 result(s)

\smallskip

\noindent\textbf{\texttt{subst}}\quad\textit{tactic}\\
\texttt{(subst '(= s t))  |  (subst '(== s t))} -- Rewrite s -> t throughout the goal, using an equation s = t -- or a quasi-equation s == t -- that is in context, in EITHER orientation (Leibniz substitution).  Cannot reach a term in OPERATOR position: use the equation as a macete instead. \\
declared in \texttt{tactics-help.scm}

\smallskip

\noindent\textbf{\texttt{succ}}\quad\textit{operator}\\
a kernel term-former \\
mentioned by 228 result(s)

\smallskip

\noindent\textbf{\texttt{succ\_ord}}\quad\textit{operator}\\
a kernel term-former \\
mentioned by 23 result(s)

\smallskip

\noindent\textbf{\texttt{sum}}\quad\textit{defined-fn}\\
defined by \texttt{sum-zero} \texttt{sum-succ} \\
mentioned by 14 result(s)

\smallskip

\noindent\textbf{\texttt{sum-ag}}\quad\textit{defined-fn}\\
defined by \texttt{sum-ag-zero} \texttt{sum-ag-succ} \\
mentioned by 26 result(s)

\smallskip

\noindent\textbf{\texttt{sum-set}}\quad\textit{operator}\\
a kernel term-former \\
mentioned by 11 result(s)

\smallskip

\noindent\textbf{\texttt{summable-weight(w)}}\quad\textit{predicate}\\
reads: \$1 is a summable sequence of positive weights \\
defined by \texttt{summable-weight} \\
mentioned by 9 result(s)

\smallskip

\noindent\textbf{\texttt{sums-to(grp, f, r)}}\quad\textit{predicate}\\
reads: \$2 sums to \$3 in \$1 \\
defined by \texttt{sums-to} \\
mentioned by 6 result(s)

\smallskip

\noindent\textbf{\texttt{sup}}\quad\textit{functoid}\\
reads: the least upper bound of \$1 \\
mentioned by 3 result(s)

\smallskip

\noindent\textbf{\texttt{sup-ord}}\quad\textit{operator}\\
a kernel term-former \\
mentioned by 9 result(s)

\smallskip

\noindent\textbf{\texttt{supp(a, m, f)}}\quad\textit{functoid}\\
reads: the support of \$3 \\
= \{x\_ in carr(m): not(f(x\_) = zero(a))\}

\smallskip


\section*{T}
\addcontentsline{toc}{section}{T}

\noindent\textbf{\texttt{ta}}\quad\textit{tactic}\\
\texttt{(ta 'name)} -- Theorem-assumption: bring the named installed theorem into context as an assumption. \\
declared in \texttt{tactics-help.scm}

\smallskip

\noindent\textbf{\texttt{taylor-differentiable(f, a, x, n)}}\quad\textit{predicate}\\
reads: \$1 is \$4 times differentiable from \$2 to \$3 \\
defined by \texttt{taylor-differentiable} \\
mentioned by 9 result(s)

\smallskip

\noindent\textbf{\texttt{taylor-differentiable-v(m, f, a, x, n)}}\quad\textit{predicate}\\
reads: \$2 is \$5 times differentiable from \$3 to \$4, as a curve in \$1 \\
defined by \texttt{taylor-differentiable-v} \\
mentioned by 5 result(s)

\smallskip

\noindent\textbf{\texttt{taylor-poly(f, a, n, x)}}\quad\textit{functoid}\\
= series-partial-sum(vnb-lambda(k, nn, (nth-deriv(f, k))(a) * (x - a) \textasciicircum{} k * recip(factorial(k))), succ(n)) \\
mentioned by 15 result(s)

\smallskip

\noindent\textbf{\texttt{taylor-poly-v}}\quad\textit{defined-fn}\\
defined by \texttt{taylor-poly-v-zero} \texttt{taylor-poly-v-succ} \\
mentioned by 10 result(s)

\smallskip

\noindent\textbf{\texttt{te}}\quad\textit{tactic}\\
\texttt{(te hyp k)} -- Tuple elim: project the k-th component out of a tuple-membership assumption. \\
declared in \texttt{tactics-help.scm}

\smallskip

\noindent\textbf{\texttt{ti}}\quad\textit{tactic}\\
\texttt{(ti)} -- Tuple intro: prove a tuple/LIST membership component-wise. \\
declared in \texttt{tactics-help.scm}

\smallskip

\noindent\textbf{\texttt{top-space}}\quad\textit{structure}\\
predicate \texttt{is-top-space}; slots: pts opens

\smallskip

\noindent\textbf{\texttt{totally-bounded(s)}}\quad\textit{predicate}\\
reads: \$1 is totally bounded \\
defined by \texttt{totally-bounded} \\
mentioned by 11 result(s)

\smallskip

\noindent\textbf{\texttt{trinum}}\quad\textit{functoid}\\
reads: the \$1-th triangular number \\
defined by \texttt{trinum-zero} \texttt{trinum-succ} \\
mentioned by 8 result(s)

\smallskip

\noindent\textbf{\texttt{tuples}}\quad\textit{operator}\\
a kernel term-former \\
mentioned by 7 result(s)

\smallskip


\section*{U}
\addcontentsline{toc}{section}{U}

\noindent\textbf{\texttt{ue}}\quad\textit{tactic}\\
\texttt{(ue hyp)} -- Union elim: split a cited (IN x (UNION ...)) membership assumption into one subgoal per set -- union-side case analysis, the dual of `ui'. \\
declared in \texttt{tactics-help.scm}

\smallskip

\noindent\textbf{\texttt{ui}}\quad\textit{tactic}\\
\texttt{(ui k)} -- Union intro: reduce a goal (IN x (UNION ...)) to membership of x in the k-th set (1-based). \\
declared in \texttt{tactics-help.scm}

\smallskip

\noindent\textbf{\texttt{union}}\quad\textit{operator}\\
a kernel term-former \\
mentioned by 39 result(s)

\smallskip

\noindent\textbf{\texttt{unitrow(a, n, i)}}\quad\textit{functoid}\\
= matof(1, n, vnb-lambda([rw, cl], cartesian(interval(1, 1), interval(1, n)), if(cl = i, one(a), zero(a)))) \\
mentioned by 7 result(s)

\smallskip


\section*{V}
\addcontentsline{toc}{section}{V}

\noindent\textbf{\texttt{vadd(s)}}\quad\textit{accessor}\\
a structure slot \\
mentioned by 268 result(s)

\smallskip

\noindent\textbf{\texttt{vec(s)}}\quad\textit{accessor}\\
a structure slot \\
mentioned by 389 result(s)

\smallskip

\noindent\textbf{\texttt{vector-space}}\quad\textit{refinement}\\
predicate \texttt{is-vector-space}; same shape as module

\smallskip

\noindent\textbf{\texttt{vlet}}\quad\textit{tactic}\\
\texttt{(vlet (n1 ...) FORMER)} -- Bind proof-object names from the current state.  FORMER is (match PATTERN) -- bind each name to its slot in a context formula matching PATTERN -- or (choice [v body]) -- eliminate an in-context existential and bind its witness. \\
declared in \texttt{tactics-help.scm}

\smallskip

\noindent\textbf{\texttt{vnb-lambda}}\quad\textit{operator}\\
a kernel term-former \\
mentioned by 217 result(s)

\smallskip

\noindent\textbf{\texttt{vneg(s)}}\quad\textit{accessor}\\
a structure slot \\
mentioned by 50 result(s)

\smallskip

\noindent\textbf{\texttt{vnrm(s)}}\quad\textit{accessor}\\
a structure slot \\
mentioned by 35 result(s)

\smallskip

\noindent\textbf{\texttt{vzero(s)}}\quad\textit{accessor}\\
a structure slot \\
mentioned by 172 result(s)

\smallskip


\section*{W}
\addcontentsline{toc}{section}{W}

\noindent\textbf{\texttt{wbc}}\quad\textit{tactic}\\
\texttt{(wbc ['name])} -- Witness-shape backchain: on an `exists v. ...' goal, cite a PSS lemma that MANUFACTURES a witness of that shape, so you can finish with inst+/grind/ew. \\
declared in \texttt{tactics-help.scm}

\smallskip

\noindent\textbf{\texttt{wk}}\quad\textit{tactic}\\
\texttt{(wk hyp)} -- Weaken: drop a cited assumption from the context to tidy the hypothesis list.  `hyp' may be a formula, a "string", or a 1-based assumption index. \\
declared in \texttt{tactics-help.scm}

\smallskip


\section*{Z}
\addcontentsline{toc}{section}{Z}

\noindent\textbf{\texttt{zero(s)}}\quad\textit{accessor}\\
a structure slot \\
mentioned by 385 result(s)

\smallskip

\noindent\textbf{\texttt{zero-ring}}\quad\textit{defined-fn}\\
defined by \texttt{zero-ring-def}

\smallskip

\noindent\textbf{\texttt{zeromat(a, m, n)}}\quad\textit{functoid}\\
= matof(m, n, vnb-lambda([i, j], cartesian(interval(1, m), interval(1, n)), zero(a))) \\
mentioned by 15 result(s)

\smallskip

\noindent\textbf{\texttt{zkept}}\quad\textit{defined-fn}\\
defined by \texttt{zkept-zero} \texttt{zkept-succ} \texttt{zkept-limit} \\
mentioned by 9 result(s)

\smallskip

\noindent\textbf{\texttt{zup}}\quad\textit{defined-fn}\\
defined by \texttt{zup-zero} \texttt{zup-succ} \texttt{zup-limit} \\
mentioned by 9 result(s)

\smallskip

\noindent\textbf{\texttt{zz-act}}\quad\textit{defined-fn}\\
defined by \texttt{zz-act-nonneg} \texttt{zz-act-neg} \\
mentioned by 52 result(s)

\smallskip

\noindent\textbf{\texttt{zz-bezout-set(a, b)}}\quad\textit{functoid}\\
= \{z in zz: forsome([x\_ in zz, y\_ in zz], z = x\_ * a + y\_ * b)\} \\
mentioned by 4 result(s)

\smallskip

\noindent\textbf{\texttt{zz-coprime(a, b)}}\quad\textit{predicate}\\
reads: \$1 and \$2 are coprime \\
defined by \texttt{zz-coprime} \\
mentioned by 2 result(s)

\smallskip

\noindent\textbf{\texttt{zz-divides(a, b)}}\quad\textit{predicate}\\
reads: \$1 divides \$2 \\
defined by \texttt{zz-divides} \\
mentioned by 5 result(s)

\smallskip

\noindent\textbf{\texttt{zz-is-gcd(a, b, d)}}\quad\textit{predicate}\\
reads: \$3 is a greatest common divisor of \$1 and \$2 \\
defined by \texttt{zz-is-gcd} \\
mentioned by 4 result(s)

\smallskip

\noindent\textbf{\texttt{zz-ring}}\quad\textit{instance}\\
an instance of ring; predicate \texttt{is-zz-ring}; same shape as euclidean-ring

\smallskip

\endgroup


%% =========================================================

\begin{thebibliography}{9}
\bibitem{IMPS}
W.~M.~Farmer, J.~D.~Guttman, F.~J.~Thayer,
\textit{IMPS: An Interactive Mathematical Proof System},
Journal of Automated Reasoning 11(2), 1993.

\bibitem{Monk1988}
L.~G.~Monk,
Inference rules using local contexts,
\textit{Journal of Automated Reasoning}, 4:445--462, 1988.
\end{thebibliography}

%% =========================================================
%% The index.  Entries are emitted by the GENERATORS (gen-tactics-ref.py,
%% gen-glossary.py), one per tactic and per glossary name, so the index
%% cannot drift from the registries the way a hand-kept one would.
%% `make' runs makeindex between the two pdflatex passes.
\printindex

\end{document}

